\documentclass[10pt,oneside]{report}
\usepackage[T1]{fontenc}
\usepackage[utf8]{inputenc}
\usepackage{lmodern}
\usepackage[letterpaper,margin=2cm,headheight=14pt]{geometry}
\usepackage{microtype}
\usepackage{xcolor}
\usepackage{array,longtable,booktabs,tabularx}
\usepackage{enumitem}
\usepackage{amssymb}
\usepackage{titlesec}
\usepackage{fancyhdr}
\usepackage{parskip}
\usepackage{multicol}
\usepackage[most]{tcolorbox}
\usepackage[colorlinks=true,linkcolor=blue!55!black,urlcolor=teal!55!black,
            bookmarksnumbered=false,bookmarksopen=true,bookmarksopenlevel=1,
            pdftitle={Learning to Build Agentic Systems: An instructional compilation guide built from public university courses},
            pdfauthor={J. A. Arroyo}]{hyperref}
\usepackage{bookmark}
\urlstyle{same}
\setcounter{secnumdepth}{-1}
\setcounter{tocdepth}{1}
\setlist{nosep,leftmargin=*}
\emergencystretch=2em
\setlength{\LTpre}{4pt}\setlength{\LTpost}{4pt}
\titleformat{\part}[display]{\centering\Huge\bfseries}{}{0pt}{}
\titleformat{\chapter}[hang]{\LARGE\bfseries\color{blue!40!black}}{}{0pt}{}
\titlespacing*{\chapter}{0pt}{0pt}{10pt}
\titleformat{\section}{\large\bfseries\color{blue!40!black}}{}{0pt}{}
\titlespacing*{\section}{0pt}{10pt}{4pt}
\titleformat{\subsection}{\normalsize\bfseries}{}{0pt}{}
\titlespacing*{\subsection}{0pt}{8pt}{2pt}
\pagestyle{fancy}
\fancyhf{}
\fancyhead[L]{\small\nouppercase{\leftmark}}
\fancyhead[R]{\small\hyperlink{toc}{Contents} $\cdot$ \hyperlink{navmatrix}{Matrix} $\cdot$ \thepage}
\renewcommand{\chaptermark}[1]{\markboth{#1}{}}
\renewcommand{\headrulewidth}{0.3pt}
\newcolumntype{P}[1]{>{\raggedright\arraybackslash}p{#1}}
\newcommand{\tg}[1]{{\setlength{\fboxsep}{1.2pt}\colorbox{tag#1}{\kern1pt\textcolor{white}{\sffamily\bfseries\fontsize{5.5}{6}\selectfont\strut #1}\kern1.5pt}}}
\newtcolorbox{infobox}[1][]{colback=blue!4!white,colframe=blue!40!black,boxrule=0.5pt,arc=1mm,left=4pt,right=4pt,top=3pt,bottom=3pt,fonttitle=\bfseries\small,#1}

\definecolor{tagFND}{HTML}{4E79A7}
\definecolor{tagARCH}{HTML}{59A14F}
\definecolor{tagTOOL}{HTML}{F28E2B}
\definecolor{tagMEM}{HTML}{B07AA1}
\definecolor{tagRSN}{HTML}{E15759}
\definecolor{tagMULTI}{HTML}{76B7B2}
\definecolor{tagCODE}{HTML}{EDC948}
\definecolor{tagWEB}{HTML}{FF9DA7}
\definecolor{tagEVAL}{HTML}{9C755F}
\definecolor{tagTRAIN}{HTML}{D37295}
\definecolor{tagSAFE}{HTML}{E03C31}
\definecolor{tagDEPLOY}{HTML}{499894}
\definecolor{tagAPP}{HTML}{86BCB6}
\definecolor{tagORG}{HTML}{BAB0AC}
\begin{document}
\begin{titlepage}\centering
\vspace*{3cm}{\Huge\bfseries Learning to Build Agentic Systems\par}\vspace{0.6cm}
{\Large An instructional compilation guide built from public university courses\par}\vspace{0.8cm}
{\large J. A. Arroyo\par}\vspace{0.8cm}
{\large 15 courses $\cdot$ 283 lectures $\cdot$ 98 assignments $\cdot$ 15 modules $\cdot$ 3 learning tracks\par}\vspace{1.5cm}
\begin{minipage}{0.8\linewidth}\centering\small
CMU 11-768 $\cdot$ Stanford CS146S (2025, 2026) $\cdot$ Northeastern Agentic AI $\cdot$ Buffalo CSE 410/510 $\cdot$ UVA CS 6501 $\cdot$ UW--Madison CS 839 $\cdot$ JHU Machine Programming $\cdot$ CU Boulder CSCI 7000 $\cdot$ Mannheim IE 685/686 $\cdot$ UChicago AI Agents for Social Science\\[4pt]
plus MCP, Anthropic Engineering, OpenClaw, Berkeley TinyAgent and Colab
\end{minipage}\vfill
{\small Snapshot compiled October 9, 2026 from the \emph{CMU-Agentic} bookmarks folder plus a targeted course search.\\ All materials belong to their respective authors. This guide only links to them.\\[6pt]
\emph{Use of AI tools.} I am the sole author. Claude (Anthropic, through Claude Code) searched for the courses, read their public pages and drafted the summaries and structure under my direction. The descriptions are summaries, each linked to its source; titles are the sources' own. I am responsible for the content.\par}
\end{titlepage}
\hypertarget{toc}{}\tableofcontents
\part{Orientation}
\chapter{How to Use This Guide}
This guide compiles 15 public university courses on LLM agents and agentic systems, plus 8 supplementary resources, into one curriculum you can navigate. Most come from your \emph{CMU-Agentic} bookmarks folder; the rest were added after a targeted search for courses strong in evaluation, training and security.
Every one of the 283 catalogued lectures and 98 assignments is classified with one shared set of topic tags, so the same subject can be followed across courses.
\begin{infobox}[title=One set of three documents]
\begin{enumerate}
\item \textbf{Learning to Build Agentic Systems} (this guide): the courses, lectures and assignments behind the practices. \emph{Use it to learn a topic.}
\item \textbf{Agentic System Design Process} (\texttt{design\_process/design\_process.pdf}): nine phases from framing a problem to running it in production, each citing the courses indexed here. \emph{Use it to design a system.}
\item \textbf{Screening Tasks for Agentic Automation} (\texttt{design\_process/task\_screening.pdf}): a one-page worksheet, filled in after Phase~1 of the design process, that feeds Phase~2 and sets up the comparison run in Phase~7. \emph{Use it to decide whether a task is worth it.}
\end{enumerate}
The modules map onto the design phases: \hyperlink{mod-ARCH}{\tg{ARCH}} for Phases 2--3, \hyperlink{mod-EVAL}{\tg{EVAL}} for Phases 4 and 7, \hyperlink{mod-TOOL}{\tg{TOOL}} and \hyperlink{mod-MEM}{\tg{MEM}} for Phase~5, \hyperlink{mod-SAFE}{\tg{SAFE}} for Phase~8 and \hyperlink{mod-DEPLOY}{\tg{DEPLOY}} for Phase~9.
\end{infobox}
\begin{infobox}[title=Four ways in]
\begin{itemize}
\item \textbf{Follow a track} (\hyperlink{tracks}{Learning tracks}): a role-based path through the modules, or the \hyperlink{plan16}{16-week self-study plan}.
\item \textbf{Study a topic} (Part~II, \hyperlink{mod-SETUP}{Modules 0--14}): each module gives learning objectives, a short conceptual primer, a \emph{curated} sequence of the best lectures across courses, an automatically generated map of \emph{every} relevant lecture, key readings and hands-on assignments.
\item \textbf{Browse a course} (Part~III, \hyperlink{catalog}{Course catalogue}): each course card is followed by its full lecture table and assignment table with all public links.
\item \textbf{Jump around} with the \hyperlink{navmatrix}{course $\times$ module matrix}. Every coloured tag box, lecture reference and assignment name is a hyperlink, and the page header always links back to Contents and the Matrix.
\end{itemize}
\end{infobox}
\section{Conventions}
\begin{itemize}
\item Lecture references look like \textbf{CMU~2} (course code and session number in that course's own schedule) and link to the row in the catalogue, which has dates, topics, readings and materials.
\item Material links are labelled by type: \texttt{[slides]}, \texttt{[video]}, \texttt{[PDF]}, \texttt{[PPTX]}, \texttt{[GitHub]}, \texttt{[arXiv:\ldots]}, \texttt{[notebook zip]}, or the site's domain.
\item In module lecture maps, $\star$ marks the lectures chosen for the curated sequence.
\item Guest talks and purely logistical sessions stay in the catalogue but are left out of the module maps.
\item In hands-on lists, ``Starter code'' means the assignment has a public repository or download, ``Spec'' links to a written specification with no code, and ``Resources'' links to benchmarks or data sources the project builds on.
\end{itemize}
\section{A note on currency}
The material spans Spring 2024 (TinyAgent) to courses still running in October 2026. Two things change fast. (1)~\textbf{MCP:} the 2026-07-28 revision made the protocol stateless (see \hyperref[sup-mcp]{the MCP supplement}), so earlier lectures teach the older handshake. (2)~\textbf{Model names and SDKs:} courses use whatever was current when they ran. Learn the concepts from the lectures and take API details from current documentation.
\chapter{Tag Legend}
Every lecture and assignment carries one or more tags. Each tag (except \tg{ORG}) corresponds to a module in Part~II. \tg{TOOL} lectures that are mainly about MCP or A2A are collected in their own module.
\begin{longtable}{l P{4cm} P{7.5cm} r}
\toprule Tag & Meaning & Module & Lectures\\ \midrule
\hyperlink{mod-FND}{\tg{FND}} & Foundations & \hyperlink{mod-FND}{Module 1: LLM Foundations and Prompting} & 26\\
\hyperlink{mod-ARCH}{\tg{ARCH}} & Architectures & \hyperlink{mod-ARCH}{Module 2: The Agent Loop and Agent Architectures} & 46\\
\hyperlink{mod-TOOL}{\tg{TOOL}} & Tools & \hyperlink{mod-TOOL}{Module 3: Tool Use and Function Calling}; \hyperlink{mod-MCP}{Module 4: MCP and Agent Interoperability} & 31\\
\hyperlink{mod-MEM}{\tg{MEM}} & Memory/RAG/Context & \hyperlink{mod-MEM}{Module 5: Memory, Retrieval and Context Engineering} & 27\\
\hyperlink{mod-RSN}{\tg{RSN}} & Reasoning/Planning & \hyperlink{mod-RSN}{Module 6: Reasoning, Planning, Reflection and Search} & 50\\
\hyperlink{mod-MULTI}{\tg{MULTI}} & Multi-agent & \hyperlink{mod-MULTI}{Module 7: Multi-Agent Systems} & 28\\
\hyperlink{mod-CODE}{\tg{CODE}} & Coding agents & \hyperlink{mod-CODE}{Module 8: Coding Agents} & 54\\
\hyperlink{mod-WEB}{\tg{WEB}} & Web/GUI/Multimodal & \hyperlink{mod-WEB}{Module 9: Web, GUI, Computer-Use and Multimodal Agents} & 11\\
\hyperlink{mod-EVAL}{\tg{EVAL}} & Evaluation & \hyperlink{mod-EVAL}{Module 10: Evaluating Agents} & 36\\
\hyperlink{mod-TRAIN}{\tg{TRAIN}} & Training & \hyperlink{mod-TRAIN}{Module 11: Training and Optimizing Agents} & 31\\
\hyperlink{mod-SAFE}{\tg{SAFE}} & Safety/Security & \hyperlink{mod-SAFE}{Module 12: Safety, Security and Oversight} & 24\\
\hyperlink{mod-DEPLOY}{\tg{DEPLOY}} & Deployment & \hyperlink{mod-DEPLOY}{Module 13: Deployment, Efficiency and Operations} & 23\\
\hyperlink{mod-APP}{\tg{APP}} & Applications & \hyperlink{mod-APP}{Module 14: Applications: Deep Research, Science and Society} & 40\\
\tg{ORG} & Logistics & -- & 62\\
\bottomrule\end{longtable}
\chapter{The Courses at a Glance}
The courses fall into a few archetypes, and knowing which kind you are reading tells you how to use it:
\begin{itemize}
\item \textbf{Systems and training} (CMU): the backbone, from harnesses to RL, with the best assignments.
\item \textbf{Builder courses} (Northeastern, UVA, Buffalo, Mannheim): framework-based, hands-on and lab-driven.
\item \textbf{Coding-agent practice} (Stanford CS146S, both years): how practitioners work with and around coding agents.
\item \textbf{Research seminars} (UW--Madison, CU Boulder): curated paper lists, for depth and for the research frontier.
\item \textbf{Program synthesis} (JHU): the PL view of code generation, from classic search to LLM coding agents.
\item \textbf{Applied social science} (UChicago): agents as simulated subjects and research instruments.
\end{itemize}
{\footnotesize\begin{longtable}{P{2.2cm} P{3.0cm} P{2.9cm} P{1.6cm} P{5.2cm}}
\toprule Course & Archetype / level & Public materials & Lectures / assignments & Best for\\ \midrule\endhead
\hyperlink{course-CMU}{\textbf{CMU}}\newline 11-768 AI Agents & Systems + training (build, evaluate, RL-train)\newline\emph{Graduate; assumes you have trained neural LMs} & Slides: L1--13\newline Video: L1--9\newline Code: A1--A3 public & 23 / 5 & The backbone course: harness engineering, evals, SFT/RL, safety, search.\\ \midrule
\hyperlink{course-STAN26}{\textbf{Stanford '26}}\newline CS146S The Modern Software Developer & Coding-agent practice (practitioner)\newline\emph{Undergrad/grad; CS111/CS161-level programming} & Slides: Weeks 1--3 so far\newline Video: none\newline Code: Weeks 1--3 assignments & 20 / 3 & How production coding agents work and how to set up a repo and a team around them.\\ \midrule
\hyperlink{course-STAN25}{\textbf{Stanford '25}}\newline CS146S The Modern Software Developer (archive) & Coding-agent practice (practitioner)\newline\emph{Undergrad/grad} & Slides: Most lectures (Google Slides)\newline Video: none\newline Code: 8 weekly assignments (branch fall2025) & 20 / 9 & Slide decks on prompting, building a coding agent, MCP servers, AI testing/security, code review, SRE.\\ \midrule
\hyperlink{course-NEU}{\textbf{Northeastern}}\newline Agentic AI & Builder (LangChain/LangGraph stack)\newline\emph{Gentle on-ramp (Python primer provided)} & Slides: 2 HTML decks\newline Video: none\newline Code: Notebooks W02b--W05b, DADA1 & 19 / 4 & Best guided, notebook-by-notebook path from an LLM API call to a deployed LangGraph agent.\\ \midrule
\hyperlink{course-UB}{\textbf{Buffalo}}\newline CSE 410/510 Introduction to Agentic Engineering & Engineering discipline (framework-agnostic)\newline\emph{Upper-undergrad / MS} & Slides: none\newline Video: none\newline Code: none (lab specs only) & 15 / 11 & Five lab specs and a capstone spec that treat an agent as an engineered, evaluated, bounded system.\\ \midrule
\hyperlink{course-UVA}{\textbf{UVA}}\newline CS 6501 Workshop on Building AI Agents & Hands-on lab workshop (small/open models)\newline\emph{Graduate workshop; Python} & Slides: PPTX per topic\newline Video: 1 + podcasts\newline Code: Starter scripts per topic & 19 / 9 & The richest set of runnable lab exercises: local LLMs, LangGraph, tools, RAG, VLMs, MCP/A2A, Tinker fine-tuning.\\ \midrule
\hyperlink{course-UW}{\textbf{UW-Madison}}\newline CS 839 AI Agents & Research seminar (paper reading)\newline\emph{Graduate research seminar} & Slides: none\newline Video: none\newline Code: Demo notebook; MIT license & 25 / 4 & A curated, critical reading list from ReAct to self-improving and evolutionary agents, plus a no-frameworks coding-agent assignment.\\ \midrule
\hyperlink{course-JHU}{\textbf{JHU}}\newline EN.601.727 Machine Programming & Program synthesis to coding agents\newline\emph{Graduate (PL + ML)} & Slides: L1--12\newline Video: none\newline Code: A1--A3 public & 28 / 5 & Ties classic program synthesis to LLM coding agents, with two MCP lectures and a strong web-app agent assignment.\\ \midrule
\hyperlink{course-CU}{\textbf{CU Boulder}}\newline CSCI 7000-005 Building AI Agents & Research seminar (agentic software engineering)\newline\emph{Graduate seminar; competent SE assumed} & Slides: 4 instructor decks\newline Video: none\newline Code: none & 28 / 4 & Research reading on agentic software engineering: refactoring agents, SWE-agent, skills, memory, MCP design choices.\\ \midrule
\hyperlink{course-MANN}{\textbf{Mannheim}}\newline IE 685/686 Large Language Models and Agents & Applied master's course + team project\newline\emph{Master's; Python + basic ML} & Slides: 8 PDF decks\newline Video: none\newline Code: 5 exercise notebooks (zip) & 12 / 15 & Compact slide decks from LLM basics to context engineering, plus nine well-specified team-project briefs.\\ \midrule
\hyperlink{course-UCHI}{\textbf{UChicago}}\newline AI Agents for Social Science and Society & Applied: agents as research instruments\newline\emph{Grad / advanced undergrad (computational social science)} & Slides: none\newline Video: none\newline Code: 9 weekly notebooks & 11 / 11 & Agents as simulated subjects and research assistants; notebooks on multi-agent simulation, LoRA, RL, steering and red-teaming.\\ \midrule
\hyperlink{course-S329A}{\textbf{Stanford 329A}}\newline CS329A Self-Improving AI Agents & Research course (self-improvement levers) + homeworks\newline\emph{Graduate; CS224N or CS229S} & Slides: none linked\newline Video: 9 recordings (Stanford Online)\newline Code: HW1--HW3 public repos (HW1 MIT, HW3 Apache-2.0) & 20 / 6 & The clearest path through test-time compute, verifiers, RL and open-ended self-improvement, with three public homeworks.\\ \midrule
\hyperlink{course-BRKAG}{\textbf{Berkeley '25}}\newline CS194/294-196 Agentic AI & Guest-lecture series + AgentBeats project (also a MOOC)\newline\emph{Undergraduate and graduate; ML/DL background} & Slides: 9 of 13 sessions\newline Video: 12 of 13 sessions\newline Code: AgentBeats tutorial repo (no license stated) & 13 / 4 & Industry and lab speakers on post-training, agent evaluation and deployment, and an evaluator-agent project (AgentBeats).\\ \midrule
\hyperlink{course-BRKADV}{\textbf{Berkeley Adv.}}\newline CS294/194-280 Advanced Large Language Model Agents & Guest-lecture series + Lean labs (also a MOOC)\newline\emph{Upper-division undergraduate and graduate} & Slides: All 12 lectures\newline Video: All 12 lectures\newline Code: 2 Lean labs in a public Drive folder (no license) & 12 / 6 & Research talks on reasoning, post-training, coding and web agents and formal math, plus a compiler-verified agent lab.\\ \midrule
\hyperlink{course-S329T}{\textbf{Stanford 329T}}\newline CS329T Trustworthy ML: Building and Evaluating Agentic Systems & Project course on agent evaluation\newline\emph{Graduate; CS229 + deep learning} & Slides: 11 decks (10 sessions)\newline Video: none (not recorded)\newline Code: HW1 handout PDF; starter branch 404 & 18 / 2 & Evaluation-first agent building: Goal-Plan-Action evaluation, LLM judges, MCP tool descriptions, GEPA and red-teaming.\\ \midrule
\end{longtable}}
\chapter{Course $\times$ Module Navigation Matrix}\hypertarget{navmatrix}{}
Each cell is the number of lectures from a course in a module. Click a number to open the module, or a course name to open its catalogue entry. Column headers are module numbers:
\hyperlink{mod-FND}{1~LLM Foundations and Prompting}; \hyperlink{mod-ARCH}{2~The Agent Loop and Agent Architectures}; \hyperlink{mod-TOOL}{3~Tool Use and Function Calling}; \hyperlink{mod-MCP}{4~MCP and Agent Interoperability}; \hyperlink{mod-MEM}{5~Memory}; \hyperlink{mod-RSN}{6~Reasoning}; \hyperlink{mod-MULTI}{7~Multi-Agent Systems}; \hyperlink{mod-CODE}{8~Coding Agents}; \hyperlink{mod-WEB}{9~Web}; \hyperlink{mod-EVAL}{10~Evaluating Agents}; \hyperlink{mod-TRAIN}{11~Training and Optimizing Agents}; \hyperlink{mod-SAFE}{12~Safety}; \hyperlink{mod-DEPLOY}{13~Deployment}; \hyperlink{mod-APP}{14~Applications}.
{\small\setlength{\tabcolsep}{4pt}\begin{longtable}{lccccccccccccccr}
\toprule Course & \hyperlink{mod-FND}{\textbf{1}} & \hyperlink{mod-ARCH}{\textbf{2}} & \hyperlink{mod-TOOL}{\textbf{3}} & \hyperlink{mod-MCP}{\textbf{4}} & \hyperlink{mod-MEM}{\textbf{5}} & \hyperlink{mod-RSN}{\textbf{6}} & \hyperlink{mod-MULTI}{\textbf{7}} & \hyperlink{mod-CODE}{\textbf{8}} & \hyperlink{mod-WEB}{\textbf{9}} & \hyperlink{mod-EVAL}{\textbf{10}} & \hyperlink{mod-TRAIN}{\textbf{11}} & \hyperlink{mod-SAFE}{\textbf{12}} & \hyperlink{mod-DEPLOY}{\textbf{13}} & \hyperlink{mod-APP}{\textbf{14}} & All\\ \midrule
\hyperlink{course-CMU}{CMU} & \textcolor{black!25}{0} & \hyperlink{mod-ARCH}{{\setlength{\fboxsep}{1.5pt}\colorbox{blue!26}{\makebox[1.1em]{3}}}} & \hyperlink{mod-TOOL}{{\setlength{\fboxsep}{1.5pt}\colorbox{blue!12}{\makebox[1.1em]{1}}}} & \textcolor{black!25}{0} & \hyperlink{mod-MEM}{{\setlength{\fboxsep}{1.5pt}\colorbox{blue!26}{\makebox[1.1em]{3}}}} & \hyperlink{mod-RSN}{{\setlength{\fboxsep}{1.5pt}\colorbox{blue!26}{\makebox[1.1em]{3}}}} & \hyperlink{mod-MULTI}{{\setlength{\fboxsep}{1.5pt}\colorbox{blue!19}{\makebox[1.1em]{2}}}} & \hyperlink{mod-CODE}{{\setlength{\fboxsep}{1.5pt}\colorbox{blue!19}{\makebox[1.1em]{2}}}} & \hyperlink{mod-WEB}{{\setlength{\fboxsep}{1.5pt}\colorbox{blue!12}{\makebox[1.1em]{1}}}} & \hyperlink{mod-EVAL}{{\setlength{\fboxsep}{1.5pt}\colorbox{blue!19}{\makebox[1.1em]{2}}}} & \hyperlink{mod-TRAIN}{{\setlength{\fboxsep}{1.5pt}\colorbox{blue!33}{\makebox[1.1em]{4}}}} & \hyperlink{mod-SAFE}{{\setlength{\fboxsep}{1.5pt}\colorbox{blue!19}{\makebox[1.1em]{2}}}} & \hyperlink{mod-DEPLOY}{{\setlength{\fboxsep}{1.5pt}\colorbox{blue!19}{\makebox[1.1em]{2}}}} & \hyperlink{mod-APP}{{\setlength{\fboxsep}{1.5pt}\colorbox{blue!26}{\makebox[1.1em]{3}}}} & 23\\
\hyperlink{course-STAN26}{Stanford '26} & \hyperlink{mod-FND}{{\setlength{\fboxsep}{1.5pt}\colorbox{blue!19}{\makebox[1.1em]{2}}}} & \hyperlink{mod-ARCH}{{\setlength{\fboxsep}{1.5pt}\colorbox{blue!26}{\makebox[1.1em]{3}}}} & \hyperlink{mod-TOOL}{{\setlength{\fboxsep}{1.5pt}\colorbox{blue!12}{\makebox[1.1em]{1}}}} & \hyperlink{mod-MCP}{{\setlength{\fboxsep}{1.5pt}\colorbox{blue!12}{\makebox[1.1em]{1}}}} & \hyperlink{mod-MEM}{{\setlength{\fboxsep}{1.5pt}\colorbox{blue!19}{\makebox[1.1em]{2}}}} & \textcolor{black!25}{0} & \hyperlink{mod-MULTI}{{\setlength{\fboxsep}{1.5pt}\colorbox{blue!12}{\makebox[1.1em]{1}}}} & \hyperlink{mod-CODE}{{\setlength{\fboxsep}{1.5pt}\colorbox{blue!40}{\makebox[1.1em]{5}}}} & \textcolor{black!25}{0} & \hyperlink{mod-EVAL}{{\setlength{\fboxsep}{1.5pt}\colorbox{blue!19}{\makebox[1.1em]{2}}}} & \textcolor{black!25}{0} & \hyperlink{mod-SAFE}{{\setlength{\fboxsep}{1.5pt}\colorbox{blue!12}{\makebox[1.1em]{1}}}} & \hyperlink{mod-DEPLOY}{{\setlength{\fboxsep}{1.5pt}\colorbox{blue!26}{\makebox[1.1em]{3}}}} & \textcolor{black!25}{0} & 20\\
\hyperlink{course-STAN25}{Stanford '25} & \hyperlink{mod-FND}{{\setlength{\fboxsep}{1.5pt}\colorbox{blue!19}{\makebox[1.1em]{2}}}} & \hyperlink{mod-ARCH}{{\setlength{\fboxsep}{1.5pt}\colorbox{blue!12}{\makebox[1.1em]{1}}}} & \hyperlink{mod-TOOL}{{\setlength{\fboxsep}{1.5pt}\colorbox{blue!12}{\makebox[1.1em]{1}}}} & \hyperlink{mod-MCP}{{\setlength{\fboxsep}{1.5pt}\colorbox{blue!12}{\makebox[1.1em]{1}}}} & \hyperlink{mod-MEM}{{\setlength{\fboxsep}{1.5pt}\colorbox{blue!12}{\makebox[1.1em]{1}}}} & \hyperlink{mod-RSN}{{\setlength{\fboxsep}{1.5pt}\colorbox{blue!12}{\makebox[1.1em]{1}}}} & \hyperlink{mod-MULTI}{{\setlength{\fboxsep}{1.5pt}\colorbox{blue!19}{\makebox[1.1em]{2}}}} & \hyperlink{mod-CODE}{{\setlength{\fboxsep}{1.5pt}\colorbox{blue!45}{\makebox[1.1em]{8}}}} & \textcolor{black!25}{0} & \hyperlink{mod-EVAL}{{\setlength{\fboxsep}{1.5pt}\colorbox{blue!19}{\makebox[1.1em]{2}}}} & \textcolor{black!25}{0} & \hyperlink{mod-SAFE}{{\setlength{\fboxsep}{1.5pt}\colorbox{blue!12}{\makebox[1.1em]{1}}}} & \hyperlink{mod-DEPLOY}{{\setlength{\fboxsep}{1.5pt}\colorbox{blue!19}{\makebox[1.1em]{2}}}} & \hyperlink{mod-APP}{{\setlength{\fboxsep}{1.5pt}\colorbox{blue!12}{\makebox[1.1em]{1}}}} & 20\\
\hyperlink{course-NEU}{Northeastern} & \hyperlink{mod-FND}{{\setlength{\fboxsep}{1.5pt}\colorbox{blue!26}{\makebox[1.1em]{3}}}} & \hyperlink{mod-ARCH}{{\setlength{\fboxsep}{1.5pt}\colorbox{blue!45}{\makebox[1.1em]{6}}}} & \hyperlink{mod-TOOL}{{\setlength{\fboxsep}{1.5pt}\colorbox{blue!19}{\makebox[1.1em]{2}}}} & \hyperlink{mod-MCP}{{\setlength{\fboxsep}{1.5pt}\colorbox{blue!19}{\makebox[1.1em]{2}}}} & \hyperlink{mod-MEM}{{\setlength{\fboxsep}{1.5pt}\colorbox{blue!45}{\makebox[1.1em]{6}}}} & \hyperlink{mod-RSN}{{\setlength{\fboxsep}{1.5pt}\colorbox{blue!33}{\makebox[1.1em]{4}}}} & \hyperlink{mod-MULTI}{{\setlength{\fboxsep}{1.5pt}\colorbox{blue!26}{\makebox[1.1em]{3}}}} & \textcolor{black!25}{0} & \textcolor{black!25}{0} & \hyperlink{mod-EVAL}{{\setlength{\fboxsep}{1.5pt}\colorbox{blue!12}{\makebox[1.1em]{1}}}} & \hyperlink{mod-TRAIN}{{\setlength{\fboxsep}{1.5pt}\colorbox{blue!12}{\makebox[1.1em]{1}}}} & \hyperlink{mod-SAFE}{{\setlength{\fboxsep}{1.5pt}\colorbox{blue!19}{\makebox[1.1em]{2}}}} & \hyperlink{mod-DEPLOY}{{\setlength{\fboxsep}{1.5pt}\colorbox{blue!40}{\makebox[1.1em]{5}}}} & \textcolor{black!25}{0} & 19\\
\hyperlink{course-UB}{Buffalo} & \hyperlink{mod-FND}{{\setlength{\fboxsep}{1.5pt}\colorbox{blue!12}{\makebox[1.1em]{1}}}} & \hyperlink{mod-ARCH}{{\setlength{\fboxsep}{1.5pt}\colorbox{blue!19}{\makebox[1.1em]{2}}}} & \hyperlink{mod-TOOL}{{\setlength{\fboxsep}{1.5pt}\colorbox{blue!12}{\makebox[1.1em]{1}}}} & \textcolor{black!25}{0} & \hyperlink{mod-MEM}{{\setlength{\fboxsep}{1.5pt}\colorbox{blue!26}{\makebox[1.1em]{3}}}} & \hyperlink{mod-RSN}{{\setlength{\fboxsep}{1.5pt}\colorbox{blue!19}{\makebox[1.1em]{2}}}} & \hyperlink{mod-MULTI}{{\setlength{\fboxsep}{1.5pt}\colorbox{blue!12}{\makebox[1.1em]{1}}}} & \hyperlink{mod-CODE}{{\setlength{\fboxsep}{1.5pt}\colorbox{blue!12}{\makebox[1.1em]{1}}}} & \textcolor{black!25}{0} & \hyperlink{mod-EVAL}{{\setlength{\fboxsep}{1.5pt}\colorbox{blue!12}{\makebox[1.1em]{1}}}} & \textcolor{black!25}{0} & \hyperlink{mod-SAFE}{{\setlength{\fboxsep}{1.5pt}\colorbox{blue!19}{\makebox[1.1em]{2}}}} & \hyperlink{mod-DEPLOY}{{\setlength{\fboxsep}{1.5pt}\colorbox{blue!19}{\makebox[1.1em]{2}}}} & \hyperlink{mod-APP}{{\setlength{\fboxsep}{1.5pt}\colorbox{blue!12}{\makebox[1.1em]{1}}}} & 15\\
\hyperlink{course-UVA}{UVA} & \hyperlink{mod-FND}{{\setlength{\fboxsep}{1.5pt}\colorbox{blue!19}{\makebox[1.1em]{2}}}} & \hyperlink{mod-ARCH}{{\setlength{\fboxsep}{1.5pt}\colorbox{blue!33}{\makebox[1.1em]{4}}}} & \hyperlink{mod-TOOL}{{\setlength{\fboxsep}{1.5pt}\colorbox{blue!26}{\makebox[1.1em]{3}}}} & \hyperlink{mod-MCP}{{\setlength{\fboxsep}{1.5pt}\colorbox{blue!12}{\makebox[1.1em]{1}}}} & \hyperlink{mod-MEM}{{\setlength{\fboxsep}{1.5pt}\colorbox{blue!19}{\makebox[1.1em]{2}}}} & \hyperlink{mod-RSN}{{\setlength{\fboxsep}{1.5pt}\colorbox{blue!33}{\makebox[1.1em]{4}}}} & \hyperlink{mod-MULTI}{{\setlength{\fboxsep}{1.5pt}\colorbox{blue!33}{\makebox[1.1em]{4}}}} & \hyperlink{mod-CODE}{{\setlength{\fboxsep}{1.5pt}\colorbox{blue!12}{\makebox[1.1em]{1}}}} & \hyperlink{mod-WEB}{{\setlength{\fboxsep}{1.5pt}\colorbox{blue!19}{\makebox[1.1em]{2}}}} & \hyperlink{mod-EVAL}{{\setlength{\fboxsep}{1.5pt}\colorbox{blue!26}{\makebox[1.1em]{3}}}} & \hyperlink{mod-TRAIN}{{\setlength{\fboxsep}{1.5pt}\colorbox{blue!19}{\makebox[1.1em]{2}}}} & \hyperlink{mod-SAFE}{{\setlength{\fboxsep}{1.5pt}\colorbox{blue!12}{\makebox[1.1em]{1}}}} & \hyperlink{mod-DEPLOY}{{\setlength{\fboxsep}{1.5pt}\colorbox{blue!12}{\makebox[1.1em]{1}}}} & \hyperlink{mod-APP}{{\setlength{\fboxsep}{1.5pt}\colorbox{blue!26}{\makebox[1.1em]{3}}}} & 19\\
\hyperlink{course-UW}{UW-Madison} & \hyperlink{mod-FND}{{\setlength{\fboxsep}{1.5pt}\colorbox{blue!12}{\makebox[1.1em]{1}}}} & \hyperlink{mod-ARCH}{{\setlength{\fboxsep}{1.5pt}\colorbox{blue!45}{\makebox[1.1em]{6}}}} & \hyperlink{mod-TOOL}{{\setlength{\fboxsep}{1.5pt}\colorbox{blue!19}{\makebox[1.1em]{2}}}} & \textcolor{black!25}{0} & \hyperlink{mod-MEM}{{\setlength{\fboxsep}{1.5pt}\colorbox{blue!26}{\makebox[1.1em]{3}}}} & \hyperlink{mod-RSN}{{\setlength{\fboxsep}{1.5pt}\colorbox{blue!40}{\makebox[1.1em]{5}}}} & \hyperlink{mod-MULTI}{{\setlength{\fboxsep}{1.5pt}\colorbox{blue!40}{\makebox[1.1em]{5}}}} & \hyperlink{mod-CODE}{{\setlength{\fboxsep}{1.5pt}\colorbox{blue!33}{\makebox[1.1em]{4}}}} & \hyperlink{mod-WEB}{{\setlength{\fboxsep}{1.5pt}\colorbox{blue!12}{\makebox[1.1em]{1}}}} & \hyperlink{mod-EVAL}{{\setlength{\fboxsep}{1.5pt}\colorbox{blue!40}{\makebox[1.1em]{5}}}} & \hyperlink{mod-TRAIN}{{\setlength{\fboxsep}{1.5pt}\colorbox{blue!19}{\makebox[1.1em]{2}}}} & \hyperlink{mod-SAFE}{{\setlength{\fboxsep}{1.5pt}\colorbox{blue!26}{\makebox[1.1em]{3}}}} & \textcolor{black!25}{0} & \hyperlink{mod-APP}{{\setlength{\fboxsep}{1.5pt}\colorbox{blue!33}{\makebox[1.1em]{4}}}} & 25\\
\hyperlink{course-JHU}{JHU} & \hyperlink{mod-FND}{{\setlength{\fboxsep}{1.5pt}\colorbox{blue!19}{\makebox[1.1em]{2}}}} & \hyperlink{mod-ARCH}{{\setlength{\fboxsep}{1.5pt}\colorbox{blue!33}{\makebox[1.1em]{4}}}} & \textcolor{black!25}{0} & \hyperlink{mod-MCP}{{\setlength{\fboxsep}{1.5pt}\colorbox{blue!19}{\makebox[1.1em]{2}}}} & \textcolor{black!25}{0} & \hyperlink{mod-RSN}{{\setlength{\fboxsep}{1.5pt}\colorbox{blue!40}{\makebox[1.1em]{5}}}} & \textcolor{black!25}{0} & \hyperlink{mod-CODE}{{\setlength{\fboxsep}{1.5pt}\colorbox{blue!45}{\makebox[1.1em]{8}}}} & \textcolor{black!25}{0} & \textcolor{black!25}{0} & \hyperlink{mod-TRAIN}{{\setlength{\fboxsep}{1.5pt}\colorbox{blue!33}{\makebox[1.1em]{4}}}} & \hyperlink{mod-SAFE}{{\setlength{\fboxsep}{1.5pt}\colorbox{blue!12}{\makebox[1.1em]{1}}}} & \textcolor{black!25}{0} & \hyperlink{mod-APP}{{\setlength{\fboxsep}{1.5pt}\colorbox{blue!45}{\makebox[1.1em]{7}}}} & 28\\
\hyperlink{course-CU}{CU Boulder} & \textcolor{black!25}{0} & \hyperlink{mod-ARCH}{{\setlength{\fboxsep}{1.5pt}\colorbox{blue!45}{\makebox[1.1em]{6}}}} & \hyperlink{mod-TOOL}{{\setlength{\fboxsep}{1.5pt}\colorbox{blue!40}{\makebox[1.1em]{5}}}} & \hyperlink{mod-MCP}{{\setlength{\fboxsep}{1.5pt}\colorbox{blue!12}{\makebox[1.1em]{1}}}} & \hyperlink{mod-MEM}{{\setlength{\fboxsep}{1.5pt}\colorbox{blue!12}{\makebox[1.1em]{1}}}} & \hyperlink{mod-RSN}{{\setlength{\fboxsep}{1.5pt}\colorbox{blue!33}{\makebox[1.1em]{4}}}} & \hyperlink{mod-MULTI}{{\setlength{\fboxsep}{1.5pt}\colorbox{blue!33}{\makebox[1.1em]{4}}}} & \hyperlink{mod-CODE}{{\setlength{\fboxsep}{1.5pt}\colorbox{blue!45}{\makebox[1.1em]{10}}}} & \hyperlink{mod-WEB}{{\setlength{\fboxsep}{1.5pt}\colorbox{blue!12}{\makebox[1.1em]{1}}}} & \hyperlink{mod-EVAL}{{\setlength{\fboxsep}{1.5pt}\colorbox{blue!26}{\makebox[1.1em]{3}}}} & \hyperlink{mod-TRAIN}{{\setlength{\fboxsep}{1.5pt}\colorbox{blue!12}{\makebox[1.1em]{1}}}} & \hyperlink{mod-SAFE}{{\setlength{\fboxsep}{1.5pt}\colorbox{blue!19}{\makebox[1.1em]{2}}}} & \hyperlink{mod-DEPLOY}{{\setlength{\fboxsep}{1.5pt}\colorbox{blue!12}{\makebox[1.1em]{1}}}} & \textcolor{black!25}{0} & 28\\
\hyperlink{course-MANN}{Mannheim} & \hyperlink{mod-FND}{{\setlength{\fboxsep}{1.5pt}\colorbox{blue!26}{\makebox[1.1em]{3}}}} & \hyperlink{mod-ARCH}{{\setlength{\fboxsep}{1.5pt}\colorbox{blue!19}{\makebox[1.1em]{2}}}} & \textcolor{black!25}{0} & \hyperlink{mod-MCP}{{\setlength{\fboxsep}{1.5pt}\colorbox{blue!12}{\makebox[1.1em]{1}}}} & \hyperlink{mod-MEM}{{\setlength{\fboxsep}{1.5pt}\colorbox{blue!19}{\makebox[1.1em]{2}}}} & \hyperlink{mod-RSN}{{\setlength{\fboxsep}{1.5pt}\colorbox{blue!12}{\makebox[1.1em]{1}}}} & \textcolor{black!25}{0} & \textcolor{black!25}{0} & \textcolor{black!25}{0} & \hyperlink{mod-EVAL}{{\setlength{\fboxsep}{1.5pt}\colorbox{blue!12}{\makebox[1.1em]{1}}}} & \hyperlink{mod-TRAIN}{{\setlength{\fboxsep}{1.5pt}\colorbox{blue!12}{\makebox[1.1em]{1}}}} & \hyperlink{mod-SAFE}{{\setlength{\fboxsep}{1.5pt}\colorbox{blue!12}{\makebox[1.1em]{1}}}} & \textcolor{black!25}{0} & \textcolor{black!25}{0} & 12\\
\hyperlink{course-UCHI}{UChicago} & \hyperlink{mod-FND}{{\setlength{\fboxsep}{1.5pt}\colorbox{blue!26}{\makebox[1.1em]{3}}}} & \textcolor{black!25}{0} & \textcolor{black!25}{0} & \hyperlink{mod-MCP}{{\setlength{\fboxsep}{1.5pt}\colorbox{blue!12}{\makebox[1.1em]{1}}}} & \hyperlink{mod-MEM}{{\setlength{\fboxsep}{1.5pt}\colorbox{blue!12}{\makebox[1.1em]{1}}}} & \hyperlink{mod-RSN}{{\setlength{\fboxsep}{1.5pt}\colorbox{blue!26}{\makebox[1.1em]{3}}}} & \hyperlink{mod-MULTI}{{\setlength{\fboxsep}{1.5pt}\colorbox{blue!19}{\makebox[1.1em]{2}}}} & \hyperlink{mod-CODE}{{\setlength{\fboxsep}{1.5pt}\colorbox{blue!12}{\makebox[1.1em]{1}}}} & \hyperlink{mod-WEB}{{\setlength{\fboxsep}{1.5pt}\colorbox{blue!12}{\makebox[1.1em]{1}}}} & \hyperlink{mod-EVAL}{{\setlength{\fboxsep}{1.5pt}\colorbox{blue!19}{\makebox[1.1em]{2}}}} & \hyperlink{mod-TRAIN}{{\setlength{\fboxsep}{1.5pt}\colorbox{blue!19}{\makebox[1.1em]{2}}}} & \hyperlink{mod-SAFE}{{\setlength{\fboxsep}{1.5pt}\colorbox{blue!19}{\makebox[1.1em]{2}}}} & \textcolor{black!25}{0} & \hyperlink{mod-APP}{{\setlength{\fboxsep}{1.5pt}\colorbox{blue!40}{\makebox[1.1em]{5}}}} & 11\\
\hyperlink{course-S329A}{Stanford 329A} & \hyperlink{mod-FND}{{\setlength{\fboxsep}{1.5pt}\colorbox{blue!12}{\makebox[1.1em]{1}}}} & \hyperlink{mod-ARCH}{{\setlength{\fboxsep}{1.5pt}\colorbox{blue!19}{\makebox[1.1em]{2}}}} & \hyperlink{mod-TOOL}{{\setlength{\fboxsep}{1.5pt}\colorbox{blue!12}{\makebox[1.1em]{1}}}} & \textcolor{black!25}{0} & \textcolor{black!25}{0} & \hyperlink{mod-RSN}{{\setlength{\fboxsep}{1.5pt}\colorbox{blue!40}{\makebox[1.1em]{5}}}} & \textcolor{black!25}{0} & \hyperlink{mod-CODE}{{\setlength{\fboxsep}{1.5pt}\colorbox{blue!33}{\makebox[1.1em]{4}}}} & \hyperlink{mod-WEB}{{\setlength{\fboxsep}{1.5pt}\colorbox{blue!12}{\makebox[1.1em]{1}}}} & \hyperlink{mod-EVAL}{{\setlength{\fboxsep}{1.5pt}\colorbox{blue!26}{\makebox[1.1em]{3}}}} & \hyperlink{mod-TRAIN}{{\setlength{\fboxsep}{1.5pt}\colorbox{blue!26}{\makebox[1.1em]{3}}}} & \textcolor{black!25}{0} & \textcolor{black!25}{0} & \hyperlink{mod-APP}{{\setlength{\fboxsep}{1.5pt}\colorbox{blue!12}{\makebox[1.1em]{1}}}} & 20\\
\hyperlink{course-BRKAG}{Berkeley '25} & \hyperlink{mod-FND}{{\setlength{\fboxsep}{1.5pt}\colorbox{blue!19}{\makebox[1.1em]{2}}}} & \hyperlink{mod-ARCH}{{\setlength{\fboxsep}{1.5pt}\colorbox{blue!12}{\makebox[1.1em]{1}}}} & \textcolor{black!25}{0} & \textcolor{black!25}{0} & \textcolor{black!25}{0} & \hyperlink{mod-RSN}{{\setlength{\fboxsep}{1.5pt}\colorbox{blue!12}{\makebox[1.1em]{1}}}} & \hyperlink{mod-MULTI}{{\setlength{\fboxsep}{1.5pt}\colorbox{blue!26}{\makebox[1.1em]{3}}}} & \hyperlink{mod-CODE}{{\setlength{\fboxsep}{1.5pt}\colorbox{blue!12}{\makebox[1.1em]{1}}}} & \hyperlink{mod-WEB}{{\setlength{\fboxsep}{1.5pt}\colorbox{blue!12}{\makebox[1.1em]{1}}}} & \hyperlink{mod-EVAL}{{\setlength{\fboxsep}{1.5pt}\colorbox{blue!33}{\makebox[1.1em]{4}}}} & \hyperlink{mod-TRAIN}{{\setlength{\fboxsep}{1.5pt}\colorbox{blue!33}{\makebox[1.1em]{4}}}} & \hyperlink{mod-SAFE}{{\setlength{\fboxsep}{1.5pt}\colorbox{blue!12}{\makebox[1.1em]{1}}}} & \hyperlink{mod-DEPLOY}{{\setlength{\fboxsep}{1.5pt}\colorbox{blue!19}{\makebox[1.1em]{2}}}} & \hyperlink{mod-APP}{{\setlength{\fboxsep}{1.5pt}\colorbox{blue!26}{\makebox[1.1em]{3}}}} & 13\\
\hyperlink{course-BRKADV}{Berkeley Adv.} & \textcolor{black!25}{0} & \hyperlink{mod-ARCH}{{\setlength{\fboxsep}{1.5pt}\colorbox{blue!19}{\makebox[1.1em]{2}}}} & \hyperlink{mod-TOOL}{{\setlength{\fboxsep}{1.5pt}\colorbox{blue!19}{\makebox[1.1em]{2}}}} & \textcolor{black!25}{0} & \hyperlink{mod-MEM}{{\setlength{\fboxsep}{1.5pt}\colorbox{blue!12}{\makebox[1.1em]{1}}}} & \hyperlink{mod-RSN}{{\setlength{\fboxsep}{1.5pt}\colorbox{blue!45}{\makebox[1.1em]{9}}}} & \textcolor{black!25}{0} & \hyperlink{mod-CODE}{{\setlength{\fboxsep}{1.5pt}\colorbox{blue!12}{\makebox[1.1em]{1}}}} & \hyperlink{mod-WEB}{{\setlength{\fboxsep}{1.5pt}\colorbox{blue!19}{\makebox[1.1em]{2}}}} & \hyperlink{mod-EVAL}{{\setlength{\fboxsep}{1.5pt}\colorbox{blue!19}{\makebox[1.1em]{2}}}} & \hyperlink{mod-TRAIN}{{\setlength{\fboxsep}{1.5pt}\colorbox{blue!26}{\makebox[1.1em]{3}}}} & \hyperlink{mod-SAFE}{{\setlength{\fboxsep}{1.5pt}\colorbox{blue!19}{\makebox[1.1em]{2}}}} & \hyperlink{mod-DEPLOY}{{\setlength{\fboxsep}{1.5pt}\colorbox{blue!12}{\makebox[1.1em]{1}}}} & \hyperlink{mod-APP}{{\setlength{\fboxsep}{1.5pt}\colorbox{blue!33}{\makebox[1.1em]{4}}}} & 12\\
\hyperlink{course-S329T}{Stanford 329T} & \hyperlink{mod-FND}{{\setlength{\fboxsep}{1.5pt}\colorbox{blue!33}{\makebox[1.1em]{4}}}} & \hyperlink{mod-ARCH}{{\setlength{\fboxsep}{1.5pt}\colorbox{blue!26}{\makebox[1.1em]{3}}}} & \hyperlink{mod-TOOL}{{\setlength{\fboxsep}{1.5pt}\colorbox{blue!12}{\makebox[1.1em]{1}}}} & \hyperlink{mod-MCP}{{\setlength{\fboxsep}{1.5pt}\colorbox{blue!12}{\makebox[1.1em]{1}}}} & \hyperlink{mod-MEM}{{\setlength{\fboxsep}{1.5pt}\colorbox{blue!12}{\makebox[1.1em]{1}}}} & \hyperlink{mod-RSN}{{\setlength{\fboxsep}{1.5pt}\colorbox{blue!12}{\makebox[1.1em]{1}}}} & \textcolor{black!25}{0} & \textcolor{black!25}{0} & \hyperlink{mod-WEB}{{\setlength{\fboxsep}{1.5pt}\colorbox{blue!12}{\makebox[1.1em]{1}}}} & \hyperlink{mod-EVAL}{{\setlength{\fboxsep}{1.5pt}\colorbox{blue!40}{\makebox[1.1em]{5}}}} & \hyperlink{mod-TRAIN}{{\setlength{\fboxsep}{1.5pt}\colorbox{blue!26}{\makebox[1.1em]{3}}}} & \hyperlink{mod-SAFE}{{\setlength{\fboxsep}{1.5pt}\colorbox{blue!12}{\makebox[1.1em]{1}}}} & \textcolor{black!25}{0} & \hyperlink{mod-APP}{{\setlength{\fboxsep}{1.5pt}\colorbox{blue!12}{\makebox[1.1em]{1}}}} & 18\\
\midrule Total & 26 & 45 & 20 & 11 & 26 & 48 & 27 & 46 & 11 & 36 & 30 & 22 & 19 & 33 & \\ \bottomrule\end{longtable}}
A lecture can appear in several modules. Logistics rows and guest talks count toward \emph{All} but not toward any module.
\chapter{Learning Tracks}\hypertarget{tracks}{}
Choose a track by goal. All tracks assume Python fluency. Track~C also assumes basic deep learning (PyTorch, transformers). Each module's curated sequence is the default path; the lecture map lists alternatives.
\section{Track A: Agent Builder (practitioner)}
\textbf{For:} You want to ship reliable tool-using agents and RAG systems.\par
\textbf{Path:} \hyperlink{mod-SETUP}{M0} $\rightarrow$ \hyperlink{mod-FND}{M1} $\rightarrow$ \hyperlink{mod-ARCH}{M2} $\rightarrow$ \hyperlink{mod-TOOL}{M3} $\rightarrow$ \hyperlink{mod-MCP}{M4} $\rightarrow$ \hyperlink{mod-MEM}{M5} $\rightarrow$ \hyperlink{mod-RSN}{M6} $\rightarrow$ \hyperlink{mod-EVAL}{M10} $\rightarrow$ \hyperlink{mod-SAFE}{M12} $\rightarrow$ \hyperlink{mod-DEPLOY}{M13}\par
\textbf{Spine courses:} Northeastern notebooks for guided builds, UVA labs for experiments, Buffalo lab specs as checkpoints.\par
\textbf{Capstone:} Buffalo's \emph{useful bounded agent} (\hyperlink{asg-UB-10}{\textbf{Buffalo: Final Team Project: Useful Bounded Agent}}), deployed like Northeastern's \hyperlink{asg-NEU-1}{\textbf{Northeastern: DADA1: The Trip Planner Agent}}.\par
\section{Track B: Coding-Agent Engineer}
\textbf{For:} You want to build or deeply customize coding agents and agent-ready workflows.\par
\textbf{Path:} \hyperlink{mod-SETUP}{M0} $\rightarrow$ \hyperlink{mod-ARCH}{M2} $\rightarrow$ \hyperlink{mod-TOOL}{M3} $\rightarrow$ \hyperlink{mod-MCP}{M4} $\rightarrow$ \hyperlink{mod-MEM}{M5} $\rightarrow$ \hyperlink{mod-CODE}{M8} $\rightarrow$ \hyperlink{mod-EVAL}{M10} $\rightarrow$ \hyperlink{mod-SAFE}{M12}\par
\textbf{Spine courses:} Stanford '26 for practice, CMU for harness internals, UW and CU Boulder for the research context.\par
\textbf{Capstone:} \hyperlink{asg-UW-1}{\textbf{UW-Madison: Build a Coding Agent}} without frameworks, then extend it to pass the \hyperlink{asg-JHU-3}{\textbf{JHU: A3: Coding LLM Agents}} tasks and evaluate it like \hyperlink{asg-JHU-2}{\textbf{JHU: A2: Evaluating Coding LLMs}}.\par
\section{Track C: Agent Researcher / Trainer}
\textbf{For:} You want to train agentic models and do research-grade evaluation.\par
\textbf{Path:} \hyperlink{mod-FND}{M1} $\rightarrow$ \hyperlink{mod-ARCH}{M2} $\rightarrow$ \hyperlink{mod-TOOL}{M3} $\rightarrow$ \hyperlink{mod-MEM}{M5} $\rightarrow$ \hyperlink{mod-RSN}{M6} $\rightarrow$ \hyperlink{mod-EVAL}{M10} $\rightarrow$ \hyperlink{mod-TRAIN}{M11} $\rightarrow$ \hyperlink{mod-MULTI}{M7} $\rightarrow$ \hyperlink{mod-SAFE}{M12} $\rightarrow$ \hyperlink{mod-APP}{M14}\par
\textbf{Spine courses:} CMU 11-768 end to end (all three assignments), with UW CS 839 and CU Boulder reading lists; Stanford CS329A for test-time compute and self-improvement, and the two Berkeley series for current research talks.\par
\textbf{Capstone:} The CMU pipeline: \hyperlink{asg-CMU-1}{\textbf{CMU: A1: Build an Agent Harness}} $\rightarrow$ \hyperlink{asg-CMU-2}{\textbf{CMU: A2: Evaluating a Data-Visualization Agent}} $\rightarrow$ \hyperlink{asg-CMU-3}{\textbf{CMU: A3: Agentic RL Training, "scottyrl"}}, then a research question framed as Northeastern requires (mechanism $\rightarrow$ manipulation $\rightarrow$ measured behavior).\par
\section{A 16-Week Self-Study Plan}\hypertarget{plan16}{}
Roughly 8--10 hours a week: watch or read the listed sessions (slides first, then papers), then build. The plan follows Track~A but includes the core of B and C. Adjust it freely.
{\footnotesize\begin{longtable}{c P{2.4cm} P{6.4cm} P{6.4cm}}
\toprule Wk & Module & Watch / read & Build\\ \midrule\endhead
1 & \hyperlink{mod-SETUP}{SETUP}, \hyperlink{mod-FND}{FND} & \hyperlink{lec-NEU-2}{Northeastern~2}, \hyperlink{lec-NEU-3}{Northeastern~3}, \hyperlink{lec-STAN25-2}{Stanford '25~2} & \hyperlink{asg-UVA-1}{\textbf{UVA: Topic 1 portfolio: Running an LLM}}; \hyperlink{asg-STAN25-1}{\textbf{Stanford '25: Week 1 --- Prompting Techniques}}\\ \midrule
2 & \hyperlink{mod-ARCH}{ARCH} & \hyperlink{lec-CMU-1}{CMU~1}, \hyperlink{lec-STAN26-1}{Stanford '26~1}, \hyperlink{lec-UW-2}{UW-Madison~2} & Raw ReAct loop from scratch (start \hyperlink{asg-UW-1}{\textbf{UW-Madison: Build a Coding Agent}})\\ \midrule
3 & \hyperlink{mod-TOOL}{TOOL} & \hyperlink{lec-CMU-2}{CMU~2}, \hyperlink{lec-UVA-5}{UVA~5}, \hyperlink{lec-UVA-7}{UVA~7} & \hyperlink{asg-UVA-3}{\textbf{UVA: Topic 3 portfolio: Tool use}}; \hyperlink{asg-UB-1}{\textbf{Buffalo: Lab 1: A Minimal Tool-Using Agent}}\\ \midrule
4 & \hyperlink{mod-MCP}{MCP} & MCP docs, \hyperlink{lec-STAN26-4}{Stanford '26~4}, \hyperlink{lec-JHU-8}{JHU~8}--\hyperlink{lec-JHU-9}{JHU~9} & \hyperlink{asg-STAN26-2}{\textbf{Stanford '26: Week 2 --- Build an MCP Server for a Third-Party API}}\\ \midrule
5 & \hyperlink{mod-MEM}{MEM (context)} & \hyperlink{lec-CMU-3}{CMU~3}, \hyperlink{lec-CMU-4}{CMU~4}, \hyperlink{lec-NEU-7}{Northeastern~7} & Add compaction and skills to your agent (\hyperlink{asg-STAN26-3}{\textbf{Stanford '26: Week 3 --- Author a Skill}})\\ \midrule
6 & \hyperlink{mod-MEM}{MEM (RAG)} & \hyperlink{lec-NEU-8}{Northeastern~8}, \hyperlink{lec-NEU-9}{Northeastern~9}, \hyperlink{lec-UVA-9}{UVA~9} & \hyperlink{asg-UVA-5}{\textbf{UVA: Topic 5: RAG experiments}}; \hyperlink{asg-UB-2}{\textbf{Buffalo: Lab 2: Retrieval and Memory}}\\ \midrule
7 & \hyperlink{mod-ARCH}{ARCH (graphs)} & \hyperlink{lec-NEU-6}{Northeastern~6}, \hyperlink{lec-UVA-3}{UVA~3}, \hyperlink{lec-CMU-15}{CMU~15} & \hyperlink{asg-UB-3}{\textbf{Buffalo: Lab 3: Workflow Graphs and State Machines}}; \hyperlink{asg-NEU-1}{\textbf{Northeastern: DADA1: The Trip Planner Agent}}\\ \midrule
8 & \hyperlink{mod-RSN}{RSN} & \hyperlink{lec-CMU-5}{CMU~5}, \hyperlink{lec-NEU-10}{Northeastern~10}, \hyperlink{lec-S329A-2}{Stanford 329A~2}; before A1: \hyperlink{lec-JHU-2}{JHU~2}--\hyperlink{lec-JHU-4}{JHU~4} & \hyperlink{asg-JHU-1}{\textbf{JHU: A1: Inductive Synthesis}}\\ \midrule
9 & \hyperlink{mod-CODE}{CODE} & \hyperlink{lec-CMU-6}{CMU~6}, \hyperlink{lec-STAN26-2}{Stanford '26~2}, \hyperlink{lec-UW-5}{UW-Madison~5}--\hyperlink{lec-UW-6}{UW-Madison~6} & \hyperlink{asg-CMU-1}{\textbf{CMU: A1: Build an Agent Harness}} or \hyperlink{asg-JHU-3}{\textbf{JHU: A3: Coding LLM Agents}}\\ \midrule
10 & \hyperlink{mod-EVAL}{EVAL} & \hyperlink{lec-NEU-18}{Northeastern~18}, \hyperlink{lec-UB-9}{Buffalo~9}, \hyperlink{lec-S329T-7}{Stanford 329T~7}, \hyperlink{lec-BRKAG-5}{Berkeley '25~5} & \hyperlink{asg-CMU-2}{\textbf{CMU: A2: Evaluating a Data-Visualization Agent}} or \hyperlink{asg-JHU-2}{\textbf{JHU: A2: Evaluating Coding LLMs}} (pick one; each fills the week. Do the other later as a multi-week project)\\ \midrule
11 & \hyperlink{mod-MULTI}{MULTI} & \hyperlink{lec-NEU-12}{Northeastern~12}--\hyperlink{lec-NEU-14}{Northeastern~14}, \hyperlink{lec-UW-20}{UW-Madison~20} & \hyperlink{asg-UVA-7}{\textbf{UVA: Topic 7: MCP + A2A}} (A2A)\\ \midrule
12 & \hyperlink{mod-WEB}{WEB} & \hyperlink{lec-CMU-7}{CMU~7}, \hyperlink{lec-UVA-10}{UVA~10} & \hyperlink{asg-UVA-6}{\textbf{UVA: Topic 6: VLM agents}}\\ \midrule
13 & \hyperlink{mod-TRAIN}{TRAIN (SFT)} & \hyperlink{lec-CMU-8}{CMU~8}, \hyperlink{lec-UVA-13}{UVA~13}, TinyAgent & \hyperlink{asg-UVA-8}{\textbf{UVA: Topic 8: Text-to-SQL fine-tuning}}\\ \midrule
14 & \hyperlink{mod-TRAIN}{TRAIN (RL)} & \hyperlink{lec-CMU-9}{CMU~9}, \hyperlink{lec-CMU-11}{CMU~11}, \hyperlink{lec-CMU-12}{CMU~12}, \hyperlink{lec-S329A-6}{Stanford 329A~6} & \hyperlink{asg-CMU-3}{\textbf{CMU: A3: Agentic RL Training, "scottyrl"}}\\ \midrule
15 & \hyperlink{mod-SAFE}{SAFE} & \hyperlink{lec-CMU-13}{CMU~13}, \hyperlink{lec-UW-11}{UW-Madison~11}--\hyperlink{lec-UW-12}{UW-Madison~12}, \hyperlink{lec-UB-10}{Buffalo~10}, \hyperlink{lec-BRKAG-13}{Berkeley '25~13} & \hyperlink{asg-UB-5}{\textbf{Buffalo: Lab 5: Human-in-the-Loop and Safety Boundaries}}\\ \midrule
16 & \hyperlink{mod-DEPLOY}{DEPLOY}, \hyperlink{mod-APP}{APP} & \hyperlink{lec-NEU-19}{Northeastern~19}, \hyperlink{lec-STAN26-15}{Stanford '26~15}, \hyperlink{lec-CMU-10}{CMU~10} & Capstone (\hyperlink{asg-UB-10}{\textbf{Buffalo: Final Team Project: Useful Bounded Agent}})\\ \midrule
\end{longtable}}
\part{Topic Modules}
\chapter{Module 0: Setup: Your Agent Workbench}\hypertarget{mod-SETUP}{}
\section{Learning objectives}
\begin{itemize}
\item Have a reproducible Python environment, API access and a sandbox to run course code in.
\item Know which courses need which accounts and services.
\end{itemize}
\section{Core concepts}
Every course here assumes a small workbench, and setting it up once saves friction later.
\begin{itemize}
\item \textbf{Python environment.} CMU and Northeastern use \texttt{uv}; Stanford '25 used Poetry/Conda. Keep one virtual environment per course repository.
\item \textbf{Models.} You need at least one hosted API key (the courses use OpenAI, Anthropic via Bedrock, Gemini and OpenRouter) and, optionally, a local runtime (Ollama is used by Northeastern, UVA and Stanford '25) for cheap, private experiments with small open models.
\item \textbf{Sandboxes.} Agents execute code. CMU runs tools in Modal sandboxes. Colab is the zero-install choice for notebooks. Never let an agent you are developing run commands directly on the machine that holds your credentials.
\item \textbf{Tracing.} Log every model call and tool call from day one (plain JSONL is enough). Every evaluation and debugging lecture assumes you have traces.
\end{itemize}
\textbf{Install policy reminder.} The courses name many packages (LangGraph, FastMCP, Tinker, Modal, \ldots). Treat each install as a deliberate decision: pin versions and check security advisories before installing.
\section{Curated learning sequence}
\begin{enumerate}
\item \hyperlink{lec-NEU-3}{Northeastern~3 \emph{LLM Foundations (II)}} --- The setup page walks through uv, \texttt{.env} keys and a first LLM API call.
\item \hyperlink{lec-UVA-1}{UVA~1 \emph{Topic 1: Running an LLM}} --- Run a small open model locally and on Colab, with and without quantization.
\item Supplement: \hyperref[sup-colab]{Google Colab} --- The runtime for every notebook-based course.
\end{enumerate}
\chapter{Module 1: LLM Foundations and Prompting}\hypertarget{mod-FND}{}
{\small Tag: \tg{FND}}
\section{Learning objectives}
\begin{itemize}
\item Explain next-token prediction, pre-training versus post-training, and why instruction-tuned models follow chat templates.
\item Use prompting techniques (few-shot, CoT, self-consistency, structured output) deliberately.
\item Run and benchmark a small open model locally.
\end{itemize}
\section{Core concepts}
An agent is only as good as the model in its loop, and most agent bugs trace back to how that model was conditioned.
\begin{itemize}
\item \textbf{Pre-training} produces a next-token predictor. \textbf{Post-training} (instruction tuning, RLHF/RL) turns it into an assistant that follows \emph{chat templates}: role-tagged messages, and for agents, tool-call and tool-result segments. CMU's tool-use lecture shows how those templates encode function calls.
\item \textbf{Prompting} is programming in natural language. Few-shot examples, chain-of-thought, self-consistency and explicit output schemas are the basic toolkit. Structured output (JSON or typed objects) is what makes a model's response machine-actionable.
\item \textbf{Cost and latency} scale with tokens, so an agent's prompt is a budget. Context-window limits and KV-cache pricing return in \hyperlink{mod-MEM}{Module~5 (Memory, Retrieval and Context Engineering)} and \hyperlink{mod-DEPLOY}{Module~13 (Deployment, Efficiency and Operations)}.
\item \textbf{Small versus large models.} UVA's first reading argues that small language models are the future of agentic AI, and the TinyAgent supplement shows a 1.1B model beating a frontier model on a narrow toolset.
\end{itemize}
\section{Curated learning sequence}
\begin{enumerate}
\item \hyperlink{lec-NEU-2}{Northeastern~2 \emph{LLM Foundations (I)}} --- Transformers, autoregressive generation, base versus instruction-tuned models (with a local model).
\item \hyperlink{lec-NEU-3}{Northeastern~3 \emph{LLM Foundations (II)}} --- InstructGPT, CoT, and the LLM API basics notebook.
\item \hyperlink{lec-STAN25-1}{Stanford '25~1 \emph{Introduction and how an LLM is made}} and \hyperlink{lec-STAN25-2}{Stanford '25~2 \emph{Power prompting for LLMs}} --- How an LLM is made, then power prompting (slides).
\item \hyperlink{lec-MANN-1}{Mannheim~1 \emph{Introduction (to LLMs)}}--\hyperlink{lec-MANN-3}{Mannheim~3} --- Compact decks on LLM basics, post-training and prompt engineering.
\item \hyperlink{lec-UVA-1}{UVA~1 \emph{Topic 1: Running an LLM}} --- Benchmark Llama-3.2-1B on MMLU across hardware and quantization settings.
\end{enumerate}
\section{Lecture map (26 lectures)}
{\footnotesize\begin{longtable}{P{1.75cm} P{0.4cm} P{5.75cm} P{4.2cm} P{2.65cm}}
\toprule Course & \# & Title (links to catalogue row) & Materials & Tags\\ \midrule\endhead
\hyperlink{course-STAN26}{Stanford '26} & 2 & \hyperlink{lec-STAN26-2}{How SOTA coding agents are designed: deep dive into system prompts} & \href{https://themodernsoftware.dev/}{[course site]} & \hyperlink{mod-CODE}{\tg{CODE}} \hyperlink{mod-FND}{\tg{FND}}\\
\hyperlink{course-STAN26}{Stanford '26} & 3 & \hyperlink{lec-STAN26-3}{Advanced prompting + agentic dev frameworks (RePPIT, spec-driven development)} & \href{https://themodernsoftware.dev/}{[course site]} & \hyperlink{mod-FND}{\tg{FND}} \hyperlink{mod-MEM}{\tg{MEM}}\\
\hyperlink{course-STAN25}{Stanford '25} & 1 & $\star$\,\hyperlink{lec-STAN25-1}{Introduction and how an LLM is made} & \href{https://docs.google.com/presentation/d/1zT2Ofy88cajLTLkd7TcuSM4BCELvF9qQdHmlz33i4t0/edit?usp=sharing}{[slides]} & \tg{ORG} \hyperlink{mod-FND}{\tg{FND}}\\
\hyperlink{course-STAN25}{Stanford '25} & 2 & $\star$\,\hyperlink{lec-STAN25-2}{Power prompting for LLMs} & \href{https://docs.google.com/presentation/d/1MIhw8p6TLGdbQ9TcxhXSs5BaPf5d_h77QY70RHNfeGs/edit?usp=drive_link}{[slides]} & \hyperlink{mod-FND}{\tg{FND}} \hyperlink{mod-RSN}{\tg{RSN}}\\
\hyperlink{course-NEU}{Northeastern} & 2 & $\star$\,\hyperlink{lec-NEU-2}{LLM Foundations (I)} & \href{https://agenticllms.github.io/week2a.html}{[agenticllms.github.io]} \href{https://www.youtube.com/watch?v=wjZofJX0v4M}{[video]} \href{https://www.youtube.com/watch?v=kCc8FmEb1nY}{[video]} & \hyperlink{mod-FND}{\tg{FND}} \hyperlink{mod-DEPLOY}{\tg{DEPLOY}}\\
\hyperlink{course-NEU}{Northeastern} & 3 & $\star$\,\hyperlink{lec-NEU-3}{LLM Foundations (II)} & \href{https://agenticllms.github.io/week2b-setup.html}{[agenticllms.github.io]} \href{https://github.com/agenticllms/lecture-files/tree/main/W02b}{[GitHub]} & \hyperlink{mod-FND}{\tg{FND}} \hyperlink{mod-TRAIN}{\tg{TRAIN}}\\
\hyperlink{course-NEU}{Northeastern} & 19 & \hyperlink{lec-NEU-19}{Agentic AI Optimization} & -- & \hyperlink{mod-DEPLOY}{\tg{DEPLOY}} \hyperlink{mod-FND}{\tg{FND}}\\
\hyperlink{course-UB}{Buffalo} & 2 & \hyperlink{lec-UB-2}{Foundation models as components} & -- & \hyperlink{mod-FND}{\tg{FND}}\\
\hyperlink{course-UVA}{UVA} & 1 & $\star$\,\hyperlink{lec-UVA-1}{Topic 1: Running an LLM} & \href{https://www.cs.virginia.edu/\%7Ermw7my/Courses/AgenticAISpring2026/Topic1Running/Running_an_LLM.html}{[cs.virginia.edu]} & \hyperlink{mod-FND}{\tg{FND}} \hyperlink{mod-EVAL}{\tg{EVAL}} \hyperlink{mod-DEPLOY}{\tg{DEPLOY}}\\
\hyperlink{course-UVA}{UVA} & 10 & \hyperlink{lec-UVA-10}{Topic 6: Vision-Language Models} & \href{https://www.cs.virginia.edu/\%7Ermw7my/Courses/AgenticAISpring2026/Topic6VLM/vlm.html}{[cs.virginia.edu]} \href{https://github.com/HenryKautz/clip_test}{[GitHub]} & \hyperlink{mod-WEB}{\tg{WEB}} \hyperlink{mod-APP}{\tg{APP}} \hyperlink{mod-FND}{\tg{FND}}\\
\hyperlink{course-UW}{UW-Madison} & 3 & \hyperlink{lec-UW-3}{Reasoning and acting I} & -- & \hyperlink{mod-RSN}{\tg{RSN}} \hyperlink{mod-FND}{\tg{FND}}\\
\hyperlink{course-JHU}{JHU} & 5 & \hyperlink{lec-JHU-5}{Language Modeling for Synthesis} & \href{https://machine-programming.github.io/lectures/lecture-5.pdf}{[PDF]} & \hyperlink{mod-FND}{\tg{FND}} \hyperlink{mod-CODE}{\tg{CODE}}\\
\hyperlink{course-JHU}{JHU} & 6 & \hyperlink{lec-JHU-6}{Decoding and Prompting for Synthesis} & \href{https://machine-programming.github.io/lectures/lecture-6.pdf}{[PDF]} & \hyperlink{mod-FND}{\tg{FND}} \hyperlink{mod-CODE}{\tg{CODE}}\\
\hyperlink{course-MANN}{Mannheim} & 1 & $\star$\,\hyperlink{lec-MANN-1}{Introduction (to LLMs)} & \href{https://www.uni-mannheim.de/media/Einrichtungen/dws/Files_Teaching/Large_Language_Models_and_Agents/FSS2026/IE685_LA_01_Introduction.pdf}{[PDF]} & \hyperlink{mod-FND}{\tg{FND}} \tg{ORG}\\
\hyperlink{course-MANN}{Mannheim} & 2 & $\star$\,\hyperlink{lec-MANN-2}{Post-training and Efficient Adaptation} & \href{https://www.uni-mannheim.de/media/Einrichtungen/dws/Files_Teaching/Large_Language_Models_and_Agents/FSS2026/IE685_LA_02_PosttrainingAndEfficientAdaptation.pdf}{[PDF]} \href{https://www.uni-mannheim.de/media/Einrichtungen/dws/Files_Teaching/Large_Language_Models_and_Agents/FSS2026/IE685_LA_EX_LLM_Introduction.ipynb.zip}{[notebook zip]} & \hyperlink{mod-TRAIN}{\tg{TRAIN}} \hyperlink{mod-FND}{\tg{FND}}\\
\hyperlink{course-MANN}{Mannheim} & 3 & $\star$\,\hyperlink{lec-MANN-3}{Prompt Engineering} & \href{https://www.uni-mannheim.de/media/Einrichtungen/dws/Files_Teaching/Large_Language_Models_and_Agents/FSS2026/IE685_LA_03_PromptEngineering.pdf}{[PDF]} \href{https://www.uni-mannheim.de/media/Einrichtungen/dws/Files_Teaching/Large_Language_Models_and_Agents/FSS2026/IE686_LA_EX_Langchain_Introduction.ipynb.zip}{[notebook zip]} & \hyperlink{mod-FND}{\tg{FND}} \hyperlink{mod-RSN}{\tg{RSN}}\\
\hyperlink{course-UCHI}{UChicago} & 1 & \hyperlink{lec-UCHI-1}{Deep Learning and Social Agents} & \href{https://github.com/KnowledgeLab/AI-Agents-for-Social-Science-and-Society-2026/blob/main/Tutorials-Homework_Notebooks/Week_1_Intro_NNs.ipynb}{[GitHub]} & \hyperlink{mod-FND}{\tg{FND}} \hyperlink{mod-APP}{\tg{APP}}\\
\hyperlink{course-UCHI}{UChicago} & 2 & \hyperlink{lec-UCHI-2}{Text Learning, Transformers, and Diffusion Models} & \href{https://github.com/KnowledgeLab/AI-Agents-for-Social-Science-and-Society-2026/blob/main/Tutorials-Homework_Notebooks/Week_2.ipynb}{[GitHub]} & \hyperlink{mod-FND}{\tg{FND}} \hyperlink{mod-RSN}{\tg{RSN}}\\
\hyperlink{course-UCHI}{UChicago} & 6 & \hyperlink{lec-UCHI-6}{Auto-Encoders, Interpretability, and Steering Agents} & \href{https://github.com/KnowledgeLab/AI-Agents-for-Social-Science-and-Society-2026/blob/main/Tutorials-Homework_Notebooks/week6_2026.ipynb}{[GitHub]} & \hyperlink{mod-SAFE}{\tg{SAFE}} \hyperlink{mod-FND}{\tg{FND}}\\
\hyperlink{course-S329A}{Stanford 329A} & 20 & \hyperlink{lec-S329A-20}{Future Research Areas} & \href{https://www.youtube.com/watch?v=AyO6wyu4DEg}{[video]} & \hyperlink{mod-FND}{\tg{FND}}\\
\hyperlink{course-BRKAG}{Berkeley '25} & 1 & \hyperlink{lec-BRKAG-1}{Introduction} & \href{https://rdi.berkeley.edu/agentic-ai/slides/introduction_25.pdf}{[slides]} & \tg{ORG} \hyperlink{mod-FND}{\tg{FND}}\\
\hyperlink{course-BRKAG}{Berkeley '25} & 2 & \hyperlink{lec-BRKAG-2}{LLM Agents Overview} & \href{https://rdi.berkeley.edu/agentic-ai/slides/lecture1.pdf}{[slides]} \href{https://www.youtube.com/watch?v=r1qZpYAmqmg}{[video]} & \hyperlink{mod-FND}{\tg{FND}} \hyperlink{mod-TRAIN}{\tg{TRAIN}}\\
\hyperlink{course-S329T}{Stanford 329T} & 1 & \hyperlink{lec-S329T-1}{Course introduction and overview} & \href{https://web.stanford.edu/class/cs329t/slides/Lecture\%201\%20-\%20CS\%20329T\%20Fall\%202025.pdf}{[slides]} & \tg{ORG} \hyperlink{mod-FND}{\tg{FND}} \hyperlink{mod-ARCH}{\tg{ARCH}}\\
\hyperlink{course-S329T}{Stanford 329T} & 2 & \hyperlink{lec-S329T-2}{Models, Prompting and RAG} & \href{https://web.stanford.edu/class/cs329t/slides/Lecture\%202\%20-\%20CS\%20329T\%20Fall\%202025.pdf}{[slides]} \href{https://web.stanford.edu/class/cs329t/homeworks/CS329T_HW1_Fall2025.pdf}{[HW1]} & \hyperlink{mod-FND}{\tg{FND}} \hyperlink{mod-MEM}{\tg{MEM}}\\
\hyperlink{course-S329T}{Stanford 329T} & 9 & \hyperlink{lec-S329T-9}{Deep dive into LLMs} & \href{https://web.stanford.edu/class/cs329t/slides/Lecture\%207\%20--\%20CS\%20329T\%20Fall\%202025.pdf}{[slides]} & \hyperlink{mod-FND}{\tg{FND}} \hyperlink{mod-TRAIN}{\tg{TRAIN}}\\
\hyperlink{course-S329T}{Stanford 329T} & 15 & \hyperlink{lec-S329T-15}{Multi-modal frontier models with agentic capabilities (guest: Deepak Ramachandran, Google DeepMind)} & -- & \hyperlink{mod-FND}{\tg{FND}} \hyperlink{mod-ARCH}{\tg{ARCH}}\\
\bottomrule\end{longtable}}
\section{Key readings}
\begin{itemize}
\item Vaswani et al., \href{https://arxiv.org/abs/1706.03762}{\emph{Attention Is All You Need}} (2017)
\item Ouyang et al., \href{https://arxiv.org/abs/2203.02155}{\emph{InstructGPT}} (2022)
\item Wei et al., \href{https://arxiv.org/abs/2201.11903}{\emph{Chain-of-Thought Prompting}} (2022)
\item Belcak et al., \href{https://arxiv.org/abs/2506.02153}{\emph{Small Language Models are the Future of Agentic AI}} (2025)
\item Karpathy, \href{https://www.youtube.com/watch?v=7xTGNNLPyMI}{\emph{Deep Dive into LLMs}} (video)
\end{itemize}
\section{Hands-on}
\begin{itemize}
\item \hyperlink{asg-STAN25-1}{\textbf{Stanford '25: Week 1 --- Prompting Techniques}} --- Fill in prompts for K-shot, chain-of-thought, tool calling, self-consistency, RAG, and reflexion until the test scripts pass, using local Ollama models (mistral-nemo:12b, llama3.1:8b). \textit{Starter code:} \href{https://github.com/mihail911/modern-software-dev-assignments/tree/fall2025/week1}{[GitHub]}
\item \hyperlink{asg-UVA-1}{\textbf{UVA: Topic 1 portfolio: Running an LLM}} --- MMLU timing/accuracy study across models, quantization and hardware; chat agent with bounded context and restartability; optional MT-Bench. \textit{Starter code:} \href{https://www.cs.virginia.edu/\%7Ermw7my/Courses/AgenticAISpring2026/Topic1Running/Running_an_LLM.html}{[cs.virginia.edu]}
\item \hyperlink{asg-MANN-1}{\textbf{Mannheim: Exercise: LLM Introduction}} --- Intro notebook for working with LLMs. \textit{Starter code:} \href{https://www.uni-mannheim.de/media/Einrichtungen/dws/Files_Teaching/Large_Language_Models_and_Agents/FSS2026/IE685_LA_EX_LLM_Introduction.ipynb.zip}{[notebook zip]}
\end{itemize}
\chapter{Module 2: The Agent Loop and Agent Architectures}\hypertarget{mod-ARCH}{}
{\small Tag: \tg{ARCH}}
\section{Learning objectives}
\begin{itemize}
\item Define an agent (model + loop + tools + environment + stopping rule) and distinguish workflows from agents.
\item Implement a ReAct-style loop from scratch, without a framework.
\item Judge when a graph framework (LangGraph) or an agent SDK earns its complexity.
\end{itemize}
\section{Core concepts}
Anthropic's \emph{Building Effective Agents}, assigned by Northeastern, Buffalo and UW, draws the key distinction. A \textbf{workflow} orchestrates LLM calls along predefined code paths (chaining, routing, parallelization, orchestrator--workers, evaluator--optimizer). An \textbf{agent} lets the model choose its own next action in a loop until a stopping condition holds.
\begin{itemize}
\item \textbf{The loop.} Build the prompt (system + history + tool specs), call the model, parse any tool calls, execute them in the environment, append the observations, and repeat until the model answers or a step or token budget runs out. That is a \emph{harness}. CMU A1, JHU A3 and the UW assignment all make you write one.
\item \textbf{State} is the message history plus anything external (files, a scratchpad, a database). Explicit state machines and graphs (Buffalo Lab~3, LangGraph in Northeastern, UVA and CMU) make control flow inspectable, resumable (checkpointing) and testable.
\item \textbf{Frameworks versus raw loops.} Write a loop by hand first. Frameworks such as LangGraph, the OpenAI Agents SDK, OpenHands and smolagents then become shortcuts you understand rather than magic. UW bans frameworks for its assignment for exactly this reason.
\item \textbf{Cognitive architectures} (CoALA) provide a vocabulary for memory, action space and decision procedure that the rest of this guide reuses.
\end{itemize}
\section{Curated learning sequence}
\begin{enumerate}
\item \hyperlink{lec-CMU-1}{CMU~1 \emph{Course Overview: What Is an Agent?}} --- What an agent is, the agent loop and the course map (slides + video).
\item \hyperlink{lec-STAN26-1}{Stanford '26~1 \emph{Course intro + build Claude Code in 200 lines}} --- Build a Claude-Code-style agent in about 200 lines, then read the completed code.
\item \hyperlink{lec-UW-2}{UW-Madison~2 \emph{Building agents}} --- Instructor walkthrough with a demo notebook. Then do the UW assignment below.
\item \hyperlink{lec-NEU-4}{Northeastern~4 \emph{Simple Agent Architecture}} and \hyperlink{lec-NEU-6}{Northeastern~6 \emph{Short-term \& Long-term Memory}} --- From a simple tools-and-memory agent to LangGraph state, nodes, edges and cycles.
\item \hyperlink{lec-UVA-3}{UVA~3 \emph{Topic 2: Agent Frameworks}} --- LangGraph lab: routing, parallel models, checkpointing and crash recovery.
\item \hyperlink{lec-CMU-14}{CMU~14 \emph{Frameworks 1: OpenHands}} and \hyperlink{lec-CMU-15}{CMU~15 \emph{Frameworks 2: LangGraph}} --- Two production frameworks compared: OpenHands and LangGraph.
\end{enumerate}
\section{Lecture map (45 lectures)}
{\footnotesize\begin{longtable}{P{1.75cm} P{0.4cm} P{5.75cm} P{4.2cm} P{2.65cm}}
\toprule Course & \# & Title (links to catalogue row) & Materials & Tags\\ \midrule\endhead
\hyperlink{course-CMU}{CMU} & 1 & $\star$\,\hyperlink{lec-CMU-1}{Course Overview: What Is an Agent?} & \href{https://www.cmu-agents.com/slides/lecture-01-agents.pdf}{[slides]} \href{https://www.cmu-agents.com/\#/schedule}{[video]} & \tg{ORG} \hyperlink{mod-ARCH}{\tg{ARCH}}\\
\hyperlink{course-CMU}{CMU} & 14 & $\star$\,\hyperlink{lec-CMU-14}{Frameworks 1: OpenHands} & -- & \hyperlink{mod-ARCH}{\tg{ARCH}} \hyperlink{mod-CODE}{\tg{CODE}}\\
\hyperlink{course-CMU}{CMU} & 15 & $\star$\,\hyperlink{lec-CMU-15}{Frameworks 2: LangGraph} & -- & \hyperlink{mod-ARCH}{\tg{ARCH}}\\
\hyperlink{course-STAN26}{Stanford '26} & 1 & $\star$\,\hyperlink{lec-STAN26-1}{Course intro + build Claude Code in 200 lines} & \href{https://themodernsoftware.dev/}{[course site]} & \hyperlink{mod-CODE}{\tg{CODE}} \hyperlink{mod-ARCH}{\tg{ARCH}}\\
\hyperlink{course-STAN26}{Stanford '26} & 15 & \hyperlink{lec-STAN26-15}{Background agents: launching tasks asynchronously} & \href{https://themodernsoftware.dev/}{[course site]} & \hyperlink{mod-DEPLOY}{\tg{DEPLOY}} \hyperlink{mod-ARCH}{\tg{ARCH}}\\
\hyperlink{course-STAN26}{Stanford '26} & 20 & \hyperlink{lec-STAN26-20}{The Software Factory: self-running, self-improving software systems} & \href{https://themodernsoftware.dev/}{[course site]} & \hyperlink{mod-ARCH}{\tg{ARCH}} \hyperlink{mod-DEPLOY}{\tg{DEPLOY}}\\
\hyperlink{course-STAN25}{Stanford '25} & 3 & \hyperlink{lec-STAN25-3}{Building a coding agent from scratch} & \href{https://docs.google.com/presentation/d/11CP26VhsjnZOmi9YFgLlonzdib9BLyAlgc4cEvC5Fps/edit?usp=sharing}{[slides]} \href{https://drive.google.com/file/d/1YtpKFVG13DHyQ2i3HOtwyVJOV90nWeL2/view?usp=drive_link}{[drive.google.com]} & \hyperlink{mod-ARCH}{\tg{ARCH}} \hyperlink{mod-TOOL}{\tg{TOOL}} \hyperlink{mod-CODE}{\tg{CODE}}\\
\hyperlink{course-NEU}{Northeastern} & 1 & \hyperlink{lec-NEU-1}{Introduction to Agentic AI} & \href{https://agenticllms.github.io/week1.html}{[agenticllms.github.io]} \href{https://colab.research.google.com/drive/1s60I7NQiTpOWg2tK-KmFyiC8ZAY24lo7?usp=sharing}{[colab.research.google.com]} & \tg{ORG} \hyperlink{mod-ARCH}{\tg{ARCH}}\\
\hyperlink{course-NEU}{Northeastern} & 4 & $\star$\,\hyperlink{lec-NEU-4}{Simple Agent Architecture} & \href{https://agenticllms.github.io/week3a.html}{[agenticllms.github.io]} \href{https://github.com/agenticllms/lecture-files/tree/main/W03a}{[GitHub]} & \hyperlink{mod-ARCH}{\tg{ARCH}} \hyperlink{mod-TOOL}{\tg{TOOL}} \hyperlink{mod-MEM}{\tg{MEM}}\\
\hyperlink{course-NEU}{Northeastern} & 6 & $\star$\,\hyperlink{lec-NEU-6}{Short-term \& Long-term Memory} & \href{https://agenticllms.github.io/week4a.html}{[agenticllms.github.io]} \href{https://agenticllms.github.io/week4a-slides.html}{[agenticllms.github.io]} \href{https://github.com/agenticllms/lecture-files/tree/main/W04a}{[GitHub]} & \hyperlink{mod-MEM}{\tg{MEM}} \hyperlink{mod-ARCH}{\tg{ARCH}} \hyperlink{mod-DEPLOY}{\tg{DEPLOY}}\\
\hyperlink{course-NEU}{Northeastern} & 7 & \hyperlink{lec-NEU-7}{Context Management} & \href{https://agenticllms.github.io/week4b.html}{[agenticllms.github.io]} \href{https://github.com/agenticllms/lecture-files/tree/main/W04b}{[GitHub]} & \hyperlink{mod-MEM}{\tg{MEM}} \hyperlink{mod-ARCH}{\tg{ARCH}}\\
\hyperlink{course-NEU}{Northeastern} & 10 & \hyperlink{lec-NEU-10}{Planning \& Task Decomposition} & -- & \hyperlink{mod-RSN}{\tg{RSN}} \hyperlink{mod-ARCH}{\tg{ARCH}}\\
\hyperlink{course-NEU}{Northeastern} & 14 & \hyperlink{lec-NEU-14}{Multi-Agent Orchestration III} & -- & \hyperlink{mod-MULTI}{\tg{MULTI}} \hyperlink{mod-ARCH}{\tg{ARCH}}\\
\hyperlink{course-UB}{Buffalo} & 1 & \hyperlink{lec-UB-1}{What is agentic engineering?} & -- & \tg{ORG} \hyperlink{mod-ARCH}{\tg{ARCH}}\\
\hyperlink{course-UB}{Buffalo} & 6 & \hyperlink{lec-UB-6}{Agent workflows} & -- & \hyperlink{mod-ARCH}{\tg{ARCH}}\\
\hyperlink{course-UVA}{UVA} & 3 & $\star$\,\hyperlink{lec-UVA-3}{Topic 2: Agent Frameworks} & \href{https://www.cs.virginia.edu/\%7Ermw7my/Courses/AgenticAISpring2026/Topic2Frameworks/Agent\%20Control\%20Flow\%20Frameworks.html}{[cs.virginia.edu]} & \hyperlink{mod-ARCH}{\tg{ARCH}} \hyperlink{mod-MULTI}{\tg{MULTI}}\\
\hyperlink{course-UVA}{UVA} & 5 & \hyperlink{lec-UVA-5}{Topic 3: Agent Tool Use (online)} & \href{https://www.cs.virginia.edu/\%7Ermw7my/Courses/AgenticAISpring2026/Topic3Tools/tools.html}{[cs.virginia.edu]} & \hyperlink{mod-TOOL}{\tg{TOOL}} \hyperlink{mod-ARCH}{\tg{ARCH}}\\
\hyperlink{course-UVA}{UVA} & 7 & \hyperlink{lec-UVA-7}{Topic 4: Exploring Tools} & \href{https://www.cs.virginia.edu/\%7Ermw7my/Courses/AgenticAISpring2026/Topic4Exploring/exploring.html}{[cs.virginia.edu]} \href{https://github.com/HenryKautz/MoltReader}{[GitHub]} & \hyperlink{mod-TOOL}{\tg{TOOL}} \hyperlink{mod-RSN}{\tg{RSN}} \hyperlink{mod-ARCH}{\tg{ARCH}}\\
\hyperlink{course-UVA}{UVA} & 17 & \hyperlink{lec-UVA-17}{Topic 11: Coding Agents} & \href{https://www.cs.virginia.edu/\%7Ermw7my/Courses/AgenticAISpring2026/Topic11Coding/coding.html}{[cs.virginia.edu]} & \hyperlink{mod-CODE}{\tg{CODE}} \hyperlink{mod-ARCH}{\tg{ARCH}}\\
\hyperlink{course-UW}{UW-Madison} & 2 & $\star$\,\hyperlink{lec-UW-2}{Building agents} & \href{https://github.com/barghouthi/ai-agents/blob/main/demos/agents.ipynb}{[GitHub]} & \hyperlink{mod-ARCH}{\tg{ARCH}} \hyperlink{mod-TOOL}{\tg{TOOL}}\\
\hyperlink{course-UW}{UW-Madison} & 13 & \hyperlink{lec-UW-13}{Optimizing agents I} & -- & \hyperlink{mod-ARCH}{\tg{ARCH}} \hyperlink{mod-TRAIN}{\tg{TRAIN}}\\
\hyperlink{course-UW}{UW-Madison} & 14 & \hyperlink{lec-UW-14}{Optimizing agents II} & -- & \hyperlink{mod-ARCH}{\tg{ARCH}} \hyperlink{mod-TRAIN}{\tg{TRAIN}}\\
\hyperlink{course-UW}{UW-Madison} & 15 & \hyperlink{lec-UW-15}{Self-improving agents I} & -- & \hyperlink{mod-ARCH}{\tg{ARCH}} \hyperlink{mod-WEB}{\tg{WEB}}\\
\hyperlink{course-UW}{UW-Madison} & 20 & \hyperlink{lec-UW-20}{Multi-agent systems II} & -- & \hyperlink{mod-MULTI}{\tg{MULTI}} \hyperlink{mod-ARCH}{\tg{ARCH}}\\
\hyperlink{course-UW}{UW-Madison} & 21 & \hyperlink{lec-UW-21}{Memory I} & -- & \hyperlink{mod-MEM}{\tg{MEM}} \hyperlink{mod-ARCH}{\tg{ARCH}}\\
\hyperlink{course-JHU}{JHU} & 7 & \hyperlink{lec-JHU-7}{Iterative Refinement and Agentic Framework} & \href{https://machine-programming.github.io/lectures/lecture-7.pdf}{[PDF]} & \hyperlink{mod-ARCH}{\tg{ARCH}} \hyperlink{mod-CODE}{\tg{CODE}}\\
\hyperlink{course-JHU}{JHU} & 8 & \hyperlink{lec-JHU-8}{Agentic Frameworks: Model Context Protocol} & \href{https://machine-programming.github.io/lectures/lecture-8.pdf}{[PDF]} & \hyperlink{mod-TOOL}{\tg{TOOL}} \hyperlink{mod-ARCH}{\tg{ARCH}}\\
\hyperlink{course-JHU}{JHU} & 9 & \hyperlink{lec-JHU-9}{Agentic Frameworks: Model Context Protocol (2)} & \href{https://machine-programming.github.io/lectures/lecture-9.pdf}{[PDF]} & \hyperlink{mod-TOOL}{\tg{TOOL}} \hyperlink{mod-ARCH}{\tg{ARCH}}\\
\hyperlink{course-JHU}{JHU} & 10 & \hyperlink{lec-JHU-10}{Agentic Frameworks: Selected Topics} & \href{https://machine-programming.github.io/lectures/lecture-10.pdf}{[PDF]} & \hyperlink{mod-ARCH}{\tg{ARCH}} \hyperlink{mod-CODE}{\tg{CODE}}\\
\hyperlink{course-CU}{CU Boulder} & 2 & \hyperlink{lec-CU-2}{Survey of agents (Dig)} & \href{https://danny.cs.colorado.edu/courses/csci7000-005_Sp26/pdfs/W1_L2_SurveyOfAgents.pdf}{[PDF]} & \hyperlink{mod-ARCH}{\tg{ARCH}} \hyperlink{mod-CODE}{\tg{CODE}}\\
\hyperlink{course-CU}{CU Boulder} & 3 & \hyperlink{lec-CU-3}{Agentic SE roadmap (Dig)} & \href{https://danny.cs.colorado.edu/courses/csci7000-005_Sp26/pdfs/W2_L2_Roadmap.pdf}{[PDF]} & \hyperlink{mod-CODE}{\tg{CODE}} \hyperlink{mod-ARCH}{\tg{ARCH}}\\
\hyperlink{course-CU}{CU Boulder} & 13 & \hyperlink{lec-CU-13}{SWE-agent \& AutoCodeRover (students)} & -- & \hyperlink{mod-CODE}{\tg{CODE}} \hyperlink{mod-ARCH}{\tg{ARCH}}\\
\hyperlink{course-CU}{CU Boulder} & 14 & \hyperlink{lec-CU-14}{AutoGen \& MetaGPT (students)} & -- & \hyperlink{mod-MULTI}{\tg{MULTI}} \hyperlink{mod-ARCH}{\tg{ARCH}}\\
\hyperlink{course-CU}{CU Boulder} & 16 & \hyperlink{lec-CU-16}{LangGraph tutorial \& OpenHands} & -- & \hyperlink{mod-ARCH}{\tg{ARCH}} \hyperlink{mod-CODE}{\tg{CODE}}\\
\hyperlink{course-CU}{CU Boulder} & 18 & \hyperlink{lec-CU-18}{RefAgent \& Agentic Skills (students)} & -- & \hyperlink{mod-CODE}{\tg{CODE}} \hyperlink{mod-ARCH}{\tg{ARCH}} \hyperlink{mod-TOOL}{\tg{TOOL}}\\
\hyperlink{course-MANN}{Mannheim} & 4 & \hyperlink{lec-MANN-4}{LLM Agents and Tool Use} & \href{https://www.uni-mannheim.de/media/Einrichtungen/dws/Files_Teaching/Large_Language_Models_and_Agents/FSS2026/IE685_LA_04_LLMAgentsAndToolUse.pdf}{[PDF]} \href{https://www.uni-mannheim.de/media/Einrichtungen/dws/Files_Teaching/Large_Language_Models_and_Agents/FSS2026/IE685_LA_EX_Agent_Programming_LangChain_MCP.ipynb.zip}{[notebook zip]} & \hyperlink{mod-ARCH}{\tg{ARCH}} \hyperlink{mod-TOOL}{\tg{TOOL}}\\
\hyperlink{course-MANN}{Mannheim} & 6 & \hyperlink{lec-MANN-6}{Context Engineering} & \href{https://www.uni-mannheim.de/media/Einrichtungen/dws/Files_Teaching/Large_Language_Models_and_Agents/FSS2026/IE685_LA_06_ContextEngineering.pdf}{[PDF]} \href{https://www.uni-mannheim.de/media/Einrichtungen/dws/Files_Teaching/Large_Language_Models_and_Agents/FSS2026/deep_agent.zip}{[zip]} & \hyperlink{mod-MEM}{\tg{MEM}} \hyperlink{mod-ARCH}{\tg{ARCH}}\\
\hyperlink{course-S329A}{Stanford 329A} & 7 & \hyperlink{lec-S329A-7}{Open-Ended Evolution of Self-Improving Agents} & -- & \hyperlink{mod-ARCH}{\tg{ARCH}} \hyperlink{mod-CODE}{\tg{CODE}} \hyperlink{mod-APP}{\tg{APP}}\\
\hyperlink{course-S329A}{Stanford 329A} & 13 & \hyperlink{lec-S329A-13}{Agentic Frameworks for Software Engineering} & -- & \hyperlink{mod-CODE}{\tg{CODE}} \hyperlink{mod-ARCH}{\tg{ARCH}}\\
\hyperlink{course-BRKAG}{Berkeley '25} & 3 & \hyperlink{lec-BRKAG-3}{Evolution of system designs from an AI engineer perspective} & \href{https://rdi.berkeley.edu/agentic-ai/slides/Yangqing.pdf}{[slides]} \href{https://www.youtube.com/watch?v=xqRAS6rAouo\&list=PLS01nW3RtgoqGkm4UeqNeZLccW-OGc1fJ\&index=2}{[video]} & \hyperlink{mod-ARCH}{\tg{ARCH}} \hyperlink{mod-DEPLOY}{\tg{DEPLOY}}\\
\hyperlink{course-BRKADV}{Berkeley Adv.} & 3 & \hyperlink{lec-BRKADV-3}{On Reasoning, Memory, and Planning of Language Agents} & \href{https://rdi.berkeley.edu/adv-llm-agents/slides/language_agents_YuSu_Berkeley.pdf}{[slides]} \href{https://www.youtube.com/live/zvI4UN2_i-w}{[video]} & \hyperlink{mod-RSN}{\tg{RSN}} \hyperlink{mod-MEM}{\tg{MEM}} \hyperlink{mod-ARCH}{\tg{ARCH}}\\
\hyperlink{course-BRKADV}{Berkeley Adv.} & 11 & \hyperlink{lec-BRKADV-11}{Abstraction and Discovery with LLM Agents} & \href{https://rdi.berkeley.edu/adv-llm-agents/slides/swarat.pdf}{[slides]} \href{https://www.youtube.com/live/IHc0TEMrEdY}{[video]} & \hyperlink{mod-RSN}{\tg{RSN}} \hyperlink{mod-ARCH}{\tg{ARCH}} \hyperlink{mod-APP}{\tg{APP}}\\
\hyperlink{course-S329T}{Stanford 329T} & 1 & \hyperlink{lec-S329T-1}{Course introduction and overview} & \href{https://web.stanford.edu/class/cs329t/slides/Lecture\%201\%20-\%20CS\%20329T\%20Fall\%202025.pdf}{[slides]} & \tg{ORG} \hyperlink{mod-FND}{\tg{FND}} \hyperlink{mod-ARCH}{\tg{ARCH}}\\
\hyperlink{course-S329T}{Stanford 329T} & 3 & \hyperlink{lec-S329T-3}{Agentic AI} & \href{https://web.stanford.edu/class/cs329t/slides/Lecture\%203\%20--\%20CS\%20329T\%20Fall\%202025.pdf}{[slides]} & \hyperlink{mod-ARCH}{\tg{ARCH}} \hyperlink{mod-EVAL}{\tg{EVAL}}\\
\hyperlink{course-S329T}{Stanford 329T} & 15 & \hyperlink{lec-S329T-15}{Multi-modal frontier models with agentic capabilities (guest: Deepak Ramachandran, Google DeepMind)} & -- & \hyperlink{mod-FND}{\tg{FND}} \hyperlink{mod-ARCH}{\tg{ARCH}}\\
\bottomrule\end{longtable}}
\section{Key readings}
\begin{itemize}
\item Anthropic, \href{https://www.anthropic.com/engineering/building-effective-agents}{\emph{Building Effective Agents}}
\item Yao et al., \href{https://arxiv.org/abs/2210.03629}{\emph{ReAct}} (2022)
\item Sumers et al., \href{https://arxiv.org/abs/2309.02427}{\emph{CoALA: Cognitive Architectures for Language Agents}} (2023)
\item Chip Huyen, \href{https://huyenchip.com/2025/01/07/agents.html}{\emph{Agents}} (2025)
\item Wang et al., \href{https://arxiv.org/abs/2308.11432}{\emph{A Survey on LLM-based Autonomous Agents}} (2024)
\end{itemize}
\section{Hands-on}
\begin{itemize}
\item \hyperlink{asg-UW-1}{\textbf{UW-Madison: Build a Coding Agent}} --- A CLI coding agent like Claude Code / Cursor CLI / Gemini CLI, using raw LLM API calls with no agent frameworks (LangChain, Strands, CrewAI banned; thin wrappers like litellm allowed). It must parse and execute tool calls, have a polished UI, a README, and one non-trivial program the agent built. \textit{Spec:} \href{https://github.com/barghouthi/ai-agents/blob/main/assignment.md}{[GitHub]}
\item \hyperlink{asg-CMU-1}{\textbf{CMU: A1: Build an Agent Harness}} --- A domain-agnostic ReAct \texttt{Agent} base class: prompt construction, tool-call loop, step limit, model-generated context compaction. On top of it, a \texttt{CodeAgent} (terminal/message tools, skill loader with progressive disclosure via \texttt{invoke\_\allowbreak{}skill}) that fixes a chess-app bug and solves SWE-bench \texttt{django\_\allowbreak{}\_\allowbreak{}django-15368} with and without compaction. Then a \texttt{ChessAgent} (play\_\allowbreak{}move, stateless simulate\_\allowbreak{}move, sandboxed run\_\allowbreak{}python, chess skill) plus an observation A/B experiment. Uses Modal sandboxes, an OpenAI-compatible LLM (default DeepSeek-V4-Flash), uv, and pytest. 100-pt rubric. \textit{Starter code:} \href{https://github.com/cmu-agents/assignment-1}{[GitHub]} \href{https://github.com/cmu-agents/assignment-1/blob/main/ASSIGNMENT.md}{[GitHub]} \href{https://github.com/cmu-agents/chess-app}{[GitHub]}
\item \hyperlink{asg-UVA-2}{\textbf{UVA: Topic 2 portfolio: LangGraph agents}} --- Routed, parallel, multi-model LangGraph chat with history and crash recovery; optional web summarizer. \textit{Starter code:} \href{https://www.cs.virginia.edu/\%7Ermw7my/Courses/AgenticAISpring2026/Topic2Frameworks/Agent\%20Control\%20Flow\%20Frameworks.html}{[cs.virginia.edu]}
\item \hyperlink{asg-UB-3}{\textbf{Buffalo: Lab 3: Workflow Graphs and State Machines}} --- Agentic workflow as a graph/state machine with explicit states, transitions, failure handling and logging. \textit{(spec only)}
\item \hyperlink{asg-UB-1}{\textbf{Buffalo: Lab 1: A Minimal Tool-Using Agent}} --- Python agent choosing among a calculator, a file reader and a search/lookup tool. Submit a design note and successful/failed runs. \textit{(spec only)}
\end{itemize}
\textbf{Related supplements:} \hyperref[sup-mcp]{Model Context Protocol docs (version 2026-07-28)}, \hyperref[sup-openclaw]{OpenClaw documentation}, \hyperref[sup-agentsec]{MIT 6.566 guest lecture: AI Agent Security (Anish Athalye, April 2026)}.
\chapter{Module 3: Tool Use and Function Calling}\hypertarget{mod-TOOL}{}
{\small Tag: \tg{TOOL} (MCP-focused lectures are in \hyperlink{mod-MCP}{Module 4})}
\section{Learning objectives}
\begin{itemize}
\item Explain how tool calls are encoded in chat templates and parsed by the harness.
\item Design agent-ergonomic tools (names, schemas, errors, granularity).
\item Implement manual tool dispatch, then framework tool nodes, and measure tool-selection accuracy.
\end{itemize}
\section{Core concepts}
Tools are how an agent acts. Three layers matter:
\begin{itemize}
\item \textbf{Encoding.} The model emits a structured call (name + JSON arguments) inside its chat template, the harness executes it, and the result returns as a tool message. Constrained decoding (e.g.\ XGrammar) can guarantee well-formed calls. Benchmarks such as BFCL measure function-calling accuracy.
\item \textbf{Design.} Tools are an interface for a non-human user. Prefer a few high-level, composable tools with clear names, typed parameters and \emph{informative errors} that tell the model how to recover. ``Code as action'' (CodeAct) replaces many tools with one interpreter. Stanford's readings include Anthropic's \emph{Writing effective tools for agents} and ``APIs don't make good MCP tools''.
\item \textbf{Selection at scale.} With dozens of tools, retrieval over tool descriptions, or a small classifier as in TinyAgent, keeps prompts short and accurate. For hundreds of tools see the code-execution pattern in \hyperlink{mod-MCP}{Module~4 (MCP and Agent Interoperability)}.
\end{itemize}
The research lineage runs from Toolformer (self-taught API calls) to Gorilla (API retrieval), PAL (program-aided reasoning), MRKL and ``What are tools anyway?''.
\section{Curated learning sequence}
\begin{enumerate}
\item \hyperlink{lec-CMU-2}{CMU~2 \emph{Capabilities 1: Tool Use}} --- Tool use end to end: APIs, chat templates, CodeAct, constrained decoding, BFCL (slides + video).
\item \hyperlink{lec-UVA-5}{UVA~5 \emph{Topic 3: Agent Tool Use (online)}} --- Lab: manual tool handling with a calculator, the OpenAI tool format, tool forcing.
\item \hyperlink{lec-UVA-7}{UVA~7 \emph{Topic 4: Exploring Tools}} --- ToolNode versus \texttt{create\_react\_agent}, parallel tool dispatch, a 2-hour mini project.
\item \hyperlink{lec-NEU-5}{Northeastern~5 \emph{Tool-Calling}} --- ReAct support agent, tools-versus-model design rules, deploy to HF Spaces.
\item \hyperlink{lec-CU-11}{CU Boulder~11 \emph{Reflexion \& Toolformer (students)}} and \hyperlink{lec-CU-12}{CU Boulder~12 \emph{Gorilla \& PAL (students)}} --- Research reading: Reflexion/Toolformer, then Gorilla/PAL.
\item Supplement: \hyperref[sup-tinyagent]{TinyAgent} --- Tool RAG and small-model function calling at the edge.
\end{enumerate}
\section{Lecture map (20 lectures)}
{\footnotesize\begin{longtable}{P{1.75cm} P{0.4cm} P{5.75cm} P{4.2cm} P{2.65cm}}
\toprule Course & \# & Title (links to catalogue row) & Materials & Tags\\ \midrule\endhead
\hyperlink{course-CMU}{CMU} & 2 & $\star$\,\hyperlink{lec-CMU-2}{Capabilities 1: Tool Use} & \href{https://www.cmu-agents.com/slides/lecture-02-tool-use.pdf}{[slides]} \href{https://www.cmu-agents.com/\#/schedule}{[video]} & \hyperlink{mod-TOOL}{\tg{TOOL}}\\
\hyperlink{course-STAN26}{Stanford '26} & 5 & \hyperlink{lec-STAN26-5}{All about agent skills (incl. web skills)} & \href{https://themodernsoftware.dev/}{[course site]} & \hyperlink{mod-MEM}{\tg{MEM}} \hyperlink{mod-TOOL}{\tg{TOOL}}\\
\hyperlink{course-STAN25}{Stanford '25} & 3 & \hyperlink{lec-STAN25-3}{Building a coding agent from scratch} & \href{https://docs.google.com/presentation/d/11CP26VhsjnZOmi9YFgLlonzdib9BLyAlgc4cEvC5Fps/edit?usp=sharing}{[slides]} \href{https://drive.google.com/file/d/1YtpKFVG13DHyQ2i3HOtwyVJOV90nWeL2/view?usp=drive_link}{[drive.google.com]} & \hyperlink{mod-ARCH}{\tg{ARCH}} \hyperlink{mod-TOOL}{\tg{TOOL}} \hyperlink{mod-CODE}{\tg{CODE}}\\
\hyperlink{course-NEU}{Northeastern} & 4 & \hyperlink{lec-NEU-4}{Simple Agent Architecture} & \href{https://agenticllms.github.io/week3a.html}{[agenticllms.github.io]} \href{https://github.com/agenticllms/lecture-files/tree/main/W03a}{[GitHub]} & \hyperlink{mod-ARCH}{\tg{ARCH}} \hyperlink{mod-TOOL}{\tg{TOOL}} \hyperlink{mod-MEM}{\tg{MEM}}\\
\hyperlink{course-NEU}{Northeastern} & 5 & $\star$\,\hyperlink{lec-NEU-5}{Tool-Calling} & \href{https://agenticllms.github.io/week3b.html}{[agenticllms.github.io]} \href{https://agenticllms.github.io/week3b-slides.html}{[agenticllms.github.io]} \href{https://github.com/agenticllms/lecture-files/tree/main/W03b}{[GitHub]} & \hyperlink{mod-TOOL}{\tg{TOOL}} \hyperlink{mod-RSN}{\tg{RSN}} \hyperlink{mod-DEPLOY}{\tg{DEPLOY}}\\
\hyperlink{course-UB}{Buffalo} & 3 & \hyperlink{lec-UB-3}{Tool/function calling and ReAct} & -- & \hyperlink{mod-TOOL}{\tg{TOOL}} \hyperlink{mod-RSN}{\tg{RSN}}\\
\hyperlink{course-UVA}{UVA} & 5 & $\star$\,\hyperlink{lec-UVA-5}{Topic 3: Agent Tool Use (online)} & \href{https://www.cs.virginia.edu/\%7Ermw7my/Courses/AgenticAISpring2026/Topic3Tools/tools.html}{[cs.virginia.edu]} & \hyperlink{mod-TOOL}{\tg{TOOL}} \hyperlink{mod-ARCH}{\tg{ARCH}}\\
\hyperlink{course-UVA}{UVA} & 6 & \hyperlink{lec-UVA-6}{Topic 3 lab continued (online)} & -- & \hyperlink{mod-TOOL}{\tg{TOOL}}\\
\hyperlink{course-UVA}{UVA} & 7 & $\star$\,\hyperlink{lec-UVA-7}{Topic 4: Exploring Tools} & \href{https://www.cs.virginia.edu/\%7Ermw7my/Courses/AgenticAISpring2026/Topic4Exploring/exploring.html}{[cs.virginia.edu]} \href{https://github.com/HenryKautz/MoltReader}{[GitHub]} & \hyperlink{mod-TOOL}{\tg{TOOL}} \hyperlink{mod-RSN}{\tg{RSN}} \hyperlink{mod-ARCH}{\tg{ARCH}}\\
\hyperlink{course-UW}{UW-Madison} & 2 & \hyperlink{lec-UW-2}{Building agents} & \href{https://github.com/barghouthi/ai-agents/blob/main/demos/agents.ipynb}{[GitHub]} & \hyperlink{mod-ARCH}{\tg{ARCH}} \hyperlink{mod-TOOL}{\tg{TOOL}}\\
\hyperlink{course-UW}{UW-Madison} & 4 & \hyperlink{lec-UW-4}{Reasoning and acting II} & -- & \hyperlink{mod-TOOL}{\tg{TOOL}} \hyperlink{mod-RSN}{\tg{RSN}}\\
\hyperlink{course-CU}{CU Boulder} & 9 & \hyperlink{lec-CU-9}{ReAct \& Web API code gen (students)} & -- & \hyperlink{mod-RSN}{\tg{RSN}} \hyperlink{mod-TOOL}{\tg{TOOL}} \hyperlink{mod-EVAL}{\tg{EVAL}}\\
\hyperlink{course-CU}{CU Boulder} & 11 & $\star$\,\hyperlink{lec-CU-11}{Reflexion \& Toolformer (students)} & -- & \hyperlink{mod-RSN}{\tg{RSN}} \hyperlink{mod-TOOL}{\tg{TOOL}} \hyperlink{mod-TRAIN}{\tg{TRAIN}}\\
\hyperlink{course-CU}{CU Boulder} & 12 & $\star$\,\hyperlink{lec-CU-12}{Gorilla \& PAL (students)} & -- & \hyperlink{mod-TOOL}{\tg{TOOL}} \hyperlink{mod-RSN}{\tg{RSN}}\\
\hyperlink{course-CU}{CU Boulder} & 18 & \hyperlink{lec-CU-18}{RefAgent \& Agentic Skills (students)} & -- & \hyperlink{mod-CODE}{\tg{CODE}} \hyperlink{mod-ARCH}{\tg{ARCH}} \hyperlink{mod-TOOL}{\tg{TOOL}}\\
\hyperlink{course-CU}{CU Boulder} & 23 & \hyperlink{lec-CU-23}{SafePro \& ToolDreamer (students)} & -- & \hyperlink{mod-SAFE}{\tg{SAFE}} \hyperlink{mod-TOOL}{\tg{TOOL}}\\
\hyperlink{course-S329A}{Stanford 329A} & 4 & \hyperlink{lec-S329A-4}{Learning from Feedback with Tools/Code} & \href{https://www.youtube.com/watch?v=Lxh9RF5S-K0}{[video]} & \hyperlink{mod-TOOL}{\tg{TOOL}} \hyperlink{mod-CODE}{\tg{CODE}} \hyperlink{mod-TRAIN}{\tg{TRAIN}}\\
\hyperlink{course-BRKADV}{Berkeley Adv.} & 5 & \hyperlink{lec-BRKADV-5}{Coding Agents and AI for Vulnerability Detection} & \href{https://rdi.berkeley.edu/adv-llm-agents/slides/Code\%20Agents\%20and\%20AI\%20for\%20Vulnerability\%20Detection.pdf}{[slides]} \href{https://www.youtube.com/live/JCk6qJtaCSU}{[video]} & \hyperlink{mod-CODE}{\tg{CODE}} \hyperlink{mod-TOOL}{\tg{TOOL}} \hyperlink{mod-SAFE}{\tg{SAFE}}\\
\hyperlink{course-BRKADV}{Berkeley Adv.} & 9 & \hyperlink{lec-BRKADV-9}{Language models for autoformalization and theorem proving} & \href{https://rdi.berkeley.edu/adv-llm-agents/slides/mathverification.pdf}{[slides]} \href{https://www.youtube.com/live/cLhWEyMQ4mQ}{[video]} & \hyperlink{mod-RSN}{\tg{RSN}} \hyperlink{mod-APP}{\tg{APP}} \hyperlink{mod-TOOL}{\tg{TOOL}}\\
\hyperlink{course-S329T}{Stanford 329T} & 10 & \hyperlink{lec-S329T-10}{Data curation and RL to optimize agents (guest: Trung Vu, Bespoke Labs)} & \href{https://web.stanford.edu/class/cs329t/slides/Lecture\%208\%20-\%20\%5B10_23_2025\%5D\%20Reinforcement\%20Learning\%20\%20for\%20Agents.pdf}{[slides]} & \hyperlink{mod-TRAIN}{\tg{TRAIN}} \hyperlink{mod-TOOL}{\tg{TOOL}} \hyperlink{mod-EVAL}{\tg{EVAL}}\\
\bottomrule\end{longtable}}
\section{Key readings}
\begin{itemize}
\item Schick et al., \href{https://arxiv.org/abs/2302.04761}{\emph{Toolformer}} (2023)
\item Patil et al., \href{https://arxiv.org/abs/2305.15334}{\emph{Gorilla}} (2023)
\item Wang et al., \href{https://arxiv.org/abs/2402.01030}{\emph{CodeAct: Executable Code Actions}} (2024)
\item Gao et al., \href{https://arxiv.org/abs/2211.10435}{\emph{PAL}} (2022)
\item \href{https://arxiv.org/abs/2403.15452}{\emph{What Are Tools Anyway?}} (2024)
\item Anthropic, \href{https://www.anthropic.com/engineering/writing-tools-for-agents}{\emph{Writing effective tools for agents}}
\end{itemize}
\section{Hands-on}
\begin{itemize}
\item \hyperlink{asg-UVA-3}{\textbf{UVA: Topic 3 portfolio: Tool use}} --- Ollama parallel timing; manual tool calling with a calculator; LangGraph multi-tool agent with tool\_\allowbreak{}map; persistent graph with Mermaid diagram. \textit{Starter code:} \href{https://www.cs.virginia.edu/\%7Ermw7my/Courses/AgenticAISpring2026/Topic3Tools/tools.html}{[cs.virginia.edu]}
\item \hyperlink{asg-UB-1}{\textbf{Buffalo: Lab 1: A Minimal Tool-Using Agent}} --- Python agent choosing among a calculator, a file reader and a search/lookup tool. Submit a design note and successful/failed runs. \textit{(spec only)}
\item \hyperlink{asg-NEU-1}{\textbf{Northeastern: DADA1: The Trip Planner Agent}} --- A conversational day-trip planner with tools: get\_\allowbreak{}weather (Open-Meteo), find\_\allowbreak{}attractions/restaurants/transit (Tavily), find\_\allowbreak{}events (Ticketmaster). It has memory-based plan revision, a structured Itinerary/DayPlan output, 8 test prompts, and a Gradio app deployed on Render. Graded on tool routing, honesty at data limits, sources and memory. Optional extra credit: a find\_\allowbreak{}performer\_\allowbreak{}events chained tool. \textit{Starter code:} \href{https://github.com/agenticllms/lecture-files/tree/main/DADA1}{[GitHub]}
\item \hyperlink{asg-UVA-4}{\textbf{UVA: Topic 4: 2-Hour Agent Project}} --- Partner-built mini agent using LangGraph/third-party tools. \textit{Starter code:} \href{https://www.cs.virginia.edu/\%7Ermw7my/Courses/AgenticAISpring2026/Topic4Exploring/exploring.html}{[cs.virginia.edu]}
\end{itemize}
\textbf{Related supplements:} \hyperref[sup-codeexec]{Anthropic Engineering: Code execution with MCP}, \hyperref[sup-tinyagent]{Berkeley BAIR: TinyAgent, Function Calling at the Edge}, \hyperref[sup-agentsec]{MIT 6.566 guest lecture: AI Agent Security (Anish Athalye, April 2026)}, \hyperref[sup-anthauto]{Anthropic: Building effective agent automations (Martin and Avilla, October 2026)}.
\chapter{Module 4: MCP and Agent Interoperability}\hypertarget{mod-MCP}{}
{\small Tag: \tg{TOOL} (lectures tagged \tg{TOOL} whose title or topics mention MCP or A2A)}
\section{Learning objectives}
\begin{itemize}
\item Describe MCP hosts, clients and servers, its primitives and transports, and what changed in the 2026-07-28 revision.
\item Build and test an MCP server with authentication.
\item Know when to switch from direct tool calls to code execution against MCP, and what A2A adds.
\end{itemize}
\section{Core concepts}
The \textbf{Model Context Protocol} standardizes how tools and context reach an agent: write one server, use it in any host. Read the \hyperref[sup-mcp]{MCP supplement} first, because \textbf{the 2026-07-28 revision made the protocol stateless}, and lectures from before mid-2026 teach the older handshake.
\begin{itemize}
\item \textbf{Primitives:} tools, resources and prompts on the server side, and elicitation on the client side. Transports are stdio (local) and Streamable HTTP (remote, OAuth).
\item \textbf{Server craft:} typed, composable tools; dry-run modes for write tools; \texttt{readOnlyHint}; structured errors; rate-limit handling; protocol-level tests (Stanford '26 week~2).
\item \textbf{Security:} an MCP server is a trust boundary. Northeastern devotes two sessions to MCP security and enterprise auth. See also \hyperlink{mod-SAFE}{Module~12 (Safety, Security and Oversight)}.
\item \textbf{Scale:} with many servers, loading every tool schema is wasteful. Anthropic's \hyperref[sup-codeexec]{code-execution pattern} and CU Boulder's MCP design-choice paper argue for progressive discovery and code-based orchestration.
\item \textbf{Agent-to-agent:} A2A (agent cards, registries) connects whole agents rather than tools. UVA runs an A2A trivia tournament.
\end{itemize}
\section{Curated learning sequence}
\begin{enumerate}
\item Supplement: \hyperref[sup-mcp]{MCP docs} --- Architecture, primitives and the 2026-07-28 changes.
\item \hyperlink{lec-STAN26-4}{Stanford '26~4 \emph{Full introduction to MCP and tool-calling}} --- Full introduction to MCP and tool calling (theory, setup, advanced tool design).
\item \hyperlink{lec-STAN25-4}{Stanford '25~4 \emph{Building a custom MCP server}} --- Building a custom MCP server (slides + completed exercise).
\item \hyperlink{lec-JHU-8}{JHU~8 \emph{Agentic Frameworks: Model Context Protocol}} and \hyperlink{lec-JHU-9}{JHU~9 \emph{Agentic Frameworks: Model Context Protocol (2)}} --- Two MCP lectures in a PL/synthesis context (slides).
\item \hyperlink{lec-NEU-16}{Northeastern~16 \emph{Model Context Protocol}} and \hyperlink{lec-NEU-17}{Northeastern~17 \emph{MCP in Agentic Systems}} --- MCP spec, the security landscape, remote MCP and agent auth.
\item \hyperlink{lec-UVA-12}{UVA~12 \emph{Topic 7: MCP and A2A}} --- MCP clients for AI2 Asta tools, then A2A agents, a registry and ngrok.
\item Supplement: \hyperref[sup-codeexec]{Code execution with MCP} --- Then \hyperlink{lec-CU-24}{CU Boulder~24 \emph{Swiss Cheese safety \& MCP design (students)}} for the empirical study.
\end{enumerate}
\section{Lecture map (11 lectures)}
{\footnotesize\begin{longtable}{P{1.75cm} P{0.4cm} P{5.75cm} P{4.2cm} P{2.65cm}}
\toprule Course & \# & Title (links to catalogue row) & Materials & Tags\\ \midrule\endhead
\hyperlink{course-STAN26}{Stanford '26} & 4 & $\star$\,\hyperlink{lec-STAN26-4}{Full introduction to MCP and tool-calling (theory, setup, advanced tool design)} & \href{https://themodernsoftware.dev/}{[course site]} & \hyperlink{mod-TOOL}{\tg{TOOL}}\\
\hyperlink{course-STAN25}{Stanford '25} & 4 & $\star$\,\hyperlink{lec-STAN25-4}{Building a custom MCP server} & \href{https://docs.google.com/presentation/d/1zSC2ra77XOUrJeyS85houg1DU7z9hq5Y4ebagTch-5o/edit?usp=drive_link}{[slides]} \href{https://drive.google.com/file/d/1J6lgZWcxPzpCpjujJSnW1aAkCYF6Yxv3/view?usp=drive_link}{[drive.google.com]} & \hyperlink{mod-TOOL}{\tg{TOOL}}\\
\hyperlink{course-NEU}{Northeastern} & 16 & $\star$\,\hyperlink{lec-NEU-16}{Model Context Protocol} & -- & \hyperlink{mod-TOOL}{\tg{TOOL}} \hyperlink{mod-SAFE}{\tg{SAFE}}\\
\hyperlink{course-NEU}{Northeastern} & 17 & $\star$\,\hyperlink{lec-NEU-17}{MCP in Agentic Systems} & -- & \hyperlink{mod-TOOL}{\tg{TOOL}} \hyperlink{mod-SAFE}{\tg{SAFE}} \hyperlink{mod-DEPLOY}{\tg{DEPLOY}}\\
\hyperlink{course-UVA}{UVA} & 12 & $\star$\,\hyperlink{lec-UVA-12}{Topic 7: MCP and A2A} & \href{https://www.cs.virginia.edu/\%7Ermw7my/Courses/AgenticAISpring2026/Topic7MCP/mcp.html}{[cs.virginia.edu]} & \hyperlink{mod-TOOL}{\tg{TOOL}} \hyperlink{mod-MULTI}{\tg{MULTI}}\\
\hyperlink{course-JHU}{JHU} & 8 & $\star$\,\hyperlink{lec-JHU-8}{Agentic Frameworks: Model Context Protocol} & \href{https://machine-programming.github.io/lectures/lecture-8.pdf}{[PDF]} & \hyperlink{mod-TOOL}{\tg{TOOL}} \hyperlink{mod-ARCH}{\tg{ARCH}}\\
\hyperlink{course-JHU}{JHU} & 9 & $\star$\,\hyperlink{lec-JHU-9}{Agentic Frameworks: Model Context Protocol (2)} & \href{https://machine-programming.github.io/lectures/lecture-9.pdf}{[PDF]} & \hyperlink{mod-TOOL}{\tg{TOOL}} \hyperlink{mod-ARCH}{\tg{ARCH}}\\
\hyperlink{course-CU}{CU Boulder} & 24 & \hyperlink{lec-CU-24}{Swiss Cheese safety \& MCP design (students)} & -- & \hyperlink{mod-SAFE}{\tg{SAFE}} \hyperlink{mod-TOOL}{\tg{TOOL}}\\
\hyperlink{course-MANN}{Mannheim} & 4 & \hyperlink{lec-MANN-4}{LLM Agents and Tool Use} & \href{https://www.uni-mannheim.de/media/Einrichtungen/dws/Files_Teaching/Large_Language_Models_and_Agents/FSS2026/IE685_LA_04_LLMAgentsAndToolUse.pdf}{[PDF]} \href{https://www.uni-mannheim.de/media/Einrichtungen/dws/Files_Teaching/Large_Language_Models_and_Agents/FSS2026/IE685_LA_EX_Agent_Programming_LangChain_MCP.ipynb.zip}{[notebook zip]} & \hyperlink{mod-ARCH}{\tg{ARCH}} \hyperlink{mod-TOOL}{\tg{TOOL}}\\
\hyperlink{course-UCHI}{UChicago} & 5 & \hyperlink{lec-UCHI-5}{Sampling, Fine Tuning, Benchmarking, and Tools} & \href{https://github.com/KnowledgeLab/AI-Agents-for-Social-Science-and-Society-2026/blob/main/Tutorials-Homework_Notebooks/Week_5_2026.ipynb}{[GitHub]} & \hyperlink{mod-TRAIN}{\tg{TRAIN}} \hyperlink{mod-EVAL}{\tg{EVAL}} \hyperlink{mod-TOOL}{\tg{TOOL}}\\
\hyperlink{course-S329T}{Stanford 329T} & 8 & \hyperlink{lec-S329T-8}{MCP, effective tool descriptions, evals (guests: Yusuf Ozuysal, Ishita Gupta)} & \href{https://web.stanford.edu/class/cs329t/slides/Model\%20Context\%20Protocol\%20\%28\%20Yusuf\%20Ozuysal\%29.pdf}{[slides (Ozuysal)]} \href{https://web.stanford.edu/class/cs329t/slides/ishita-mcp.pdf}{[slides (Gupta)]} & \hyperlink{mod-TOOL}{\tg{TOOL}} \hyperlink{mod-EVAL}{\tg{EVAL}}\\
\bottomrule\end{longtable}}
\section{Key readings}
\begin{itemize}
\item MCP specification and \href{https://modelcontextprotocol.io/docs/2026-07-28/tutorials/security/security_best_practices}{security best practices}
\item Anthropic, \href{https://www.anthropic.com/engineering/code-execution-with-mcp}{\emph{Code execution with MCP}} (2025)
\item Felendler et al., \href{https://arxiv.org/abs/2602.15945}{\emph{From Tool Orchestration to Code Execution: MCP Design Choices}} (2026)
\item \href{https://arxiv.org/abs/2503.23278}{\emph{MCP Landscape \& Security}} (arXiv 2503.23278)
\end{itemize}
\section{Hands-on}
\begin{itemize}
\item \hyperlink{asg-STAN26-2}{\textbf{Stanford '26: Week 2 --- Build an MCP Server for a Third-Party API}} --- FastMCP server over an OAuth 2.0 API with at least 3 composing tools, a write tool with dry-run, structured errors, readOnlyHint, full auth-code flow with refresh, and a protocol-level test. \textit{Starter code:} \href{https://github.com/mihail911/modern-software-dev-assignments/blob/master/week2/assignment.md}{[GitHub]}
\item \hyperlink{asg-STAN25-3}{\textbf{Stanford '25: Week 3 --- Build a Custom MCP Server}} --- MCP server wrapping a real external API: at least 2 typed tools, robust error and rate-limit handling, STDIO or remote HTTP, optional OAuth2. \textit{Starter code:} \href{https://github.com/mihail911/modern-software-dev-assignments/tree/fall2025/week3}{[GitHub]}
\item \hyperlink{asg-UVA-7}{\textbf{UVA: Topic 7: MCP + A2A}} --- MCP client for AI2 Asta tools; A2A agents, registry and trivia tournament over ngrok. \textit{Starter code:} \href{https://www.cs.virginia.edu/\%7Ermw7my/Courses/AgenticAISpring2026/Topic7MCP/mcp.html}{[cs.virginia.edu]}
\item \hyperlink{asg-MANN-3}{\textbf{Mannheim: Exercise: Agent Programming with LangChain + MCP}} --- Tool-using agents with LangChain and MCP. \textit{Starter code:} \href{https://www.uni-mannheim.de/media/Einrichtungen/dws/Files_Teaching/Large_Language_Models_and_Agents/FSS2026/IE685_LA_EX_Agent_Programming_LangChain_MCP.ipynb.zip}{[notebook zip]}
\end{itemize}
\textbf{Related supplements:} \hyperref[sup-mcp]{Model Context Protocol docs (version 2026-07-28)}, \hyperref[sup-codeexec]{Anthropic Engineering: Code execution with MCP}, \hyperref[sup-tinyagent]{Berkeley BAIR: TinyAgent, Function Calling at the Edge}, \hyperref[sup-agentsec]{MIT 6.566 guest lecture: AI Agent Security (Anish Athalye, April 2026)}, \hyperref[sup-anthauto]{Anthropic: Building effective agent automations (Martin and Avilla, October 2026)}.
\chapter{Module 5: Memory, Retrieval and Context Engineering}\hypertarget{mod-MEM}{}
{\small Tag: \tg{MEM}}
\section{Learning objectives}
\begin{itemize}
\item Manage a finite context window: selection, compression/compaction, offloading and persistence.
\item Build and evaluate a RAG pipeline (chunking, embeddings, top-K, reranking, unanswerable questions).
\item Implement long-term memory and skills with progressive disclosure.
\end{itemize}
\section{Core concepts}
Apart from what the model learned in training, everything an agent can use at step $t$ is in its context: the information supplied for that inference. \textbf{Context engineering} decides what goes in.
\begin{itemize}
\item \textbf{Limits.} Models degrade on long inputs (``lost in the middle'', context rot), and every token costs money and latency. KV-cache reuse makes a stable prompt prefix cheap. CMU's long-context lecture uses real provider pricing and harness source code.
\item \textbf{Compaction.} Summarize or fold old turns (model-generated compaction in CMU A1; AgentFold and context-folding in UVA and UW), and offload large observations to files.
\item \textbf{Retrieval (RAG).} Chunk, embed, index, retrieve top-K, optionally rerank, then generate with citations. Treat chunk size, overlap, K, query rewriting and score thresholds as experimental variables (UVA Topic~5 runs 11 controlled experiments). Agentic RAG adds self-grading, correction and routing (Self-RAG, CRAG, Adaptive-RAG).
\item \textbf{Long-term memory and skills.} OS-style memory tiers (MemGPT), reflection-based memories (Generative Agents, Reflexion), workflow and skill induction (AWM, ASI). \emph{Skills} are SKILL.md bundles loaded on demand, the same progressive-disclosure idea Stanford '26 week~3 and CMU A1 teach.
\end{itemize}
\section{Curated learning sequence}
\begin{enumerate}
\item \hyperlink{lec-CMU-3}{CMU~3 \emph{Capabilities 2: Context Management for Long-Context Agents}} --- Context management for long-context agents: limits, compaction, KV-cache economics.
\item \hyperlink{lec-CMU-4}{CMU~4 \emph{Capabilities 3: Skills and Memory}} --- Skills and memory: MemGPT, AWM, ASI, ReasoningBank, SkillsBench.
\item \hyperlink{lec-NEU-7}{Northeastern~7 \emph{Context Management}}--\hyperlink{lec-NEU-9}{Northeastern~9} --- Context engineering, then RAG from scratch to Pinecone, then a self-grading retrieval loop.
\item \hyperlink{lec-UVA-9}{UVA~9 \emph{Topic 5: Retrieval Augmented Generation}} --- The RAG experiment lab (chunking, top-K, unanswerable questions, small versus large model).
\item \hyperlink{lec-MANN-6}{Mannheim~6 \emph{Context Engineering}} --- Context-engineering deck and the deep-agent exercise.
\item \hyperlink{lec-STAN26-5}{Stanford '26~5 \emph{All about agent skills (incl. web skills)}} --- Agent skills, including web skills (then author one in the Week~3 assignment).
\item \hyperlink{lec-UW-21}{UW-Madison~21 \emph{Memory I}} and \hyperlink{lec-UW-22}{UW-Madison~22 \emph{Memory \& other topics II}} --- Research reading: MemGPT, CoALA, A-Mem, context folding.
\item Supplement: \hyperref[sup-notion]{Notion Data Scout} --- Curated, owned context for a production data agent.
\end{enumerate}
\section{Lecture map (26 lectures)}
{\footnotesize\begin{longtable}{P{1.75cm} P{0.4cm} P{5.75cm} P{4.2cm} P{2.65cm}}
\toprule Course & \# & Title (links to catalogue row) & Materials & Tags\\ \midrule\endhead
\hyperlink{course-CMU}{CMU} & 3 & $\star$\,\hyperlink{lec-CMU-3}{Capabilities 2: Context Management for Long-Context Agents} & \href{https://www.cmu-agents.com/slides/lecture-03-long-context.pdf}{[slides]} \href{https://www.cmu-agents.com/\#/schedule}{[video]} & \hyperlink{mod-MEM}{\tg{MEM}}\\
\hyperlink{course-CMU}{CMU} & 4 & $\star$\,\hyperlink{lec-CMU-4}{Capabilities 3: Skills and Memory} & \href{https://www.cmu-agents.com/slides/lecture-04-memory-and-skills.pdf}{[slides]} \href{https://www.cmu-agents.com/\#/schedule}{[video]} & \hyperlink{mod-MEM}{\tg{MEM}}\\
\hyperlink{course-CMU}{CMU} & 10 & \hyperlink{lec-CMU-10}{Domains 3: Deep Research Agents} & \href{https://www.cmu-agents.com/slides/lecture-10-deep-research-agents.pdf}{[slides]} & \hyperlink{mod-APP}{\tg{APP}} \hyperlink{mod-MEM}{\tg{MEM}}\\
\hyperlink{course-STAN26}{Stanford '26} & 3 & \hyperlink{lec-STAN26-3}{Advanced prompting + agentic dev frameworks (RePPIT, spec-driven development)} & \href{https://themodernsoftware.dev/}{[course site]} & \hyperlink{mod-FND}{\tg{FND}} \hyperlink{mod-MEM}{\tg{MEM}}\\
\hyperlink{course-STAN26}{Stanford '26} & 5 & $\star$\,\hyperlink{lec-STAN26-5}{All about agent skills (incl. web skills)} & \href{https://themodernsoftware.dev/}{[course site]} & \hyperlink{mod-MEM}{\tg{MEM}} \hyperlink{mod-TOOL}{\tg{TOOL}}\\
\hyperlink{course-STAN25}{Stanford '25} & 5 & \hyperlink{lec-STAN25-5}{From first prompt to optimal IDE setup} & \href{https://docs.google.com/presentation/d/11pQNCde_mmRnImBat0Zymnp8TCS_cT_1up7zbcj6Sjg/edit?usp=sharing}{[slides]} \href{https://drive.google.com/file/d/1MZ0Qx68Vzw4x5x_XcV8XiPLp7fFDe1LJ/view?usp=drive_link}{[drive.google.com]} & \hyperlink{mod-MEM}{\tg{MEM}} \hyperlink{mod-CODE}{\tg{CODE}}\\
\hyperlink{course-NEU}{Northeastern} & 4 & \hyperlink{lec-NEU-4}{Simple Agent Architecture} & \href{https://agenticllms.github.io/week3a.html}{[agenticllms.github.io]} \href{https://github.com/agenticllms/lecture-files/tree/main/W03a}{[GitHub]} & \hyperlink{mod-ARCH}{\tg{ARCH}} \hyperlink{mod-TOOL}{\tg{TOOL}} \hyperlink{mod-MEM}{\tg{MEM}}\\
\hyperlink{course-NEU}{Northeastern} & 6 & \hyperlink{lec-NEU-6}{Short-term \& Long-term Memory} & \href{https://agenticllms.github.io/week4a.html}{[agenticllms.github.io]} \href{https://agenticllms.github.io/week4a-slides.html}{[agenticllms.github.io]} \href{https://github.com/agenticllms/lecture-files/tree/main/W04a}{[GitHub]} & \hyperlink{mod-MEM}{\tg{MEM}} \hyperlink{mod-ARCH}{\tg{ARCH}} \hyperlink{mod-DEPLOY}{\tg{DEPLOY}}\\
\hyperlink{course-NEU}{Northeastern} & 7 & $\star$\,\hyperlink{lec-NEU-7}{Context Management} & \href{https://agenticllms.github.io/week4b.html}{[agenticllms.github.io]} \href{https://github.com/agenticllms/lecture-files/tree/main/W04b}{[GitHub]} & \hyperlink{mod-MEM}{\tg{MEM}} \hyperlink{mod-ARCH}{\tg{ARCH}}\\
\hyperlink{course-NEU}{Northeastern} & 8 & $\star$\,\hyperlink{lec-NEU-8}{Retrieval-Augmented Generation} & \href{https://agenticllms.github.io/week5a.html}{[agenticllms.github.io]} \href{https://github.com/agenticllms/lecture-files/tree/main/W05a}{[GitHub]} & \hyperlink{mod-MEM}{\tg{MEM}}\\
\hyperlink{course-NEU}{Northeastern} & 9 & $\star$\,\hyperlink{lec-NEU-9}{Reflection \& Self-critique} & \href{https://agenticllms.github.io/week5b.html}{[agenticllms.github.io]} \href{https://github.com/agenticllms/lecture-files/tree/main/W05b}{[GitHub]} & \hyperlink{mod-RSN}{\tg{RSN}} \hyperlink{mod-MEM}{\tg{MEM}}\\
\hyperlink{course-NEU}{Northeastern} & 11 & \hyperlink{lec-NEU-11}{Agentic RAG} & -- & \hyperlink{mod-MEM}{\tg{MEM}} \hyperlink{mod-RSN}{\tg{RSN}}\\
\hyperlink{course-UB}{Buffalo} & 4 & \hyperlink{lec-UB-4}{Retrieval-augmented generation and context construction} & -- & \hyperlink{mod-MEM}{\tg{MEM}}\\
\hyperlink{course-UB}{Buffalo} & 5 & \hyperlink{lec-UB-5}{Embeddings, vector search, reranking} & -- & \hyperlink{mod-MEM}{\tg{MEM}}\\
\hyperlink{course-UB}{Buffalo} & 8 & \hyperlink{lec-UB-8}{Human-agent interaction and memory} & -- & \hyperlink{mod-MEM}{\tg{MEM}} \hyperlink{mod-SAFE}{\tg{SAFE}}\\
\hyperlink{course-UVA}{UVA} & 9 & $\star$\,\hyperlink{lec-UVA-9}{Topic 5: Retrieval Augmented Generation} & \href{https://www.cs.virginia.edu/\%7Ermw7my/Courses/AgenticAISpring2026/Topic5RAG/rag.html}{[cs.virginia.edu]} \href{https://github.com/HenryKautz/pdf_ocr}{[GitHub]} & \hyperlink{mod-MEM}{\tg{MEM}} \hyperlink{mod-EVAL}{\tg{EVAL}}\\
\hyperlink{course-UVA}{UVA} & 15 & \hyperlink{lec-UVA-15}{Topic 9: Context Management} & \href{https://www.cs.virginia.edu/\%7Ermw7my/Courses/AgenticAISpring2026/Topic9Context/context.html}{[cs.virginia.edu]} & \hyperlink{mod-MEM}{\tg{MEM}}\\
\hyperlink{course-UW}{UW-Madison} & 8 & \hyperlink{lec-UW-8}{Deep research agents II} & -- & \hyperlink{mod-APP}{\tg{APP}} \hyperlink{mod-MEM}{\tg{MEM}} \hyperlink{mod-EVAL}{\tg{EVAL}}\\
\hyperlink{course-UW}{UW-Madison} & 21 & $\star$\,\hyperlink{lec-UW-21}{Memory I} & -- & \hyperlink{mod-MEM}{\tg{MEM}} \hyperlink{mod-ARCH}{\tg{ARCH}}\\
\hyperlink{course-UW}{UW-Madison} & 22 & $\star$\,\hyperlink{lec-UW-22}{Memory \& other topics II} & -- & \hyperlink{mod-MEM}{\tg{MEM}} \hyperlink{mod-MULTI}{\tg{MULTI}} \hyperlink{mod-EVAL}{\tg{EVAL}}\\
\hyperlink{course-CU}{CU Boulder} & 21 & \hyperlink{lec-CU-21}{MemSkill \& More with Less (students)} & -- & \hyperlink{mod-MEM}{\tg{MEM}} \hyperlink{mod-DEPLOY}{\tg{DEPLOY}}\\
\hyperlink{course-MANN}{Mannheim} & 5 & \hyperlink{lec-MANN-5}{Retrieval Augmented Generation} & \href{https://www.uni-mannheim.de/media/Einrichtungen/dws/Files_Teaching/Large_Language_Models_and_Agents/FSS2026/IE685_LA_05_RetrievalAugmentedGeneration.pdf}{[PDF]} \href{https://www.uni-mannheim.de/media/Einrichtungen/dws/Files_Teaching/Large_Language_Models_and_Agents/FSS2026/IE685_LA_EX_Retrieval_Augmented_Generation.zip}{[zip]} & \hyperlink{mod-MEM}{\tg{MEM}}\\
\hyperlink{course-MANN}{Mannheim} & 6 & $\star$\,\hyperlink{lec-MANN-6}{Context Engineering} & \href{https://www.uni-mannheim.de/media/Einrichtungen/dws/Files_Teaching/Large_Language_Models_and_Agents/FSS2026/IE685_LA_06_ContextEngineering.pdf}{[PDF]} \href{https://www.uni-mannheim.de/media/Einrichtungen/dws/Files_Teaching/Large_Language_Models_and_Agents/FSS2026/deep_agent.zip}{[zip]} & \hyperlink{mod-MEM}{\tg{MEM}} \hyperlink{mod-ARCH}{\tg{ARCH}}\\
\hyperlink{course-UCHI}{UChicago} & 3 & \hyperlink{lec-UCHI-3}{LLMs Prompted for Multi-Agent Simulation} & \href{https://github.com/KnowledgeLab/AI-Agents-for-Social-Science-and-Society-2026/blob/main/Tutorials-Homework_Notebooks/Week_3.ipynb}{[GitHub]} & \hyperlink{mod-MULTI}{\tg{MULTI}} \hyperlink{mod-MEM}{\tg{MEM}} \hyperlink{mod-APP}{\tg{APP}}\\
\hyperlink{course-BRKADV}{Berkeley Adv.} & 3 & \hyperlink{lec-BRKADV-3}{On Reasoning, Memory, and Planning of Language Agents} & \href{https://rdi.berkeley.edu/adv-llm-agents/slides/language_agents_YuSu_Berkeley.pdf}{[slides]} \href{https://www.youtube.com/live/zvI4UN2_i-w}{[video]} & \hyperlink{mod-RSN}{\tg{RSN}} \hyperlink{mod-MEM}{\tg{MEM}} \hyperlink{mod-ARCH}{\tg{ARCH}}\\
\hyperlink{course-S329T}{Stanford 329T} & 2 & \hyperlink{lec-S329T-2}{Models, Prompting and RAG} & \href{https://web.stanford.edu/class/cs329t/slides/Lecture\%202\%20-\%20CS\%20329T\%20Fall\%202025.pdf}{[slides]} \href{https://web.stanford.edu/class/cs329t/homeworks/CS329T_HW1_Fall2025.pdf}{[HW1]} & \hyperlink{mod-FND}{\tg{FND}} \hyperlink{mod-MEM}{\tg{MEM}}\\
\bottomrule\end{longtable}}
\section{Key readings}
\begin{itemize}
\item Lewis et al., \href{https://arxiv.org/abs/2005.11401}{\emph{Retrieval-Augmented Generation}} (2020)
\item Liu et al., \href{https://arxiv.org/abs/2307.03172}{\emph{Lost in the Middle}} (2023)
\item Packer et al., \href{https://arxiv.org/abs/2310.08560}{\emph{MemGPT}} (2023)
\item Park et al., \href{https://arxiv.org/abs/2304.03442}{\emph{Generative Agents}} (2023)
\item Anthropic, \href{https://www.anthropic.com/engineering/effective-context-engineering-for-ai-agents}{\emph{Effective context engineering for AI agents}}
\item Asai et al., \href{https://arxiv.org/abs/2310.11511}{\emph{Self-RAG}} (2023); Yan et al., \href{https://arxiv.org/abs/2401.15884}{\emph{Corrective RAG}} (2024)
\end{itemize}
\section{Hands-on}
\begin{itemize}
\item \hyperlink{asg-UVA-5}{\textbf{UVA: Topic 5: RAG experiments}} --- 11 controlled RAG experiments (RAG vs no-RAG, top-K, chunking, prompts, unanswerable questions). \textit{Starter code:} \href{https://www.cs.virginia.edu/\%7Ermw7my/Courses/AgenticAISpring2026/Topic5RAG/rag.html}{[cs.virginia.edu]}
\item \hyperlink{asg-UB-2}{\textbf{Buffalo: Lab 2: Retrieval and Memory}} --- Small RAG pipeline over a course-selected corpus. Compare keyword/dense/hybrid retrieval, compute cosine similarity, do error analysis. \textit{(spec only)}
\item \hyperlink{asg-STAN26-3}{\textbf{Stanford '26: Week 3 --- Author a Skill}} --- A SKILL.md (Agent Skills standard) for a repeatable workflow on a real OSS repo (htmx, dayjs, yt-dlp, etc.), with trigger tests including a near-miss and an end-to-end run. \textit{Starter code:} \href{https://github.com/mihail911/modern-software-dev-assignments/blob/master/week3/assignment.md}{[GitHub]}
\item \hyperlink{asg-NEU-2}{\textbf{Northeastern: Husky Tech Customer Support Agent}} --- A support agent with 4 tools and ticket filing. It is built as a notebook, then a project skeleton, then deployed to HF Spaces, then rewritten in LangGraph, then given RAG over husky\_\allowbreak{}policy.md, then a reflection/self-grading loop. \textit{Starter code:} \href{https://github.com/agenticllms/lecture-files/tree/main/W03b}{[GitHub]}
\item \hyperlink{asg-MANN-4}{\textbf{Mannheim: Exercise: RAG}} --- RAG pipeline notebook. \textit{Starter code:} \href{https://www.uni-mannheim.de/media/Einrichtungen/dws/Files_Teaching/Large_Language_Models_and_Agents/FSS2026/IE685_LA_EX_Retrieval_Augmented_Generation.zip}{[zip]}
\end{itemize}
\textbf{Related supplements:} \hyperref[sup-codeexec]{Anthropic Engineering: Code execution with MCP}, \hyperref[sup-notion]{Notion: How we built a personal data scientist for every Notion employee (David Thomas, October 2026)}.
\chapter{Module 6: Reasoning, Planning, Reflection and Search}\hypertarget{mod-RSN}{}
{\small Tag: \tg{RSN}}
\section{Learning objectives}
\begin{itemize}
\item Choose between ReAct-style interleaving, plan-then-execute and decomposition strategies.
\item Add reflection/self-critique loops and measure whether they help.
\item Use search (best-of-N, critics/verifiers, tree search) to trade compute for reliability.
\end{itemize}
\section{Core concepts}
\begin{itemize}
\item \textbf{Thinking in tokens.} CoT and self-consistency let the model compute before answering. UVA's CoT lecture looks at why this works and at latent (hidden) reasoning.
\item \textbf{Planning.} Plan-and-execute, least-to-most and decomposed prompting split long tasks. Code-as-policies and DAG plans (LLMCompiler, used in TinyAgent) make plans executable and parallel. CMU's planning lecture asks \emph{when} to plan (``Thinking vs.\ Doing'', ``Calibrate-Then-Act'').
\item \textbf{Reflection.} Reflexion and Self-Refine critique and retry. Northeastern's self-grading RAG loop is a concrete build. Reflection only helps when the feedback signal carries real information (tests, tools, retrieval scores), so measure it.
\item \textbf{Search.} Sample many candidates, then pick with a verifier or critic (CMU Search~1), or explore a tree of states (ToT, LATS, CMU Search~2). AlphaCode, AlphaEvolve and ShinkaEvolve (UW) scale this to evolutionary program search. JHU starts from the classic version: enumerative program synthesis.
\end{itemize}
\section{Curated learning sequence}
\begin{enumerate}
\item \hyperlink{lec-CMU-5}{CMU~5 \emph{Capabilities 4: Planning, Task Decomposition, Multi-Agent Coordination}} --- Planning, task decomposition and multi-agent coordination (slides + video).
\item \hyperlink{lec-UW-3}{UW-Madison~3 \emph{Reasoning and acting I}} and \hyperlink{lec-UW-4}{UW-Madison~4 \emph{Reasoning and acting II}} --- CoT/ReAct, then Toolformer/ToT, with supplementary papers.
\item \hyperlink{lec-NEU-9}{Northeastern~9 \emph{Reflection \& Self-critique}} and \hyperlink{lec-NEU-10}{Northeastern~10 \emph{Planning \& Task Decomposition}} --- Reflection and self-critique, then planning with skill libraries (Voyager) and ToT.
\item \hyperlink{lec-UVA-16}{UVA~16 \emph{Topic 10: Understanding Chain of Thought}} --- Understanding chain of thought, including latent CoT.
\item \hyperlink{lec-CMU-20}{CMU~20 \emph{Search 1: Reranking and Critic Models}} and \hyperlink{lec-CMU-21}{CMU~21 \emph{Search 2: Tree Search}} --- Reranking and critic models, then tree search.
\item \hyperlink{lec-UW-17}{UW-Madison~17 \emph{Searching the solution space I}} and \hyperlink{lec-UW-18}{UW-Madison~18 \emph{Searching the solution space II}} --- Searching the solution space: AlphaCode, AlphaGeometry, evolutionary search.
\item \hyperlink{lec-JHU-2}{JHU~2 \emph{Syntax, Semantics, and Bottom-up Inductive Synthesis}}--\hyperlink{lec-JHU-4}{JHU~4} --- Program synthesis as search (the foundation for A1).
\end{enumerate}
\section{Lecture map (48 lectures)}
{\footnotesize\begin{longtable}{P{1.75cm} P{0.4cm} P{5.75cm} P{4.2cm} P{2.65cm}}
\toprule Course & \# & Title (links to catalogue row) & Materials & Tags\\ \midrule\endhead
\hyperlink{course-CMU}{CMU} & 5 & $\star$\,\hyperlink{lec-CMU-5}{Capabilities 4: Planning, Task Decomposition, Multi-Agent Coordination} & \href{https://www.cmu-agents.com/slides/lecture-05-planning.pdf}{[slides]} \href{https://www.cmu-agents.com/\#/schedule}{[video]} & \hyperlink{mod-RSN}{\tg{RSN}} \hyperlink{mod-MULTI}{\tg{MULTI}}\\
\hyperlink{course-CMU}{CMU} & 20 & $\star$\,\hyperlink{lec-CMU-20}{Search 1: Reranking and Critic Models} & -- & \hyperlink{mod-RSN}{\tg{RSN}} \hyperlink{mod-EVAL}{\tg{EVAL}}\\
\hyperlink{course-CMU}{CMU} & 21 & $\star$\,\hyperlink{lec-CMU-21}{Search 2: Tree Search} & -- & \hyperlink{mod-RSN}{\tg{RSN}}\\
\hyperlink{course-STAN25}{Stanford '25} & 2 & \hyperlink{lec-STAN25-2}{Power prompting for LLMs} & \href{https://docs.google.com/presentation/d/1MIhw8p6TLGdbQ9TcxhXSs5BaPf5d_h77QY70RHNfeGs/edit?usp=drive_link}{[slides]} & \hyperlink{mod-FND}{\tg{FND}} \hyperlink{mod-RSN}{\tg{RSN}}\\
\hyperlink{course-NEU}{Northeastern} & 5 & \hyperlink{lec-NEU-5}{Tool-Calling} & \href{https://agenticllms.github.io/week3b.html}{[agenticllms.github.io]} \href{https://agenticllms.github.io/week3b-slides.html}{[agenticllms.github.io]} \href{https://github.com/agenticllms/lecture-files/tree/main/W03b}{[GitHub]} & \hyperlink{mod-TOOL}{\tg{TOOL}} \hyperlink{mod-RSN}{\tg{RSN}} \hyperlink{mod-DEPLOY}{\tg{DEPLOY}}\\
\hyperlink{course-NEU}{Northeastern} & 9 & $\star$\,\hyperlink{lec-NEU-9}{Reflection \& Self-critique} & \href{https://agenticllms.github.io/week5b.html}{[agenticllms.github.io]} \href{https://github.com/agenticllms/lecture-files/tree/main/W05b}{[GitHub]} & \hyperlink{mod-RSN}{\tg{RSN}} \hyperlink{mod-MEM}{\tg{MEM}}\\
\hyperlink{course-NEU}{Northeastern} & 10 & $\star$\,\hyperlink{lec-NEU-10}{Planning \& Task Decomposition} & -- & \hyperlink{mod-RSN}{\tg{RSN}} \hyperlink{mod-ARCH}{\tg{ARCH}}\\
\hyperlink{course-NEU}{Northeastern} & 11 & \hyperlink{lec-NEU-11}{Agentic RAG} & -- & \hyperlink{mod-MEM}{\tg{MEM}} \hyperlink{mod-RSN}{\tg{RSN}}\\
\hyperlink{course-UB}{Buffalo} & 3 & \hyperlink{lec-UB-3}{Tool/function calling and ReAct} & -- & \hyperlink{mod-TOOL}{\tg{TOOL}} \hyperlink{mod-RSN}{\tg{RSN}}\\
\hyperlink{course-UB}{Buffalo} & 7 & \hyperlink{lec-UB-7}{Planning and decomposition} & -- & \hyperlink{mod-RSN}{\tg{RSN}}\\
\hyperlink{course-UVA}{UVA} & 7 & \hyperlink{lec-UVA-7}{Topic 4: Exploring Tools} & \href{https://www.cs.virginia.edu/\%7Ermw7my/Courses/AgenticAISpring2026/Topic4Exploring/exploring.html}{[cs.virginia.edu]} \href{https://github.com/HenryKautz/MoltReader}{[GitHub]} & \hyperlink{mod-TOOL}{\tg{TOOL}} \hyperlink{mod-RSN}{\tg{RSN}} \hyperlink{mod-ARCH}{\tg{ARCH}}\\
\hyperlink{course-UVA}{UVA} & 16 & $\star$\,\hyperlink{lec-UVA-16}{Topic 10: Understanding Chain of Thought} & \href{https://www.cs.virginia.edu/\%7Ermw7my/Courses/AgenticAISpring2026/Topic10COT/cot.html}{[cs.virginia.edu]} & \hyperlink{mod-RSN}{\tg{RSN}}\\
\hyperlink{course-UVA}{UVA} & 20 & \hyperlink{lec-UVA-20}{Topic 14: Agentic AI for Mathematics and Verification (guest: Prof. Kevin Sullivan)} & \href{https://www.cs.virginia.edu/\%7Ermw7my/Courses/AgenticAISpring2026/Topic14Math/math.html}{[cs.virginia.edu]} & \hyperlink{mod-RSN}{\tg{RSN}} \hyperlink{mod-APP}{\tg{APP}}\\
\hyperlink{course-UVA}{UVA} & 21 & \hyperlink{lec-UVA-21}{Topic 15: Visual Reasoning} & \href{https://www.cs.virginia.edu/\%7Ermw7my/Courses/AgenticAISpring2026/Topic15Visual/visual.html}{[cs.virginia.edu]} \href{https://www.youtube.com/watch?v=Gs9TPZmzo-s}{[video]} & \hyperlink{mod-RSN}{\tg{RSN}} \hyperlink{mod-EVAL}{\tg{EVAL}}\\
\hyperlink{course-UW}{UW-Madison} & 3 & $\star$\,\hyperlink{lec-UW-3}{Reasoning and acting I} & -- & \hyperlink{mod-RSN}{\tg{RSN}} \hyperlink{mod-FND}{\tg{FND}}\\
\hyperlink{course-UW}{UW-Madison} & 4 & $\star$\,\hyperlink{lec-UW-4}{Reasoning and acting II} & -- & \hyperlink{mod-TOOL}{\tg{TOOL}} \hyperlink{mod-RSN}{\tg{RSN}}\\
\hyperlink{course-UW}{UW-Madison} & 16 & \hyperlink{lec-UW-16}{Self-improving agents II} & -- & \hyperlink{mod-CODE}{\tg{CODE}} \hyperlink{mod-RSN}{\tg{RSN}}\\
\hyperlink{course-UW}{UW-Madison} & 17 & $\star$\,\hyperlink{lec-UW-17}{Searching the solution space I} & -- & \hyperlink{mod-RSN}{\tg{RSN}} \hyperlink{mod-CODE}{\tg{CODE}}\\
\hyperlink{course-UW}{UW-Madison} & 18 & $\star$\,\hyperlink{lec-UW-18}{Searching the solution space II} & -- & \hyperlink{mod-RSN}{\tg{RSN}}\\
\hyperlink{course-JHU}{JHU} & 2 & $\star$\,\hyperlink{lec-JHU-2}{Syntax, Semantics, and Bottom-up Inductive Synthesis} & \href{https://machine-programming.github.io/lectures/lecture-2.pdf}{[PDF]} & \hyperlink{mod-RSN}{\tg{RSN}}\\
\hyperlink{course-JHU}{JHU} & 3 & $\star$\,\hyperlink{lec-JHU-3}{Type Systems and Top-down Enumerative Synthesis} & \href{https://machine-programming.github.io/lectures/lecture-3.pdf}{[PDF]} & \hyperlink{mod-RSN}{\tg{RSN}}\\
\hyperlink{course-JHU}{JHU} & 4 & $\star$\,\hyperlink{lec-JHU-4}{Functional Specifications and Synthesis} & \href{https://machine-programming.github.io/lectures/lecture-4.pdf}{[PDF]} & \hyperlink{mod-RSN}{\tg{RSN}}\\
\hyperlink{course-JHU}{JHU} & 26 & \hyperlink{lec-JHU-26}{Special Topics on Neurosymbolic Programming (1)} & -- & \hyperlink{mod-RSN}{\tg{RSN}} \hyperlink{mod-APP}{\tg{APP}}\\
\hyperlink{course-JHU}{JHU} & 27 & \hyperlink{lec-JHU-27}{Special Topics on Neurosymbolic Programming (2)} & -- & \hyperlink{mod-RSN}{\tg{RSN}} \hyperlink{mod-APP}{\tg{APP}}\\
\hyperlink{course-CU}{CU Boulder} & 9 & \hyperlink{lec-CU-9}{ReAct \& Web API code gen (students)} & -- & \hyperlink{mod-RSN}{\tg{RSN}} \hyperlink{mod-TOOL}{\tg{TOOL}} \hyperlink{mod-EVAL}{\tg{EVAL}}\\
\hyperlink{course-CU}{CU Boulder} & 11 & \hyperlink{lec-CU-11}{Reflexion \& Toolformer (students)} & -- & \hyperlink{mod-RSN}{\tg{RSN}} \hyperlink{mod-TOOL}{\tg{TOOL}} \hyperlink{mod-TRAIN}{\tg{TRAIN}}\\
\hyperlink{course-CU}{CU Boulder} & 12 & \hyperlink{lec-CU-12}{Gorilla \& PAL (students)} & -- & \hyperlink{mod-TOOL}{\tg{TOOL}} \hyperlink{mod-RSN}{\tg{RSN}}\\
\hyperlink{course-CU}{CU Boulder} & 15 & \hyperlink{lec-CU-15}{Multiagent debate \& Agentless (students)} & -- & \hyperlink{mod-MULTI}{\tg{MULTI}} \hyperlink{mod-RSN}{\tg{RSN}} \hyperlink{mod-CODE}{\tg{CODE}}\\
\hyperlink{course-MANN}{Mannheim} & 3 & \hyperlink{lec-MANN-3}{Prompt Engineering} & \href{https://www.uni-mannheim.de/media/Einrichtungen/dws/Files_Teaching/Large_Language_Models_and_Agents/FSS2026/IE685_LA_03_PromptEngineering.pdf}{[PDF]} \href{https://www.uni-mannheim.de/media/Einrichtungen/dws/Files_Teaching/Large_Language_Models_and_Agents/FSS2026/IE686_LA_EX_Langchain_Introduction.ipynb.zip}{[notebook zip]} & \hyperlink{mod-FND}{\tg{FND}} \hyperlink{mod-RSN}{\tg{RSN}}\\
\hyperlink{course-UCHI}{UChicago} & 2 & \hyperlink{lec-UCHI-2}{Text Learning, Transformers, and Diffusion Models} & \href{https://github.com/KnowledgeLab/AI-Agents-for-Social-Science-and-Society-2026/blob/main/Tutorials-Homework_Notebooks/Week_2.ipynb}{[GitHub]} & \hyperlink{mod-FND}{\tg{FND}} \hyperlink{mod-RSN}{\tg{RSN}}\\
\hyperlink{course-UCHI}{UChicago} & 7 & \hyperlink{lec-UCHI-7}{Reinforcement Learning to Optimize AI Agents and Institutions} & \href{https://github.com/KnowledgeLab/AI-Agents-for-Social-Science-and-Society-2026/blob/main/Tutorials-Homework_Notebooks/Week_7_RL.ipynb}{[GitHub]} & \hyperlink{mod-TRAIN}{\tg{TRAIN}} \hyperlink{mod-RSN}{\tg{RSN}} \hyperlink{mod-MULTI}{\tg{MULTI}}\\
\hyperlink{course-UCHI}{UChicago} & 8 & \hyperlink{lec-UCHI-8}{Multi-Modal and Embodied Agents} & \href{https://github.com/KnowledgeLab/AI-Agents-for-Social-Science-and-Society-2026/blob/main/Tutorials-Homework_Notebooks/Week_8_Multimodal.ipynb}{[GitHub]} & \hyperlink{mod-WEB}{\tg{WEB}} \hyperlink{mod-RSN}{\tg{RSN}} \hyperlink{mod-APP}{\tg{APP}}\\
\hyperlink{course-S329A}{Stanford 329A} & 2 & \hyperlink{lec-S329A-2}{Test-time Compute Scaling} & \href{https://www.youtube.com/watch?v=-Ggc37xLj_Y}{[video]} & \hyperlink{mod-RSN}{\tg{RSN}} \hyperlink{mod-EVAL}{\tg{EVAL}}\\
\hyperlink{course-S329A}{Stanford 329A} & 3 & \hyperlink{lec-S329A-3}{Robust Verification} & \href{https://www.youtube.com/watch?v=p7TdPUcPoik}{[video]} & \hyperlink{mod-RSN}{\tg{RSN}} \hyperlink{mod-EVAL}{\tg{EVAL}}\\
\hyperlink{course-S329A}{Stanford 329A} & 5 & \hyperlink{lec-S329A-5}{Multi-step Reasoning/Planning} & \href{https://www.youtube.com/watch?v=Ml_fp9XkB8Y}{[video]} & \hyperlink{mod-RSN}{\tg{RSN}} \hyperlink{mod-TRAIN}{\tg{TRAIN}}\\
\hyperlink{course-S329A}{Stanford 329A} & 6 & \hyperlink{lec-S329A-6}{Train Time Scaling/Scaling RL} & \href{https://www.youtube.com/watch?v=yVnmHSAy3ck}{[video]} & \hyperlink{mod-TRAIN}{\tg{TRAIN}} \hyperlink{mod-RSN}{\tg{RSN}}\\
\hyperlink{course-S329A}{Stanford 329A} & 8 & \hyperlink{lec-S329A-8}{Self-Improvement with Search \& Deep Research Agents} & \href{https://www.youtube.com/watch?v=Uni9dqyuuDM}{[video]} & \hyperlink{mod-WEB}{\tg{WEB}} \hyperlink{mod-RSN}{\tg{RSN}} \hyperlink{mod-CODE}{\tg{CODE}}\\
\hyperlink{course-BRKAG}{Berkeley '25} & 7 & \hyperlink{lec-BRKAG-7}{Multi-Agent AI} & \href{https://rdi.berkeley.edu/agentic-ai/slides/multiagent_berkeley_2025.pptx.pdf}{[slides]} \href{https://www.youtube.com/watch?v=SrLcGdVOb9w}{[video]} & \hyperlink{mod-MULTI}{\tg{MULTI}} \hyperlink{mod-RSN}{\tg{RSN}}\\
\hyperlink{course-BRKADV}{Berkeley Adv.} & 1 & \hyperlink{lec-BRKADV-1}{Inference-Time Techniques for LLM Reasoning} & \href{https://rdi.berkeley.edu/adv-llm-agents/slides/inference_time_techniques_lecture_sp25.pdf}{[slides]} \href{https://rdi.berkeley.edu/adv-llm-agents/slides/llm-agents-berkeley-intro-sp25.pdf}{[intro]} \href{https://www.youtube.com/live/g0Dwtf3BH-0}{[video]} & \hyperlink{mod-RSN}{\tg{RSN}} \tg{ORG}\\
\hyperlink{course-BRKADV}{Berkeley Adv.} & 2 & \hyperlink{lec-BRKADV-2}{Learning to reason with LLMs} & \href{https://rdi.berkeley.edu/adv-llm-agents/slides/Jason-Weston-Reasoning-Alignment-Berkeley-Talk.pdf}{[slides]} \href{https://www.youtube.com/live/_MNlLhU33H0}{[video]} & \hyperlink{mod-TRAIN}{\tg{TRAIN}} \hyperlink{mod-RSN}{\tg{RSN}}\\
\hyperlink{course-BRKADV}{Berkeley Adv.} & 3 & \hyperlink{lec-BRKADV-3}{On Reasoning, Memory, and Planning of Language Agents} & \href{https://rdi.berkeley.edu/adv-llm-agents/slides/language_agents_YuSu_Berkeley.pdf}{[slides]} \href{https://www.youtube.com/live/zvI4UN2_i-w}{[video]} & \hyperlink{mod-RSN}{\tg{RSN}} \hyperlink{mod-MEM}{\tg{MEM}} \hyperlink{mod-ARCH}{\tg{ARCH}}\\
\hyperlink{course-BRKADV}{Berkeley Adv.} & 4 & \hyperlink{lec-BRKADV-4}{Open Training Recipes for Reasoning in Language Models} & \href{https://rdi.berkeley.edu/adv-llm-agents/slides/OLMo-Tulu-Reasoning-Hanna.pdf}{[slides]} \href{https://www.youtube.com/live/cMiu3A7YBks}{[video]} & \hyperlink{mod-TRAIN}{\tg{TRAIN}} \hyperlink{mod-RSN}{\tg{RSN}}\\
\hyperlink{course-BRKADV}{Berkeley Adv.} & 6 & \hyperlink{lec-BRKADV-6}{Multimodal Autonomous AI Agents} & \href{https://rdi.berkeley.edu/adv-llm-agents/slides/ruslan-multimodal.pdf}{[slides]} \href{https://www.youtube.com/live/RPINOYM12RU}{[video]} & \hyperlink{mod-WEB}{\tg{WEB}} \hyperlink{mod-EVAL}{\tg{EVAL}} \hyperlink{mod-RSN}{\tg{RSN}}\\
\hyperlink{course-BRKADV}{Berkeley Adv.} & 8 & \hyperlink{lec-BRKADV-8}{AlphaProof: when reinforcement learning meets formal mathematics} & \href{https://rdi.berkeley.edu/adv-llm-agents/slides/alphaproof.pdf}{[slides]} \href{https://www.youtube.com/live/3gaEMscOMAU}{[video]} & \hyperlink{mod-RSN}{\tg{RSN}} \hyperlink{mod-TRAIN}{\tg{TRAIN}} \hyperlink{mod-APP}{\tg{APP}}\\
\hyperlink{course-BRKADV}{Berkeley Adv.} & 9 & \hyperlink{lec-BRKADV-9}{Language models for autoformalization and theorem proving} & \href{https://rdi.berkeley.edu/adv-llm-agents/slides/mathverification.pdf}{[slides]} \href{https://www.youtube.com/live/cLhWEyMQ4mQ}{[video]} & \hyperlink{mod-RSN}{\tg{RSN}} \hyperlink{mod-APP}{\tg{APP}} \hyperlink{mod-TOOL}{\tg{TOOL}}\\
\hyperlink{course-BRKADV}{Berkeley Adv.} & 10 & \hyperlink{lec-BRKADV-10}{Bridging Informal and Formal Mathematical Reasoning (Berkeley page: "Advanced topics in theorem proving")} & \href{https://rdi.berkeley.edu/adv-llm-agents/slides/welleck2025_berkeley_bridging.pdf}{[slides]} \href{https://www.youtube.com/live/Gy5Nm17l9oo}{[video]} & \hyperlink{mod-RSN}{\tg{RSN}} \hyperlink{mod-APP}{\tg{APP}}\\
\hyperlink{course-BRKADV}{Berkeley Adv.} & 11 & \hyperlink{lec-BRKADV-11}{Abstraction and Discovery with LLM Agents} & \href{https://rdi.berkeley.edu/adv-llm-agents/slides/swarat.pdf}{[slides]} \href{https://www.youtube.com/live/IHc0TEMrEdY}{[video]} & \hyperlink{mod-RSN}{\tg{RSN}} \hyperlink{mod-ARCH}{\tg{ARCH}} \hyperlink{mod-APP}{\tg{APP}}\\
\hyperlink{course-S329T}{Stanford 329T} & 11 & \hyperlink{lec-S329T-11}{GEPA (guest: Lakshya Agrawal)} & \href{https://web.stanford.edu/class/cs329t/slides/Lecture\%209\%20-\%2020251028\%20GEPA_Talk.pdf}{[slides]} & \hyperlink{mod-TRAIN}{\tg{TRAIN}} \hyperlink{mod-RSN}{\tg{RSN}}\\
\bottomrule\end{longtable}}
\section{Key readings}
\begin{itemize}
\item Wang et al., \href{https://arxiv.org/abs/2203.11171}{\emph{Self-Consistency}} (2022)
\item Yao et al., \href{https://arxiv.org/abs/2305.10601}{\emph{Tree of Thoughts}} (2023)
\item Shinn et al., \href{https://arxiv.org/abs/2303.11366}{\emph{Reflexion}} (2023); Madaan et al., \href{https://arxiv.org/abs/2303.17651}{\emph{Self-Refine}} (2023)
\item Zhou et al., \href{https://arxiv.org/abs/2205.10625}{\emph{Least-to-Most Prompting}} (2022)
\item Zhou et al., \href{https://arxiv.org/abs/2310.04406}{\emph{LATS}} (2023)
\item Kim et al., \href{https://arxiv.org/abs/2312.04511}{\emph{LLMCompiler}} (2023)
\end{itemize}
\section{Hands-on}
\begin{itemize}
\item \hyperlink{asg-JHU-1}{\textbf{JHU: A1: Inductive Synthesis}} --- A bottom-up enumerative synthesizer for a geometric-shape DSL (grow / eliminate observational equivalents / search); your own string DSL and synthesizer; then prompts plus output parsing so Claude Sonnet 4.6 (via Bedrock) solves at least 80\% of string tasks. \textit{Starter code:} \href{https://github.com/machine-programming/assignment-1}{[GitHub]}
\item \hyperlink{asg-STAN25-1}{\textbf{Stanford '25: Week 1 --- Prompting Techniques}} --- Fill in prompts for K-shot, chain-of-thought, tool calling, self-consistency, RAG, and reflexion until the test scripts pass, using local Ollama models (mistral-nemo:12b, llama3.1:8b). \textit{Starter code:} \href{https://github.com/mihail911/modern-software-dev-assignments/tree/fall2025/week1}{[GitHub]}
\item \hyperlink{asg-NEU-2}{\textbf{Northeastern: Husky Tech Customer Support Agent}} --- A support agent with 4 tools and ticket filing. It is built as a notebook, then a project skeleton, then deployed to HF Spaces, then rewritten in LangGraph, then given RAG over husky\_\allowbreak{}policy.md, then a reflection/self-grading loop. \textit{Starter code:} \href{https://github.com/agenticllms/lecture-files/tree/main/W03b}{[GitHub]}
\end{itemize}
\chapter{Module 7: Multi-Agent Systems}\hypertarget{mod-MULTI}{}
{\small Tag: \tg{MULTI}}
\section{Learning objectives}
\begin{itemize}
\item Implement handoffs, routers, orchestrator--worker and debate patterns.
\item Diagnose multi-agent failure modes and decide when one agent is enough.
\end{itemize}
\section{Core concepts}
Multiple agents can parallelize work, isolate context (sub-agents with fresh windows) and specialize roles. They also multiply cost and add new failure modes: role confusion, lost information at handoffs, and endless loops.
\begin{itemize}
\item \textbf{Patterns:} router/handoff, orchestrator--workers (Anthropic's multi-agent research system), role-based SOPs (MetaGPT, CAMEL), conversation frameworks (AutoGen), debate, and planner/implementer/reviewer sub-agents in coding tools (Stanford '26 week~4).
\item \textbf{Evidence:} \emph{Why do multi-agent LLM systems fail?} (Northeastern reads it twice) and \emph{Towards a Science of Scaling Agent Systems} (UW) set this against OpenHands' ``don't sleep on single-agent systems''. Start with one agent and add more only when a measurement says so.
\item \textbf{Simulation:} populations of agents as social-science instruments (Generative Agents, Concordia, UChicago week~3).
\end{itemize}
\section{Curated learning sequence}
\begin{enumerate}
\item \hyperlink{lec-NEU-12}{Northeastern~12 \emph{Multi-Agent Orchestration I}}--\hyperlink{lec-NEU-14}{Northeastern~14} --- Three sessions: failure modes and handoffs, role-based MAS, the LangGraph runtime.
\item \hyperlink{lec-CMU-18}{CMU~18 \emph{Interaction 1: Multi-Agent Interaction}} --- Multi-agent interaction (guest lecture).
\item \hyperlink{lec-UW-19}{UW-Madison~19 \emph{Multi-agent systems I}} and \hyperlink{lec-UW-20}{UW-Madison~20 \emph{Multi-agent systems II}} --- Generative Agents/CAMEL, then AutoGen and the scaling science of agent systems.
\item \hyperlink{lec-CU-14}{CU Boulder~14 \emph{AutoGen \& MetaGPT (students)}} and \hyperlink{lec-CU-15}{CU Boulder~15 \emph{Multiagent debate \& Agentless (students)}} --- AutoGen/MetaGPT, then debate and Agentless.
\item \hyperlink{lec-UCHI-3}{UChicago~3 \emph{LLMs Prompted for Multi-Agent Simulation}} --- Prompted multi-agent social simulation (notebook with AutoGen and Concordia).
\end{enumerate}
\section{Lecture map (27 lectures)}
{\footnotesize\begin{longtable}{P{1.75cm} P{0.4cm} P{5.75cm} P{4.2cm} P{2.65cm}}
\toprule Course & \# & Title (links to catalogue row) & Materials & Tags\\ \midrule\endhead
\hyperlink{course-CMU}{CMU} & 5 & \hyperlink{lec-CMU-5}{Capabilities 4: Planning, Task Decomposition, Multi-Agent Coordination} & \href{https://www.cmu-agents.com/slides/lecture-05-planning.pdf}{[slides]} \href{https://www.cmu-agents.com/\#/schedule}{[video]} & \hyperlink{mod-RSN}{\tg{RSN}} \hyperlink{mod-MULTI}{\tg{MULTI}}\\
\hyperlink{course-CMU}{CMU} & 18 & $\star$\,\hyperlink{lec-CMU-18}{Interaction 1: Multi-Agent Interaction} & -- & \hyperlink{mod-MULTI}{\tg{MULTI}}\\
\hyperlink{course-STAN26}{Stanford '26} & 7 & \hyperlink{lec-STAN26-7}{CLAUDE.md, AGENTS.md, hooks; subagent patterns (planner/implementer/reviewer)} & \href{https://themodernsoftware.dev/}{[course site]} & \hyperlink{mod-CODE}{\tg{CODE}} \hyperlink{mod-MULTI}{\tg{MULTI}}\\
\hyperlink{course-STAN25}{Stanford '25} & 7 & \hyperlink{lec-STAN25-7}{How to be an agent manager} & \href{https://docs.google.com/presentation/d/19mgkwAnJDc7JuJy0zhhoY0ZC15DiNpxL8kchPDnRkRQ/edit?usp=sharing}{[slides]} & \hyperlink{mod-CODE}{\tg{CODE}} \hyperlink{mod-MULTI}{\tg{MULTI}}\\
\hyperlink{course-STAN25}{Stanford '25} & 17 & \hyperlink{lec-STAN25-17}{Incident response and DevOps} & \href{https://docs.google.com/presentation/d/1Mfe-auWAsg9URCujneKnHr0AbO8O-_U4QXBVOlO4qp0/edit?usp=sharing}{[slides]} & \hyperlink{mod-DEPLOY}{\tg{DEPLOY}} \hyperlink{mod-MULTI}{\tg{MULTI}}\\
\hyperlink{course-NEU}{Northeastern} & 12 & $\star$\,\hyperlink{lec-NEU-12}{Multi-Agent Orchestration I} & -- & \hyperlink{mod-MULTI}{\tg{MULTI}}\\
\hyperlink{course-NEU}{Northeastern} & 13 & $\star$\,\hyperlink{lec-NEU-13}{Multi-Agent Orchestration II} & -- & \hyperlink{mod-MULTI}{\tg{MULTI}}\\
\hyperlink{course-NEU}{Northeastern} & 14 & $\star$\,\hyperlink{lec-NEU-14}{Multi-Agent Orchestration III} & -- & \hyperlink{mod-MULTI}{\tg{MULTI}} \hyperlink{mod-ARCH}{\tg{ARCH}}\\
\hyperlink{course-UB}{Buffalo} & 11 & \hyperlink{lec-UB-11}{Multi-agent systems and coordination} & -- & \hyperlink{mod-MULTI}{\tg{MULTI}}\\
\hyperlink{course-UVA}{UVA} & 3 & \hyperlink{lec-UVA-3}{Topic 2: Agent Frameworks} & \href{https://www.cs.virginia.edu/\%7Ermw7my/Courses/AgenticAISpring2026/Topic2Frameworks/Agent\%20Control\%20Flow\%20Frameworks.html}{[cs.virginia.edu]} & \hyperlink{mod-ARCH}{\tg{ARCH}} \hyperlink{mod-MULTI}{\tg{MULTI}}\\
\hyperlink{course-UVA}{UVA} & 12 & \hyperlink{lec-UVA-12}{Topic 7: MCP and A2A} & \href{https://www.cs.virginia.edu/\%7Ermw7my/Courses/AgenticAISpring2026/Topic7MCP/mcp.html}{[cs.virginia.edu]} & \hyperlink{mod-TOOL}{\tg{TOOL}} \hyperlink{mod-MULTI}{\tg{MULTI}}\\
\hyperlink{course-UVA}{UVA} & 13 & \hyperlink{lec-UVA-13}{Topic 7 cont. (A2A) + Topic 8: LLM Fine-Tuning} & \href{https://www.cs.virginia.edu/\%7Ermw7my/Courses/AgenticAISpring2026/Topic8FineTuning/finetuning.html}{[cs.virginia.edu]} & \hyperlink{mod-TRAIN}{\tg{TRAIN}} \hyperlink{mod-MULTI}{\tg{MULTI}}\\
\hyperlink{course-UVA}{UVA} & 19 & \hyperlink{lec-UVA-19}{Topic 13: Agentic AI for Science} & \href{https://www.cs.virginia.edu/\%7Ermw7my/Courses/AgenticAISpring2026/Topic13Science/science.html}{[cs.virginia.edu]} \href{https://github.com/tonymenzo/heptapod}{[GitHub]} & \hyperlink{mod-APP}{\tg{APP}} \hyperlink{mod-MULTI}{\tg{MULTI}}\\
\hyperlink{course-UW}{UW-Madison} & 7 & \hyperlink{lec-UW-7}{Deep research agents I} & -- & \hyperlink{mod-APP}{\tg{APP}} \hyperlink{mod-MULTI}{\tg{MULTI}}\\
\hyperlink{course-UW}{UW-Madison} & 10 & \hyperlink{lec-UW-10}{Agents for science II} & -- & \hyperlink{mod-APP}{\tg{APP}} \hyperlink{mod-MULTI}{\tg{MULTI}}\\
\hyperlink{course-UW}{UW-Madison} & 19 & $\star$\,\hyperlink{lec-UW-19}{Multi-agent systems I} & -- & \hyperlink{mod-MULTI}{\tg{MULTI}}\\
\hyperlink{course-UW}{UW-Madison} & 20 & $\star$\,\hyperlink{lec-UW-20}{Multi-agent systems II} & -- & \hyperlink{mod-MULTI}{\tg{MULTI}} \hyperlink{mod-ARCH}{\tg{ARCH}}\\
\hyperlink{course-UW}{UW-Madison} & 22 & \hyperlink{lec-UW-22}{Memory \& other topics II} & -- & \hyperlink{mod-MEM}{\tg{MEM}} \hyperlink{mod-MULTI}{\tg{MULTI}} \hyperlink{mod-EVAL}{\tg{EVAL}}\\
\hyperlink{course-CU}{CU Boulder} & 4 & \hyperlink{lec-CU-4}{Refactoring agents (students)} & -- & \hyperlink{mod-CODE}{\tg{CODE}} \hyperlink{mod-MULTI}{\tg{MULTI}} \hyperlink{mod-EVAL}{\tg{EVAL}}\\
\hyperlink{course-CU}{CU Boulder} & 5 & \hyperlink{lec-CU-5}{MM-Assist \& MANTRA (Dig)} & \href{https://danny.cs.colorado.edu/courses/csci7000-005_Sp26/pdfs/W3_L2_MM-Mantra.pdf}{[PDF]} & \hyperlink{mod-CODE}{\tg{CODE}} \hyperlink{mod-MULTI}{\tg{MULTI}}\\
\hyperlink{course-CU}{CU Boulder} & 14 & $\star$\,\hyperlink{lec-CU-14}{AutoGen \& MetaGPT (students)} & -- & \hyperlink{mod-MULTI}{\tg{MULTI}} \hyperlink{mod-ARCH}{\tg{ARCH}}\\
\hyperlink{course-CU}{CU Boulder} & 15 & $\star$\,\hyperlink{lec-CU-15}{Multiagent debate \& Agentless (students)} & -- & \hyperlink{mod-MULTI}{\tg{MULTI}} \hyperlink{mod-RSN}{\tg{RSN}} \hyperlink{mod-CODE}{\tg{CODE}}\\
\hyperlink{course-UCHI}{UChicago} & 3 & $\star$\,\hyperlink{lec-UCHI-3}{LLMs Prompted for Multi-Agent Simulation} & \href{https://github.com/KnowledgeLab/AI-Agents-for-Social-Science-and-Society-2026/blob/main/Tutorials-Homework_Notebooks/Week_3.ipynb}{[GitHub]} & \hyperlink{mod-MULTI}{\tg{MULTI}} \hyperlink{mod-MEM}{\tg{MEM}} \hyperlink{mod-APP}{\tg{APP}}\\
\hyperlink{course-UCHI}{UChicago} & 7 & \hyperlink{lec-UCHI-7}{Reinforcement Learning to Optimize AI Agents and Institutions} & \href{https://github.com/KnowledgeLab/AI-Agents-for-Social-Science-and-Society-2026/blob/main/Tutorials-Homework_Notebooks/Week_7_RL.ipynb}{[GitHub]} & \hyperlink{mod-TRAIN}{\tg{TRAIN}} \hyperlink{mod-RSN}{\tg{RSN}} \hyperlink{mod-MULTI}{\tg{MULTI}}\\
\hyperlink{course-BRKAG}{Berkeley '25} & 7 & \hyperlink{lec-BRKAG-7}{Multi-Agent AI} & \href{https://rdi.berkeley.edu/agentic-ai/slides/multiagent_berkeley_2025.pptx.pdf}{[slides]} \href{https://www.youtube.com/watch?v=SrLcGdVOb9w}{[video]} & \hyperlink{mod-MULTI}{\tg{MULTI}} \hyperlink{mod-RSN}{\tg{RSN}}\\
\hyperlink{course-BRKAG}{Berkeley '25} & 9 & \hyperlink{lec-BRKAG-9}{AI Agents to Automate Scientific Discoveries} & \href{https://www.youtube.com/watch?v=yqPIsTTdUkc}{[video]} & \hyperlink{mod-APP}{\tg{APP}} \hyperlink{mod-MULTI}{\tg{MULTI}}\\
\hyperlink{course-BRKAG}{Berkeley '25} & 11 & \hyperlink{lec-BRKAG-11}{Multi-Agent Systems in the Era of LLMs} & \href{https://www.youtube.com/watch?v=ntjOxjZMaac\&t=2s}{[video]} & \hyperlink{mod-MULTI}{\tg{MULTI}}\\
\bottomrule\end{longtable}}
\section{Key readings}
\begin{itemize}
\item Wu et al., \href{https://arxiv.org/abs/2308.08155}{\emph{AutoGen}} (2023)
\item Hong et al., \href{https://arxiv.org/abs/2308.00352}{\emph{MetaGPT}} (2023)
\item Li et al., \href{https://arxiv.org/abs/2303.17760}{\emph{CAMEL}} (2023)
\item Du et al., \href{https://arxiv.org/abs/2305.14325}{\emph{Multiagent Debate}} (2023)
\item Cemri et al., \href{https://arxiv.org/abs/2503.13657}{\emph{Why Do Multi-Agent LLM Systems Fail?}} (2025)
\item Anthropic, \href{https://www.anthropic.com/engineering/multi-agent-research-system}{\emph{How we built our multi-agent research system}}
\end{itemize}
\section{Hands-on}
\begin{itemize}
\item \hyperlink{asg-UVA-7}{\textbf{UVA: Topic 7: MCP + A2A}} --- MCP client for AI2 Asta tools; A2A agents, registry and trivia tournament over ngrok. \textit{Starter code:} \href{https://www.cs.virginia.edu/\%7Ermw7my/Courses/AgenticAISpring2026/Topic7MCP/mcp.html}{[cs.virginia.edu]}
\item \hyperlink{asg-UCHI-3}{\textbf{UChicago: Week 3 notebook}} --- Prompting (0-shot, few-shot, chain-of-reasoning) via the OpenAI API; simulated survey and social-media respondents compared with human data; AutoGen multi-agent dialogues with reflection; Concordia situated simulation; RAG over 5-10 documents vs. a vanilla model. \textit{Starter code:} \href{https://github.com/KnowledgeLab/AI-Agents-for-Social-Science-and-Society-2026/blob/main/Tutorials-Homework_Notebooks/Week_3.ipynb}{[GitHub]}
\item \hyperlink{asg-MANN-6}{\textbf{Mannheim: Project topic 1: Multi-Agent Gaming}} --- Agent teams that cooperate in, or compete at, games (Codenames worked well before). \textit{Resources:} \href{https://github.com/THUDM/AgentBench}{[GitHub]} \href{https://github.com/microsoft/SmartPlay}{[GitHub]}
\item \hyperlink{asg-STAN25-5}{\textbf{Stanford '25: Week 5 --- Agentic Development with Warp}} --- Warp Drive prompts, rules, or MCP integrations, plus multi-agent concurrent sessions working from a task list. \textit{Starter code:} \href{https://github.com/mihail911/modern-software-dev-assignments/tree/fall2025/week5}{[GitHub]}
\end{itemize}
\chapter{Module 8: Coding Agents}\hypertarget{mod-CODE}{}
{\small Tag: \tg{CODE}}
\section{Learning objectives}
\begin{itemize}
\item Explain the architecture of production coding agents (system prompt, core tools, sandbox, context management, sub-agents).
\item Build a coding agent and evaluate it on repository-level tasks.
\item Make repositories and teams agent-ready (AGENTS.md/CLAUDE.md, hooks, tests, review).
\end{itemize}
\section{Core concepts}
Coding is the domain where agents work best today, because the environment provides dense, verifiable feedback (compilers, tests, linters).
\begin{itemize}
\item \textbf{Anatomy.} A long system prompt, a small core toolset (read, write, edit, bash, search), a sandbox, compaction, TODO/plan tools, and optionally sub-agents. Stanford '26 dissects real system prompts and traces. JHU A3 makes you implement path-validated file tools and TODO tools.
\item \textbf{Interfaces matter.} SWE-agent showed that agent--computer interface design (viewer windows, linting edits) moves benchmark scores. Agentless showed that a fixed pipeline can be competitive, so keep a simple baseline.
\item \textbf{Benchmarks.} SWE-bench (and SWE-Gym, R2E-Gym for training), Terminal-Bench, BaxBench and RedCode (security). See \hyperlink{mod-EVAL}{Module~10 (Evaluating Agents)}.
\item \textbf{Practice.} Spec-driven development, context files (AGENTS.md/CLAUDE.md), hooks as lint/test gates, agentic code review, background agents, and ``agent readiness'' of a codebase (Stanford '26 weeks~2--8).
\item \textbf{Research frontier.} Refactoring agents (CU Boulder), coding-agent configuration studies, and evolutionary coding agents (AlphaEvolve, Live-SWE-agent).
\end{itemize}
\section{Curated learning sequence}
\begin{enumerate}
\item \hyperlink{lec-CMU-6}{CMU~6 \emph{Domains 1: Coding Agents}} --- Agents for coding and software development (slides + video).
\item \hyperlink{lec-STAN26-1}{Stanford '26~1 \emph{Course intro + build Claude Code in 200 lines}} and \hyperlink{lec-STAN26-2}{Stanford '26~2 \emph{How SOTA coding agents are designed: deep dive into system prompts}} --- Build a mini Claude Code, then a deep dive into production system prompts.
\item \hyperlink{lec-UW-5}{UW-Madison~5 \emph{Coding agents I}} and \hyperlink{lec-UW-6}{UW-Madison~6 \emph{Coding agents II}} --- SWE-bench, SWE-agent, OpenHands, BaxBench.
\item \hyperlink{lec-CU-13}{CU Boulder~13 \emph{SWE-agent \& AutoCodeRover (students)}} --- SWE-agent and AutoCodeRover (agent--computer interfaces, repository repair).
\item \hyperlink{lec-JHU-7}{JHU~7 \emph{Iterative Refinement and Agentic Framework}} and \hyperlink{lec-JHU-10}{JHU~10 \emph{Agentic Frameworks: Selected Topics}} --- Iterative refinement and agentic frameworks for code (slides).
\item \hyperlink{lec-STAN26-7}{Stanford '26~7 \emph{CLAUDE.md, AGENTS.md, hooks; subagent patterns}}, \hyperlink{lec-STAN26-9}{Stanford '26~9}, \hyperlink{lec-STAN26-11}{Stanford '26~11}, \hyperlink{lec-STAN26-15}{Stanford '26~15} --- Customizing the agent, agent-ready repos, code review, background agents.
\item \hyperlink{lec-UVA-17}{UVA~17 \emph{Topic 11: Coding Agents}} --- History, Claude Code architecture, skills files and evolutionary programming.
\end{enumerate}
\section{Lecture map (46 lectures)}
{\footnotesize\begin{longtable}{P{1.75cm} P{0.4cm} P{5.75cm} P{4.2cm} P{2.65cm}}
\toprule Course & \# & Title (links to catalogue row) & Materials & Tags\\ \midrule\endhead
\hyperlink{course-CMU}{CMU} & 6 & $\star$\,\hyperlink{lec-CMU-6}{Domains 1: Coding Agents} & \href{https://www.cmu-agents.com/slides/lecture-06-coding-agents.pdf}{[slides]} \href{https://www.cmu-agents.com/\#/schedule}{[video]} & \hyperlink{mod-CODE}{\tg{CODE}}\\
\hyperlink{course-CMU}{CMU} & 14 & \hyperlink{lec-CMU-14}{Frameworks 1: OpenHands} & -- & \hyperlink{mod-ARCH}{\tg{ARCH}} \hyperlink{mod-CODE}{\tg{CODE}}\\
\hyperlink{course-STAN26}{Stanford '26} & 1 & $\star$\,\hyperlink{lec-STAN26-1}{Course intro + build Claude Code in 200 lines} & \href{https://themodernsoftware.dev/}{[course site]} & \hyperlink{mod-CODE}{\tg{CODE}} \hyperlink{mod-ARCH}{\tg{ARCH}}\\
\hyperlink{course-STAN26}{Stanford '26} & 2 & $\star$\,\hyperlink{lec-STAN26-2}{How SOTA coding agents are designed: deep dive into system prompts} & \href{https://themodernsoftware.dev/}{[course site]} & \hyperlink{mod-CODE}{\tg{CODE}} \hyperlink{mod-FND}{\tg{FND}}\\
\hyperlink{course-STAN26}{Stanford '26} & 7 & $\star$\,\hyperlink{lec-STAN26-7}{CLAUDE.md, AGENTS.md, hooks; subagent patterns (planner/implementer/reviewer)} & \href{https://themodernsoftware.dev/}{[course site]} & \hyperlink{mod-CODE}{\tg{CODE}} \hyperlink{mod-MULTI}{\tg{MULTI}}\\
\hyperlink{course-STAN26}{Stanford '26} & 9 & $\star$\,\hyperlink{lec-STAN26-9}{Agent readiness: structure, docs, tests, checks} & \href{https://themodernsoftware.dev/}{[course site]} & \hyperlink{mod-CODE}{\tg{CODE}} \hyperlink{mod-EVAL}{\tg{EVAL}}\\
\hyperlink{course-STAN26}{Stanford '26} & 11 & $\star$\,\hyperlink{lec-STAN26-11}{Agentic code review: best practices and architectures} & \href{https://themodernsoftware.dev/}{[course site]} & \hyperlink{mod-CODE}{\tg{CODE}} \hyperlink{mod-EVAL}{\tg{EVAL}}\\
\hyperlink{course-STAN25}{Stanford '25} & 3 & \hyperlink{lec-STAN25-3}{Building a coding agent from scratch} & \href{https://docs.google.com/presentation/d/11CP26VhsjnZOmi9YFgLlonzdib9BLyAlgc4cEvC5Fps/edit?usp=sharing}{[slides]} \href{https://drive.google.com/file/d/1YtpKFVG13DHyQ2i3HOtwyVJOV90nWeL2/view?usp=drive_link}{[drive.google.com]} & \hyperlink{mod-ARCH}{\tg{ARCH}} \hyperlink{mod-TOOL}{\tg{TOOL}} \hyperlink{mod-CODE}{\tg{CODE}}\\
\hyperlink{course-STAN25}{Stanford '25} & 5 & \hyperlink{lec-STAN25-5}{From first prompt to optimal IDE setup} & \href{https://docs.google.com/presentation/d/11pQNCde_mmRnImBat0Zymnp8TCS_cT_1up7zbcj6Sjg/edit?usp=sharing}{[slides]} \href{https://drive.google.com/file/d/1MZ0Qx68Vzw4x5x_XcV8XiPLp7fFDe1LJ/view?usp=drive_link}{[drive.google.com]} & \hyperlink{mod-MEM}{\tg{MEM}} \hyperlink{mod-CODE}{\tg{CODE}}\\
\hyperlink{course-STAN25}{Stanford '25} & 7 & \hyperlink{lec-STAN25-7}{How to be an agent manager} & \href{https://docs.google.com/presentation/d/19mgkwAnJDc7JuJy0zhhoY0ZC15DiNpxL8kchPDnRkRQ/edit?usp=sharing}{[slides]} & \hyperlink{mod-CODE}{\tg{CODE}} \hyperlink{mod-MULTI}{\tg{MULTI}}\\
\hyperlink{course-STAN25}{Stanford '25} & 9 & \hyperlink{lec-STAN25-9}{How to build a breakthrough AI developer product} & \href{https://docs.google.com/presentation/d/1Djd4eBLBbRkma8rFnJAWMT0ptct_UGB8hipmoqFVkxQ/edit?usp=sharing}{[slides]} & \hyperlink{mod-CODE}{\tg{CODE}} \hyperlink{mod-DEPLOY}{\tg{DEPLOY}}\\
\hyperlink{course-STAN25}{Stanford '25} & 11 & \hyperlink{lec-STAN25-11}{AI QA, SAST, DAST, and beyond} & \href{https://docs.google.com/presentation/d/1C05bCLasMDigBbkwdWbiz4WrXibzi6ua4hQQbTod_8c/edit?usp=sharing}{[slides]} & \hyperlink{mod-SAFE}{\tg{SAFE}} \hyperlink{mod-EVAL}{\tg{EVAL}} \hyperlink{mod-CODE}{\tg{CODE}}\\
\hyperlink{course-STAN25}{Stanford '25} & 13 & \hyperlink{lec-STAN25-13}{AI code review} & \href{https://docs.google.com/presentation/d/1NkPzpuSQt6Esbnr2-EnxM9007TL6ebSPFwITyVY-QxU/edit?usp=sharing}{[slides]} & \hyperlink{mod-CODE}{\tg{CODE}} \hyperlink{mod-EVAL}{\tg{EVAL}}\\
\hyperlink{course-STAN25}{Stanford '25} & 15 & \hyperlink{lec-STAN25-15}{End-to-end apps from a single prompt} & \href{https://docs.google.com/presentation/d/1GrVLsfMFIXMiGjIW9D7EJIyLYh_-3ReHHNd_vRfZUoo/edit?usp=sharing}{[slides]} & \hyperlink{mod-CODE}{\tg{CODE}} \hyperlink{mod-APP}{\tg{APP}}\\
\hyperlink{course-STAN25}{Stanford '25} & 19 & \hyperlink{lec-STAN25-19}{Software development in 10 years} & -- & \hyperlink{mod-CODE}{\tg{CODE}} \tg{ORG}\\
\hyperlink{course-UB}{Buffalo} & 12 & \hyperlink{lec-UB-12}{Coding agents, testing, CI, observability} & -- & \hyperlink{mod-CODE}{\tg{CODE}} \hyperlink{mod-DEPLOY}{\tg{DEPLOY}}\\
\hyperlink{course-UVA}{UVA} & 17 & $\star$\,\hyperlink{lec-UVA-17}{Topic 11: Coding Agents} & \href{https://www.cs.virginia.edu/\%7Ermw7my/Courses/AgenticAISpring2026/Topic11Coding/coding.html}{[cs.virginia.edu]} & \hyperlink{mod-CODE}{\tg{CODE}} \hyperlink{mod-ARCH}{\tg{ARCH}}\\
\hyperlink{course-UW}{UW-Madison} & 5 & $\star$\,\hyperlink{lec-UW-5}{Coding agents I} & -- & \hyperlink{mod-CODE}{\tg{CODE}} \hyperlink{mod-EVAL}{\tg{EVAL}}\\
\hyperlink{course-UW}{UW-Madison} & 6 & $\star$\,\hyperlink{lec-UW-6}{Coding agents II} & -- & \hyperlink{mod-CODE}{\tg{CODE}} \hyperlink{mod-SAFE}{\tg{SAFE}} \hyperlink{mod-EVAL}{\tg{EVAL}}\\
\hyperlink{course-UW}{UW-Madison} & 16 & \hyperlink{lec-UW-16}{Self-improving agents II} & -- & \hyperlink{mod-CODE}{\tg{CODE}} \hyperlink{mod-RSN}{\tg{RSN}}\\
\hyperlink{course-UW}{UW-Madison} & 17 & \hyperlink{lec-UW-17}{Searching the solution space I} & -- & \hyperlink{mod-RSN}{\tg{RSN}} \hyperlink{mod-CODE}{\tg{CODE}}\\
\hyperlink{course-JHU}{JHU} & 5 & \hyperlink{lec-JHU-5}{Language Modeling for Synthesis} & \href{https://machine-programming.github.io/lectures/lecture-5.pdf}{[PDF]} & \hyperlink{mod-FND}{\tg{FND}} \hyperlink{mod-CODE}{\tg{CODE}}\\
\hyperlink{course-JHU}{JHU} & 6 & \hyperlink{lec-JHU-6}{Decoding and Prompting for Synthesis} & \href{https://machine-programming.github.io/lectures/lecture-6.pdf}{[PDF]} & \hyperlink{mod-FND}{\tg{FND}} \hyperlink{mod-CODE}{\tg{CODE}}\\
\hyperlink{course-JHU}{JHU} & 7 & $\star$\,\hyperlink{lec-JHU-7}{Iterative Refinement and Agentic Framework} & \href{https://machine-programming.github.io/lectures/lecture-7.pdf}{[PDF]} & \hyperlink{mod-ARCH}{\tg{ARCH}} \hyperlink{mod-CODE}{\tg{CODE}}\\
\hyperlink{course-JHU}{JHU} & 10 & $\star$\,\hyperlink{lec-JHU-10}{Agentic Frameworks: Selected Topics} & \href{https://machine-programming.github.io/lectures/lecture-10.pdf}{[PDF]} & \hyperlink{mod-ARCH}{\tg{ARCH}} \hyperlink{mod-CODE}{\tg{CODE}}\\
\hyperlink{course-JHU}{JHU} & 11 & \hyperlink{lec-JHU-11}{Pre-training of Coding Language Models (1)} & \href{https://machine-programming.github.io/lectures/lecture-11.pdf}{[PDF]} & \hyperlink{mod-TRAIN}{\tg{TRAIN}} \hyperlink{mod-CODE}{\tg{CODE}}\\
\hyperlink{course-JHU}{JHU} & 12 & \hyperlink{lec-JHU-12}{Pre-training of Coding Language Models (2)} & \href{https://machine-programming.github.io/lectures/lecture-12.pdf}{[PDF]} & \hyperlink{mod-TRAIN}{\tg{TRAIN}} \hyperlink{mod-CODE}{\tg{CODE}}\\
\hyperlink{course-JHU}{JHU} & 13 & \hyperlink{lec-JHU-13}{Post-training of Coding Language Models} & -- & \hyperlink{mod-TRAIN}{\tg{TRAIN}} \hyperlink{mod-CODE}{\tg{CODE}}\\
\hyperlink{course-JHU}{JHU} & 15 & \hyperlink{lec-JHU-15}{PL Kaleidoscope: Formal Verification and Theorem Proving} & -- & \hyperlink{mod-APP}{\tg{APP}} \hyperlink{mod-CODE}{\tg{CODE}}\\
\hyperlink{course-CU}{CU Boulder} & 2 & \hyperlink{lec-CU-2}{Survey of agents (Dig)} & \href{https://danny.cs.colorado.edu/courses/csci7000-005_Sp26/pdfs/W1_L2_SurveyOfAgents.pdf}{[PDF]} & \hyperlink{mod-ARCH}{\tg{ARCH}} \hyperlink{mod-CODE}{\tg{CODE}}\\
\hyperlink{course-CU}{CU Boulder} & 3 & \hyperlink{lec-CU-3}{Agentic SE roadmap (Dig)} & \href{https://danny.cs.colorado.edu/courses/csci7000-005_Sp26/pdfs/W2_L2_Roadmap.pdf}{[PDF]} & \hyperlink{mod-CODE}{\tg{CODE}} \hyperlink{mod-ARCH}{\tg{ARCH}}\\
\hyperlink{course-CU}{CU Boulder} & 4 & \hyperlink{lec-CU-4}{Refactoring agents (students)} & -- & \hyperlink{mod-CODE}{\tg{CODE}} \hyperlink{mod-MULTI}{\tg{MULTI}} \hyperlink{mod-EVAL}{\tg{EVAL}}\\
\hyperlink{course-CU}{CU Boulder} & 5 & \hyperlink{lec-CU-5}{MM-Assist \& MANTRA (Dig)} & \href{https://danny.cs.colorado.edu/courses/csci7000-005_Sp26/pdfs/W3_L2_MM-Mantra.pdf}{[PDF]} & \hyperlink{mod-CODE}{\tg{CODE}} \hyperlink{mod-MULTI}{\tg{MULTI}}\\
\hyperlink{course-CU}{CU Boulder} & 13 & $\star$\,\hyperlink{lec-CU-13}{SWE-agent \& AutoCodeRover (students)} & -- & \hyperlink{mod-CODE}{\tg{CODE}} \hyperlink{mod-ARCH}{\tg{ARCH}}\\
\hyperlink{course-CU}{CU Boulder} & 15 & \hyperlink{lec-CU-15}{Multiagent debate \& Agentless (students)} & -- & \hyperlink{mod-MULTI}{\tg{MULTI}} \hyperlink{mod-RSN}{\tg{RSN}} \hyperlink{mod-CODE}{\tg{CODE}}\\
\hyperlink{course-CU}{CU Boulder} & 16 & \hyperlink{lec-CU-16}{LangGraph tutorial \& OpenHands} & -- & \hyperlink{mod-ARCH}{\tg{ARCH}} \hyperlink{mod-CODE}{\tg{CODE}}\\
\hyperlink{course-CU}{CU Boulder} & 18 & \hyperlink{lec-CU-18}{RefAgent \& Agentic Skills (students)} & -- & \hyperlink{mod-CODE}{\tg{CODE}} \hyperlink{mod-ARCH}{\tg{ARCH}} \hyperlink{mod-TOOL}{\tg{TOOL}}\\
\hyperlink{course-CU}{CU Boulder} & 19 & \hyperlink{lec-CU-19}{AppAgent \& coding-agent configuration (students)} & -- & \hyperlink{mod-WEB}{\tg{WEB}} \hyperlink{mod-CODE}{\tg{CODE}}\\
\hyperlink{course-CU}{CU Boulder} & 22 & \hyperlink{lec-CU-22}{MetaClaw \& code smells (students)} & -- & \hyperlink{mod-CODE}{\tg{CODE}} \hyperlink{mod-EVAL}{\tg{EVAL}}\\
\hyperlink{course-UCHI}{UChicago} & 0 & \hyperlink{lec-UCHI-0}{Introduction / Review} & -- & \tg{ORG} \hyperlink{mod-CODE}{\tg{CODE}}\\
\hyperlink{course-S329A}{Stanford 329A} & 4 & \hyperlink{lec-S329A-4}{Learning from Feedback with Tools/Code} & \href{https://www.youtube.com/watch?v=Lxh9RF5S-K0}{[video]} & \hyperlink{mod-TOOL}{\tg{TOOL}} \hyperlink{mod-CODE}{\tg{CODE}} \hyperlink{mod-TRAIN}{\tg{TRAIN}}\\
\hyperlink{course-S329A}{Stanford 329A} & 7 & \hyperlink{lec-S329A-7}{Open-Ended Evolution of Self-Improving Agents} & -- & \hyperlink{mod-ARCH}{\tg{ARCH}} \hyperlink{mod-CODE}{\tg{CODE}} \hyperlink{mod-APP}{\tg{APP}}\\
\hyperlink{course-S329A}{Stanford 329A} & 8 & \hyperlink{lec-S329A-8}{Self-Improvement with Search \& Deep Research Agents} & \href{https://www.youtube.com/watch?v=Uni9dqyuuDM}{[video]} & \hyperlink{mod-WEB}{\tg{WEB}} \hyperlink{mod-RSN}{\tg{RSN}} \hyperlink{mod-CODE}{\tg{CODE}}\\
\hyperlink{course-S329A}{Stanford 329A} & 13 & \hyperlink{lec-S329A-13}{Agentic Frameworks for Software Engineering} & -- & \hyperlink{mod-CODE}{\tg{CODE}} \hyperlink{mod-ARCH}{\tg{ARCH}}\\
\hyperlink{course-BRKAG}{Berkeley '25} & 4 & \hyperlink{lec-BRKAG-4}{Post-Training Verifiable Agents} & \href{https://rdi.berkeley.edu/agentic-ai/slides/AgenticAI_Jiantao\%20Jiao.pdf}{[slides]} \href{https://www.youtube.com/watch?v=3l0Zxus34es\&list=PLS01nW3RtgoqGkm4UeqNeZLccW-OGc1fJ\&index=3}{[video]} & \hyperlink{mod-TRAIN}{\tg{TRAIN}} \hyperlink{mod-EVAL}{\tg{EVAL}} \hyperlink{mod-CODE}{\tg{CODE}} \hyperlink{mod-WEB}{\tg{WEB}}\\
\hyperlink{course-BRKADV}{Berkeley Adv.} & 5 & \hyperlink{lec-BRKADV-5}{Coding Agents and AI for Vulnerability Detection} & \href{https://rdi.berkeley.edu/adv-llm-agents/slides/Code\%20Agents\%20and\%20AI\%20for\%20Vulnerability\%20Detection.pdf}{[slides]} \href{https://www.youtube.com/live/JCk6qJtaCSU}{[video]} & \hyperlink{mod-CODE}{\tg{CODE}} \hyperlink{mod-TOOL}{\tg{TOOL}} \hyperlink{mod-SAFE}{\tg{SAFE}}\\
\bottomrule\end{longtable}}
\section{Key readings}
\begin{itemize}
\item Jimenez et al., \href{https://arxiv.org/abs/2310.06770}{\emph{SWE-bench}} (2023)
\item Yang et al., \href{https://arxiv.org/abs/2405.15793}{\emph{SWE-agent}} (2024)
\item Wang et al., \href{https://arxiv.org/abs/2407.16741}{\emph{OpenHands}} (2024)
\item Xia et al., \href{https://arxiv.org/abs/2407.01489}{\emph{Agentless}} (2024)
\item Zhang et al., \href{https://arxiv.org/abs/2404.05427}{\emph{AutoCodeRover}} (2024)
\item Hassan et al., \href{https://arxiv.org/abs/2509.06216}{\emph{Agentic Software Engineering}} roadmap (2025)
\item Anthropic, \href{https://www.anthropic.com/engineering/claude-code-best-practices}{\emph{Claude Code best practices}}
\end{itemize}
\section{Hands-on}
\begin{itemize}
\item \hyperlink{asg-UW-1}{\textbf{UW-Madison: Build a Coding Agent}} --- A CLI coding agent like Claude Code / Cursor CLI / Gemini CLI, using raw LLM API calls with no agent frameworks (LangChain, Strands, CrewAI banned; thin wrappers like litellm allowed). It must parse and execute tool calls, have a polished UI, a README, and one non-trivial program the agent built. \textit{Spec:} \href{https://github.com/barghouthi/ai-agents/blob/main/assignment.md}{[GitHub]}
\item \hyperlink{asg-JHU-3}{\textbf{JHU: A3: Coding LLM Agents}} --- "WAA" (Web-App Agent): a Python tool-calling loop (tool registry, system prompt, LLM query, parse, execute, terminate or hit the turn limit), file-system tools with path validation and protected files, TODO tools; the agent then builds a personal site, an Express chat room and a creative app; extra credit adds a database plus DB tools. \textit{Starter code:} \href{https://github.com/machine-programming/assignment-3}{[GitHub]}
\item \hyperlink{asg-CMU-1}{\textbf{CMU: A1: Build an Agent Harness}} --- A domain-agnostic ReAct \texttt{Agent} base class: prompt construction, tool-call loop, step limit, model-generated context compaction. On top of it, a \texttt{CodeAgent} (terminal/message tools, skill loader with progressive disclosure via \texttt{invoke\_\allowbreak{}skill}) that fixes a chess-app bug and solves SWE-bench \texttt{django\_\allowbreak{}\_\allowbreak{}django-15368} with and without compaction. Then a \texttt{ChessAgent} (play\_\allowbreak{}move, stateless simulate\_\allowbreak{}move, sandboxed run\_\allowbreak{}python, chess skill) plus an observation A/B experiment. Uses Modal sandboxes, an OpenAI-compatible LLM (default DeepSeek-V4-Flash), uv, and pytest. 100-pt rubric. \textit{Starter code:} \href{https://github.com/cmu-agents/assignment-1}{[GitHub]} \href{https://github.com/cmu-agents/assignment-1/blob/main/ASSIGNMENT.md}{[GitHub]} \href{https://github.com/cmu-agents/chess-app}{[GitHub]}
\item \hyperlink{asg-STAN26-1}{\textbf{Stanford '26: Week 1 --- Trace Dissection of a Real Claude Code Session}} --- Capture a Claude Code session through a mitmproxy reverse proxy, then annotate the system prompt, \texttt{<system-reminder>}s, tool inventory and schemas, error recovery, planning, subagents, and context management. \textit{Starter code:} \href{https://github.com/mihail911/modern-software-dev-assignments/blob/master/week1/assignment.md}{[GitHub]}
\item \hyperlink{asg-STAN25-4}{\textbf{Stanford '25: Week 4 --- The Autonomous Coding Agent IRL}} --- At least 2 Claude Code automations (slash commands, CLAUDE.md, subagents, MCP) used to extend a FastAPI app. \textit{Starter code:} \href{https://github.com/mihail911/modern-software-dev-assignments/tree/fall2025/week4}{[GitHub]}
\item \hyperlink{asg-STAN25-7}{\textbf{Stanford '25: Week 7 --- Exploring AI Code Review Using Graphite}} --- Implement 4 tasks with one-shot agent prompts, open PRs, and compare manual review with Graphite Diamond AI review. \textit{Starter code:} \href{https://github.com/mihail911/modern-software-dev-assignments/tree/fall2025/week7}{[GitHub]}
\end{itemize}
\chapter{Module 9: Web, GUI, Computer-Use and Multimodal Agents}\hypertarget{mod-WEB}{}
{\small Tag: \tg{WEB}}
\section{Learning objectives}
\begin{itemize}
\item Explain observation and action spaces for browser/GUI agents (DOM, accessibility tree, screenshots, coordinates).
\item Build a vision-language agent and know the main web/OS benchmarks.
\end{itemize}
\section{Core concepts}
When there is no API, the agent operates the interface a human would. Observations are HTML/DOM, accessibility trees or screenshots (set-of-marks). Actions are clicks, typing and navigation, or full OS control. These tasks have long horizons, partial observability and irreversible actions, which makes them hard and makes safety and human confirmation central.
Benchmarks progress from MiniWoB and WebShop through Mind2Web and WebArena/VisualWebArena to OSWorld. Training environments include CUA-Gym and OpenCUA. Mobile agents (AppAgent) and embodied agents (SayCan, Voyager) extend the same loop. Personal agents such as the OpenClaw family control real accounts, so read UVA's risk discussion together with \hyperlink{mod-SAFE}{Module~12 (Safety, Security and Oversight)}.
\section{Curated learning sequence}
\begin{enumerate}
\item \hyperlink{lec-CMU-7}{CMU~7 \emph{Computer Use Agents}} --- Computer-use agents: WebArena, VisualWebArena, OSWorld, OpenCUA (slides + video).
\item \hyperlink{lec-UVA-10}{UVA~10 \emph{Topic 6: Vision-Language Models}} --- Vision-language models and a VLM chat agent and video-surveillance agent lab.
\item \hyperlink{lec-CU-19}{CU Boulder~19 \emph{AppAgent \& coding-agent configuration (students)}} --- AppAgent (smartphone GUI) and coding-agent configuration.
\item \hyperlink{lec-UVA-18}{UVA~18 \emph{Topic 12: Personal Agents and Distilling}} --- Personal agents (OpenClaw family, Claude Cowork) and their risks; distillation.
\item \hyperlink{lec-UCHI-8}{UChicago~8 \emph{Multi-Modal and Embodied Agents}} --- Multimodal and embodied agents (notebook).
\end{enumerate}
\section{Lecture map (11 lectures)}
{\footnotesize\begin{longtable}{P{1.75cm} P{0.4cm} P{5.75cm} P{4.2cm} P{2.65cm}}
\toprule Course & \# & Title (links to catalogue row) & Materials & Tags\\ \midrule\endhead
\hyperlink{course-CMU}{CMU} & 7 & $\star$\,\hyperlink{lec-CMU-7}{Computer Use Agents} & \href{https://www.cmu-agents.com/slides/lecture-07-computer-use-agents.pdf}{[slides]} \href{https://www.cmu-agents.com/\#/schedule}{[video]} & \hyperlink{mod-WEB}{\tg{WEB}}\\
\hyperlink{course-UVA}{UVA} & 10 & $\star$\,\hyperlink{lec-UVA-10}{Topic 6: Vision-Language Models} & \href{https://www.cs.virginia.edu/\%7Ermw7my/Courses/AgenticAISpring2026/Topic6VLM/vlm.html}{[cs.virginia.edu]} \href{https://github.com/HenryKautz/clip_test}{[GitHub]} & \hyperlink{mod-WEB}{\tg{WEB}} \hyperlink{mod-APP}{\tg{APP}} \hyperlink{mod-FND}{\tg{FND}}\\
\hyperlink{course-UVA}{UVA} & 18 & $\star$\,\hyperlink{lec-UVA-18}{Topic 12: Personal Agents and Distilling} & \href{https://www.cs.virginia.edu/\%7Ermw7my/Courses/AgenticAISpring2026/Topic12Personal/personal.html}{[cs.virginia.edu]} & \hyperlink{mod-WEB}{\tg{WEB}} \hyperlink{mod-SAFE}{\tg{SAFE}} \hyperlink{mod-TRAIN}{\tg{TRAIN}}\\
\hyperlink{course-UW}{UW-Madison} & 15 & \hyperlink{lec-UW-15}{Self-improving agents I} & -- & \hyperlink{mod-ARCH}{\tg{ARCH}} \hyperlink{mod-WEB}{\tg{WEB}}\\
\hyperlink{course-CU}{CU Boulder} & 19 & $\star$\,\hyperlink{lec-CU-19}{AppAgent \& coding-agent configuration (students)} & -- & \hyperlink{mod-WEB}{\tg{WEB}} \hyperlink{mod-CODE}{\tg{CODE}}\\
\hyperlink{course-UCHI}{UChicago} & 8 & $\star$\,\hyperlink{lec-UCHI-8}{Multi-Modal and Embodied Agents} & \href{https://github.com/KnowledgeLab/AI-Agents-for-Social-Science-and-Society-2026/blob/main/Tutorials-Homework_Notebooks/Week_8_Multimodal.ipynb}{[GitHub]} & \hyperlink{mod-WEB}{\tg{WEB}} \hyperlink{mod-RSN}{\tg{RSN}} \hyperlink{mod-APP}{\tg{APP}}\\
\hyperlink{course-S329A}{Stanford 329A} & 8 & \hyperlink{lec-S329A-8}{Self-Improvement with Search \& Deep Research Agents} & \href{https://www.youtube.com/watch?v=Uni9dqyuuDM}{[video]} & \hyperlink{mod-WEB}{\tg{WEB}} \hyperlink{mod-RSN}{\tg{RSN}} \hyperlink{mod-CODE}{\tg{CODE}}\\
\hyperlink{course-BRKAG}{Berkeley '25} & 4 & \hyperlink{lec-BRKAG-4}{Post-Training Verifiable Agents} & \href{https://rdi.berkeley.edu/agentic-ai/slides/AgenticAI_Jiantao\%20Jiao.pdf}{[slides]} \href{https://www.youtube.com/watch?v=3l0Zxus34es\&list=PLS01nW3RtgoqGkm4UeqNeZLccW-OGc1fJ\&index=3}{[video]} & \hyperlink{mod-TRAIN}{\tg{TRAIN}} \hyperlink{mod-EVAL}{\tg{EVAL}} \hyperlink{mod-CODE}{\tg{CODE}} \hyperlink{mod-WEB}{\tg{WEB}}\\
\hyperlink{course-BRKADV}{Berkeley Adv.} & 6 & \hyperlink{lec-BRKADV-6}{Multimodal Autonomous AI Agents} & \href{https://rdi.berkeley.edu/adv-llm-agents/slides/ruslan-multimodal.pdf}{[slides]} \href{https://www.youtube.com/live/RPINOYM12RU}{[video]} & \hyperlink{mod-WEB}{\tg{WEB}} \hyperlink{mod-EVAL}{\tg{EVAL}} \hyperlink{mod-RSN}{\tg{RSN}}\\
\hyperlink{course-BRKADV}{Berkeley Adv.} & 7 & \hyperlink{lec-BRKADV-7}{Multimodal Agents -- From Perception to Action} & \href{https://rdi.berkeley.edu/adv-llm-agents/slides/Multimodal_Agent_caiming.pdf}{[slides]} \href{https://www.youtube.com/live/n__Tim8K2IY}{[video]} & \hyperlink{mod-WEB}{\tg{WEB}} \hyperlink{mod-EVAL}{\tg{EVAL}}\\
\hyperlink{course-S329T}{Stanford 329T} & 4 & \hyperlink{lec-S329T-4}{Recap, further discussion, and planning for projects} & \href{https://web.stanford.edu/class/cs329t/slides/Lecture\%204\%20--\%20CS\%20329T\%20Fall\%202025.pdf}{[slides]} & \tg{ORG} \hyperlink{mod-WEB}{\tg{WEB}}\\
\bottomrule\end{longtable}}
\section{Key readings}
\begin{itemize}
\item Zhou et al., \href{https://arxiv.org/abs/2307.13854}{\emph{WebArena}} (2023); Koh et al., \href{https://arxiv.org/abs/2401.13649}{\emph{VisualWebArena}} (2024)
\item Xie et al., \href{https://arxiv.org/abs/2404.07972}{\emph{OSWorld}} (2024)
\item Deng et al., \href{https://arxiv.org/abs/2306.06070}{\emph{Mind2Web}} (2023)
\item Yao et al., \href{https://arxiv.org/abs/2207.01206}{\emph{WebShop}} (2022)
\item Zhang et al., \href{https://arxiv.org/abs/2312.13771}{\emph{AppAgent}} (2023)
\item Wang et al., \href{https://arxiv.org/abs/2305.16291}{\emph{Voyager}} (2023)
\end{itemize}
\section{Hands-on}
\begin{itemize}
\item \hyperlink{asg-UVA-6}{\textbf{UVA: Topic 6: VLM agents}} --- Image-chat LangGraph agent; video-surveillance agent with OpenCV + LLaVA (optional TTS intruder alert). \textit{Starter code:} \href{https://www.cs.virginia.edu/\%7Ermw7my/Courses/AgenticAISpring2026/Topic6VLM/vlm.html}{[cs.virginia.edu]}
\item \hyperlink{asg-MANN-9}{\textbf{Mannheim: Project topic 4: Online Shopping Assistant}} --- Product search/recommendation agent with structured feature comparison. \textit{Resources:} \href{https://webarena.dev/}{[webarena.dev]} \href{https://webshop-pnlp.github.io/}{[webshop-pnlp.github.io]}
\item \hyperlink{asg-MANN-7}{\textbf{Mannheim: Project topic 2: QA Agents}} --- RAG (multi-)agent QA over web search that explains its evidence. \textit{Resources:} \href{https://openai.com/index/browsecomp/}{[openai.com]} \href{https://www.aicrowd.com/challenges/meta-comprehensive-rag-benchmark-kdd-cup-2024}{[aicrowd.com]}
\end{itemize}
\chapter{Module 10: Evaluating Agents}\hypertarget{mod-EVAL}{}
{\small Tag: \tg{EVAL}}
\section{Learning objectives}
\begin{itemize}
\item Design task-level evals with success criteria, failure taxonomies and reproducible scripts.
\item Build and validate an LLM-as-judge against human labels.
\item Report results honestly (pass@k, variance, cost) and avoid benchmark pitfalls.
\end{itemize}
\section{Core concepts}
Evaluation separates agent engineering from demo-building, and every serious course in this collection weights it heavily.
\begin{itemize}
\item \textbf{Levels.} Unit tests for tools, trajectory and intermediate-step checks, end-to-end task success, and human judgments of usefulness.
\item \textbf{Verifiers.} Deterministic tests where possible (SWE-bench, Harbor tasks with reference solutions and \emph{mutants} in CMU A2). LLM-as-judge where not, \emph{validated against human labels} (CMU A2 scores judges by MCC).
\item \textbf{Statistics and cost.} Agents are stochastic: report pass@k over repeated runs (JHU A2), confidence intervals and dollars per task. \emph{AI Agents That Matter} criticizes leaderboards that ignore cost and overfitting.
\item \textbf{Error analysis.} A failure taxonomy (Buffalo Lab~4; CMU A2's four error families) turns logs into engineering decisions.
\item \textbf{Benchmarks:} SWE-bench, $\tau$-bench, AgentBench, WebArena, AgentDojo (security), BrowseComp and DeepResearch Bench (research agents), and METR's time-horizon measurements.
\end{itemize}
\section{Curated learning sequence}
\begin{enumerate}
\item \hyperlink{asg-CMU-2}{\textbf{CMU: A2: Evaluating a Data-Visualization Agent}} --- The single best exercise here: author tasks, label runs, build and score a judge.
\item \hyperlink{lec-NEU-18}{Northeastern~18 \emph{Agentic AI Evaluation}} --- SWE-bench, $\tau$-bench, LLM-as-a-judge, LangSmith evaluation.
\item \hyperlink{lec-UB-9}{Buffalo~9 \emph{Evaluation of agents}} --- From single-prompt to task-level evals (then Lab~4).
\item \hyperlink{asg-JHU-2}{\textbf{JHU: A2: Evaluating Coding LLMs}} --- A rigorous pass@k study across languages, models and prompting strategies.
\item \hyperlink{lec-CMU-20}{CMU~20 \emph{Search 1: Reranking and Critic Models}} --- Critic models: evaluation used at inference time.
\item \hyperlink{lec-UW-22}{UW-Madison~22 \emph{Memory \& other topics II}} --- Includes METR's \emph{Measuring AI Ability to Complete Long Tasks}.
\end{enumerate}
\section{Lecture map (36 lectures)}
{\footnotesize\begin{longtable}{P{1.75cm} P{0.4cm} P{5.75cm} P{4.2cm} P{2.65cm}}
\toprule Course & \# & Title (links to catalogue row) & Materials & Tags\\ \midrule\endhead
\hyperlink{course-CMU}{CMU} & 19 & \hyperlink{lec-CMU-19}{Interaction 2: Human-Agent Interaction} & -- & \hyperlink{mod-APP}{\tg{APP}} \hyperlink{mod-EVAL}{\tg{EVAL}}\\
\hyperlink{course-CMU}{CMU} & 20 & $\star$\,\hyperlink{lec-CMU-20}{Search 1: Reranking and Critic Models} & -- & \hyperlink{mod-RSN}{\tg{RSN}} \hyperlink{mod-EVAL}{\tg{EVAL}}\\
\hyperlink{course-STAN26}{Stanford '26} & 9 & \hyperlink{lec-STAN26-9}{Agent readiness: structure, docs, tests, checks} & \href{https://themodernsoftware.dev/}{[course site]} & \hyperlink{mod-CODE}{\tg{CODE}} \hyperlink{mod-EVAL}{\tg{EVAL}}\\
\hyperlink{course-STAN26}{Stanford '26} & 11 & \hyperlink{lec-STAN26-11}{Agentic code review: best practices and architectures} & \href{https://themodernsoftware.dev/}{[course site]} & \hyperlink{mod-CODE}{\tg{CODE}} \hyperlink{mod-EVAL}{\tg{EVAL}}\\
\hyperlink{course-STAN25}{Stanford '25} & 11 & \hyperlink{lec-STAN25-11}{AI QA, SAST, DAST, and beyond} & \href{https://docs.google.com/presentation/d/1C05bCLasMDigBbkwdWbiz4WrXibzi6ua4hQQbTod_8c/edit?usp=sharing}{[slides]} & \hyperlink{mod-SAFE}{\tg{SAFE}} \hyperlink{mod-EVAL}{\tg{EVAL}} \hyperlink{mod-CODE}{\tg{CODE}}\\
\hyperlink{course-STAN25}{Stanford '25} & 13 & \hyperlink{lec-STAN25-13}{AI code review} & \href{https://docs.google.com/presentation/d/1NkPzpuSQt6Esbnr2-EnxM9007TL6ebSPFwITyVY-QxU/edit?usp=sharing}{[slides]} & \hyperlink{mod-CODE}{\tg{CODE}} \hyperlink{mod-EVAL}{\tg{EVAL}}\\
\hyperlink{course-NEU}{Northeastern} & 18 & $\star$\,\hyperlink{lec-NEU-18}{Agentic AI Evaluation} & -- & \hyperlink{mod-EVAL}{\tg{EVAL}}\\
\hyperlink{course-UB}{Buffalo} & 9 & $\star$\,\hyperlink{lec-UB-9}{Evaluation of agents} & -- & \hyperlink{mod-EVAL}{\tg{EVAL}}\\
\hyperlink{course-UVA}{UVA} & 1 & \hyperlink{lec-UVA-1}{Topic 1: Running an LLM} & \href{https://www.cs.virginia.edu/\%7Ermw7my/Courses/AgenticAISpring2026/Topic1Running/Running_an_LLM.html}{[cs.virginia.edu]} & \hyperlink{mod-FND}{\tg{FND}} \hyperlink{mod-EVAL}{\tg{EVAL}} \hyperlink{mod-DEPLOY}{\tg{DEPLOY}}\\
\hyperlink{course-UVA}{UVA} & 9 & \hyperlink{lec-UVA-9}{Topic 5: Retrieval Augmented Generation} & \href{https://www.cs.virginia.edu/\%7Ermw7my/Courses/AgenticAISpring2026/Topic5RAG/rag.html}{[cs.virginia.edu]} \href{https://github.com/HenryKautz/pdf_ocr}{[GitHub]} & \hyperlink{mod-MEM}{\tg{MEM}} \hyperlink{mod-EVAL}{\tg{EVAL}}\\
\hyperlink{course-UVA}{UVA} & 21 & \hyperlink{lec-UVA-21}{Topic 15: Visual Reasoning} & \href{https://www.cs.virginia.edu/\%7Ermw7my/Courses/AgenticAISpring2026/Topic15Visual/visual.html}{[cs.virginia.edu]} \href{https://www.youtube.com/watch?v=Gs9TPZmzo-s}{[video]} & \hyperlink{mod-RSN}{\tg{RSN}} \hyperlink{mod-EVAL}{\tg{EVAL}}\\
\hyperlink{course-UW}{UW-Madison} & 5 & \hyperlink{lec-UW-5}{Coding agents I} & -- & \hyperlink{mod-CODE}{\tg{CODE}} \hyperlink{mod-EVAL}{\tg{EVAL}}\\
\hyperlink{course-UW}{UW-Madison} & 6 & \hyperlink{lec-UW-6}{Coding agents II} & -- & \hyperlink{mod-CODE}{\tg{CODE}} \hyperlink{mod-SAFE}{\tg{SAFE}} \hyperlink{mod-EVAL}{\tg{EVAL}}\\
\hyperlink{course-UW}{UW-Madison} & 8 & \hyperlink{lec-UW-8}{Deep research agents II} & -- & \hyperlink{mod-APP}{\tg{APP}} \hyperlink{mod-MEM}{\tg{MEM}} \hyperlink{mod-EVAL}{\tg{EVAL}}\\
\hyperlink{course-UW}{UW-Madison} & 11 & \hyperlink{lec-UW-11}{Security \& prompt injection I} & -- & \hyperlink{mod-SAFE}{\tg{SAFE}} \hyperlink{mod-EVAL}{\tg{EVAL}}\\
\hyperlink{course-UW}{UW-Madison} & 22 & $\star$\,\hyperlink{lec-UW-22}{Memory \& other topics II} & -- & \hyperlink{mod-MEM}{\tg{MEM}} \hyperlink{mod-MULTI}{\tg{MULTI}} \hyperlink{mod-EVAL}{\tg{EVAL}}\\
\hyperlink{course-CU}{CU Boulder} & 4 & \hyperlink{lec-CU-4}{Refactoring agents (students)} & -- & \hyperlink{mod-CODE}{\tg{CODE}} \hyperlink{mod-MULTI}{\tg{MULTI}} \hyperlink{mod-EVAL}{\tg{EVAL}}\\
\hyperlink{course-CU}{CU Boulder} & 9 & \hyperlink{lec-CU-9}{ReAct \& Web API code gen (students)} & -- & \hyperlink{mod-RSN}{\tg{RSN}} \hyperlink{mod-TOOL}{\tg{TOOL}} \hyperlink{mod-EVAL}{\tg{EVAL}}\\
\hyperlink{course-CU}{CU Boulder} & 22 & \hyperlink{lec-CU-22}{MetaClaw \& code smells (students)} & -- & \hyperlink{mod-CODE}{\tg{CODE}} \hyperlink{mod-EVAL}{\tg{EVAL}}\\
\hyperlink{course-MANN}{Mannheim} & x1 & \hyperlink{lec-MANN-x1}{Evaluation of LLMs and LLM agents} & -- & \hyperlink{mod-EVAL}{\tg{EVAL}}\\
\hyperlink{course-UCHI}{UChicago} & 4 & \hyperlink{lec-UCHI-4}{Experimental Designs with AI Agent and Human Interactions} & \href{https://github.com/KnowledgeLab/AI-Agents-for-Social-Science-and-Society-2026/blob/main/Tutorials-Homework_Notebooks/week_4_2026.ipynb}{[GitHub]} & \hyperlink{mod-APP}{\tg{APP}} \hyperlink{mod-EVAL}{\tg{EVAL}}\\
\hyperlink{course-UCHI}{UChicago} & 5 & \hyperlink{lec-UCHI-5}{Sampling, Fine Tuning, Benchmarking, and Tools} & \href{https://github.com/KnowledgeLab/AI-Agents-for-Social-Science-and-Society-2026/blob/main/Tutorials-Homework_Notebooks/Week_5_2026.ipynb}{[GitHub]} & \hyperlink{mod-TRAIN}{\tg{TRAIN}} \hyperlink{mod-EVAL}{\tg{EVAL}} \hyperlink{mod-TOOL}{\tg{TOOL}}\\
\hyperlink{course-S329A}{Stanford 329A} & 2 & \hyperlink{lec-S329A-2}{Test-time Compute Scaling} & \href{https://www.youtube.com/watch?v=-Ggc37xLj_Y}{[video]} & \hyperlink{mod-RSN}{\tg{RSN}} \hyperlink{mod-EVAL}{\tg{EVAL}}\\
\hyperlink{course-S329A}{Stanford 329A} & 3 & \hyperlink{lec-S329A-3}{Robust Verification} & \href{https://www.youtube.com/watch?v=p7TdPUcPoik}{[video]} & \hyperlink{mod-RSN}{\tg{RSN}} \hyperlink{mod-EVAL}{\tg{EVAL}}\\
\hyperlink{course-S329A}{Stanford 329A} & 17 & \hyperlink{lec-S329A-17}{Agentic Evaluations \& Long-Horizon Tasks} & \href{https://www.youtube.com/watch?v=8JAqLnTaZu4}{[video]} & \hyperlink{mod-EVAL}{\tg{EVAL}}\\
\hyperlink{course-BRKAG}{Berkeley '25} & 4 & \hyperlink{lec-BRKAG-4}{Post-Training Verifiable Agents} & \href{https://rdi.berkeley.edu/agentic-ai/slides/AgenticAI_Jiantao\%20Jiao.pdf}{[slides]} \href{https://www.youtube.com/watch?v=3l0Zxus34es\&list=PLS01nW3RtgoqGkm4UeqNeZLccW-OGc1fJ\&index=3}{[video]} & \hyperlink{mod-TRAIN}{\tg{TRAIN}} \hyperlink{mod-EVAL}{\tg{EVAL}} \hyperlink{mod-CODE}{\tg{CODE}} \hyperlink{mod-WEB}{\tg{WEB}}\\
\hyperlink{course-BRKAG}{Berkeley '25} & 5 & \hyperlink{lec-BRKAG-5}{Agent Evaluation \& Project Overview} & \href{https://rdi.berkeley.edu/agentic-ai/slides/LLM\%20Agent\%20Evaluations\%20\%26\%20Project\%20Overview.pdf}{[slides]} \href{https://www.youtube.com/watch?v=VfOA2a0dj4w}{[video]} & \hyperlink{mod-EVAL}{\tg{EVAL}} \tg{ORG}\\
\hyperlink{course-BRKAG}{Berkeley '25} & 8 & \hyperlink{lec-BRKAG-8}{Predictable Noise in LLM} & \href{https://rdi.berkeley.edu/agentic-ai/slides/PredEval.pdf}{[slides]} \href{https://www.youtube.com/watch?v=HV8pugcFVO0}{[video]} & \hyperlink{mod-EVAL}{\tg{EVAL}}\\
\hyperlink{course-BRKAG}{Berkeley '25} & 10 & \hyperlink{lec-BRKAG-10}{Practical Lessons from Deploying Real-World AI Agents} & \href{https://rdi.berkeley.edu/agentic-ai/slides/Practical_Lessons_AI_Agents.pdf}{[slides]} \href{https://www.youtube.com/watch?v=sfJM4LaiYsM}{[video]} & \hyperlink{mod-DEPLOY}{\tg{DEPLOY}} \hyperlink{mod-EVAL}{\tg{EVAL}} \hyperlink{mod-APP}{\tg{APP}}\\
\hyperlink{course-BRKADV}{Berkeley Adv.} & 6 & \hyperlink{lec-BRKADV-6}{Multimodal Autonomous AI Agents} & \href{https://rdi.berkeley.edu/adv-llm-agents/slides/ruslan-multimodal.pdf}{[slides]} \href{https://www.youtube.com/live/RPINOYM12RU}{[video]} & \hyperlink{mod-WEB}{\tg{WEB}} \hyperlink{mod-EVAL}{\tg{EVAL}} \hyperlink{mod-RSN}{\tg{RSN}}\\
\hyperlink{course-BRKADV}{Berkeley Adv.} & 7 & \hyperlink{lec-BRKADV-7}{Multimodal Agents -- From Perception to Action} & \href{https://rdi.berkeley.edu/adv-llm-agents/slides/Multimodal_Agent_caiming.pdf}{[slides]} \href{https://www.youtube.com/live/n__Tim8K2IY}{[video]} & \hyperlink{mod-WEB}{\tg{WEB}} \hyperlink{mod-EVAL}{\tg{EVAL}}\\
\hyperlink{course-S329T}{Stanford 329T} & 3 & \hyperlink{lec-S329T-3}{Agentic AI} & \href{https://web.stanford.edu/class/cs329t/slides/Lecture\%203\%20--\%20CS\%20329T\%20Fall\%202025.pdf}{[slides]} & \hyperlink{mod-ARCH}{\tg{ARCH}} \hyperlink{mod-EVAL}{\tg{EVAL}}\\
\hyperlink{course-S329T}{Stanford 329T} & 7 & \hyperlink{lec-S329T-7}{Systematic evaluation of agentic systems (guest: Allison Jia)} & \href{https://web.stanford.edu/class/cs329t/slides/Allison\%20Jia\%20-\%20slides.pdf}{[slides]} & \hyperlink{mod-EVAL}{\tg{EVAL}}\\
\hyperlink{course-S329T}{Stanford 329T} & 8 & \hyperlink{lec-S329T-8}{MCP, effective tool descriptions, evals (guests: Yusuf Ozuysal, Ishita Gupta)} & \href{https://web.stanford.edu/class/cs329t/slides/Model\%20Context\%20Protocol\%20\%28\%20Yusuf\%20Ozuysal\%29.pdf}{[slides (Ozuysal)]} \href{https://web.stanford.edu/class/cs329t/slides/ishita-mcp.pdf}{[slides (Gupta)]} & \hyperlink{mod-TOOL}{\tg{TOOL}} \hyperlink{mod-EVAL}{\tg{EVAL}}\\
\hyperlink{course-S329T}{Stanford 329T} & 10 & \hyperlink{lec-S329T-10}{Data curation and RL to optimize agents (guest: Trung Vu, Bespoke Labs)} & \href{https://web.stanford.edu/class/cs329t/slides/Lecture\%208\%20-\%20\%5B10_23_2025\%5D\%20Reinforcement\%20Learning\%20\%20for\%20Agents.pdf}{[slides]} & \hyperlink{mod-TRAIN}{\tg{TRAIN}} \hyperlink{mod-TOOL}{\tg{TOOL}} \hyperlink{mod-EVAL}{\tg{EVAL}}\\
\hyperlink{course-S329T}{Stanford 329T} & 12 & \hyperlink{lec-S329T-12}{Security for agents: adversarial attacks and guardrails (guests: Matt Fredrikson, Andy Zou)} & \href{https://web.stanford.edu/class/cs329t/slides/lecture\%2010\%20-\%2020251030\%20-\%20security\%20for\%20agents\%20-\%20andy\%20zou.pdf}{[slides]} & \hyperlink{mod-SAFE}{\tg{SAFE}} \hyperlink{mod-EVAL}{\tg{EVAL}}\\
\bottomrule\end{longtable}}
\section{Key readings}
\begin{itemize}
\item Zheng et al., \href{https://arxiv.org/abs/2306.05685}{\emph{Judging LLM-as-a-Judge}} (2023)
\item Yao et al., \href{https://arxiv.org/abs/2406.12045}{\emph{$\tau$-bench}} (2024)
\item Liu et al., \href{https://arxiv.org/abs/2308.03688}{\emph{AgentBench}} (2023)
\item Kapoor et al., \href{https://arxiv.org/abs/2407.01502}{\emph{AI Agents That Matter}} (2024)
\item Mohammadi et al., \href{https://arxiv.org/abs/2507.21504}{\emph{Evaluation and Benchmarking of LLM Agents}} (KDD 2025)
\item METR, \href{https://arxiv.org/abs/2503.14499}{\emph{Measuring AI Ability to Complete Long Tasks}} (2025)
\end{itemize}
\section{Hands-on}
\begin{itemize}
\item \hyperlink{asg-CMU-2}{\textbf{CMU: A2: Evaluating a Data-Visualization Agent}} --- A \texttt{validate(run)} LLM-as-judge validator that flags four error families (execution\_\allowbreak{}failure, wrong\_\allowbreak{}data, wrong\_\allowbreak{}chart, hard\_\allowbreak{}to\_\allowbreak{}read), scored by per-family and macro MCC against human labels. Students author 5+ new tasks, run 3 fixed agents (15+ runs), human-label the runs, and package the tasks as Harbor environments with deterministic tests, a reference solution, and two mutants. The judge is fixed to Qwen3-VL-30B on Modal. Grading: report 40\%, private-set MCC 30\%, quiz 20\%, code 10\%. \textit{Starter code:} \href{https://github.com/cmu-agents/assignment-2}{[GitHub]}
\item \hyperlink{asg-JHU-2}{\textbf{JHU: A2: Evaluating Coding LLMs}} --- Execution harnesses for Python/Rust/OCaml; 4 prompting strategies (zero-shot, two-step CoT, iterative refinement with feedback up to 3 rounds, a custom method of about 5 calls or fewer); pass@1/2/3 pipeline over 24 experiments (Claude Haiku vs Sonnet x 3 languages x 4 strategies) on a 16-item xCodeEval/MBPP set; plots. \textit{Starter code:} \href{https://github.com/machine-programming/assignment-2}{[GitHub]}
\item \hyperlink{asg-UB-4}{\textbf{Buffalo: Lab 4: Evaluation Harness}} --- Eval suite for a bounded task: test cases, success criteria, metrics, failure taxonomy, reproducible script. \textit{(spec only)}
\item \hyperlink{asg-UVA-5}{\textbf{UVA: Topic 5: RAG experiments}} --- 11 controlled RAG experiments (RAG vs no-RAG, top-K, chunking, prompts, unanswerable questions). \textit{Starter code:} \href{https://www.cs.virginia.edu/\%7Ermw7my/Courses/AgenticAISpring2026/Topic5RAG/rag.html}{[cs.virginia.edu]}
\end{itemize}
\textbf{Related supplements:} \hyperref[sup-notion]{Notion: How we built a personal data scientist for every Notion employee (David Thomas, October 2026)}.
\chapter{Module 11: Training and Optimizing Agents}\hypertarget{mod-TRAIN}{}
{\small Tag: \tg{TRAIN}}
\section{Learning objectives}
\begin{itemize}
\item Fine-tune a model on agent trajectories (SFT, LoRA/QLoRA) and know when fine-tuning beats prompting or RAG.
\item Implement policy-gradient RL (REINFORCE, GRPO) for multi-turn agents, including async rollouts.
\item Optimize prompts and programs automatically (DSPy, TextGrad, GEPA).
\end{itemize}
\section{Core concepts}
\begin{itemize}
\item \textbf{SFT and distillation.} Collect or synthesize high-quality trajectories from a strong model, filter them with verifiers, and fine-tune a smaller model (TinyAgent, CMU Training~1, UVA's Tinker text-to-SQL lab). LoRA and QLoRA make this cheap.
\item \textbf{RL from verifiable rewards.} The policy generates trajectories in an environment and receives a reward (tests passed, task solved). Advantages are computed by REINFORCE baselines, PPO with GAE, or GRPO's group-relative normalization (DeepSeekMath, DeepSeek-R1, Dr.~GRPO). Multi-turn agent RL adds credit assignment across tool calls and requires masking tool-output tokens.
\item \textbf{RL systems.} Rollouts dominate cost, so asynchronous producer--consumer training with off-policy corrections (importance weighting) is standard. CMU A3 makes you build both sync and async GRPO on Tinker with a SkyRL server.
\item \textbf{Optimizing without gradients.} DSPy compiles LM programs. TextGrad, GEPA, ADAS and Promptbreeder evolve prompts or whole agent designs (UW weeks on optimizing and self-improving agents).
\end{itemize}
\section{Curated learning sequence}
\begin{enumerate}
\item \hyperlink{lec-CMU-8}{CMU~8 \emph{Training 1: SFT}} --- Supervised fine-tuning for agents (slides + video).
\item \hyperlink{lec-CMU-9}{CMU~9 \emph{Training 2: RL Basics}} --- RL basics: ReST, ReST-EM, DeepSeekMath/GRPO, Dr.~GRPO (slides + video).
\item \hyperlink{lec-CMU-11}{CMU~11 \emph{Training 3: Advanced RL Algorithms}} and \hyperlink{lec-CMU-12}{CMU~12 \emph{Training 4: RL Systems}} --- Advanced RL algorithms, then RL systems (slides).
\item \hyperlink{lec-UVA-13}{UVA~13 \emph{Topic 7 cont. (A2A) + Topic 8: LLM Fine-Tuning}} --- Fine-tuning versus RAG, QLoRA, the Tinker pipeline, text-to-SQL.
\item \hyperlink{lec-JHU-11}{JHU~11 \emph{Pre-training of Coding Language Models (1)}}--\hyperlink{lec-JHU-13}{JHU~13} --- Pre- and post-training of code LLMs.
\item \hyperlink{lec-UW-13}{UW-Madison~13 \emph{Optimizing agents I}} and \hyperlink{lec-UW-14}{UW-Madison~14 \emph{Optimizing agents II}} --- DSPy/TextGrad, then ADAS/GEPA.
\item Supplement: \hyperref[sup-tinyagent]{TinyAgent} --- An end-to-end distillation case study.
\end{enumerate}
\section{Lecture map (30 lectures)}
{\footnotesize\begin{longtable}{P{1.75cm} P{0.4cm} P{5.75cm} P{4.2cm} P{2.65cm}}
\toprule Course & \# & Title (links to catalogue row) & Materials & Tags\\ \midrule\endhead
\hyperlink{course-CMU}{CMU} & 8 & $\star$\,\hyperlink{lec-CMU-8}{Training 1: SFT} & \href{https://www.cmu-agents.com/slides/lecture-08-sft.pdf}{[slides]} \href{https://www.cmu-agents.com/\#/schedule}{[video]} & \hyperlink{mod-TRAIN}{\tg{TRAIN}}\\
\hyperlink{course-CMU}{CMU} & 9 & $\star$\,\hyperlink{lec-CMU-9}{Training 2: RL Basics} & \href{https://www.cmu-agents.com/slides/lecture-09-rl-basics.pdf}{[slides]} \href{https://www.cmu-agents.com/\#/schedule}{[video]} & \hyperlink{mod-TRAIN}{\tg{TRAIN}}\\
\hyperlink{course-CMU}{CMU} & 11 & $\star$\,\hyperlink{lec-CMU-11}{Training 3: Advanced RL Algorithms} & \href{https://www.cmu-agents.com/slides/lecture-11-rl-advanced.pdf}{[slides]} & \hyperlink{mod-TRAIN}{\tg{TRAIN}}\\
\hyperlink{course-CMU}{CMU} & 12 & $\star$\,\hyperlink{lec-CMU-12}{Training 4: RL Systems} & \href{https://www.cmu-agents.com/slides/lecture-12-rl-systems.pdf}{[slides]} & \hyperlink{mod-TRAIN}{\tg{TRAIN}} \hyperlink{mod-DEPLOY}{\tg{DEPLOY}}\\
\hyperlink{course-NEU}{Northeastern} & 3 & \hyperlink{lec-NEU-3}{LLM Foundations (II)} & \href{https://agenticllms.github.io/week2b-setup.html}{[agenticllms.github.io]} \href{https://github.com/agenticllms/lecture-files/tree/main/W02b}{[GitHub]} & \hyperlink{mod-FND}{\tg{FND}} \hyperlink{mod-TRAIN}{\tg{TRAIN}}\\
\hyperlink{course-UVA}{UVA} & 13 & $\star$\,\hyperlink{lec-UVA-13}{Topic 7 cont. (A2A) + Topic 8: LLM Fine-Tuning} & \href{https://www.cs.virginia.edu/\%7Ermw7my/Courses/AgenticAISpring2026/Topic8FineTuning/finetuning.html}{[cs.virginia.edu]} & \hyperlink{mod-TRAIN}{\tg{TRAIN}} \hyperlink{mod-MULTI}{\tg{MULTI}}\\
\hyperlink{course-UVA}{UVA} & 18 & \hyperlink{lec-UVA-18}{Topic 12: Personal Agents and Distilling} & \href{https://www.cs.virginia.edu/\%7Ermw7my/Courses/AgenticAISpring2026/Topic12Personal/personal.html}{[cs.virginia.edu]} & \hyperlink{mod-WEB}{\tg{WEB}} \hyperlink{mod-SAFE}{\tg{SAFE}} \hyperlink{mod-TRAIN}{\tg{TRAIN}}\\
\hyperlink{course-UW}{UW-Madison} & 13 & $\star$\,\hyperlink{lec-UW-13}{Optimizing agents I} & -- & \hyperlink{mod-ARCH}{\tg{ARCH}} \hyperlink{mod-TRAIN}{\tg{TRAIN}}\\
\hyperlink{course-UW}{UW-Madison} & 14 & $\star$\,\hyperlink{lec-UW-14}{Optimizing agents II} & -- & \hyperlink{mod-ARCH}{\tg{ARCH}} \hyperlink{mod-TRAIN}{\tg{TRAIN}}\\
\hyperlink{course-JHU}{JHU} & 11 & $\star$\,\hyperlink{lec-JHU-11}{Pre-training of Coding Language Models (1)} & \href{https://machine-programming.github.io/lectures/lecture-11.pdf}{[PDF]} & \hyperlink{mod-TRAIN}{\tg{TRAIN}} \hyperlink{mod-CODE}{\tg{CODE}}\\
\hyperlink{course-JHU}{JHU} & 12 & $\star$\,\hyperlink{lec-JHU-12}{Pre-training of Coding Language Models (2)} & \href{https://machine-programming.github.io/lectures/lecture-12.pdf}{[PDF]} & \hyperlink{mod-TRAIN}{\tg{TRAIN}} \hyperlink{mod-CODE}{\tg{CODE}}\\
\hyperlink{course-JHU}{JHU} & 13 & $\star$\,\hyperlink{lec-JHU-13}{Post-training of Coding Language Models} & -- & \hyperlink{mod-TRAIN}{\tg{TRAIN}} \hyperlink{mod-CODE}{\tg{CODE}}\\
\hyperlink{course-JHU}{JHU} & 14 & \hyperlink{lec-JHU-14}{Special Topics in Training Language Models} & -- & \hyperlink{mod-TRAIN}{\tg{TRAIN}}\\
\hyperlink{course-CU}{CU Boulder} & 11 & \hyperlink{lec-CU-11}{Reflexion \& Toolformer (students)} & -- & \hyperlink{mod-RSN}{\tg{RSN}} \hyperlink{mod-TOOL}{\tg{TOOL}} \hyperlink{mod-TRAIN}{\tg{TRAIN}}\\
\hyperlink{course-MANN}{Mannheim} & 2 & \hyperlink{lec-MANN-2}{Post-training and Efficient Adaptation} & \href{https://www.uni-mannheim.de/media/Einrichtungen/dws/Files_Teaching/Large_Language_Models_and_Agents/FSS2026/IE685_LA_02_PosttrainingAndEfficientAdaptation.pdf}{[PDF]} \href{https://www.uni-mannheim.de/media/Einrichtungen/dws/Files_Teaching/Large_Language_Models_and_Agents/FSS2026/IE685_LA_EX_LLM_Introduction.ipynb.zip}{[notebook zip]} & \hyperlink{mod-TRAIN}{\tg{TRAIN}} \hyperlink{mod-FND}{\tg{FND}}\\
\hyperlink{course-UCHI}{UChicago} & 5 & \hyperlink{lec-UCHI-5}{Sampling, Fine Tuning, Benchmarking, and Tools} & \href{https://github.com/KnowledgeLab/AI-Agents-for-Social-Science-and-Society-2026/blob/main/Tutorials-Homework_Notebooks/Week_5_2026.ipynb}{[GitHub]} & \hyperlink{mod-TRAIN}{\tg{TRAIN}} \hyperlink{mod-EVAL}{\tg{EVAL}} \hyperlink{mod-TOOL}{\tg{TOOL}}\\
\hyperlink{course-UCHI}{UChicago} & 7 & \hyperlink{lec-UCHI-7}{Reinforcement Learning to Optimize AI Agents and Institutions} & \href{https://github.com/KnowledgeLab/AI-Agents-for-Social-Science-and-Society-2026/blob/main/Tutorials-Homework_Notebooks/Week_7_RL.ipynb}{[GitHub]} & \hyperlink{mod-TRAIN}{\tg{TRAIN}} \hyperlink{mod-RSN}{\tg{RSN}} \hyperlink{mod-MULTI}{\tg{MULTI}}\\
\hyperlink{course-S329A}{Stanford 329A} & 4 & \hyperlink{lec-S329A-4}{Learning from Feedback with Tools/Code} & \href{https://www.youtube.com/watch?v=Lxh9RF5S-K0}{[video]} & \hyperlink{mod-TOOL}{\tg{TOOL}} \hyperlink{mod-CODE}{\tg{CODE}} \hyperlink{mod-TRAIN}{\tg{TRAIN}}\\
\hyperlink{course-S329A}{Stanford 329A} & 5 & \hyperlink{lec-S329A-5}{Multi-step Reasoning/Planning} & \href{https://www.youtube.com/watch?v=Ml_fp9XkB8Y}{[video]} & \hyperlink{mod-RSN}{\tg{RSN}} \hyperlink{mod-TRAIN}{\tg{TRAIN}}\\
\hyperlink{course-S329A}{Stanford 329A} & 6 & \hyperlink{lec-S329A-6}{Train Time Scaling/Scaling RL} & \href{https://www.youtube.com/watch?v=yVnmHSAy3ck}{[video]} & \hyperlink{mod-TRAIN}{\tg{TRAIN}} \hyperlink{mod-RSN}{\tg{RSN}}\\
\hyperlink{course-BRKAG}{Berkeley '25} & 2 & \hyperlink{lec-BRKAG-2}{LLM Agents Overview} & \href{https://rdi.berkeley.edu/agentic-ai/slides/lecture1.pdf}{[slides]} \href{https://www.youtube.com/watch?v=r1qZpYAmqmg}{[video]} & \hyperlink{mod-FND}{\tg{FND}} \hyperlink{mod-TRAIN}{\tg{TRAIN}}\\
\hyperlink{course-BRKAG}{Berkeley '25} & 4 & \hyperlink{lec-BRKAG-4}{Post-Training Verifiable Agents} & \href{https://rdi.berkeley.edu/agentic-ai/slides/AgenticAI_Jiantao\%20Jiao.pdf}{[slides]} \href{https://www.youtube.com/watch?v=3l0Zxus34es\&list=PLS01nW3RtgoqGkm4UeqNeZLccW-OGc1fJ\&index=3}{[video]} & \hyperlink{mod-TRAIN}{\tg{TRAIN}} \hyperlink{mod-EVAL}{\tg{EVAL}} \hyperlink{mod-CODE}{\tg{CODE}} \hyperlink{mod-WEB}{\tg{WEB}}\\
\hyperlink{course-BRKAG}{Berkeley '25} & 6 & \hyperlink{lec-BRKAG-6}{Some Challenges and Lessons from Training Agentic Models} & \href{https://rdi.berkeley.edu/agentic-ai/slides/weizhu.pdf}{[slides]} \href{https://www.youtube.com/watch?v=xNxrBHZPDvM}{[video]} & \hyperlink{mod-TRAIN}{\tg{TRAIN}}\\
\hyperlink{course-BRKAG}{Berkeley '25} & 12 & \hyperlink{lec-BRKAG-12}{Autonomous Agents: Embodiment, Interaction, and Learning} & \href{https://www.youtube.com/watch?v=iDhzzugMOLA}{[video]} & \hyperlink{mod-APP}{\tg{APP}} \hyperlink{mod-TRAIN}{\tg{TRAIN}}\\
\hyperlink{course-BRKADV}{Berkeley Adv.} & 2 & \hyperlink{lec-BRKADV-2}{Learning to reason with LLMs} & \href{https://rdi.berkeley.edu/adv-llm-agents/slides/Jason-Weston-Reasoning-Alignment-Berkeley-Talk.pdf}{[slides]} \href{https://www.youtube.com/live/_MNlLhU33H0}{[video]} & \hyperlink{mod-TRAIN}{\tg{TRAIN}} \hyperlink{mod-RSN}{\tg{RSN}}\\
\hyperlink{course-BRKADV}{Berkeley Adv.} & 4 & \hyperlink{lec-BRKADV-4}{Open Training Recipes for Reasoning in Language Models} & \href{https://rdi.berkeley.edu/adv-llm-agents/slides/OLMo-Tulu-Reasoning-Hanna.pdf}{[slides]} \href{https://www.youtube.com/live/cMiu3A7YBks}{[video]} & \hyperlink{mod-TRAIN}{\tg{TRAIN}} \hyperlink{mod-RSN}{\tg{RSN}}\\
\hyperlink{course-BRKADV}{Berkeley Adv.} & 8 & \hyperlink{lec-BRKADV-8}{AlphaProof: when reinforcement learning meets formal mathematics} & \href{https://rdi.berkeley.edu/adv-llm-agents/slides/alphaproof.pdf}{[slides]} \href{https://www.youtube.com/live/3gaEMscOMAU}{[video]} & \hyperlink{mod-RSN}{\tg{RSN}} \hyperlink{mod-TRAIN}{\tg{TRAIN}} \hyperlink{mod-APP}{\tg{APP}}\\
\hyperlink{course-S329T}{Stanford 329T} & 9 & \hyperlink{lec-S329T-9}{Deep dive into LLMs} & \href{https://web.stanford.edu/class/cs329t/slides/Lecture\%207\%20--\%20CS\%20329T\%20Fall\%202025.pdf}{[slides]} & \hyperlink{mod-FND}{\tg{FND}} \hyperlink{mod-TRAIN}{\tg{TRAIN}}\\
\hyperlink{course-S329T}{Stanford 329T} & 10 & \hyperlink{lec-S329T-10}{Data curation and RL to optimize agents (guest: Trung Vu, Bespoke Labs)} & \href{https://web.stanford.edu/class/cs329t/slides/Lecture\%208\%20-\%20\%5B10_23_2025\%5D\%20Reinforcement\%20Learning\%20\%20for\%20Agents.pdf}{[slides]} & \hyperlink{mod-TRAIN}{\tg{TRAIN}} \hyperlink{mod-TOOL}{\tg{TOOL}} \hyperlink{mod-EVAL}{\tg{EVAL}}\\
\hyperlink{course-S329T}{Stanford 329T} & 11 & \hyperlink{lec-S329T-11}{GEPA (guest: Lakshya Agrawal)} & \href{https://web.stanford.edu/class/cs329t/slides/Lecture\%209\%20-\%2020251028\%20GEPA_Talk.pdf}{[slides]} & \hyperlink{mod-TRAIN}{\tg{TRAIN}} \hyperlink{mod-RSN}{\tg{RSN}}\\
\bottomrule\end{longtable}}
\section{Key readings}
\begin{itemize}
\item Ouyang et al., \href{https://arxiv.org/abs/2203.02155}{\emph{InstructGPT}} (2022)
\item Shao et al., \href{https://arxiv.org/abs/2402.03300}{\emph{DeepSeekMath (GRPO)}} (2024)
\item DeepSeek-AI, \href{https://arxiv.org/abs/2501.12948}{\emph{DeepSeek-R1}} (2025)
\item Singh et al., \href{https://arxiv.org/abs/2312.06585}{\emph{ReST-EM}} (2023)
\item Hu et al., \href{https://arxiv.org/abs/2106.09685}{\emph{LoRA}} (2021); Dettmers et al., \href{https://arxiv.org/abs/2305.14314}{\emph{QLoRA}} (2023)
\item Khattab et al., \href{https://arxiv.org/abs/2310.03714}{\emph{DSPy}} (2023); Agrawal et al., \href{https://arxiv.org/abs/2507.19457}{\emph{GEPA}} (2025)
\end{itemize}
\section{Hands-on}
\begin{itemize}
\item \hyperlink{asg-CMU-3}{\textbf{CMU: A3: Agentic RL Training, "scottyrl"}} --- Fill in TODOs for REINFORCE/GRPO advantages, trajectory to masked token sequences, grouped rollouts, a sync trainer, and an async producer-consumer GRPO trainer. Tasks: Number Search (sync and async), ablations (group size 2, stale-4 async with and without importance weighting), and async GRPO on TextCraft-Synth. The final adapter is exported to Hugging Face. Uses Tinker 0.24.0 with a SkyRL Tinker server on Modal and vLLM/SGLang. \textit{Starter code:} \href{https://github.com/cmu-agents/assignment-3}{[GitHub]}
\item \hyperlink{asg-UVA-8}{\textbf{UVA: Topic 8: Text-to-SQL fine-tuning}} --- Fine-tune Llama on text-to-SQL with Tinker and score semantic SQL matches. \textit{Starter code:} \href{https://www.cs.virginia.edu/\%7Ermw7my/Courses/AgenticAISpring2026/Topic8FineTuning/finetuning.html}{[cs.virginia.edu]}
\item \hyperlink{asg-UCHI-5}{\textbf{UChicago: Week 5 notebook}} --- 3 of 4 modules: sampling; LoRA/QLoRA fine-tuning with varied ranks plus CKA comparison; designing an agent benchmark; a tool-using agent and an MCP server. \textit{Starter code:} \href{https://github.com/KnowledgeLab/AI-Agents-for-Social-Science-and-Society-2026/blob/main/Tutorials-Homework_Notebooks/Week_5_2026.ipynb}{[GitHub]}
\item \hyperlink{asg-UCHI-7}{\textbf{UChicago: Week 7 notebook}} --- Bandits (epsilon-greedy, annealing, UCB), RL for alignment, then RL for reasoning or multi-agent RL; memo pilot. \textit{Starter code:} \href{https://github.com/KnowledgeLab/AI-Agents-for-Social-Science-and-Society-2026/blob/main/Tutorials-Homework_Notebooks/Week_7_RL.ipynb}{[GitHub]}
\end{itemize}
\textbf{Related supplements:} \hyperref[sup-tinyagent]{Berkeley BAIR: TinyAgent, Function Calling at the Edge}.
\chapter{Module 12: Safety, Security and Oversight}\hypertarget{mod-SAFE}{}
{\small Tag: \tg{SAFE}}
\section{Learning objectives}
\begin{itemize}
\item Threat-model an agent: prompt injection, tool misuse, data exfiltration, supply chain.
\item Apply defenses: sandboxing, least privilege, capability separation, human approval, monitoring.
\item Evaluate security with benchmarks such as AgentDojo.
\end{itemize}
\section{Core concepts}
An agent with tools is a confused-deputy risk. Simon Willison's \textbf{lethal trifecta} (UW reading) names the dangerous combination: access to private data, exposure to untrusted content, and an exfiltration channel.
\begin{itemize}
\item \textbf{Attacks.} Indirect prompt injection through web pages, documents, tool outputs or MCP servers. Jailbreaks. Malicious or compromised tools and packages. CMU's safety lecture works through real incidents: the OpenClaw incident report, a database-deletion incident, CVE-2025-32711, and the Cline supply-chain compromise.
\item \textbf{Defenses.} Sandboxed execution and credential isolation; least-privilege tools and permission modes; human-in-the-loop approval for irreversible actions (Buffalo Lab~5); capability-based designs that separate control flow from untrusted data (CaMeL); guardrails; layered ``Swiss cheese'' safety.
\item \textbf{Oversight.} Tracing, monitoring and anomaly detection (CMU Safety~2). Static analysis of agent-written code (Semgrep, Stanford weeks on security).
\item \textbf{Alignment.} Sleeper agents, misalignment after narrow fine-tuning, representation engineering (UChicago week~9, UVA project ideas).
\end{itemize}
\section{Curated learning sequence}
\begin{enumerate}
\item \hyperlink{lec-CMU-13}{CMU~13 \emph{Safety 1: Agent Safety}} --- Agent safety through real incidents (slides).
\item \hyperlink{lec-UW-11}{UW-Madison~11 \emph{Security \& prompt injection I}} and \hyperlink{lec-UW-12}{UW-Madison~12 \emph{Security \& prompt injection II}} --- Indirect prompt injection and AgentDojo, then CaMeL and browser-agent security.
\item \hyperlink{lec-UB-10}{Buffalo~10 \emph{Security and privacy}} --- Prompt injection, tool misuse, sandboxing, privacy (then Lab~5).
\item \hyperlink{lec-STAN25-11}{Stanford '25~11 \emph{AI QA, SAST, DAST, and beyond}} --- AI QA, SAST/DAST and prompt-injection RCE case studies (slides).
\item \hyperlink{lec-CMU-16}{CMU~16 \emph{Safety 2: Observability and Monitoring}} --- Observability and monitoring.
\item \hyperlink{lec-NEU-16}{Northeastern~16 \emph{Model Context Protocol}} and \hyperlink{lec-NEU-17}{Northeastern~17 \emph{MCP in Agentic Systems}} --- MCP security and enterprise auth.
\item \hyperlink{lec-MANN-9}{Mannheim~9 \emph{LLM Security and Safety (+ project coaching)}} --- LLM security and safety (full-stack safety survey).
\end{enumerate}
\section{Lecture map (22 lectures)}
{\footnotesize\begin{longtable}{P{1.75cm} P{0.4cm} P{5.75cm} P{4.2cm} P{2.65cm}}
\toprule Course & \# & Title (links to catalogue row) & Materials & Tags\\ \midrule\endhead
\hyperlink{course-CMU}{CMU} & 13 & $\star$\,\hyperlink{lec-CMU-13}{Safety 1: Agent Safety} & \href{https://www.cmu-agents.com/slides/lecture-13-agent-safety.pdf}{[slides]} & \hyperlink{mod-SAFE}{\tg{SAFE}}\\
\hyperlink{course-CMU}{CMU} & 16 & $\star$\,\hyperlink{lec-CMU-16}{Safety 2: Observability and Monitoring} & -- & \hyperlink{mod-SAFE}{\tg{SAFE}} \hyperlink{mod-DEPLOY}{\tg{DEPLOY}}\\
\hyperlink{course-STAN26}{Stanford '26} & 13 & \hyperlink{lec-STAN26-13}{Security in AI codebases} & \href{https://themodernsoftware.dev/}{[course site]} & \hyperlink{mod-SAFE}{\tg{SAFE}}\\
\hyperlink{course-STAN25}{Stanford '25} & 11 & $\star$\,\hyperlink{lec-STAN25-11}{AI QA, SAST, DAST, and beyond} & \href{https://docs.google.com/presentation/d/1C05bCLasMDigBbkwdWbiz4WrXibzi6ua4hQQbTod_8c/edit?usp=sharing}{[slides]} & \hyperlink{mod-SAFE}{\tg{SAFE}} \hyperlink{mod-EVAL}{\tg{EVAL}} \hyperlink{mod-CODE}{\tg{CODE}}\\
\hyperlink{course-NEU}{Northeastern} & 16 & $\star$\,\hyperlink{lec-NEU-16}{Model Context Protocol} & -- & \hyperlink{mod-TOOL}{\tg{TOOL}} \hyperlink{mod-SAFE}{\tg{SAFE}}\\
\hyperlink{course-NEU}{Northeastern} & 17 & $\star$\,\hyperlink{lec-NEU-17}{MCP in Agentic Systems} & -- & \hyperlink{mod-TOOL}{\tg{TOOL}} \hyperlink{mod-SAFE}{\tg{SAFE}} \hyperlink{mod-DEPLOY}{\tg{DEPLOY}}\\
\hyperlink{course-UB}{Buffalo} & 8 & \hyperlink{lec-UB-8}{Human-agent interaction and memory} & -- & \hyperlink{mod-MEM}{\tg{MEM}} \hyperlink{mod-SAFE}{\tg{SAFE}}\\
\hyperlink{course-UB}{Buffalo} & 10 & $\star$\,\hyperlink{lec-UB-10}{Security and privacy} & -- & \hyperlink{mod-SAFE}{\tg{SAFE}}\\
\hyperlink{course-UVA}{UVA} & 18 & \hyperlink{lec-UVA-18}{Topic 12: Personal Agents and Distilling} & \href{https://www.cs.virginia.edu/\%7Ermw7my/Courses/AgenticAISpring2026/Topic12Personal/personal.html}{[cs.virginia.edu]} & \hyperlink{mod-WEB}{\tg{WEB}} \hyperlink{mod-SAFE}{\tg{SAFE}} \hyperlink{mod-TRAIN}{\tg{TRAIN}}\\
\hyperlink{course-UW}{UW-Madison} & 6 & \hyperlink{lec-UW-6}{Coding agents II} & -- & \hyperlink{mod-CODE}{\tg{CODE}} \hyperlink{mod-SAFE}{\tg{SAFE}} \hyperlink{mod-EVAL}{\tg{EVAL}}\\
\hyperlink{course-UW}{UW-Madison} & 11 & $\star$\,\hyperlink{lec-UW-11}{Security \& prompt injection I} & -- & \hyperlink{mod-SAFE}{\tg{SAFE}} \hyperlink{mod-EVAL}{\tg{EVAL}}\\
\hyperlink{course-UW}{UW-Madison} & 12 & $\star$\,\hyperlink{lec-UW-12}{Security \& prompt injection II} & -- & \hyperlink{mod-SAFE}{\tg{SAFE}}\\
\hyperlink{course-JHU}{JHU} & 17 & \hyperlink{lec-JHU-17}{PL Kaleidoscope: Security and Privacy} & -- & \hyperlink{mod-SAFE}{\tg{SAFE}} \hyperlink{mod-APP}{\tg{APP}}\\
\hyperlink{course-CU}{CU Boulder} & 23 & \hyperlink{lec-CU-23}{SafePro \& ToolDreamer (students)} & -- & \hyperlink{mod-SAFE}{\tg{SAFE}} \hyperlink{mod-TOOL}{\tg{TOOL}}\\
\hyperlink{course-CU}{CU Boulder} & 24 & \hyperlink{lec-CU-24}{Swiss Cheese safety \& MCP design (students)} & -- & \hyperlink{mod-SAFE}{\tg{SAFE}} \hyperlink{mod-TOOL}{\tg{TOOL}}\\
\hyperlink{course-MANN}{Mannheim} & 9 & $\star$\,\hyperlink{lec-MANN-9}{LLM Security and Safety (+ project coaching)} & \href{https://www.uni-mannheim.de/media/Einrichtungen/dws/Files_Teaching/Large_Language_Models_and_Agents/FSS2026/IE685_LA_08_SecurityAndSafety.pdf}{[PDF]} & \hyperlink{mod-SAFE}{\tg{SAFE}}\\
\hyperlink{course-UCHI}{UChicago} & 6 & \hyperlink{lec-UCHI-6}{Auto-Encoders, Interpretability, and Steering Agents} & \href{https://github.com/KnowledgeLab/AI-Agents-for-Social-Science-and-Society-2026/blob/main/Tutorials-Homework_Notebooks/week6_2026.ipynb}{[GitHub]} & \hyperlink{mod-SAFE}{\tg{SAFE}} \hyperlink{mod-FND}{\tg{FND}}\\
\hyperlink{course-UCHI}{UChicago} & 9 & \hyperlink{lec-UCHI-9}{Alignment, Ethics, Safety, and Novelty} & \href{https://github.com/KnowledgeLab/AI-Agents-for-Social-Science-and-Society-2026/blob/main/Tutorials-Homework_Notebooks/Week_9.ipynb}{[GitHub]} & \hyperlink{mod-SAFE}{\tg{SAFE}} \hyperlink{mod-APP}{\tg{APP}}\\
\hyperlink{course-BRKAG}{Berkeley '25} & 13 & \hyperlink{lec-BRKAG-13}{Agentic AI Safety \& Security} & \href{https://youtube.com/live/CvZDJxd4LKM?feature=share}{[video]} & \hyperlink{mod-SAFE}{\tg{SAFE}}\\
\hyperlink{course-BRKADV}{Berkeley Adv.} & 5 & \hyperlink{lec-BRKADV-5}{Coding Agents and AI for Vulnerability Detection} & \href{https://rdi.berkeley.edu/adv-llm-agents/slides/Code\%20Agents\%20and\%20AI\%20for\%20Vulnerability\%20Detection.pdf}{[slides]} \href{https://www.youtube.com/live/JCk6qJtaCSU}{[video]} & \hyperlink{mod-CODE}{\tg{CODE}} \hyperlink{mod-TOOL}{\tg{TOOL}} \hyperlink{mod-SAFE}{\tg{SAFE}}\\
\hyperlink{course-BRKADV}{Berkeley Adv.} & 12 & \hyperlink{lec-BRKADV-12}{Towards building safe and secure agentic AI} & \href{https://rdi.berkeley.edu/adv-llm-agents/slides/dawn-agentic-ai.pdf}{[slides]} \href{https://www.youtube.com/live/ti6yPE2VPZc}{[video]} & \hyperlink{mod-SAFE}{\tg{SAFE}} \hyperlink{mod-DEPLOY}{\tg{DEPLOY}}\\
\hyperlink{course-S329T}{Stanford 329T} & 12 & \hyperlink{lec-S329T-12}{Security for agents: adversarial attacks and guardrails (guests: Matt Fredrikson, Andy Zou)} & \href{https://web.stanford.edu/class/cs329t/slides/lecture\%2010\%20-\%2020251030\%20-\%20security\%20for\%20agents\%20-\%20andy\%20zou.pdf}{[slides]} & \hyperlink{mod-SAFE}{\tg{SAFE}} \hyperlink{mod-EVAL}{\tg{EVAL}}\\
\bottomrule\end{longtable}}
\section{Key readings}
\begin{itemize}
\item Greshake et al., \href{https://arxiv.org/abs/2302.12173}{\emph{Indirect Prompt Injection}} (2023)
\item Debenedetti et al., \href{https://arxiv.org/abs/2406.13352}{\emph{AgentDojo}} (2024)
\item Debenedetti et al., \href{https://arxiv.org/abs/2503.18813}{\emph{CaMeL: Defeating Prompt Injections by Design}} (2025)
\item Willison, \href{https://simonwillison.net/2025/Jun/16/the-lethal-trifecta/}{\emph{The lethal trifecta}}
\item Wang et al., \href{https://arxiv.org/abs/2504.15585}{\emph{Full Stack Safety survey}} (2025)
\item Hubinger et al., \href{https://arxiv.org/abs/2401.05566}{\emph{Sleeper Agents}} (2024)
\end{itemize}
\section{Hands-on}
\begin{itemize}
\item \hyperlink{asg-UB-5}{\textbf{Buffalo: Lab 5: Human-in-the-Loop and Safety Boundaries}} --- Add human approval, structured outputs, guardrails or permission boundaries. Analyze where autonomy should be restricted. \textit{(spec only)}
\item \hyperlink{asg-STAN25-6}{\textbf{Stanford '25: Week 6 --- Scan and Fix Vulnerabilities with Semgrep}} --- Run Semgrep on the app, triage findings, fix at least 3 issues (SQLi, unsafe DOM, CORS, etc.). \textit{Starter code:} \href{https://github.com/mihail911/modern-software-dev-assignments/tree/fall2025/week6}{[GitHub]}
\item \hyperlink{asg-UCHI-9}{\textbf{UChicago: Week 9 notebook}} --- Any 4 of 7 modules: safety judge, Constitutional AI, deception detection, moral geometry, automated red-teaming, AGrail guardrails; 150-300-word reflections (not AI-delegable). \textit{Starter code:} \href{https://github.com/KnowledgeLab/AI-Agents-for-Social-Science-and-Society-2026/blob/main/Tutorials-Homework_Notebooks/Week_9.ipynb}{[GitHub]}
\item \hyperlink{asg-STAN26-2}{\textbf{Stanford '26: Week 2 --- Build an MCP Server for a Third-Party API}} --- FastMCP server over an OAuth 2.0 API with at least 3 composing tools, a write tool with dry-run, structured errors, readOnlyHint, full auth-code flow with refresh, and a protocol-level test. \textit{Starter code:} \href{https://github.com/mihail911/modern-software-dev-assignments/blob/master/week2/assignment.md}{[GitHub]}
\end{itemize}
\textbf{Related supplements:} \hyperref[sup-openclaw]{OpenClaw documentation}, \hyperref[sup-agentsec]{MIT 6.566 guest lecture: AI Agent Security (Anish Athalye, April 2026)}, \hyperref[sup-anthauto]{Anthropic: Building effective agent automations (Martin and Avilla, October 2026)}.
\chapter{Module 13: Deployment, Efficiency and Operations}\hypertarget{mod-DEPLOY}{}
{\small Tag: \tg{DEPLOY}}
\section{Learning objectives}
\begin{itemize}
\item Deploy an agent as a service and operate it (logging, tracing, cost tracking, fallbacks).
\item Reduce cost and latency: caching, model routing, small or quantized models, code execution.
\item Understand long-running and background agent architectures.
\end{itemize}
\section{Core concepts}
\begin{itemize}
\item \textbf{Shipping.} Northeastern deploys agents to Hugging Face Spaces and Render. Buffalo adds CI and observability. Stanford '25 covers SRE and incident response with agents.
\item \textbf{Cost and latency.} Prompt and KV caching (keep a stable prefix), model routing and gateways (Stanford '26 week~10), fallbacks and spend policies (Northeastern's optimization session), shorter contexts through tool retrieval or code execution, and small or quantized models on device (TinyAgent, UVA Topic~1).
\item \textbf{Long-running agents.} Background and asynchronous agents (Stanford '26 week~8) and always-on gateways with channels, cron jobs and hooks (OpenClaw) need durable state, idempotent tools and resumability (checkpointing in \hyperlink{mod-ARCH}{Module~2 (The Agent Loop and Agent Architectures)}).
\item \textbf{Observability.} Traces of every step are the raw material for evaluation (\hyperlink{mod-EVAL}{Module~10 (Evaluating Agents)}), debugging and safety monitoring (\hyperlink{mod-SAFE}{Module~12 (Safety, Security and Oversight)}).
\end{itemize}
\section{Curated learning sequence}
\begin{enumerate}
\item \hyperlink{lec-NEU-5}{Northeastern~5 \emph{Tool-Calling}} and \hyperlink{lec-NEU-6}{Northeastern~6 \emph{Short-term \& Long-term Memory}} --- Deploy to Hugging Face Spaces, then Render.
\item \hyperlink{lec-NEU-19}{Northeastern~19 \emph{Agentic AI Optimization}} --- DSPy plus cost tracking, caching, fallbacks and spend policies.
\item \hyperlink{lec-STAN26-15}{Stanford '26~15 \emph{Background agents: launching tasks asynchronously}} and \hyperlink{lec-STAN26-19}{Stanford '26~19 \emph{Coding agents in big teams}} --- Background agents, then coding agents in big teams (gateways, routing, cost).
\item \hyperlink{lec-STAN25-17}{Stanford '25~17 \emph{Incident response and DevOps}} --- Incident response and DevOps with agents (slides).
\item \hyperlink{lec-CMU-12}{CMU~12 \emph{Training 4: RL Systems}} --- RL systems: the systems side of agent training.
\item Supplements: \hyperref[sup-openclaw]{OpenClaw}, \hyperref[sup-tinyagent]{TinyAgent} --- Always-on gateway architecture; edge deployment.
\item Supplements: \hyperref[sup-anthauto]{Anthropic agent automations}, \hyperref[sup-notion]{Notion Data Scout} --- Durable state and guardrails for a scheduled agent; operating a production data agent.
\end{enumerate}
\section{Lecture map (19 lectures)}
{\footnotesize\begin{longtable}{P{1.75cm} P{0.4cm} P{5.75cm} P{4.2cm} P{2.65cm}}
\toprule Course & \# & Title (links to catalogue row) & Materials & Tags\\ \midrule\endhead
\hyperlink{course-CMU}{CMU} & 12 & $\star$\,\hyperlink{lec-CMU-12}{Training 4: RL Systems} & \href{https://www.cmu-agents.com/slides/lecture-12-rl-systems.pdf}{[slides]} & \hyperlink{mod-TRAIN}{\tg{TRAIN}} \hyperlink{mod-DEPLOY}{\tg{DEPLOY}}\\
\hyperlink{course-CMU}{CMU} & 16 & \hyperlink{lec-CMU-16}{Safety 2: Observability and Monitoring} & -- & \hyperlink{mod-SAFE}{\tg{SAFE}} \hyperlink{mod-DEPLOY}{\tg{DEPLOY}}\\
\hyperlink{course-STAN26}{Stanford '26} & 15 & $\star$\,\hyperlink{lec-STAN26-15}{Background agents: launching tasks asynchronously} & \href{https://themodernsoftware.dev/}{[course site]} & \hyperlink{mod-DEPLOY}{\tg{DEPLOY}} \hyperlink{mod-ARCH}{\tg{ARCH}}\\
\hyperlink{course-STAN26}{Stanford '26} & 19 & $\star$\,\hyperlink{lec-STAN26-19}{Coding agents in big teams (MCP portals, LLM gateways, cost optimization, model routing)} & \href{https://themodernsoftware.dev/}{[course site]} & \hyperlink{mod-DEPLOY}{\tg{DEPLOY}}\\
\hyperlink{course-STAN26}{Stanford '26} & 20 & \hyperlink{lec-STAN26-20}{The Software Factory: self-running, self-improving software systems} & \href{https://themodernsoftware.dev/}{[course site]} & \hyperlink{mod-ARCH}{\tg{ARCH}} \hyperlink{mod-DEPLOY}{\tg{DEPLOY}}\\
\hyperlink{course-STAN25}{Stanford '25} & 9 & \hyperlink{lec-STAN25-9}{How to build a breakthrough AI developer product} & \href{https://docs.google.com/presentation/d/1Djd4eBLBbRkma8rFnJAWMT0ptct_UGB8hipmoqFVkxQ/edit?usp=sharing}{[slides]} & \hyperlink{mod-CODE}{\tg{CODE}} \hyperlink{mod-DEPLOY}{\tg{DEPLOY}}\\
\hyperlink{course-STAN25}{Stanford '25} & 17 & $\star$\,\hyperlink{lec-STAN25-17}{Incident response and DevOps} & \href{https://docs.google.com/presentation/d/1Mfe-auWAsg9URCujneKnHr0AbO8O-_U4QXBVOlO4qp0/edit?usp=sharing}{[slides]} & \hyperlink{mod-DEPLOY}{\tg{DEPLOY}} \hyperlink{mod-MULTI}{\tg{MULTI}}\\
\hyperlink{course-NEU}{Northeastern} & 2 & \hyperlink{lec-NEU-2}{LLM Foundations (I)} & \href{https://agenticllms.github.io/week2a.html}{[agenticllms.github.io]} \href{https://www.youtube.com/watch?v=wjZofJX0v4M}{[video]} \href{https://www.youtube.com/watch?v=kCc8FmEb1nY}{[video]} & \hyperlink{mod-FND}{\tg{FND}} \hyperlink{mod-DEPLOY}{\tg{DEPLOY}}\\
\hyperlink{course-NEU}{Northeastern} & 5 & $\star$\,\hyperlink{lec-NEU-5}{Tool-Calling} & \href{https://agenticllms.github.io/week3b.html}{[agenticllms.github.io]} \href{https://agenticllms.github.io/week3b-slides.html}{[agenticllms.github.io]} \href{https://github.com/agenticllms/lecture-files/tree/main/W03b}{[GitHub]} & \hyperlink{mod-TOOL}{\tg{TOOL}} \hyperlink{mod-RSN}{\tg{RSN}} \hyperlink{mod-DEPLOY}{\tg{DEPLOY}}\\
\hyperlink{course-NEU}{Northeastern} & 6 & $\star$\,\hyperlink{lec-NEU-6}{Short-term \& Long-term Memory} & \href{https://agenticllms.github.io/week4a.html}{[agenticllms.github.io]} \href{https://agenticllms.github.io/week4a-slides.html}{[agenticllms.github.io]} \href{https://github.com/agenticllms/lecture-files/tree/main/W04a}{[GitHub]} & \hyperlink{mod-MEM}{\tg{MEM}} \hyperlink{mod-ARCH}{\tg{ARCH}} \hyperlink{mod-DEPLOY}{\tg{DEPLOY}}\\
\hyperlink{course-NEU}{Northeastern} & 17 & \hyperlink{lec-NEU-17}{MCP in Agentic Systems} & -- & \hyperlink{mod-TOOL}{\tg{TOOL}} \hyperlink{mod-SAFE}{\tg{SAFE}} \hyperlink{mod-DEPLOY}{\tg{DEPLOY}}\\
\hyperlink{course-NEU}{Northeastern} & 19 & $\star$\,\hyperlink{lec-NEU-19}{Agentic AI Optimization} & -- & \hyperlink{mod-DEPLOY}{\tg{DEPLOY}} \hyperlink{mod-FND}{\tg{FND}}\\
\hyperlink{course-UB}{Buffalo} & 12 & \hyperlink{lec-UB-12}{Coding agents, testing, CI, observability} & -- & \hyperlink{mod-CODE}{\tg{CODE}} \hyperlink{mod-DEPLOY}{\tg{DEPLOY}}\\
\hyperlink{course-UB}{Buffalo} & 13 & \hyperlink{lec-UB-13}{Applied agents and responsible deployment} & -- & \hyperlink{mod-APP}{\tg{APP}} \hyperlink{mod-DEPLOY}{\tg{DEPLOY}}\\
\hyperlink{course-UVA}{UVA} & 1 & \hyperlink{lec-UVA-1}{Topic 1: Running an LLM} & \href{https://www.cs.virginia.edu/\%7Ermw7my/Courses/AgenticAISpring2026/Topic1Running/Running_an_LLM.html}{[cs.virginia.edu]} & \hyperlink{mod-FND}{\tg{FND}} \hyperlink{mod-EVAL}{\tg{EVAL}} \hyperlink{mod-DEPLOY}{\tg{DEPLOY}}\\
\hyperlink{course-CU}{CU Boulder} & 21 & \hyperlink{lec-CU-21}{MemSkill \& More with Less (students)} & -- & \hyperlink{mod-MEM}{\tg{MEM}} \hyperlink{mod-DEPLOY}{\tg{DEPLOY}}\\
\hyperlink{course-BRKAG}{Berkeley '25} & 3 & \hyperlink{lec-BRKAG-3}{Evolution of system designs from an AI engineer perspective} & \href{https://rdi.berkeley.edu/agentic-ai/slides/Yangqing.pdf}{[slides]} \href{https://www.youtube.com/watch?v=xqRAS6rAouo\&list=PLS01nW3RtgoqGkm4UeqNeZLccW-OGc1fJ\&index=2}{[video]} & \hyperlink{mod-ARCH}{\tg{ARCH}} \hyperlink{mod-DEPLOY}{\tg{DEPLOY}}\\
\hyperlink{course-BRKAG}{Berkeley '25} & 10 & \hyperlink{lec-BRKAG-10}{Practical Lessons from Deploying Real-World AI Agents} & \href{https://rdi.berkeley.edu/agentic-ai/slides/Practical_Lessons_AI_Agents.pdf}{[slides]} \href{https://www.youtube.com/watch?v=sfJM4LaiYsM}{[video]} & \hyperlink{mod-DEPLOY}{\tg{DEPLOY}} \hyperlink{mod-EVAL}{\tg{EVAL}} \hyperlink{mod-APP}{\tg{APP}}\\
\hyperlink{course-BRKADV}{Berkeley Adv.} & 12 & \hyperlink{lec-BRKADV-12}{Towards building safe and secure agentic AI} & \href{https://rdi.berkeley.edu/adv-llm-agents/slides/dawn-agentic-ai.pdf}{[slides]} \href{https://www.youtube.com/live/ti6yPE2VPZc}{[video]} & \hyperlink{mod-SAFE}{\tg{SAFE}} \hyperlink{mod-DEPLOY}{\tg{DEPLOY}}\\
\bottomrule\end{longtable}}
\section{Key readings}
\begin{itemize}
\item Google, \href{https://sre.google/sre-book/introduction/}{\emph{SRE book, introduction}}
\item Anthropic, \href{https://www.anthropic.com/engineering/code-execution-with-mcp}{\emph{Code execution with MCP}} (2025)
\item Erdogan et al., \href{https://arxiv.org/abs/2409.00608}{\emph{TinyAgent}} (2024)
\item Gao \& Peng, \href{https://arxiv.org/abs/2510.16786}{\emph{More with Less}} (2025)
\item Anthropic, \href{https://claude.dev/blog/building-effective-agent-automations/}{\emph{Building effective agent automations}} (2026)
\item Notion, \href{https://www.notion.com/blog/how-we-built-a-personal-data-scientist-for-every-notion-employee}{\emph{How we built a personal data scientist for every Notion employee}} (2026)
\end{itemize}
\section{Hands-on}
\begin{itemize}
\item \hyperlink{asg-NEU-1}{\textbf{Northeastern: DADA1: The Trip Planner Agent}} --- A conversational day-trip planner with tools: get\_\allowbreak{}weather (Open-Meteo), find\_\allowbreak{}attractions/restaurants/transit (Tavily), find\_\allowbreak{}events (Ticketmaster). It has memory-based plan revision, a structured Itinerary/DayPlan output, 8 test prompts, and a Gradio app deployed on Render. Graded on tool routing, honesty at data limits, sources and memory. Optional extra credit: a find\_\allowbreak{}performer\_\allowbreak{}events chained tool. \textit{Starter code:} \href{https://github.com/agenticllms/lecture-files/tree/main/DADA1}{[GitHub]}
\item \hyperlink{asg-UVA-1}{\textbf{UVA: Topic 1 portfolio: Running an LLM}} --- MMLU timing/accuracy study across models, quantization and hardware; chat agent with bounded context and restartability; optional MT-Bench. \textit{Starter code:} \href{https://www.cs.virginia.edu/\%7Ermw7my/Courses/AgenticAISpring2026/Topic1Running/Running_an_LLM.html}{[cs.virginia.edu]}
\end{itemize}
\textbf{Related supplements:} \hyperref[sup-codeexec]{Anthropic Engineering: Code execution with MCP}, \hyperref[sup-openclaw]{OpenClaw documentation}, \hyperref[sup-tinyagent]{Berkeley BAIR: TinyAgent, Function Calling at the Edge}, \hyperref[sup-anthauto]{Anthropic: Building effective agent automations (Martin and Avilla, October 2026)}, \hyperref[sup-notion]{Notion: How we built a personal data scientist for every Notion employee (David Thomas, October 2026)}, \hyperref[sup-colab]{Google Colab}.
\chapter{Module 14: Applications: Deep Research, Science and Society}\hypertarget{mod-APP}{}
{\small Tag: \tg{APP}}
\section{Learning objectives}
\begin{itemize}
\item Understand how deep-research and scientific-discovery agents are built and evaluated.
\item Scope an application agent for a bounded domain with a measurable outcome.
\end{itemize}
\section{Core concepts}
\begin{itemize}
\item \textbf{Deep research agents} combine search, browsing, reading and synthesis over many steps (STORM, PaperQA, OpenScholar, Tongyi DeepResearch, DR Tulu). They are evaluated on BrowseComp, DeepResearch Bench and LiveResearchBench.
\item \textbf{Agents for science} range from AI Scientist and autonomous chemistry to the AI co-scientist, sequential falsification, HEPTAPOD for particle physics, and Lean-based mathematics (UVA's guest lecture).
\item \textbf{Society.} Simulated subjects and digital doubles, mixed human--AI experiments, persuasion, the future of work (CMU), and human--agent interaction design (CMU Interaction~2).
\item \textbf{Scoping your own.} Mannheim's nine project briefs and Buffalo's ``useful bounded agent'' spec are good templates: a narrow domain, explicit tools, a benchmark and a risk assessment.
\end{itemize}
\section{Curated learning sequence}
\begin{enumerate}
\item \hyperlink{lec-CMU-10}{CMU~10 \emph{Domains 3: Deep Research Agents}} --- Deep research agents (slides).
\item \hyperlink{lec-UW-7}{UW-Madison~7 \emph{Deep research agents I}}--\hyperlink{lec-UW-10}{UW-Madison~10} --- Deep research, then agents for science (four sessions).
\item \hyperlink{lec-UVA-19}{UVA~19 \emph{Topic 13: Agentic AI for Science}} and \hyperlink{lec-UVA-20}{UVA~20 \emph{Topic 14: Agentic AI for Mathematics and Verification}} --- Agentic AI for science, then mathematics and verification (Lean).
\item \hyperlink{lec-CMU-17}{CMU~17 \emph{Agents and the Future of Work}} and \hyperlink{lec-CMU-19}{CMU~19 \emph{Interaction 2: Human-Agent Interaction}} --- The future of work; human--agent interaction.
\item \hyperlink{course-UCHI}{UChicago AI Agents for Social Science and Society} --- The whole course, for agents as research instruments.
\end{enumerate}
\section{Lecture map (33 lectures)}
{\footnotesize\begin{longtable}{P{1.75cm} P{0.4cm} P{5.75cm} P{4.2cm} P{2.65cm}}
\toprule Course & \# & Title (links to catalogue row) & Materials & Tags\\ \midrule\endhead
\hyperlink{course-CMU}{CMU} & 10 & $\star$\,\hyperlink{lec-CMU-10}{Domains 3: Deep Research Agents} & \href{https://www.cmu-agents.com/slides/lecture-10-deep-research-agents.pdf}{[slides]} & \hyperlink{mod-APP}{\tg{APP}} \hyperlink{mod-MEM}{\tg{MEM}}\\
\hyperlink{course-CMU}{CMU} & 17 & $\star$\,\hyperlink{lec-CMU-17}{Agents and the Future of Work} & -- & \hyperlink{mod-APP}{\tg{APP}}\\
\hyperlink{course-CMU}{CMU} & 19 & $\star$\,\hyperlink{lec-CMU-19}{Interaction 2: Human-Agent Interaction} & -- & \hyperlink{mod-APP}{\tg{APP}} \hyperlink{mod-EVAL}{\tg{EVAL}}\\
\hyperlink{course-STAN25}{Stanford '25} & 15 & \hyperlink{lec-STAN25-15}{End-to-end apps from a single prompt} & \href{https://docs.google.com/presentation/d/1GrVLsfMFIXMiGjIW9D7EJIyLYh_-3ReHHNd_vRfZUoo/edit?usp=sharing}{[slides]} & \hyperlink{mod-CODE}{\tg{CODE}} \hyperlink{mod-APP}{\tg{APP}}\\
\hyperlink{course-UB}{Buffalo} & 13 & \hyperlink{lec-UB-13}{Applied agents and responsible deployment} & -- & \hyperlink{mod-APP}{\tg{APP}} \hyperlink{mod-DEPLOY}{\tg{DEPLOY}}\\
\hyperlink{course-UVA}{UVA} & 10 & \hyperlink{lec-UVA-10}{Topic 6: Vision-Language Models} & \href{https://www.cs.virginia.edu/\%7Ermw7my/Courses/AgenticAISpring2026/Topic6VLM/vlm.html}{[cs.virginia.edu]} \href{https://github.com/HenryKautz/clip_test}{[GitHub]} & \hyperlink{mod-WEB}{\tg{WEB}} \hyperlink{mod-APP}{\tg{APP}} \hyperlink{mod-FND}{\tg{FND}}\\
\hyperlink{course-UVA}{UVA} & 19 & $\star$\,\hyperlink{lec-UVA-19}{Topic 13: Agentic AI for Science} & \href{https://www.cs.virginia.edu/\%7Ermw7my/Courses/AgenticAISpring2026/Topic13Science/science.html}{[cs.virginia.edu]} \href{https://github.com/tonymenzo/heptapod}{[GitHub]} & \hyperlink{mod-APP}{\tg{APP}} \hyperlink{mod-MULTI}{\tg{MULTI}}\\
\hyperlink{course-UVA}{UVA} & 20 & $\star$\,\hyperlink{lec-UVA-20}{Topic 14: Agentic AI for Mathematics and Verification (guest: Prof. Kevin Sullivan)} & \href{https://www.cs.virginia.edu/\%7Ermw7my/Courses/AgenticAISpring2026/Topic14Math/math.html}{[cs.virginia.edu]} & \hyperlink{mod-RSN}{\tg{RSN}} \hyperlink{mod-APP}{\tg{APP}}\\
\hyperlink{course-UW}{UW-Madison} & 7 & $\star$\,\hyperlink{lec-UW-7}{Deep research agents I} & -- & \hyperlink{mod-APP}{\tg{APP}} \hyperlink{mod-MULTI}{\tg{MULTI}}\\
\hyperlink{course-UW}{UW-Madison} & 8 & $\star$\,\hyperlink{lec-UW-8}{Deep research agents II} & -- & \hyperlink{mod-APP}{\tg{APP}} \hyperlink{mod-MEM}{\tg{MEM}} \hyperlink{mod-EVAL}{\tg{EVAL}}\\
\hyperlink{course-UW}{UW-Madison} & 9 & $\star$\,\hyperlink{lec-UW-9}{Agents for science I} & -- & \hyperlink{mod-APP}{\tg{APP}}\\
\hyperlink{course-UW}{UW-Madison} & 10 & $\star$\,\hyperlink{lec-UW-10}{Agents for science II} & -- & \hyperlink{mod-APP}{\tg{APP}} \hyperlink{mod-MULTI}{\tg{MULTI}}\\
\hyperlink{course-JHU}{JHU} & 15 & \hyperlink{lec-JHU-15}{PL Kaleidoscope: Formal Verification and Theorem Proving} & -- & \hyperlink{mod-APP}{\tg{APP}} \hyperlink{mod-CODE}{\tg{CODE}}\\
\hyperlink{course-JHU}{JHU} & 16 & \hyperlink{lec-JHU-16}{PL Kaleidoscope: Differentiable Programming and Machine Learning} & -- & \hyperlink{mod-APP}{\tg{APP}}\\
\hyperlink{course-JHU}{JHU} & 17 & \hyperlink{lec-JHU-17}{PL Kaleidoscope: Security and Privacy} & -- & \hyperlink{mod-SAFE}{\tg{SAFE}} \hyperlink{mod-APP}{\tg{APP}}\\
\hyperlink{course-JHU}{JHU} & 18 & \hyperlink{lec-JHU-18}{PL Kaleidoscope: Hardware and High-Performance Computing} & -- & \hyperlink{mod-APP}{\tg{APP}}\\
\hyperlink{course-JHU}{JHU} & 19 & \hyperlink{lec-JHU-19}{PL Kaleidoscope: Design, Interactive, and Visual Programming} & -- & \hyperlink{mod-APP}{\tg{APP}}\\
\hyperlink{course-JHU}{JHU} & 26 & \hyperlink{lec-JHU-26}{Special Topics on Neurosymbolic Programming (1)} & -- & \hyperlink{mod-RSN}{\tg{RSN}} \hyperlink{mod-APP}{\tg{APP}}\\
\hyperlink{course-JHU}{JHU} & 27 & \hyperlink{lec-JHU-27}{Special Topics on Neurosymbolic Programming (2)} & -- & \hyperlink{mod-RSN}{\tg{RSN}} \hyperlink{mod-APP}{\tg{APP}}\\
\hyperlink{course-UCHI}{UChicago} & 1 & \hyperlink{lec-UCHI-1}{Deep Learning and Social Agents} & \href{https://github.com/KnowledgeLab/AI-Agents-for-Social-Science-and-Society-2026/blob/main/Tutorials-Homework_Notebooks/Week_1_Intro_NNs.ipynb}{[GitHub]} & \hyperlink{mod-FND}{\tg{FND}} \hyperlink{mod-APP}{\tg{APP}}\\
\hyperlink{course-UCHI}{UChicago} & 3 & \hyperlink{lec-UCHI-3}{LLMs Prompted for Multi-Agent Simulation} & \href{https://github.com/KnowledgeLab/AI-Agents-for-Social-Science-and-Society-2026/blob/main/Tutorials-Homework_Notebooks/Week_3.ipynb}{[GitHub]} & \hyperlink{mod-MULTI}{\tg{MULTI}} \hyperlink{mod-MEM}{\tg{MEM}} \hyperlink{mod-APP}{\tg{APP}}\\
\hyperlink{course-UCHI}{UChicago} & 4 & \hyperlink{lec-UCHI-4}{Experimental Designs with AI Agent and Human Interactions} & \href{https://github.com/KnowledgeLab/AI-Agents-for-Social-Science-and-Society-2026/blob/main/Tutorials-Homework_Notebooks/week_4_2026.ipynb}{[GitHub]} & \hyperlink{mod-APP}{\tg{APP}} \hyperlink{mod-EVAL}{\tg{EVAL}}\\
\hyperlink{course-UCHI}{UChicago} & 8 & \hyperlink{lec-UCHI-8}{Multi-Modal and Embodied Agents} & \href{https://github.com/KnowledgeLab/AI-Agents-for-Social-Science-and-Society-2026/blob/main/Tutorials-Homework_Notebooks/Week_8_Multimodal.ipynb}{[GitHub]} & \hyperlink{mod-WEB}{\tg{WEB}} \hyperlink{mod-RSN}{\tg{RSN}} \hyperlink{mod-APP}{\tg{APP}}\\
\hyperlink{course-UCHI}{UChicago} & 9 & \hyperlink{lec-UCHI-9}{Alignment, Ethics, Safety, and Novelty} & \href{https://github.com/KnowledgeLab/AI-Agents-for-Social-Science-and-Society-2026/blob/main/Tutorials-Homework_Notebooks/Week_9.ipynb}{[GitHub]} & \hyperlink{mod-SAFE}{\tg{SAFE}} \hyperlink{mod-APP}{\tg{APP}}\\
\hyperlink{course-S329A}{Stanford 329A} & 7 & \hyperlink{lec-S329A-7}{Open-Ended Evolution of Self-Improving Agents} & -- & \hyperlink{mod-ARCH}{\tg{ARCH}} \hyperlink{mod-CODE}{\tg{CODE}} \hyperlink{mod-APP}{\tg{APP}}\\
\hyperlink{course-BRKAG}{Berkeley '25} & 9 & \hyperlink{lec-BRKAG-9}{AI Agents to Automate Scientific Discoveries} & \href{https://www.youtube.com/watch?v=yqPIsTTdUkc}{[video]} & \hyperlink{mod-APP}{\tg{APP}} \hyperlink{mod-MULTI}{\tg{MULTI}}\\
\hyperlink{course-BRKAG}{Berkeley '25} & 10 & \hyperlink{lec-BRKAG-10}{Practical Lessons from Deploying Real-World AI Agents} & \href{https://rdi.berkeley.edu/agentic-ai/slides/Practical_Lessons_AI_Agents.pdf}{[slides]} \href{https://www.youtube.com/watch?v=sfJM4LaiYsM}{[video]} & \hyperlink{mod-DEPLOY}{\tg{DEPLOY}} \hyperlink{mod-EVAL}{\tg{EVAL}} \hyperlink{mod-APP}{\tg{APP}}\\
\hyperlink{course-BRKAG}{Berkeley '25} & 12 & \hyperlink{lec-BRKAG-12}{Autonomous Agents: Embodiment, Interaction, and Learning} & \href{https://www.youtube.com/watch?v=iDhzzugMOLA}{[video]} & \hyperlink{mod-APP}{\tg{APP}} \hyperlink{mod-TRAIN}{\tg{TRAIN}}\\
\hyperlink{course-BRKADV}{Berkeley Adv.} & 8 & \hyperlink{lec-BRKADV-8}{AlphaProof: when reinforcement learning meets formal mathematics} & \href{https://rdi.berkeley.edu/adv-llm-agents/slides/alphaproof.pdf}{[slides]} \href{https://www.youtube.com/live/3gaEMscOMAU}{[video]} & \hyperlink{mod-RSN}{\tg{RSN}} \hyperlink{mod-TRAIN}{\tg{TRAIN}} \hyperlink{mod-APP}{\tg{APP}}\\
\hyperlink{course-BRKADV}{Berkeley Adv.} & 9 & \hyperlink{lec-BRKADV-9}{Language models for autoformalization and theorem proving} & \href{https://rdi.berkeley.edu/adv-llm-agents/slides/mathverification.pdf}{[slides]} \href{https://www.youtube.com/live/cLhWEyMQ4mQ}{[video]} & \hyperlink{mod-RSN}{\tg{RSN}} \hyperlink{mod-APP}{\tg{APP}} \hyperlink{mod-TOOL}{\tg{TOOL}}\\
\hyperlink{course-BRKADV}{Berkeley Adv.} & 10 & \hyperlink{lec-BRKADV-10}{Bridging Informal and Formal Mathematical Reasoning (Berkeley page: "Advanced topics in theorem proving")} & \href{https://rdi.berkeley.edu/adv-llm-agents/slides/welleck2025_berkeley_bridging.pdf}{[slides]} \href{https://www.youtube.com/live/Gy5Nm17l9oo}{[video]} & \hyperlink{mod-RSN}{\tg{RSN}} \hyperlink{mod-APP}{\tg{APP}}\\
\hyperlink{course-BRKADV}{Berkeley Adv.} & 11 & \hyperlink{lec-BRKADV-11}{Abstraction and Discovery with LLM Agents} & \href{https://rdi.berkeley.edu/adv-llm-agents/slides/swarat.pdf}{[slides]} \href{https://www.youtube.com/live/IHc0TEMrEdY}{[video]} & \hyperlink{mod-RSN}{\tg{RSN}} \hyperlink{mod-ARCH}{\tg{ARCH}} \hyperlink{mod-APP}{\tg{APP}}\\
\hyperlink{course-S329T}{Stanford 329T} & 18 & \hyperlink{lec-S329T-18}{Final Project Fair} & -- & \tg{ORG} \hyperlink{mod-APP}{\tg{APP}}\\
\bottomrule\end{longtable}}
\section{Key readings}
\begin{itemize}
\item Shao et al., \href{https://arxiv.org/abs/2402.14207}{\emph{STORM}} (2024)
\item L\'ala et al., \href{https://arxiv.org/abs/2312.07559}{\emph{PaperQA}} (2023)
\item Asai et al., \href{https://arxiv.org/abs/2411.14199}{\emph{OpenScholar}} (2024)
\item Lu et al., \href{https://arxiv.org/abs/2408.06292}{\emph{The AI Scientist}} (2024)
\item Gottweis et al., \href{https://arxiv.org/abs/2502.18864}{\emph{Towards an AI Co-Scientist}} (2025)
\item Park et al., \href{https://arxiv.org/abs/2304.03442}{\emph{Generative Agents}} (2023)
\end{itemize}
\section{Hands-on}
\begin{itemize}
\item \hyperlink{asg-MANN-7}{\textbf{Mannheim: Project topic 2: QA Agents}} --- RAG (multi-)agent QA over web search that explains its evidence. \textit{Resources:} \href{https://openai.com/index/browsecomp/}{[openai.com]} \href{https://www.aicrowd.com/challenges/meta-comprehensive-rag-benchmark-kdd-cup-2024}{[aicrowd.com]}
\item \hyperlink{asg-MANN-13}{\textbf{Mannheim: Project topic 8: Wikipedia Article Generation}} --- RAG/agent system to generate Wikipedia articles, evaluated against real articles with rubrics. \textit{(spec only)}
\item \hyperlink{asg-UB-10}{\textbf{Buffalo: Final Team Project: Useful Bounded Agent}} --- Working agent for a bounded task. Deliverables: problem statement, environment/tool inventory, architecture diagram, repo, eval dataset/script, results and error analysis, security/privacy/ethics risk assessment, demo, final report. \textit{(spec only)}
\item \hyperlink{asg-UCHI-11}{\textbf{UChicago: Final project}} --- Agent project for social insight, solo or teams of up to 4: at least 3 data types and 3 model types from 3 weeks, qualitatively validated; Substack/Medium post (3,600 words and 13 visuals solo; 5,000 and 17 for groups) with an open repo, plus an Ignite talk. \textit{(spec only)}
\end{itemize}
\textbf{Related supplements:} \hyperref[sup-notion]{Notion: How we built a personal data scientist for every Notion employee (David Thomas, October 2026)}.
\part{Course Catalogue}\hypertarget{catalog}{}
\chapter{CMU: 11-768 AI Agents}\hypertarget{course-CMU}{}
{\small\begin{tabularx}{\linewidth}{@{}>{\bfseries}l X@{}}\toprule
Institution & Carnegie Mellon University (LTI)\\
Term & Fall 2026\\
Instructors & Graham Neubig, Daniel Fried\\
Website & \href{https://www.cmu-agents.com/\#/schedule}{https://www.cmu-agents.com/\#/schedule}\\
Level & Graduate; assumes you have trained neural LMs\\
Archetype & Systems + training (build, evaluate, RL-train)\\
Public materials & Slides: L1--13; video: L1--9; code: A1--A3 public\\
Best for & The backbone course: harness engineering, evals, SFT/RL, safety, search.\\
\bottomrule\end{tabularx}}

The most complete public agents course in the collection, and the one this guide is built around.
It goes from capabilities (tool use, long-context management, skills and memory, planning) to domains (coding, computer use, deep research), then to training (SFT, RL basics, advanced RL, RL systems), safety and observability, frameworks (OpenHands, LangGraph), interaction and search.
The readings are taken from production systems: harness source code, provider KV-cache pricing pages and the MCP spec.
The three assignments form one pipeline: build a harness (A1), evaluate an agent with an LLM judge (A2), then RL-train an agentic policy (A3).
Slides for lectures 1--13 and recordings for lectures 1--9 are public on the schedule page.
\section{Lectures}
{\scriptsize\setlength{\tabcolsep}{2.5pt}\begin{longtable}{P{0.45cm} P{1.2cm} P{2.5cm} P{3.6cm} P{3.6cm} P{3.3cm} P{1.5cm}}
\toprule \# & Date & Title & Key topics & Readings & Materials & Tags\\ \midrule\endhead
\hypertarget{lec-CMU-1}{}1 & Tue Aug 25 & \textbf{Course Overview: What Is an Agent?} & Module: Intro \& Capabilities & Toolformer; ReAct; Mini-SWE-Agent & \href{https://www.cmu-agents.com/slides/lecture-01-agents.pdf}{slides} \href{https://www.cmu-agents.com/\#/schedule}{video} & \tg{ORG} \hyperlink{mod-ARCH}{\tg{ARCH}}\\ \midrule
\hypertarget{lec-CMU-2}{}2 & Thu Aug 27 & \textbf{Capabilities 1: Tool Use} & Module: Intro \& Capabilities & What Are Tools Anyway?; CodeAct; Toolformer; XGrammar (refs: OpenAI Chat Completions, HF chat templates, BFCL V4) & \href{https://www.cmu-agents.com/slides/lecture-02-tool-use.pdf}{slides} \href{https://www.cmu-agents.com/\#/schedule}{video} & \hyperlink{mod-TOOL}{\tg{TOOL}}\\ \midrule
\hypertarget{lec-CMU-3}{}3 & Tue Sep 1 & \textbf{Capabilities 2: Context Management for Long-Context Agents} & Module: Intro \& Capabilities & (refs: DistServe, Needle in a Haystack, Longformer) & \href{https://www.cmu-agents.com/slides/lecture-03-long-context.pdf}{slides} \href{https://www.cmu-agents.com/\#/schedule}{video} & \hyperlink{mod-MEM}{\tg{MEM}}\\ \midrule
\hypertarget{lec-CMU-4}{}4 & Thu Sep 3 & \textbf{Capabilities 3: Skills and Memory} & Module: Intro \& Capabilities & OpenHands Agent Skills; SkillsBench; MemGPT; AWM; ASI; ReasoningBank (refs: Mem0, Reflexion, ExpeL) & \href{https://www.cmu-agents.com/slides/lecture-04-memory-and-skills.pdf}{slides} \href{https://www.cmu-agents.com/\#/schedule}{video} & \hyperlink{mod-MEM}{\tg{MEM}}\\ \midrule
\hypertarget{lec-CMU-5}{}5 & Tue Sep 8 & \textbf{Capabilities 4: Planning, Task Decomposition, Multi-Agent Coordination} & Module: Intro \& Capabilities & Cursor Plan Mode; Least-to-Most; Decomposed Prompting; Code as Policies; SayCan; Plan-and-Act; Thinking vs. Doing; Calibrate-Then-Act; RAO; MACU & \href{https://www.cmu-agents.com/slides/lecture-05-planning.pdf}{slides} \href{https://www.cmu-agents.com/\#/schedule}{video} & \hyperlink{mod-RSN}{\tg{RSN}} \hyperlink{mod-MULTI}{\tg{MULTI}}\\ \midrule
\hypertarget{lec-CMU-6}{}6 & Thu Sep 10 & \textbf{Domains 1: Coding Agents} & Module: Domains & (refs: OLMo, StarCoder, Code Llama) & \href{https://www.cmu-agents.com/slides/lecture-06-coding-agents.pdf}{slides} \href{https://www.cmu-agents.com/\#/schedule}{video} & \hyperlink{mod-CODE}{\tg{CODE}}\\ \midrule
\hypertarget{lec-CMU-7}{}7 & Tue Sep 15 & \textbf{Computer Use Agents} & Module: Domains; lecturer: JY Koh & WebArena; VisualWebArena; OSWorld; OpenCUA; CUA-Gym; MACU & \href{https://www.cmu-agents.com/slides/lecture-07-computer-use-agents.pdf}{slides} \href{https://www.cmu-agents.com/\#/schedule}{video} & \hyperlink{mod-WEB}{\tg{WEB}}\\ \midrule
\hypertarget{lec-CMU-8}{}8 & Thu Sep 17 & \textbf{Training 1: SFT} & Module: Training; lecturer: Yueqi Song & (refs: DeepSeek-R1, Kimi, Nemotron) & \href{https://www.cmu-agents.com/slides/lecture-08-sft.pdf}{slides} \href{https://www.cmu-agents.com/\#/schedule}{video} & \hyperlink{mod-TRAIN}{\tg{TRAIN}}\\ \midrule
\hypertarget{lec-CMU-9}{}9 & Tue Sep 22 & \textbf{Training 2: RL Basics} & Module: Training & ReST; ReST-EM; DeepSeekMath; DrGRPO & \href{https://www.cmu-agents.com/slides/lecture-09-rl-basics.pdf}{slides} \href{https://www.cmu-agents.com/\#/schedule}{video} & \hyperlink{mod-TRAIN}{\tg{TRAIN}}\\ \midrule
\hypertarget{lec-CMU-10}{}10 & Thu Sep 24 & \textbf{Domains 3: Deep Research Agents} & Module: Domains; lecturer: Akari Asai & OpenScholar; BrowseComp; DeepResearch Bench; Tongyi DeepResearch; DR Tulu; AgentIR & \href{https://www.cmu-agents.com/slides/lecture-10-deep-research-agents.pdf}{slides} & \hyperlink{mod-APP}{\tg{APP}} \hyperlink{mod-MEM}{\tg{MEM}}\\ \midrule
\hypertarget{lec-CMU-11}{}11 & Tue Sep 29 & \textbf{Training 3: Advanced RL Algorithms} & Module: Training & (refs: InstructGPT, GAE, Let's Verify Step by Step, AReaL) & \href{https://www.cmu-agents.com/slides/lecture-11-rl-advanced.pdf}{slides} & \hyperlink{mod-TRAIN}{\tg{TRAIN}}\\ \midrule
\hypertarget{lec-CMU-12}{}12 & Thu Oct 1 & \textbf{Training 4: RL Systems} & Module: Training; lecturer: Apurva Gandhi & -- & \href{https://www.cmu-agents.com/slides/lecture-12-rl-systems.pdf}{slides} & \hyperlink{mod-TRAIN}{\tg{TRAIN}} \hyperlink{mod-DEPLOY}{\tg{DEPLOY}}\\ \midrule
\hypertarget{lec-CMU-13}{}13 & Tue Oct 6 & \textbf{Safety 1: Agent Safety} & Module: Agent Safety & (refs: OpenClaw incident report, CVE-2025-32711, Cline supply-chain postmortem) & \href{https://www.cmu-agents.com/slides/lecture-13-agent-safety.pdf}{slides} & \hyperlink{mod-SAFE}{\tg{SAFE}}\\ \midrule
\hypertarget{lec-CMU-14}{}14 & Thu Oct 8 & \textbf{Frameworks 1: OpenHands} & Module: Agent Frameworks & -- & -- & \hyperlink{mod-ARCH}{\tg{ARCH}} \hyperlink{mod-CODE}{\tg{CODE}}\\ \midrule
\hypertarget{lec-CMU-15}{}15 & Tue Oct 20 & \textbf{Frameworks 2: LangGraph} & Module: Agent Frameworks & -- & -- & \hyperlink{mod-ARCH}{\tg{ARCH}}\\ \midrule
\hypertarget{lec-CMU-16}{}16 & Thu Oct 22 & \textbf{Safety 2: Observability and Monitoring} & Module: Frameworks \& Safety; lecturer: Eric Wallace & -- & -- & \hyperlink{mod-SAFE}{\tg{SAFE}} \hyperlink{mod-DEPLOY}{\tg{DEPLOY}}\\ \midrule
\hypertarget{lec-CMU-17}{}17 & Tue Oct 27 & \textbf{Agents and the Future of Work} & Module: Interaction \& Projects; lecturer: Zora Wang & -- & -- & \hyperlink{mod-APP}{\tg{APP}}\\ \midrule
\hypertarget{lec-CMU-18}{}18 & Thu Oct 29 & \textbf{Interaction 1: Multi-Agent Interaction} & Module: Interaction \& Projects; lecturer: Saujas Vaduguru & -- & -- & \hyperlink{mod-MULTI}{\tg{MULTI}}\\ \midrule
\hypertarget{lec-CMU-19}{}19 & Tue Nov 10 & \textbf{Interaction 2: Human-Agent Interaction} & Module: Interaction \& Projects; lecturer: Valerie Chen & -- & -- & \hyperlink{mod-APP}{\tg{APP}} \hyperlink{mod-EVAL}{\tg{EVAL}}\\ \midrule
\hypertarget{lec-CMU-20}{}20 & Thu Nov 12 & \textbf{Search 1: Reranking and Critic Models} & Module: Search \& Inference & -- & -- & \hyperlink{mod-RSN}{\tg{RSN}} \hyperlink{mod-EVAL}{\tg{EVAL}}\\ \midrule
\hypertarget{lec-CMU-21}{}21 & Tue Nov 17 & \textbf{Search 2: Tree Search} & Module: Search \& Inference; lecturer: JY Koh & -- & -- & \hyperlink{mod-RSN}{\tg{RSN}}\\ \midrule
\hypertarget{lec-CMU-22}{}22 & Thu Nov 19 & \textbf{Guest Lecture} & Module: Advanced Topics; lecturer: Karthik Narasimhan & -- & -- & --\\ \midrule
\hypertarget{lec-CMU-23}{}23 & Tue Nov 24 & \textbf{Guest Lecture} & Module: Advanced Topics; lecturer: Sasha Rush & -- & -- & --\\ \midrule
\end{longtable}}
\section{Assignments and projects}
{\scriptsize\setlength{\tabcolsep}{2.5pt}\begin{longtable}{P{2.8cm} P{5.4cm} P{3.0cm} P{3.5cm} P{1.5cm}}
\toprule Name & What you build & Skills & Starter code / resources & Tags\\ \midrule\endhead
\hypertarget{asg-CMU-1}{}\textbf{A1: Build an Agent Harness (due Sep 14; 10\%)} & A domain-agnostic ReAct \texttt{Agent} base class: prompt construction, tool-call loop, step limit, model-generated context compaction. On top of it, a \texttt{CodeAgent} (terminal/message tools, skill loader with progressive disclosure via \texttt{invoke\_\allowbreak{}skill}) that fixes a chess-app bug and solves SWE-bench \texttt{django\_\allowbreak{}\_\allowbreak{}django-15368} with and without compaction. Then a \texttt{ChessAgent} (play\_\allowbreak{}move, stateless simulate\_\allowbreak{}move, sandboxed run\_\allowbreak{}python, chess skill) plus an observation A/B experiment. Uses Modal sandboxes, an OpenAI-compatible LLM (default DeepSeek-V4-Flash), uv, and pytest. 100-pt rubric. & ReAct loops, tool dispatch, sandboxing, skills, context compaction, token accounting & \href{https://github.com/cmu-agents/assignment-1}{[GitHub]} (handout: \href{https://github.com/cmu-agents/assignment-1/blob/main/ASSIGNMENT.md}{[GitHub]}); chess app: \href{https://github.com/cmu-agents/chess-app}{[GitHub]} & \hyperlink{mod-ARCH}{\tg{ARCH}} \hyperlink{mod-TOOL}{\tg{TOOL}} \hyperlink{mod-CODE}{\tg{CODE}}\\ \midrule
\hypertarget{asg-CMU-2}{}\textbf{A2: Evaluating a Data-Visualization Agent (due Sep 24; 15\%)} & A \texttt{validate(run)} LLM-as-judge validator that flags four error families (execution\_\allowbreak{}failure, wrong\_\allowbreak{}data, wrong\_\allowbreak{}chart, hard\_\allowbreak{}to\_\allowbreak{}read), scored by per-family and macro MCC against human labels. Students author 5+ new tasks, run 3 fixed agents (15+ runs), human-label the runs, and package the tasks as Harbor environments with deterministic tests, a reference solution, and two mutants. The judge is fixed to Qwen3-VL-30B on Modal. Grading: report 40\%, private-set MCC 30\%, quiz 20\%, code 10\%. & Eval design, LLM-as-judge, human labeling, verifier robustness, benchmark packaging & \href{https://github.com/cmu-agents/assignment-2}{[GitHub]} & \hyperlink{mod-EVAL}{\tg{EVAL}}\\ \midrule
\hypertarget{asg-CMU-3}{}\textbf{A3: Agentic RL Training, "scottyrl" (due Oct 22; 15\%)} & Fill in TODOs for REINFORCE/GRPO advantages, trajectory to masked token sequences, grouped rollouts, a sync trainer, and an async producer-consumer GRPO trainer. Tasks: Number Search (sync and async), ablations (group size 2, stale-4 async with and without importance weighting), and async GRPO on TextCraft-Synth. The final adapter is exported to Hugging Face. Uses Tinker 0.24.0 with a SkyRL Tinker server on Modal and vLLM/SGLang. & Policy-gradient RL for multi-turn agents, async RL systems, ablation methodology & \href{https://github.com/cmu-agents/assignment-3}{[GitHub]} & \hyperlink{mod-TRAIN}{\tg{TRAIN}}\\ \midrule
\hypertarget{asg-CMU-4}{}\textbf{Team research project (50\%)} & Open-ended agent research: proposal, check-in, poster (Dec 1/3), and final report (30\%) & Research design, experimentation, writing & -- & \tg{ORG}\\ \midrule
\hypertarget{asg-CMU-5}{}\textbf{Lecture highlights (10\%)} & Short written highlights of lectures, done without AI and submitted within 24 h (best 20 of 22 count) & Synthesis & -- & \tg{ORG}\\ \midrule
\end{longtable}}
\chapter{Stanford '26: CS146S The Modern Software Developer}\hypertarget{course-STAN26}{}
{\small\begin{tabularx}{\linewidth}{@{}>{\bfseries}l X@{}}\toprule
Institution & Stanford University\\
Term & Fall 2026 (current)\\
Instructors & Mihail Eric\\
Website & \href{https://themodernsoftware.dev/}{https://themodernsoftware.dev/}\\
Level & Undergrad/grad; CS111/CS161-level programming\\
Archetype & Coding-agent practice (practitioner)\\
Public materials & Slides: Weeks 1--3 so far; video: none; code: Weeks 1--3 assignments\\
Best for & How production coding agents work and how to set up a repo and a team around them.\\
\bottomrule\end{tabularx}}

A practitioner course about building software with coding agents, with an industry guest in most weeks.
The Fall 2026 version focuses on agents. Week~1 builds ``Claude Code in 200 lines'' and dissects production system prompts. Later weeks cover context engineering and MCP, agent skills, CLAUDE.md/AGENTS.md and hooks, agent-ready codebases, agentic code review, security, background agents, and the ``software factory''.
The assignments are short and practical: dissect a real Claude Code trace, build an OAuth MCP server, and author an Agent Skill.
\section{Lectures}
{\scriptsize\setlength{\tabcolsep}{2.5pt}\begin{longtable}{P{0.45cm} P{1.2cm} P{2.5cm} P{3.6cm} P{3.6cm} P{3.3cm} P{1.5cm}}
\toprule \# & Date & Title & Key topics & Readings & Materials & Tags\\ \midrule\endhead
\hypertarget{lec-STAN26-1}{}1 & Tue 9/22 & \textbf{Course intro + build Claude Code in 200 lines} & Week 1: The Internals of Coding Agents & Building a Coding Agent; Lessons from Building Claude Code; Dive into Claude Code & \href{https://themodernsoftware.dev/}{course site} & \hyperlink{mod-CODE}{\tg{CODE}} \hyperlink{mod-ARCH}{\tg{ARCH}}\\ \midrule
\hypertarget{lec-STAN26-2}{}2 & Thu 9/24 & \textbf{How SOTA coding agents are designed: deep dive into system prompts} & Week 1: The Internals of Coding Agents & -- & \href{https://themodernsoftware.dev/}{course site} & \hyperlink{mod-CODE}{\tg{CODE}} \hyperlink{mod-FND}{\tg{FND}}\\ \midrule
\hypertarget{lec-STAN26-3}{}3 & Tue 9/29 & \textbf{Advanced prompting + agentic dev frameworks (RePPIT, spec-driven development)} & Week 2: Advanced Context Engineering & -- & \href{https://themodernsoftware.dev/}{course site} & \hyperlink{mod-FND}{\tg{FND}} \hyperlink{mod-MEM}{\tg{MEM}}\\ \midrule
\hypertarget{lec-STAN26-4}{}4 & Thu 10/1 & \textbf{Full introduction to MCP and tool-calling (theory, setup, advanced tool design)} & Week 2: Advanced Context Engineering & -- & \href{https://themodernsoftware.dev/}{course site} & \hyperlink{mod-TOOL}{\tg{TOOL}}\\ \midrule
\hypertarget{lec-STAN26-5}{}5 & Tue 10/6 & \textbf{All about agent skills (incl. web skills)} & Week 3: Agent Skills and CLI & pstack; Matt Pocock's Skills; browse.sh & \href{https://themodernsoftware.dev/}{course site} & \hyperlink{mod-MEM}{\tg{MEM}} \hyperlink{mod-TOOL}{\tg{TOOL}}\\ \midrule
\hypertarget{lec-STAN26-6}{}6 & Thu 10/8 & \textbf{Guest: Lee Robinson} & Week 3: Agent Skills and CLI & -- & \href{https://themodernsoftware.dev/}{course site} & \hyperlink{mod-APP}{\tg{APP}}\\ \midrule
\hypertarget{lec-STAN26-7}{}7 & Tue 10/13 & \textbf{CLAUDE.md, AGENTS.md, hooks; subagent patterns (planner/implementer/reviewer)} & Week 4: Customizing Your Agent and Repository & -- & \href{https://themodernsoftware.dev/}{course site} & \hyperlink{mod-CODE}{\tg{CODE}} \hyperlink{mod-MULTI}{\tg{MULTI}}\\ \midrule
\hypertarget{lec-STAN26-8}{}8 & Thu 10/15 & \textbf{Guest: Boris Cherny (Claude Code, Anthropic)} & Week 4: Customizing Your Agent and Repository & -- & \href{https://themodernsoftware.dev/}{course site} & \hyperlink{mod-CODE}{\tg{CODE}}\\ \midrule
\hypertarget{lec-STAN26-9}{}9 & Tue 10/20 & \textbf{Agent readiness: structure, docs, tests, checks} & Week 5: Agent-Ready Codebases & -- & \href{https://themodernsoftware.dev/}{course site} & \hyperlink{mod-CODE}{\tg{CODE}} \hyperlink{mod-EVAL}{\tg{EVAL}}\\ \midrule
\hypertarget{lec-STAN26-10}{}10 & Thu 10/22 & \textbf{Guest: Eno Reyes (Factory)} & Week 5: Agent-Ready Codebases & -- & \href{https://themodernsoftware.dev/}{course site} & \hyperlink{mod-CODE}{\tg{CODE}}\\ \midrule
\hypertarget{lec-STAN26-11}{}11 & Tue 10/27 & \textbf{Agentic code review: best practices and architectures} & Week 6: Agentic Code Review & -- & \href{https://themodernsoftware.dev/}{course site} & \hyperlink{mod-CODE}{\tg{CODE}} \hyperlink{mod-EVAL}{\tg{EVAL}}\\ \midrule
\hypertarget{lec-STAN26-12}{}12 & Thu 10/29 & \textbf{Guest: Silas Alberti (Cognition)} & Week 6: Agentic Code Review & -- & \href{https://themodernsoftware.dev/}{course site} & \hyperlink{mod-CODE}{\tg{CODE}}\\ \midrule
\hypertarget{lec-STAN26-13}{}13 & Tue 11/3 & \textbf{Security in AI codebases} & Week 7: Security & -- & \href{https://themodernsoftware.dev/}{course site} & \hyperlink{mod-SAFE}{\tg{SAFE}}\\ \midrule
\hypertarget{lec-STAN26-14}{}14 & Thu 11/5 & \textbf{Guest: Isaac Evans (Semgrep)} & Week 7: Security & -- & \href{https://themodernsoftware.dev/}{course site} & \hyperlink{mod-SAFE}{\tg{SAFE}}\\ \midrule
\hypertarget{lec-STAN26-15}{}15 & Tue 11/10 & \textbf{Background agents: launching tasks asynchronously} & Week 8: Background Agents & -- & \href{https://themodernsoftware.dev/}{course site} & \hyperlink{mod-DEPLOY}{\tg{DEPLOY}} \hyperlink{mod-ARCH}{\tg{ARCH}}\\ \midrule
\hypertarget{lec-STAN26-16}{}16 & Thu 11/12 & \textbf{Guest: Rajesh Bhatia (Cloudflare)} & Week 8: Background Agents & -- & \href{https://themodernsoftware.dev/}{course site} & \hyperlink{mod-DEPLOY}{\tg{DEPLOY}}\\ \midrule
\hypertarget{lec-STAN26-17}{}17 & Tue 11/17 & \textbf{Guest: Elad Gil} & Week 9: Building an AI-Native Team & -- & \href{https://themodernsoftware.dev/}{course site} & \hyperlink{mod-APP}{\tg{APP}}\\ \midrule
\hypertarget{lec-STAN26-18}{}18 & Thu 11/19 & \textbf{Guest: Amjad Masad (Replit)} & Week 9: Building an AI-Native Team & -- & \href{https://themodernsoftware.dev/}{course site} & \hyperlink{mod-APP}{\tg{APP}}\\ \midrule
\hypertarget{lec-STAN26-19}{}19 & Tue 12/1 & \textbf{Coding agents in big teams (MCP portals, LLM gateways, cost optimization, model routing)} & Week 10: Software Factory + Future & -- & \href{https://themodernsoftware.dev/}{course site} & \hyperlink{mod-DEPLOY}{\tg{DEPLOY}}\\ \midrule
\hypertarget{lec-STAN26-20}{}20 & Thu 12/3 & \textbf{The Software Factory: self-running, self-improving software systems} & Week 10: Software Factory + Future & -- & \href{https://themodernsoftware.dev/}{course site} & \hyperlink{mod-ARCH}{\tg{ARCH}} \hyperlink{mod-DEPLOY}{\tg{DEPLOY}}\\ \midrule
\end{longtable}}
\section{Assignments and projects}
{\scriptsize\setlength{\tabcolsep}{2.5pt}\begin{longtable}{P{2.8cm} P{5.4cm} P{3.0cm} P{3.5cm} P{1.5cm}}
\toprule Name & What you build & Skills & Starter code / resources & Tags\\ \midrule\endhead
\hypertarget{asg-STAN26-1}{}\textbf{Week 1 --- Trace Dissection of a Real Claude Code Session} & Capture a Claude Code session through a mitmproxy reverse proxy, then annotate the system prompt, \texttt{<system-reminder>}s, tool inventory and schemas, error recovery, planning, subagents, and context management & Reverse-engineering a production coding agent, prompt and tool design analysis & \href{https://github.com/mihail911/modern-software-dev-assignments/blob/master/week1/assignment.md}{[GitHub]} & \hyperlink{mod-CODE}{\tg{CODE}} \hyperlink{mod-ARCH}{\tg{ARCH}} \hyperlink{mod-MEM}{\tg{MEM}}\\ \midrule
\hypertarget{asg-STAN26-2}{}\textbf{Week 2 --- Build an MCP Server for a Third-Party API} & FastMCP server over an OAuth 2.0 API with at least 3 composing tools, a write tool with dry-run, structured errors, readOnlyHint, full auth-code flow with refresh, and a protocol-level test & Agent-ergonomic tool design, OAuth, MCP & \href{https://github.com/mihail911/modern-software-dev-assignments/blob/master/week2/assignment.md}{[GitHub]} & \hyperlink{mod-TOOL}{\tg{TOOL}} \hyperlink{mod-SAFE}{\tg{SAFE}}\\ \midrule
\hypertarget{asg-STAN26-3}{}\textbf{Week 3 --- Author a Skill} & A SKILL.md (Agent Skills standard) for a repeatable workflow on a real OSS repo (htmx, dayjs, yt-dlp, etc.), with trigger tests including a near-miss and an end-to-end run & Skill authoring, trigger description design, eval of triggering & \href{https://github.com/mihail911/modern-software-dev-assignments/blob/master/week3/assignment.md}{[GitHub]} & \hyperlink{mod-CODE}{\tg{CODE}} \hyperlink{mod-MEM}{\tg{MEM}} \hyperlink{mod-EVAL}{\tg{EVAL}}\\ \midrule
\end{longtable}}
\chapter{Stanford '25: CS146S The Modern Software Developer (archive)}\hypertarget{course-STAN25}{}
{\small\begin{tabularx}{\linewidth}{@{}>{\bfseries}l X@{}}\toprule
Institution & Stanford University\\
Term & Fall 2025\\
Instructors & Mihail Eric\\
Website & \href{https://themodernsoftware.dev/fall2025}{https://themodernsoftware.dev/fall2025}\\
Level & Undergrad/grad\\
Archetype & Coding-agent practice (practitioner)\\
Public materials & Slides: Most lectures (Google Slides); video: none; code: 8 weekly assignments (branch fall2025)\\
Best for & Slide decks on prompting, building a coding agent, MCP servers, AI testing/security, code review, SRE.\\
\bottomrule\end{tabularx}}

The previous offering of CS146S. Its value here is that nearly every lecture has a public slide deck, and the eight assignments (prompting playground, AI-IDE workflow, custom MCP server, Claude Code automations, Warp, Semgrep, Graphite review, multi-stack build) live on the \texttt{fall2025} branch of the assignments repository.
The lecture dates come from an authorized translation and may be off by a day, so treat them as approximate.
\section{Lectures}
{\scriptsize\setlength{\tabcolsep}{2.5pt}\begin{longtable}{P{0.45cm} P{1.2cm} P{2.5cm} P{3.6cm} P{3.6cm} P{3.3cm} P{1.5cm}}
\toprule \# & Date & Title & Key topics & Readings & Materials & Tags\\ \midrule\endhead
\hypertarget{lec-STAN25-1}{}1 & 2025-09-22 & \textbf{Introduction and how an LLM is made} & Course logistics; what LLMs are & Karpathy "Deep Dive into LLMs" \href{https://www.youtube.com/watch?v=7xTGNNLPyMI}{[video]} ; "How OpenAI Uses Codex" \href{https://cdn.openai.com/pdf/6a2631dc-783e-479b-b1a4-af0cfbd38630/how-openai-uses-codex.pdf}{[PDF]} & Slides \href{https://docs.google.com/presentation/d/1zT2Ofy88cajLTLkd7TcuSM4BCELvF9qQdHmlz33i4t0/edit?usp=sharing}{[slides]} & \tg{ORG} \hyperlink{mod-FND}{\tg{FND}}\\ \midrule
\hypertarget{lec-STAN25-2}{}2 & 2025-09-26 & \textbf{Power prompting for LLMs} & Effective prompting techniques & Google "What is prompt engineering" \href{https://cloud.google.com/discover/what-is-prompt-engineering}{[cloud.google.com]} ; Prompting Guide \href{https://www.promptingguide.ai/techniques}{[promptingguide.ai]} ; "AI Prompt Engineering: A Deep Dive" \href{https://www.youtube.com/watch?v=T9aRN5JkmL8}{[video]} & Slides \href{https://docs.google.com/presentation/d/1MIhw8p6TLGdbQ9TcxhXSs5BaPf5d_h77QY70RHNfeGs/edit?usp=drive_link}{[slides]} & \hyperlink{mod-FND}{\tg{FND}} \hyperlink{mod-RSN}{\tg{RSN}}\\ \midrule
\hypertarget{lec-STAN25-3}{}3 & 2025-09-29 & \textbf{Building a coding agent from scratch} & Agent architecture and components; tool use / function calling & Week 2 readings (MCP, below) & Slides \href{https://docs.google.com/presentation/d/11CP26VhsjnZOmi9YFgLlonzdib9BLyAlgc4cEvC5Fps/edit?usp=sharing}{[slides]} ; completed exercise \href{https://drive.google.com/file/d/1YtpKFVG13DHyQ2i3HOtwyVJOV90nWeL2/view?usp=drive_link}{[drive.google.com]} & \hyperlink{mod-ARCH}{\tg{ARCH}} \hyperlink{mod-TOOL}{\tg{TOOL}} \hyperlink{mod-CODE}{\tg{CODE}}\\ \midrule
\hypertarget{lec-STAN25-4}{}4 & 2025-10-03 & \textbf{Building a custom MCP server} & Model Context Protocol & MCP intro \href{https://stytch.com/blog/model-context-protocol-introduction/}{[stytch.com]} ; MCP servers \href{https://github.com/modelcontextprotocol/servers}{[GitHub]} ; Cloudflare remote MCP auth \href{https://developers.cloudflare.com/agents/guides/remote-mcp-server/\#add-authentication}{[developers.cloudflare.com]} ; MCP TS SDK \href{https://github.com/modelcontextprotocol/typescript-sdk/tree/main?tab=readme-ov-file\#server}{[GitHub]} ; MCP Registry \href{https://blog.modelcontextprotocol.io/posts/2025-09-08-mcp-registry-preview/}{[blog.modelcontextprotocol.io]} ; "APIs don't make good MCP tools" \href{https://www.reillywood.com/blog/apis-dont-make-good-mcp-tools/}{[reillywood.com]} & Slides \href{https://docs.google.com/presentation/d/1zSC2ra77XOUrJeyS85houg1DU7z9hq5Y4ebagTch-5o/edit?usp=drive_link}{[slides]} ; completed exercise \href{https://drive.google.com/file/d/1J6lgZWcxPzpCpjujJSnW1aAkCYF6Yxv3/view?usp=drive_link}{[drive.google.com]} & \hyperlink{mod-TOOL}{\tg{TOOL}}\\ \midrule
\hypertarget{lec-STAN25-5}{}5 & 2025-10-06 & \textbf{From first prompt to optimal IDE setup} & Context management and code understanding; PRDs/specs for agents; IDE integrations & "Specs Are the New Source Code" \href{https://blog.ravi-mehta.com/p/specs-are-the-new-source-code}{[blog.ravi-mehta.com]} ; "How Long Contexts Fail" \href{https://www.dbreunig.com/2025/06/22/how-contexts-fail-and-how-to-fix-them.html}{[dbreunig.com]} ; Devin Agents 101 \href{https://devin.ai/agents101\#introduction}{[devin.ai]} ; HumanLayer ACE-FCA \href{https://github.com/humanlayer/advanced-context-engineering-for-coding-agents/blob/main/ace-fca.md}{[GitHub]} ; "How FAANG Vibe Codes" \href{https://x.com/rohanpaul_ai/status/1959414096589422619}{[x.com]} ; Anthropic "Writing effective tools for agents" \href{https://www.anthropic.com/engineering/writing-tools-for-agents}{[anthropic.com]} & Slides \href{https://docs.google.com/presentation/d/11pQNCde_mmRnImBat0Zymnp8TCS_cT_1up7zbcj6Sjg/edit?usp=sharing}{[slides]} ; design doc template \href{https://drive.google.com/file/d/1MZ0Qx68Vzw4x5x_XcV8XiPLp7fFDe1LJ/view?usp=drive_link}{[drive.google.com]} & \hyperlink{mod-MEM}{\tg{MEM}} \hyperlink{mod-CODE}{\tg{CODE}}\\ \midrule
\hypertarget{lec-STAN25-6}{}6 & 2025-10-10 & \textbf{Guest: Silas Alberti (Cognition)} & Coding agents in industry (Devin) & (week 3 readings) & Slides \href{https://docs.google.com/presentation/d/1i0pRttHf72lgz8C-n7DSegcLBgncYZe_ppU7dB9zhUA/edit?usp=sharing}{[slides]} & \hyperlink{mod-CODE}{\tg{CODE}}\\ \midrule
\hypertarget{lec-STAN25-7}{}7 & 2025-10-13 & \textbf{How to be an agent manager} & Managing levels of agent autonomy; human-agent collaboration patterns & "How Anthropic Uses Claude Code" \href{https://www-cdn.anthropic.com/58284b19e702b49db9302d5b6f135ad8871e7658.pdf}{[PDF]} ; Claude Code best practices \href{https://www.anthropic.com/engineering/claude-code-best-practices}{[anthropic.com]} ; awesome-claude-agents \href{https://github.com/vijaythecoder/awesome-claude-agents}{[GitHub]} ; SuperClaude \href{https://github.com/SuperClaude-Org/SuperClaude_Framework}{[GitHub]} ; "Good Context Good Code" \href{https://blog.stockapp.com/good-context-good-code/}{[blog.stockapp.com]} ; "Peeking Under the Hood of Claude Code" \href{https://medium.com/@outsightai/peeking-under-the-hood-of-claude-code-70f5a94a9a62}{[medium.com]} & Slides \href{https://docs.google.com/presentation/d/19mgkwAnJDc7JuJy0zhhoY0ZC15DiNpxL8kchPDnRkRQ/edit?usp=sharing}{[slides]} & \hyperlink{mod-CODE}{\tg{CODE}} \hyperlink{mod-MULTI}{\tg{MULTI}}\\ \midrule
\hypertarget{lec-STAN25-8}{}8 & 2025-10-17 & \textbf{Guest: Boris Cherny (creator of Claude Code, Anthropic)} & Claude Code & (week 4 readings) & Slides \href{https://docs.google.com/presentation/d/1bv7Zozn6z45CAh-IyX99dMPMyXCHC7zj95UfwErBYQ8/edit?usp=sharing}{[slides]} & \hyperlink{mod-CODE}{\tg{CODE}}\\ \midrule
\hypertarget{lec-STAN25-9}{}9 & 2025-10-20 & \textbf{How to build a breakthrough AI developer product} & AI-enhanced CLIs; terminal automation and scripting & Warp University \href{https://www.warp.dev/university?slug=university}{[warp.dev]} ; Warp vs Claude Code \href{https://www.warp.dev/university/getting-started/warp-vs-claude-code}{[warp.dev]} ; "How Warp uses Warp to build Warp" \href{https://notion.warp.dev/How-Warp-uses-Warp-to-build-Warp-21643263616d81a6b9e3e63fd8a7380c}{[notion.warp.dev]} & Slides \href{https://docs.google.com/presentation/d/1Djd4eBLBbRkma8rFnJAWMT0ptct_UGB8hipmoqFVkxQ/edit?usp=sharing}{[slides]} & \hyperlink{mod-CODE}{\tg{CODE}} \hyperlink{mod-DEPLOY}{\tg{DEPLOY}}\\ \midrule
\hypertarget{lec-STAN25-10}{}10 & 2025-10-24 & \textbf{Guest: Zach Lloyd (CEO, Warp)} & The agentic terminal & (week 5 readings) & Slides (Figma) \href{https://www.figma.com/slides/kwbcmtqTFQMfUhiMH8BiEx/Warp---Stanford--Copy-?node-id=9-116\&t=oBWBCk8mjg2l2NR5-1}{[figma.com]} & \hyperlink{mod-CODE}{\tg{CODE}}\\ \midrule
\hypertarget{lec-STAN25-11}{}11 & 2025-10-27 & \textbf{AI QA, SAST, DAST, and beyond} & Secure vibe coding; history of vulnerability detection; AI-generated test suites & SAST vs DAST \href{https://www.splunk.com/en_us/blog/learn/sast-vs-dast.html}{[splunk.com]} ; Copilot RCE via prompt injection \href{https://embracethered.com/blog/posts/2025/github-copilot-remote-code-execution-via-prompt-injection/}{[embracethered.com]} ; Semgrep, finding vulns with Claude Code and Codex \href{https://semgrep.dev/blog/2025/finding-vulnerabilities-in-modern-web-apps-using-claude-code-and-openai-codex/}{[semgrep.dev]} ; Unit42 agentic AI threats \href{https://unit42.paloaltonetworks.com/agentic-ai-threats/}{[unit42.paloaltonetworks.com]} ; OWASP Top Ten \href{https://owasp.org/www-project-top-ten/}{[owasp.org]} ; Chroma "Context Rot" \href{https://research.trychroma.com/context-rot}{[research.trychroma.com]} ; o3 CVE-2025-37899 prompt \href{https://github.com/SeanHeelan/o3_finds_cve-2025-37899/blob/master/system_prompt_uafs.prompt}{[GitHub]} & Slides \href{https://docs.google.com/presentation/d/1C05bCLasMDigBbkwdWbiz4WrXibzi6ua4hQQbTod_8c/edit?usp=sharing}{[slides]} & \hyperlink{mod-SAFE}{\tg{SAFE}} \hyperlink{mod-EVAL}{\tg{EVAL}} \hyperlink{mod-CODE}{\tg{CODE}}\\ \midrule
\hypertarget{lec-STAN25-12}{}12 & 2025-10-31 & \textbf{Guest: Isaac Evans (CEO, Semgrep)} & Static analysis and AI security & (week 6 readings) & n/a (no slides listed) & \hyperlink{mod-SAFE}{\tg{SAFE}}\\ \midrule
\hypertarget{lec-STAN25-13}{}13 & 2025-11-03 & \textbf{AI code review} & Which AI code systems to trust; debugging and diagnosis; generated documentation & "Code Reviews: Just Do It" \href{https://blog.codinghorror.com/code-reviews-just-do-it/}{[blog.codinghorror.com]} ; GitHub "How to review code effectively" \href{https://github.blog/developer-skills/github/how-to-review-code-effectively-a-github-staff-engineers-philosophy/}{[github.blog]} ; "AI-Assisted Assessment of Coding Practices in Modern Code Review" \href{https://arxiv.org/pdf/2405.13565}{[arXiv:2405.13565]} ; Graphite AI code review best practices \href{https://graphite.dev/guides/ai-code-review-implementation-best-practices}{[graphite.dev]} ; "Code Review Essentials" \href{https://blakesmith.me/2015/02/09/code-review-essentials-for-software-teams.html}{[blakesmith.me]} ; "Lessons from millions of AI code reviews" \href{https://www.youtube.com/watch?v=TswQeKftnaw}{[video]} & Slides \href{https://docs.google.com/presentation/d/1NkPzpuSQt6Esbnr2-EnxM9007TL6ebSPFwITyVY-QxU/edit?usp=sharing}{[slides]} & \hyperlink{mod-CODE}{\tg{CODE}} \hyperlink{mod-EVAL}{\tg{EVAL}}\\ \midrule
\hypertarget{lec-STAN25-14}{}14 & 2025-11-07 & \textbf{Guest: Tomas Reimers (CPO, Graphite)} & AI code review in practice & (week 7 readings) & Slides \href{https://drive.google.com/file/d/1hwF-RIkOJ_OFy17BKhzFyCtxSS7Pcf7p/view?usp=drive_link}{[drive.google.com]} & \hyperlink{mod-CODE}{\tg{CODE}}\\ \midrule
\hypertarget{lec-STAN25-15}{}15 & 2025-11-10 & \textbf{End-to-end apps from a single prompt} & Design and frontend for everyone; rapid UI/UX prototyping & n/a (none listed) & Slides \href{https://docs.google.com/presentation/d/1GrVLsfMFIXMiGjIW9D7EJIyLYh_-3ReHHNd_vRfZUoo/edit?usp=sharing}{[slides]} & \hyperlink{mod-CODE}{\tg{CODE}} \hyperlink{mod-APP}{\tg{APP}}\\ \midrule
\hypertarget{lec-STAN25-16}{}16 & 2025-11-14 & \textbf{Guest: Gaspar Garcia (Head of AI Research, Vercel)} & AI app generation (v0) & -- & Slides \href{https://docs.google.com/presentation/d/1Jf2aN5zIChd5tT86rZWWqY-iDWbxgR-uynKJxBR7E9E/edit?usp=sharing}{[slides]} & \hyperlink{mod-CODE}{\tg{CODE}} \hyperlink{mod-APP}{\tg{APP}}\\ \midrule
\hypertarget{lec-STAN25-17}{}17 & 2025-11-17 & \textbf{Incident response and DevOps} & Monitoring and observability for AI systems; automated incident response; triage and debugging & Google SRE book intro \href{https://sre.google/sre-book/introduction/}{[sre.google]} ; observability basics \href{https://last9.io/blog/traces-spans-observability-basics/}{[last9.io]} ; Resolve.ai posts \href{https://resolve.ai/blog/kubernetes-troubleshooting-in-resolve-ai}{[resolve.ai]} , \href{https://resolve.ai/blog/product-deep-dive}{[resolve.ai]} , \href{https://resolve.ai/blog/role-of-multi-agent-systems-AI-native-engineering}{[resolve.ai]} , \href{https://resolve.ai/blog/Top-5-Benefits}{[resolve.ai]} & Slides \href{https://docs.google.com/presentation/d/1Mfe-auWAsg9URCujneKnHr0AbO8O-_U4QXBVOlO4qp0/edit?usp=sharing}{[slides]} & \hyperlink{mod-DEPLOY}{\tg{DEPLOY}} \hyperlink{mod-MULTI}{\tg{MULTI}}\\ \midrule
\hypertarget{lec-STAN25-18}{}18 & 2025-11-21 & \textbf{Guests: Mayank Agarwal (CTO) and Milind Ganjoo, Resolve.ai} & Agentic on-call/SRE & (week 9 readings) & Slides \href{https://drive.google.com/file/d/11WnEbMGc9kny_WBpMN10I8oP8XsiQOnM/view?usp=sharing}{[drive.google.com]} & \hyperlink{mod-DEPLOY}{\tg{DEPLOY}}\\ \midrule
\hypertarget{lec-STAN25-19}{}19 & 2025-12-01 & \textbf{Software development in 10 years} & The future of developer roles; emerging AI coding paradigms & -- & -- & \hyperlink{mod-CODE}{\tg{CODE}} \tg{ORG}\\ \midrule
\hypertarget{lec-STAN25-20}{}20 & 2025-12-05 & \textbf{Guest: Martin Casado (General Partner, a16z)} & Industry trends and predictions & -- & -- & \hyperlink{mod-APP}{\tg{APP}} \tg{ORG}\\ \midrule
\end{longtable}}
\section{Assignments and projects}
{\scriptsize\setlength{\tabcolsep}{2.5pt}\begin{longtable}{P{2.8cm} P{5.4cm} P{3.0cm} P{3.5cm} P{1.5cm}}
\toprule Name & What you build & Skills & Starter code / resources & Tags\\ \midrule\endhead
\hypertarget{asg-STAN25-1}{}\textbf{Week 1 --- Prompting Techniques ("LLM Prompting Playground")} & Fill in prompts for K-shot, chain-of-thought, tool calling, self-consistency, RAG, and reflexion until the test scripts pass, using local Ollama models (mistral-nemo:12b, llama3.1:8b) & Prompt design, iterative refinement & \href{https://github.com/mihail911/modern-software-dev-assignments/tree/fall2025/week1}{[GitHub]} & \hyperlink{mod-FND}{\tg{FND}} \hyperlink{mod-RSN}{\tg{RSN}} \hyperlink{mod-TOOL}{\tg{TOOL}}\\ \midrule
\hypertarget{asg-STAN25-2}{}\textbf{Week 2 --- Action Item Extractor ("First Steps in the AI IDE")} & Extend a FastAPI and SQLite notes app with LLM structured-output extraction, tests, a refactor, frontend endpoints, and a generated README, using Cursor & AI-IDE workflow, LLM integration, testing & \href{https://github.com/mihail911/modern-software-dev-assignments/tree/fall2025/week2}{[GitHub]} & \hyperlink{mod-CODE}{\tg{CODE}} \hyperlink{mod-FND}{\tg{FND}}\\ \midrule
\hypertarget{asg-STAN25-3}{}\textbf{Week 3 --- Build a Custom MCP Server} & MCP server wrapping a real external API: at least 2 typed tools, robust error and rate-limit handling, STDIO or remote HTTP, optional OAuth2 & MCP tools/resources/prompts, transport, auth & \href{https://github.com/mihail911/modern-software-dev-assignments/tree/fall2025/week3}{[GitHub]} & \hyperlink{mod-TOOL}{\tg{TOOL}}\\ \midrule
\hypertarget{asg-STAN25-4}{}\textbf{Week 4 --- The Autonomous Coding Agent IRL ("Coding with Claude Code")} & At least 2 Claude Code automations (slash commands, CLAUDE.md, subagents, MCP) used to extend a FastAPI app & Claude Code customization, subagents & \href{https://github.com/mihail911/modern-software-dev-assignments/tree/fall2025/week4}{[GitHub]} & \hyperlink{mod-CODE}{\tg{CODE}} \hyperlink{mod-MULTI}{\tg{MULTI}}\\ \midrule
\hypertarget{asg-STAN25-5}{}\textbf{Week 5 --- Agentic Development with Warp} & Warp Drive prompts, rules, or MCP integrations, plus multi-agent concurrent sessions working from a task list & Agentic terminal workflows, parallel agents & \href{https://github.com/mihail911/modern-software-dev-assignments/tree/fall2025/week5}{[GitHub]} & \hyperlink{mod-CODE}{\tg{CODE}} \hyperlink{mod-MULTI}{\tg{MULTI}}\\ \midrule
\hypertarget{asg-STAN25-6}{}\textbf{Week 6 --- Scan and Fix Vulnerabilities with Semgrep ("Writing Secure AI Code")} & Run Semgrep on the app, triage findings, fix at least 3 issues (SQLi, unsafe DOM, CORS, etc.) & SAST, secrets and dependency scanning, remediation & \href{https://github.com/mihail911/modern-software-dev-assignments/tree/fall2025/week6}{[GitHub]} & \hyperlink{mod-SAFE}{\tg{SAFE}} \hyperlink{mod-CODE}{\tg{CODE}}\\ \midrule
\hypertarget{asg-STAN25-7}{}\textbf{Week 7 --- Exploring AI Code Review Using Graphite ("Code Review Reps")} & Implement 4 tasks with one-shot agent prompts, open PRs, and compare manual review with Graphite Diamond AI review & AI code review, PR hygiene & \href{https://github.com/mihail911/modern-software-dev-assignments/tree/fall2025/week7}{[GitHub]} & \hyperlink{mod-CODE}{\tg{CODE}} \hyperlink{mod-EVAL}{\tg{EVAL}}\\ \midrule
\hypertarget{asg-STAN25-8}{}\textbf{Week 8 --- Multi-Stack AI-Accelerated Web App Build} & Build the same CRUD app in 3 stacks, at least one with bolt.new and one with a non-JS language & AI app generation, stack comparison & \href{https://github.com/mihail911/modern-software-dev-assignments/tree/fall2025/week8}{[GitHub]} & \hyperlink{mod-CODE}{\tg{CODE}} \hyperlink{mod-APP}{\tg{APP}}\\ \midrule
\hypertarget{asg-STAN25-9}{}\textbf{Final project} & "Final project showcasing modern development practices" (details not public) & -- & -- & \hyperlink{mod-CODE}{\tg{CODE}} \tg{ORG}\\ \midrule
\end{longtable}}
\chapter{Northeastern: Agentic AI}\hypertarget{course-NEU}{}
{\small\begin{tabularx}{\linewidth}{@{}>{\bfseries}l X@{}}\toprule
Institution & Northeastern University (Khoury)\\
Term & Fall 2026 (in progress)\\
Instructors & Caglar Yildirim\\
Website & \href{https://agenticllms.github.io/}{https://agenticllms.github.io/}\\
Level & Gentle on-ramp (Python primer provided)\\
Archetype & Builder (LangChain/LangGraph stack)\\
Public materials & Slides: 2 HTML decks; video: none; code: Notebooks W02b--W05b, DADA1\\
Best for & Best guided, notebook-by-notebook path from an LLM API call to a deployed LangGraph agent.\\
\bottomrule\end{tabularx}}

Pairs one seminal paper with the matching LangChain/LangGraph documentation in every session.
One running example, the ``Husky Tech'' customer-support agent, grows week by week: a notebook prototype, then a project skeleton, then a Hugging Face Spaces deployment, a LangGraph rewrite, RAG, and finally a self-grading retrieval loop.
Later sessions cover planning, agentic RAG, three multi-agent sessions, MCP and its security, evaluation, and optimization with DSPy and cost controls.
Sessions after F~11/13 were not posted when this guide was compiled.
\section{Lectures}
{\scriptsize\setlength{\tabcolsep}{2.5pt}\begin{longtable}{P{0.45cm} P{1.2cm} P{2.5cm} P{3.6cm} P{3.6cm} P{3.3cm} P{1.5cm}}
\toprule \# & Date & Title & Key topics & Readings & Materials & Tags\\ \midrule\endhead
\hypertarget{lec-NEU-1}{}1 & F 9/11 & \textbf{Introduction to Agentic AI} & Workflows vs agents; course setup (VS Code, Python primer) & Building Effective Agents (Anthropic) [core]; Rise and Potential of LLM Agents (arXiv 2309.07864); LangGraph "Workflows and agents"; LangChain "Philosophy" & \href{https://agenticllms.github.io/week1.html}{[agenticllms.github.io]} ; Python primer \href{https://colab.research.google.com/drive/1s60I7NQiTpOWg2tK-KmFyiC8ZAY24lo7?usp=sharing}{[colab.research.google.com]} & \tg{ORG} \hyperlink{mod-ARCH}{\tg{ARCH}}\\ \midrule
\hypertarget{lec-NEU-2}{}2 & T 9/15 & \textbf{LLM Foundations (I)} & Transformers, autoregressive generation; run a local open-weight model (Ollama, Qwen 3.5); base vs instruction-tuned comparison & Attention Is All You Need [core]; GPT-1 "Improving Language Understanding"; videos: 3Blue1Brown LLMs visualized, Karpathy "Let's build GPT" & \href{https://agenticllms.github.io/week2a.html}{[agenticllms.github.io]} ; \href{https://www.youtube.com/watch?v=wjZofJX0v4M}{[video]} ; \href{https://www.youtube.com/watch?v=kCc8FmEb1nY}{[video]} & \hyperlink{mod-FND}{\tg{FND}} \hyperlink{mod-DEPLOY}{\tg{DEPLOY}}\\ \midrule
\hypertarget{lec-NEU-3}{}3 & F 9/18 & \textbf{LLM Foundations (II)} & Instruction tuning/RLHF, CoT; LLM API basics; environment setup (uv, .env keys) & InstructGPT [core]; Chain-of-Thought (2201.11903); LangChain Models, Messages docs & \href{https://agenticllms.github.io/week2b-setup.html}{[agenticllms.github.io]} ; code W02b (W02b.1\_\allowbreak{}LLM\_\allowbreak{}API\_\allowbreak{}Basics.ipynb, W02b.2\_\allowbreak{}AgentBasics.ipynb) \href{https://github.com/agenticllms/lecture-files/tree/main/W02b}{[GitHub]} & \hyperlink{mod-FND}{\tg{FND}} \hyperlink{mod-TRAIN}{\tg{TRAIN}}\\ \midrule
\hypertarget{lec-NEU-4}{}4 & T 9/22 & \textbf{Simple Agent Architecture} & Tools and agent loop; short-term memory, structured output; agentui.py chat UI & Toolformer [core]; MRKL (2205.00445); LangChain Tools, Agents docs & \href{https://agenticllms.github.io/week3a.html}{[agenticllms.github.io]} ; code W03a (W03a.1\_\allowbreak{}ShortTerm\_\allowbreak{}Memory.ipynb, W03a.2\_\allowbreak{}Structured\_\allowbreak{}Output.ipynb, agentui.py) \href{https://github.com/agenticllms/lecture-files/tree/main/W03a}{[GitHub]} & \hyperlink{mod-ARCH}{\tg{ARCH}} \hyperlink{mod-TOOL}{\tg{TOOL}} \hyperlink{mod-MEM}{\tg{MEM}}\\ \midrule
\hypertarget{lec-NEU-5}{}5 & F 9/25 & \textbf{Tool-Calling} & ReAct; Husky Tech customer support agent; notebook to project skeleton; deploy to HF Spaces (Gradio); tools-vs-model design rules & ReAct [core]; What Are Tools Anyway? (2403.15452); LangChain short-term memory, structured output docs & \href{https://agenticllms.github.io/week3b.html}{[agenticllms.github.io]} ; slides \href{https://agenticllms.github.io/week3b-slides.html}{[agenticllms.github.io]} ; code W03b \href{https://github.com/agenticllms/lecture-files/tree/main/W03b}{[GitHub]} & \hyperlink{mod-TOOL}{\tg{TOOL}} \hyperlink{mod-RSN}{\tg{RSN}} \hyperlink{mod-DEPLOY}{\tg{DEPLOY}}\\ \midrule
\hypertarget{lec-NEU-6}{}6 & T 9/29 & \textbf{Short-term \& Long-term Memory} & Memory in agents; LangGraph basics (state, nodes, edges, reducers, cycles); deploy to Render & Generative Agents [core]; MemGPT (2310.08560); Thinking in LangGraph, Graph API, Persistence docs & \href{https://agenticllms.github.io/week4a.html}{[agenticllms.github.io]} ; slides "Why LangGraph?" \href{https://agenticllms.github.io/week4a-slides.html}{[agenticllms.github.io]} ; code W04a \href{https://github.com/agenticllms/lecture-files/tree/main/W04a}{[GitHub]} & \hyperlink{mod-MEM}{\tg{MEM}} \hyperlink{mod-ARCH}{\tg{ARCH}} \hyperlink{mod-DEPLOY}{\tg{DEPLOY}}\\ \midrule
\hypertarget{lec-NEU-7}{}7 & F 10/2 & \textbf{Context Management} & Lost-in-the-middle, context engineering, middleware; LangGraph messages/tools/state; support agent in LangGraph & Lost in the Middle [core]; Effective Context Engineering (Anthropic); LangChain context engineering, middleware docs & \href{https://agenticllms.github.io/week4b.html}{[agenticllms.github.io]} ; code W04b \href{https://github.com/agenticllms/lecture-files/tree/main/W04b}{[GitHub]} & \hyperlink{mod-MEM}{\tg{MEM}} \hyperlink{mod-ARCH}{\tg{ARCH}}\\ \midrule
\hypertarget{lec-NEU-8}{}8 & T 10/6 & \textbf{Retrieval-Augmented Generation} & Chunking, embeddings, cosine similarity, vector store (dict, then Pinecone) & RAG (2005.11401) [core]; RAG Survey sec. 2-3 (2312.10997); LangChain Retrieval, Semantic search docs & \href{https://agenticllms.github.io/week5a.html}{[agenticllms.github.io]} ; code W05a \href{https://github.com/agenticllms/lecture-files/tree/main/W05a}{[GitHub]} & \hyperlink{mod-MEM}{\tg{MEM}}\\ \midrule
\hypertarget{lec-NEU-9}{}9 & F 10/9 & \textbf{Reflection \& Self-critique} & Self-grading retrieval loop, query rewriting; RAG inside the support agent; rate limits/cost & Reflexion [core]; Self-RAG; Self-Refine; LangGraph Graph API loops, unit testing docs & \href{https://agenticllms.github.io/week5b.html}{[agenticllms.github.io]} ; code W05b \href{https://github.com/agenticllms/lecture-files/tree/main/W05b}{[GitHub]} & \hyperlink{mod-RSN}{\tg{RSN}} \hyperlink{mod-MEM}{\tg{MEM}}\\ \midrule
\hypertarget{lec-NEU-10}{}10 & T 10/13 & \textbf{Planning \& Task Decomposition} & Skill libraries, ToT, task decomposition; subgraphs, skills & Voyager [core]; Tree of Thoughts; HuggingGPT (skim); LangGraph Subgraphs, Skills docs & -- & \hyperlink{mod-RSN}{\tg{RSN}} \hyperlink{mod-ARCH}{\tg{ARCH}}\\ \midrule
\hypertarget{lec-NEU-11}{}11 & F 10/16 & \textbf{Agentic RAG} & Corrective/adaptive RAG, multi-source routing & Corrective RAG (2401.15884) [core]; Adaptive-RAG (2403.14403); Custom RAG agent, router knowledge base docs & -- & \hyperlink{mod-MEM}{\tg{MEM}} \hyperlink{mod-RSN}{\tg{RSN}}\\ \midrule
\hypertarget{lec-NEU-12}{}12 & T 10/20 & \textbf{Multi-Agent Orchestration I} & MAS failure modes; handoffs, routers & Why Do Multi-Agent LLM Systems Fail? (2503.13657) [core]; AutoGen; Anthropic multi-agent research system; handoffs/router docs & -- & \hyperlink{mod-MULTI}{\tg{MULTI}}\\ \midrule
\hypertarget{lec-NEU-13}{}13 & F 10/23 & \textbf{Multi-Agent Orchestration II} & Role-based MAS; subagents, custom workflows & MetaGPT [core]; CAMEL (skim); Subagents, Custom workflow docs & -- & \hyperlink{mod-MULTI}{\tg{MULTI}}\\ \midrule
\hypertarget{lec-NEU-14}{}14 & T 10/27 & \textbf{Multi-Agent Orchestration III} & MAS failures revisited; LangGraph runtime (Pregel), Studio & Why MAS Fail sec. 4-5 (re-read) [core]; Agents Are Not Enough (2412.16241, optional); LangGraph runtime, Studio docs & -- & \hyperlink{mod-MULTI}{\tg{MULTI}} \hyperlink{mod-ARCH}{\tg{ARCH}}\\ \midrule
\hypertarget{lec-NEU-15}{}15 & F 10/30 & \textbf{Midterm Exam} & Exam & -- & -- & \tg{ORG}\\ \midrule
\hypertarget{lec-NEU-16}{}16 & T 11/3 & \textbf{Model Context Protocol} & MCP spec, MCP security landscape & MCP intro/spec [core]; MCP Landscape \& Security (2503.23278); LangChain MCP, Agent Server MCP docs & -- & \hyperlink{mod-TOOL}{\tg{TOOL}} \hyperlink{mod-SAFE}{\tg{SAFE}}\\ \midrule
\hypertarget{lec-NEU-17}{}17 & F 11/6 & \textbf{MCP in Agentic Systems} & Remote MCP, agent auth, enterprise MCP security & Enterprise MCP Security (2504.08623, skim) [core]; LangSmith Remote MCP, Agent Auth docs & -- & \hyperlink{mod-TOOL}{\tg{TOOL}} \hyperlink{mod-SAFE}{\tg{SAFE}} \hyperlink{mod-DEPLOY}{\tg{DEPLOY}}\\ \midrule
\hypertarget{lec-NEU-18}{}18 & T 11/10 & \textbf{Agentic AI Evaluation} & Benchmarks, trajectory/intermediate-step eval, LLM-as-judge & SWE-bench [core]; tau-bench (2406.12045); LLM-as-a-Judge (2306.05685); AgentBench (skim); LangSmith evaluation docs & -- & \hyperlink{mod-EVAL}{\tg{EVAL}}\\ \midrule
\hypertarget{lec-NEU-19}{}19 & F 11/13 & \textbf{Agentic AI Optimization} & Prompt/program optimization (DSPy); cost tracking, caching, fallbacks, spend policies & DSPy (2310.03714) [core]; LangSmith cost tracking, caching, fallbacks, spend policies, insights docs & -- & \hyperlink{mod-DEPLOY}{\tg{DEPLOY}} \hyperlink{mod-FND}{\tg{FND}}\\ \midrule
\end{longtable}}
\section{Assignments and projects}
{\scriptsize\setlength{\tabcolsep}{2.5pt}\begin{longtable}{P{2.8cm} P{5.4cm} P{3.0cm} P{3.5cm} P{1.5cm}}
\toprule Name & What you build & Skills & Starter code / resources & Tags\\ \midrule\endhead
\hypertarget{asg-NEU-1}{}\textbf{DADA1: The Trip Planner Agent} & A conversational day-trip planner with tools: get\_\allowbreak{}weather (Open-Meteo), find\_\allowbreak{}attractions/restaurants/transit (Tavily), find\_\allowbreak{}events (Ticketmaster). It has memory-based plan revision, a structured Itinerary/DayPlan output, 8 test prompts, and a Gradio app deployed on Render. Graded on tool routing, honesty at data limits, sources and memory. Optional extra credit: a find\_\allowbreak{}performer\_\allowbreak{}events chained tool & Tool design, prompting, memory, structured output, deployment & \href{https://github.com/agenticllms/lecture-files/tree/main/DADA1}{[GitHub]} & \hyperlink{mod-TOOL}{\tg{TOOL}} \hyperlink{mod-MEM}{\tg{MEM}} \hyperlink{mod-DEPLOY}{\tg{DEPLOY}}\\ \midrule
\hypertarget{asg-NEU-2}{}\textbf{Husky Tech Customer Support Agent (running in-class build, W03b-W05b)} & A support agent with 4 tools and ticket filing. It is built as a notebook, then a project skeleton, then deployed to HF Spaces, then rewritten in LangGraph, then given RAG over husky\_\allowbreak{}policy.md, then a reflection/self-grading loop & Tool calling, LangGraph, RAG, reflection, deployment & \href{https://github.com/agenticllms/lecture-files/tree/main/W03b}{[GitHub]} (also W04b, W05a, W05b) & \hyperlink{mod-TOOL}{\tg{TOOL}} \hyperlink{mod-ARCH}{\tg{ARCH}} \hyperlink{mod-MEM}{\tg{MEM}}\\ \midrule
\hypertarget{asg-NEU-3}{}\textbf{Final Project (research)} & A research contribution: agentic mechanism -> experimental manipulation -> measured behavior (e.g., memory policy -> stale-action rate). Suggested scope: controlled experiments on benchmarks, LoRA fine-tuning of 4-9B open models, trajectory analysis, fixed-token-budget architecture comparisons. Anchor papers cover eval, fine-tuning, safety, tool use/MCP and memory & Research design, experimentation, evaluation & n/a (resources: \href{https://agenticllms.github.io/research.html}{[agenticllms.github.io]}) & \hyperlink{mod-EVAL}{\tg{EVAL}} \hyperlink{mod-TRAIN}{\tg{TRAIN}} \tg{ORG}\\ \midrule
\hypertarget{asg-NEU-4}{}\textbf{Midterm Exam (F 10/30)} & Exam & -- & -- & \tg{ORG}\\ \midrule
\end{longtable}}
\chapter{Buffalo: CSE 410/510 Introduction to Agentic Engineering}\hypertarget{course-UB}{}
{\small\begin{tabularx}{\linewidth}{@{}>{\bfseries}l X@{}}\toprule
Institution & University at Buffalo\\
Term & Fall 2026\\
Instructors & Joaquin Carbonara, David Doermann\\
Website & \href{https://cse.buffalo.edu/courses/cse410/2026_fall/cse410ae/}{https://cse.buffalo.edu/courses/cse410/2026\_fall/cse410ae/}\\
Level & Upper-undergrad / MS\\
Archetype & Engineering discipline (framework-agnostic)\\
Public materials & Slides: none; video: none; code: none (lab specs only)\\
Best for & Five lab specs and a capstone spec that treat an agent as an engineered, evaluated, bounded system.\\
\bottomrule\end{tabularx}}

Treats an agent as an engineered system ``not a magic chatbot'': a bounded environment, designed tools, explicit evaluation and safety constraints.
No slides or code are public. What makes it useful is the lab sequence: a minimal tool-using agent, retrieval and memory, workflow graphs, an evaluation harness, and human-in-the-loop safety boundaries. The capstone spec, a useful bounded agent with an architecture document, an eval suite and a risk assessment, is a good template for your own portfolio project.
\section{Lectures}
{\scriptsize\setlength{\tabcolsep}{2.5pt}\begin{longtable}{P{0.45cm} P{1.2cm} P{2.5cm} P{3.6cm} P{3.6cm} P{3.3cm} P{1.5cm}}
\toprule \# & Date & Title & Key topics & Readings & Materials & Tags\\ \midrule\endhead
\hypertarget{lec-UB-1}{}1 & Aug 24-28 & \textbf{What is agentic engineering?} & Agents, environments, tools, goals, constraints & Reading Note 1 (specific reading not listed) & -- & \tg{ORG} \hyperlink{mod-ARCH}{\tg{ARCH}}\\ \midrule
\hypertarget{lec-UB-2}{}2 & Aug 31-Sep 4 & \textbf{Foundation models as components} & Prompting, context, structured outputs & Reading Note 2 & -- & \hyperlink{mod-FND}{\tg{FND}}\\ \midrule
\hypertarget{lec-UB-3}{}3 & Sep 8-11 & \textbf{Tool/function calling and ReAct} & Tool calling, ReAct; studio on linear algebra and cosine similarity & Reading Note 3 & -- & \hyperlink{mod-TOOL}{\tg{TOOL}} \hyperlink{mod-RSN}{\tg{RSN}}\\ \midrule
\hypertarget{lec-UB-4}{}4 & Sep 14-18 & \textbf{Retrieval-augmented generation and context construction} & RAG, context building & Reading Note 4 & -- & \hyperlink{mod-MEM}{\tg{MEM}}\\ \midrule
\hypertarget{lec-UB-5}{}5 & Sep 21-25 & \textbf{Embeddings, vector search, reranking} & Embeddings, vector DBs, reranking & -- & -- & \hyperlink{mod-MEM}{\tg{MEM}}\\ \midrule
\hypertarget{lec-UB-6}{}6 & Sep 28-Oct 2 & \textbf{Agent workflows} & Graphs, state machines, LangGraph or custom orchestration & -- & -- & \hyperlink{mod-ARCH}{\tg{ARCH}}\\ \midrule
\hypertarget{lec-UB-7}{}7 & Oct 5-9 & \textbf{Planning and decomposition} & When should an agent plan? & -- & -- & \hyperlink{mod-RSN}{\tg{RSN}}\\ \midrule
\hypertarget{lec-UB-8}{}8 & Oct 14-16 & \textbf{Human-agent interaction and memory} & Human-in-the-loop, memory architectures & -- & -- & \hyperlink{mod-MEM}{\tg{MEM}} \hyperlink{mod-SAFE}{\tg{SAFE}}\\ \midrule
\hypertarget{lec-UB-9}{}9 & Oct 19-23 & \textbf{Evaluation of agents} & From single-prompt to task-level evals & -- & -- & \hyperlink{mod-EVAL}{\tg{EVAL}}\\ \midrule
\hypertarget{lec-UB-10}{}10 & Oct 26-30 & \textbf{Security and privacy} & Prompt injection, tool misuse, sandboxing, privacy & -- & -- & \hyperlink{mod-SAFE}{\tg{SAFE}}\\ \midrule
\hypertarget{lec-UB-11}{}11 & Nov 2-6 & \textbf{Multi-agent systems and coordination} & MAS coordination & -- & -- & \hyperlink{mod-MULTI}{\tg{MULTI}}\\ \midrule
\hypertarget{lec-UB-12}{}12 & Nov 9-13 & \textbf{Coding agents, testing, CI, observability} & Coding agents, testing, CI, observability & -- & -- & \hyperlink{mod-CODE}{\tg{CODE}} \hyperlink{mod-DEPLOY}{\tg{DEPLOY}}\\ \midrule
\hypertarget{lec-UB-13}{}13 & Nov 16-20 & \textbf{Applied agents and responsible deployment} & Data science, discovery and education agents; responsible deployment & -- & -- & \hyperlink{mod-APP}{\tg{APP}} \hyperlink{mod-DEPLOY}{\tg{DEPLOY}}\\ \midrule
\hypertarget{lec-UB-14}{}14 & Nov 23-28 & \textbf{No class (Thanksgiving)} & Independent work on report and demo & -- & -- & \tg{ORG}\\ \midrule
\hypertarget{lec-UB-15}{}15 & Nov 30-Dec 7 (Wed Dec 2) & \textbf{Final demos, critique, synthesis} & Team demos & -- & -- & \tg{ORG}\\ \midrule
\end{longtable}}
\section{Assignments and projects}
{\scriptsize\setlength{\tabcolsep}{2.5pt}\begin{longtable}{P{2.8cm} P{5.4cm} P{3.0cm} P{3.5cm} P{1.5cm}}
\toprule Name & What you build & Skills & Starter code / resources & Tags\\ \midrule\endhead
\hypertarget{asg-UB-1}{}\textbf{Lab 1: A Minimal Tool-Using Agent (Wk 2-4)} & Python agent choosing among a calculator, a file reader and a search/lookup tool. Submit a design note and successful/failed runs & Tool calling, agent loop, failure analysis & -- & \hyperlink{mod-TOOL}{\tg{TOOL}} \hyperlink{mod-ARCH}{\tg{ARCH}}\\ \midrule
\hypertarget{asg-UB-2}{}\textbf{Lab 2: Retrieval and Memory (Wk 4-6)} & Small RAG pipeline over a course-selected corpus. Compare keyword/dense/hybrid retrieval, compute cosine similarity, do error analysis & RAG, embeddings, retrieval eval & -- & \hyperlink{mod-MEM}{\tg{MEM}} \hyperlink{mod-EVAL}{\tg{EVAL}}\\ \midrule
\hypertarget{asg-UB-3}{}\textbf{Lab 3: Workflow Graphs and State Machines (Wk 6-9)} & Agentic workflow as a graph/state machine with explicit states, transitions, failure handling and logging & Orchestration, state design, logging & -- & \hyperlink{mod-ARCH}{\tg{ARCH}}\\ \midrule
\hypertarget{asg-UB-4}{}\textbf{Lab 4: Evaluation Harness (Wk 9-11)} & Eval suite for a bounded task: test cases, success criteria, metrics, failure taxonomy, reproducible script & Agent evaluation & -- & \hyperlink{mod-EVAL}{\tg{EVAL}}\\ \midrule
\hypertarget{asg-UB-5}{}\textbf{Lab 5: Human-in-the-Loop and Safety Boundaries (Wk 11-13)} & Add human approval, structured outputs, guardrails or permission boundaries. Analyze where autonomy should be restricted & HITL, guardrails, permissions & -- & \hyperlink{mod-SAFE}{\tg{SAFE}}\\ \midrule
\hypertarget{asg-UB-6}{}\textbf{Math mini-assignment 1 (Wk 3-5)} & Linear algebra for ML: vectors, cosine similarity & Math for retrieval & -- & \hyperlink{mod-FND}{\tg{FND}} \hyperlink{mod-MEM}{\tg{MEM}}\\ \midrule
\hypertarget{asg-UB-7}{}\textbf{Math mini-assignment 2 (Wk 5-8)} & Content not specified & -- & -- & \hyperlink{mod-FND}{\tg{FND}}\\ \midrule
\hypertarget{asg-UB-8}{}\textbf{Reading/design notes 1-4 (Wk 1-4)} & Weekly notes on readings (content not specified) & Reading, design writing & -- & \tg{ORG}\\ \midrule
\hypertarget{asg-UB-9}{}\textbf{Project proposal (Wk 7), Midterm architecture review (Wk 8-10), Project checkpoint (Wk 12), Draft report (before Wk 15)} & Team project milestones & Design documentation & -- & \tg{ORG} \hyperlink{mod-ARCH}{\tg{ARCH}}\\ \midrule
\hypertarget{asg-UB-10}{}\textbf{Final Team Project: Useful Bounded Agent} & Working agent for a bounded task. Deliverables: problem statement, environment/tool inventory, architecture diagram, repo, eval dataset/script, results and error analysis, security/privacy/ethics risk assessment, demo, final report & End-to-end agent engineering & -- & \hyperlink{mod-ARCH}{\tg{ARCH}} \hyperlink{mod-EVAL}{\tg{EVAL}} \hyperlink{mod-SAFE}{\tg{SAFE}}\\ \midrule
\hypertarget{asg-UB-11}{}\textbf{Final oral exam (10\%)} & Oral exam & -- & -- & \tg{ORG}\\ \midrule
\end{longtable}}
\chapter{UVA: CS 6501 Workshop on Building AI Agents}\hypertarget{course-UVA}{}
{\small\begin{tabularx}{\linewidth}{@{}>{\bfseries}l X@{}}\toprule
Institution & University of Virginia\\
Term & Spring 2026\\
Instructors & Henry Kautz\\
Website & \href{https://www.cs.virginia.edu/\%7Ermw7my/Courses/AgenticAISpring2026/agenticAI.html}{https://www.cs.virginia.edu/\textasciitilde{}rmw7my/Courses/AgenticAISpring2026/agenticAI.html}\\
Level & Graduate workshop; Python\\
Archetype & Hands-on lab workshop (small/open models)\\
Public materials & Slides: PPTX per topic; video: 1 + podcasts; code: Starter scripts per topic\\
Best for & The richest set of runnable lab exercises: local LLMs, LangGraph, tools, RAG, VLMs, MCP/A2A, Tinker fine-tuning.\\
\bottomrule\end{tabularx}}

A workshop where most of the programming happens in class.
Each topic page has learning goals, an exercise sheet of measurable experiments, PPTX slides and starter scripts.
The first half is labs: running an LLM, LangGraph control flow, tool calling, the tool ecosystem, RAG, vision-language agents, MCP and A2A, and QLoRA-style fine-tuning with Tinker.
The second half is state-of-the-art lectures: context management, chain-of-thought internals, coding agents, personal agents and distillation, and agents for science, mathematics and visual reasoning.
Row numbers in the catalogue are calendar slots. The topic number appears in each title.
\section{Lectures}
{\scriptsize\setlength{\tabcolsep}{2.5pt}\begin{longtable}{P{0.45cm} P{1.2cm} P{2.5cm} P{3.6cm} P{3.6cm} P{3.3cm} P{1.5cm}}
\toprule \# & Date & Title & Key topics & Readings & Materials & Tags\\ \midrule\endhead
\hypertarget{lec-UVA-1}{}1 & Tue 13 Jan 2026 & \textbf{Topic 1: Running an LLM} & LLM overview, tokenizers, open models, HF benchmark datasets; run Llama 3.2-1B MMLU eval with/without quantization on CPU/GPU/Colab; build a simple chat agent with context management & Belcak et al., Small Language Models are the Future of Agentic AI (\href{https://arxiv.org/pdf/2506.02153}{[arXiv:2506.02153]}); Illustrated GPT-2 (\href{https://jalammar.github.io/illustrated-gpt2/}{[jalammar.github.io]}) & Page: \href{https://www.cs.virginia.edu/\%7Ermw7my/Courses/AgenticAISpring2026/Topic1Running/Running_an_LLM.html}{[cs.virginia.edu]} ; slides \texttt{Topic1Running/\allowbreak{}01 Running an LLM.\allowbreak{}pptx}; code \texttt{llama\_\allowbreak{}mmlu\_\allowbreak{}eval.\allowbreak{}py}, \texttt{simple\_\allowbreak{}chat\_\allowbreak{}agent.\allowbreak{}py}; NoteGPT podcast MP3 & \hyperlink{mod-FND}{\tg{FND}} \hyperlink{mod-EVAL}{\tg{EVAL}} \hyperlink{mod-DEPLOY}{\tg{DEPLOY}}\\ \midrule
\hypertarget{lec-UVA-3}{}3 & Tue 20 Jan 2026 & \textbf{Topic 2: Agent Frameworks} & LangGraph nodes/edges, conditional routing, parallel models (Llama + Qwen), chat history, multi-agent chat, checkpointing/crash recovery; optional web-search summarizer & LangGraph Graph API overview (\href{https://docs.langchain.com/oss/python/langgraph/graph-api}{[docs.langchain.com]}) & Page: \href{https://www.cs.virginia.edu/\%7Ermw7my/Courses/AgenticAISpring2026/Topic2Frameworks/Agent\%20Control\%20Flow\%20Frameworks.html}{[cs.virginia.edu]} ; slides \texttt{02 Agent Frameworks.\allowbreak{}pptx}; code \texttt{langgraph\_\allowbreak{}simple\_\allowbreak{}llama\_\allowbreak{}agent.\allowbreak{}py} & \hyperlink{mod-ARCH}{\tg{ARCH}} \hyperlink{mod-MULTI}{\tg{MULTI}}\\ \midrule
\hypertarget{lec-UVA-5}{}5 & Tue 27 Jan 2026 & \textbf{Topic 3: Agent Tool Use (online)} & Ollama server, parallel calls, OpenAI API/GPT-4o mini, OpenAI tool-call message format, manual tool handling (calculator), LangGraph multi-tool agent, tool forcing, persistent conversation graph & Schick et al., Toolformer (\href{https://arxiv.org/abs/2302.04761}{[arXiv:2302.04761]}) & Page: \href{https://www.cs.virginia.edu/\%7Ermw7my/Courses/AgenticAISpring2026/Topic3Tools/tools.html}{[cs.virginia.edu]} ; slides \texttt{03 Tools.\allowbreak{}pptx}; code \texttt{manual-tool-handling.\allowbreak{}py}, \texttt{langgraph-tool-handling.\allowbreak{}py}, \texttt{gpt4omini\_\allowbreak{}test.\allowbreak{}py}; guides \texttt{ollama\_\allowbreak{}guide.\allowbreak{}html}, \texttt{tool-forcing-strategies.\allowbreak{}html}; NotebookLM podcast \texttt{toolformer.\allowbreak{}m4a} & \hyperlink{mod-TOOL}{\tg{TOOL}} \hyperlink{mod-ARCH}{\tg{ARCH}}\\ \midrule
\hypertarget{lec-UVA-6}{}6 & Thu 29 Jan 2026 & \textbf{Topic 3 lab continued (online)} & Breakout rooms continuing Topic 3 & -- & -- & \hyperlink{mod-TOOL}{\tg{TOOL}}\\ \midrule
\hypertarget{lec-UVA-7}{}7 & Tue 03 Feb 2026 & \textbf{Topic 4: Exploring Tools} & ToolNode vs create\_\allowbreak{}react\_\allowbreak{}agent, parallel tool dispatch, tools on remote servers, LangGraph built-in and third-party tool library; "2-Hour Agent Project" with a partner & -- & Page: \href{https://www.cs.virginia.edu/\%7Ermw7my/Courses/AgenticAISpring2026/Topic4Exploring/exploring.html}{[cs.virginia.edu]} ; slides \texttt{exploring.\allowbreak{}pptx}; code \texttt{toolnode\_\allowbreak{}example.\allowbreak{}py}, \texttt{react\_\allowbreak{}agent\_\allowbreak{}example.\allowbreak{}py}; \texttt{agent\_\allowbreak{}tools\_\allowbreak{}reference.\allowbreak{}html}; MoltReader \href{https://github.com/HenryKautz/MoltReader}{[GitHub]} & \hyperlink{mod-TOOL}{\tg{TOOL}} \hyperlink{mod-RSN}{\tg{RSN}} \hyperlink{mod-ARCH}{\tg{ARCH}}\\ \midrule
\hypertarget{lec-UVA-9}{}9 & Tue 10 Feb 2026 & \textbf{Topic 5: Retrieval Augmented Generation} & Vector DBs, RAG pipeline, chunk size/overlap, top-K, query phrasing, score thresholds, prompt templates, unanswerable questions, cross-document synthesis (Qwen 2.5 1.5B vs GPT-4o mini) & -- & Page: \href{https://www.cs.virginia.edu/\%7Ermw7my/Courses/AgenticAISpring2026/Topic5RAG/rag.html}{[cs.virginia.edu]} ; slides \texttt{rag.\allowbreak{}pptx}; notebook \texttt{manual\_\allowbreak{}rag\_\allowbreak{}pipeline\_\allowbreak{}universal.\allowbreak{}ipynb}; \texttt{Corpora.\allowbreak{}zip}; pdf\_\allowbreak{}ocr \href{https://github.com/HenryKautz/pdf_ocr}{[GitHub]} & \hyperlink{mod-MEM}{\tg{MEM}} \hyperlink{mod-EVAL}{\tg{EVAL}}\\ \midrule
\hypertarget{lec-UVA-10}{}10 & Thu 12 Feb 2026 & \textbf{Topic 6: Vision-Language Models} & CLIP, contrastive alignment, projection layers, image embeddings, image generation; VLM chat agent in LangGraph; video-surveillance agent (OpenCV keyframes + LLaVA) & Illustrated Stable Diffusion (\href{https://jalammar.github.io/illustrated-stable-diffusion/}{[jalammar.github.io]}) & Page: \href{https://www.cs.virginia.edu/\%7Ermw7my/Courses/AgenticAISpring2026/Topic6VLM/vlm.html}{[cs.virginia.edu]} ; slides \texttt{vlm\_\allowbreak{}part1.\allowbreak{}pptx}, \texttt{vlm\_\allowbreak{}part2.\allowbreak{}pdf}, \texttt{vlm\_\allowbreak{}part3.\allowbreak{}pdf}; \texttt{vlm\_\allowbreak{}guide.\allowbreak{}html}, \texttt{image\_\allowbreak{}generation\_\allowbreak{}guide.\allowbreak{}html}; clip\_\allowbreak{}test \href{https://github.com/HenryKautz/clip_test}{[GitHub]} & \hyperlink{mod-WEB}{\tg{WEB}} \hyperlink{mod-APP}{\tg{APP}} \hyperlink{mod-FND}{\tg{FND}}\\ \midrule
\hypertarget{lec-UVA-12}{}12 & Tue 10 Mar 2026 & \textbf{Topic 7: MCP and A2A} & Model Context Protocol, tool discovery, remote Asta tools (AI2), Agent-to-Agent protocol, agent cards, registry, ngrok tunneling; quiz-tournament exercise & -- & Page: \href{https://www.cs.virginia.edu/\%7Ermw7my/Courses/AgenticAISpring2026/Topic7MCP/mcp.html}{[cs.virginia.edu]} ; slides \texttt{mcp.\allowbreak{}pptx}; \texttt{mcp\_\allowbreak{}asta\_\allowbreak{}lesson\_\allowbreak{}plan.\allowbreak{}html}, \texttt{a2a\_\allowbreak{}lesson\_\allowbreak{}plan.\allowbreak{}html}; code \texttt{a2a\_\allowbreak{}agent\_\allowbreak{}template.\allowbreak{}py}, \texttt{a2a\_\allowbreak{}test.\allowbreak{}py}, \texttt{a2a\_\allowbreak{}trivia.\allowbreak{}py}, \texttt{A2a\_\allowbreak{}registry.\allowbreak{}py}, \texttt{semantic\_\allowbreak{}similarity.\allowbreak{}py}, \texttt{test\_\allowbreak{}ngrok.\allowbreak{}sh} & \hyperlink{mod-TOOL}{\tg{TOOL}} \hyperlink{mod-MULTI}{\tg{MULTI}}\\ \midrule
\hypertarget{lec-UVA-13}{}13 & Tue 17 Mar 2026 & \textbf{Topic 7 cont. (A2A) + Topic 8: LLM Fine-Tuning} & Kinds of fine-tuning, fine-tuning vs RAG, QLoRA, Tinker pipeline; text-to-SQL fine-tuning of Llama & QLoRA (\href{https://arxiv.org/abs/2305.14314}{[arXiv:2305.14314]}), listed on the home page & Page: \href{https://www.cs.virginia.edu/\%7Ermw7my/Courses/AgenticAISpring2026/Topic8FineTuning/finetuning.html}{[cs.virginia.edu]} ; slides \texttt{finetuning.\allowbreak{}pptx}; \texttt{finetuning\_\allowbreak{}intro\_\allowbreak{}lesson\_\allowbreak{}plan.\allowbreak{}html}, \texttt{sql\_\allowbreak{}finetuning\_\allowbreak{}lesson\_\allowbreak{}plan.\allowbreak{}html}; data \texttt{sql\_\allowbreak{}create\_\allowbreak{}context\_\allowbreak{}v4.\allowbreak{}json}; \texttt{sql\_\allowbreak{}matches.\allowbreak{}py} & \hyperlink{mod-TRAIN}{\tg{TRAIN}} \hyperlink{mod-MULTI}{\tg{MULTI}}\\ \midrule
\hypertarget{lec-UVA-15}{}15 & Tue 24 Mar 2026 & \textbf{Topic 9: Context Management} & Compression, selection, persistence; how frontier models manage long contexts & AgentFold (\href{https://arxiv.org/abs/2510.24699}{[arXiv:2510.24699]}) and Anthropic "Effective context engineering" are listed on the home page; the topic page itself lists none & Page: \href{https://www.cs.virginia.edu/\%7Ermw7my/Courses/AgenticAISpring2026/Topic9Context/context.html}{[cs.virginia.edu]} ; slides \texttt{context.\allowbreak{}pptx}; \texttt{context\_\allowbreak{}management\_\allowbreak{}lesson\_\allowbreak{}plan.\allowbreak{}html} & \hyperlink{mod-MEM}{\tg{MEM}}\\ \midrule
\hypertarget{lec-UVA-16}{}16 & Thu 26 Mar 2026 & \textbf{Topic 10: Understanding Chain of Thought} & CoT prompt varieties, why CoT works, tokenized vs hidden-layer (latent) CoT, hidden thought tokens & CoT (\href{https://arxiv.org/abs/2201.11903}{[arXiv:2201.11903]}) on the home page; topic page lists none & Page: \href{https://www.cs.virginia.edu/\%7Ermw7my/Courses/AgenticAISpring2026/Topic10COT/cot.html}{[cs.virginia.edu]} ; slides \texttt{chain\_\allowbreak{}of\_\allowbreak{}thought.\allowbreak{}pptx}; \texttt{chain\_\allowbreak{}of\_\allowbreak{}thought\_\allowbreak{}lesson\_\allowbreak{}plan.\allowbreak{}html} & \hyperlink{mod-RSN}{\tg{RSN}}\\ \midrule
\hypertarget{lec-UVA-17}{}17 & Tue 31 Mar 2026 & \textbf{Topic 11: Coding Agents} & History from statistical LMs, frontier coding agent architecture (Claude Code), skills files, evolutionary programming & AlphaEvolve (\href{https://doi.org/10.48550/arXiv.2506.13131}{[doi.org]}); AVO (\href{https://doi.org/10.48550/arXiv.2603.24517}{[doi.org]}); EvoSkill (\href{https://doi.org/10.48550/arXiv.2603.02766}{[doi.org]}) & Page: \href{https://www.cs.virginia.edu/\%7Ermw7my/Courses/AgenticAISpring2026/Topic11Coding/coding.html}{[cs.virginia.edu]} ; slides \texttt{coding\_\allowbreak{}1\_\allowbreak{}history.\allowbreak{}pptx}, \texttt{coding\_\allowbreak{}2\_\allowbreak{}claude.\allowbreak{}pptx}, \texttt{coding\_\allowbreak{}3\_\allowbreak{}evolution.\allowbreak{}pptx} & \hyperlink{mod-CODE}{\tg{CODE}} \hyperlink{mod-ARCH}{\tg{ARCH}}\\ \midrule
\hypertarget{lec-UVA-18}{}18 & Thu 02 Apr 2026 & \textbf{Topic 12: Personal Agents and Distilling} & Clawdbot to OpenClaw/NullClaw/CoPaw/NemoClaw, Claude Cowork; risks of agents controlling accounts; OpenClaw-RL; knowledge distillation & n/a (resource links only) & Page: \href{https://www.cs.virginia.edu/\%7Ermw7my/Courses/AgenticAISpring2026/Topic12Personal/personal.html}{[cs.virginia.edu]} ; slides \texttt{personal.\allowbreak{}pptx}, \texttt{Knowledge Distillation Slides.\allowbreak{}pptx}; \texttt{personal\_\allowbreak{}lesson\_\allowbreak{}plan.\allowbreak{}html}, \texttt{distill.\allowbreak{}html} & \hyperlink{mod-WEB}{\tg{WEB}} \hyperlink{mod-SAFE}{\tg{SAFE}} \hyperlink{mod-TRAIN}{\tg{TRAIN}}\\ \midrule
\hypertarget{lec-UVA-19}{}19 & Tue 07 Apr 2026 & \textbf{Topic 13: Agentic AI for Science} & Research as an agent loop, scientific agent stack, AI2 tools (Semantic Scholar, OpenScholar, ScholarQA, Asta), JFC HEP case study, Orchestral AI, HEPTAPOD & AutoDiscovery (\href{https://arxiv.org/abs/2507.00310}{[arXiv:2507.00310]}); AI agents perform experimental HEP (\href{https://inspirehep.net/literature/3132489}{[inspirehep.net]}); HEPTAPOD (\href{https://arxiv.org/abs/2512.15867}{[arXiv:2512.15867]}) & Page: \href{https://www.cs.virginia.edu/\%7Ermw7my/Courses/AgenticAISpring2026/Topic13Science/science.html}{[cs.virginia.edu]} ; slides \texttt{science.\allowbreak{}pptx}; HEPTAPOD code \href{https://github.com/tonymenzo/heptapod}{[GitHub]} & \hyperlink{mod-APP}{\tg{APP}} \hyperlink{mod-MULTI}{\tg{MULTI}}\\ \midrule
\hypertarget{lec-UVA-20}{}20 & Thu 09 Apr 2026 & \textbf{Topic 14: Agentic AI for Mathematics and Verification (guest: Prof. Kevin Sullivan)} & Lean as lingua franca; AlphaProof, Lean Copilot, APOLLO, CLEVER; formal methods & Lean4CS1 notes (\href{https://kevinsullivan.github.io/Lean4CS1/}{[kevinsullivan.github.io]}) & Page: \href{https://www.cs.virginia.edu/\%7Ermw7my/Courses/AgenticAISpring2026/Topic14Math/math.html}{[cs.virginia.edu]} ; \texttt{Kevin Sullivan Agents-Summary.\allowbreak{}html}, \texttt{Kevin Sullivan Agents-Class Note.\allowbreak{}html} & \hyperlink{mod-RSN}{\tg{RSN}} \hyperlink{mod-APP}{\tg{APP}}\\ \midrule
\hypertarget{lec-UVA-21}{}21 & Tue 14 Apr 2026 & \textbf{Topic 15: Visual Reasoning} & Video as reasoning substrate, diffusion transformers, Chain-of-Frames vs Chain-of-Steps, emergent working memory and self-correction, MIRAGE critique of VLM evaluation & Demystifying Video Reasoning (\href{https://arxiv.org/abs/2603.16870}{[arXiv:2603.16870]}); Chain-of-Frames (\href{https://arxiv.org/abs/2506.00318}{[arXiv:2506.00318]}); MIRAGE (\href{https://arxiv.org/abs/2603.21687}{[arXiv:2603.21687]}) & Page: \href{https://www.cs.virginia.edu/\%7Ermw7my/Courses/AgenticAISpring2026/Topic15Visual/visual.html}{[cs.virginia.edu]} ; slides \texttt{visual.\allowbreak{}pptx}, \texttt{mirage.\allowbreak{}pptx}; video \href{https://www.youtube.com/watch?v=Gs9TPZmzo-s}{[video]} & \hyperlink{mod-RSN}{\tg{RSN}} \hyperlink{mod-EVAL}{\tg{EVAL}}\\ \midrule
\hypertarget{lec-UVA-22}{}22 & Thu 16 Apr 2026 & \textbf{Final project presentations begin} & 5-minute videos due at noon & -- & -- & \tg{ORG}\\ \midrule
\hypertarget{lec-UVA-23}{}23 & Tue 21 Apr 2026 & \textbf{Final project presentations} & -- & -- & -- & \tg{ORG}\\ \midrule
\hypertarget{lec-UVA-24}{}24 & Tue 28 Apr 2026 & \textbf{Final project presentations} & -- & -- & \texttt{Spin/\allowbreak{}index.\allowbreak{}html} (questioner selector) & \tg{ORG}\\ \midrule
\end{longtable}}
\section{Assignments and projects}
{\scriptsize\setlength{\tabcolsep}{2.5pt}\begin{longtable}{P{2.8cm} P{5.4cm} P{3.0cm} P{3.5cm} P{1.5cm}}
\toprule Name & What you build & Skills & Starter code / resources & Tags\\ \midrule\endhead
\hypertarget{asg-UVA-1}{}\textbf{Topic 1 portfolio: Running an LLM} & MMLU timing/accuracy study across models, quantization and hardware; chat agent with bounded context and restartability; optional MT-Bench & HF transformers, quantization, benchmarking, context management & \texttt{llama\_\allowbreak{}mmlu\_\allowbreak{}eval.\allowbreak{}py}, \texttt{simple\_\allowbreak{}chat\_\allowbreak{}agent.\allowbreak{}py} via \href{https://www.cs.virginia.edu/\%7Ermw7my/Courses/AgenticAISpring2026/Topic1Running/Running_an_LLM.html}{[cs.virginia.edu]} & \hyperlink{mod-FND}{\tg{FND}} \hyperlink{mod-EVAL}{\tg{EVAL}} \hyperlink{mod-MEM}{\tg{MEM}}\\ \midrule
\hypertarget{asg-UVA-2}{}\textbf{Topic 2 portfolio: LangGraph agents} & Routed, parallel, multi-model LangGraph chat with history and crash recovery; optional web summarizer & LangGraph graphs, routing, checkpointing & \texttt{langgraph\_\allowbreak{}simple\_\allowbreak{}llama\_\allowbreak{}agent.\allowbreak{}py} via \href{https://www.cs.virginia.edu/\%7Ermw7my/Courses/AgenticAISpring2026/Topic2Frameworks/Agent\%20Control\%20Flow\%20Frameworks.html}{[cs.virginia.edu]} & \hyperlink{mod-ARCH}{\tg{ARCH}} \hyperlink{mod-MULTI}{\tg{MULTI}}\\ \midrule
\hypertarget{asg-UVA-3}{}\textbf{Topic 3 portfolio: Tool use} & Ollama parallel timing; manual tool calling with a calculator; LangGraph multi-tool agent with tool\_\allowbreak{}map; persistent graph with Mermaid diagram & OpenAI tool format, tool forcing, checkpointing & \texttt{manual-tool-handling.\allowbreak{}py}, \texttt{langgraph-tool-handling.\allowbreak{}py} via \href{https://www.cs.virginia.edu/\%7Ermw7my/Courses/AgenticAISpring2026/Topic3Tools/tools.html}{[cs.virginia.edu]} & \hyperlink{mod-TOOL}{\tg{TOOL}} \hyperlink{mod-ARCH}{\tg{ARCH}}\\ \midrule
\hypertarget{asg-UVA-4}{}\textbf{Topic 4: 2-Hour Agent Project} & Partner-built mini agent using LangGraph/third-party tools & ToolNode, ReAct agents, tool ecosystem & \texttt{toolnode\_\allowbreak{}example.\allowbreak{}py}, \texttt{react\_\allowbreak{}agent\_\allowbreak{}example.\allowbreak{}py} via \href{https://www.cs.virginia.edu/\%7Ermw7my/Courses/AgenticAISpring2026/Topic4Exploring/exploring.html}{[cs.virginia.edu]} & \hyperlink{mod-TOOL}{\tg{TOOL}} \hyperlink{mod-ARCH}{\tg{ARCH}}\\ \midrule
\hypertarget{asg-UVA-5}{}\textbf{Topic 5: RAG experiments} & 11 controlled RAG experiments (RAG vs no-RAG, top-K, chunking, prompts, unanswerable questions) & Vector retrieval, RAG evaluation & \texttt{manual\_\allowbreak{}rag\_\allowbreak{}pipeline\_\allowbreak{}universal.\allowbreak{}ipynb}, \texttt{Corpora.\allowbreak{}zip} via \href{https://www.cs.virginia.edu/\%7Ermw7my/Courses/AgenticAISpring2026/Topic5RAG/rag.html}{[cs.virginia.edu]} & \hyperlink{mod-MEM}{\tg{MEM}} \hyperlink{mod-EVAL}{\tg{EVAL}}\\ \midrule
\hypertarget{asg-UVA-6}{}\textbf{Topic 6: VLM agents} & Image-chat LangGraph agent; video-surveillance agent with OpenCV + LLaVA (optional TTS intruder alert) & VLMs, Ollama llava, Gradio & \texttt{gradio\_\allowbreak{}quickstart.\allowbreak{}py} via \href{https://www.cs.virginia.edu/\%7Ermw7my/Courses/AgenticAISpring2026/Topic6VLM/vlm.html}{[cs.virginia.edu]} & \hyperlink{mod-WEB}{\tg{WEB}} \hyperlink{mod-APP}{\tg{APP}}\\ \midrule
\hypertarget{asg-UVA-7}{}\textbf{Topic 7: MCP + A2A} & MCP client for AI2 Asta tools; A2A agents, registry and trivia tournament over ngrok & MCP, A2A, networking & \texttt{a2a\_\allowbreak{}agent\_\allowbreak{}template.\allowbreak{}py} etc. via \href{https://www.cs.virginia.edu/\%7Ermw7my/Courses/AgenticAISpring2026/Topic7MCP/mcp.html}{[cs.virginia.edu]} & \hyperlink{mod-TOOL}{\tg{TOOL}} \hyperlink{mod-MULTI}{\tg{MULTI}}\\ \midrule
\hypertarget{asg-UVA-8}{}\textbf{Topic 8: Text-to-SQL fine-tuning} & Fine-tune Llama on text-to-SQL with Tinker and score semantic SQL matches & LoRA/QLoRA, Tinker, eval & \texttt{sql\_\allowbreak{}create\_\allowbreak{}context\_\allowbreak{}v4.\allowbreak{}json}, \texttt{sql\_\allowbreak{}matches.\allowbreak{}py} via \href{https://www.cs.virginia.edu/\%7Ermw7my/Courses/AgenticAISpring2026/Topic8FineTuning/finetuning.html}{[cs.virginia.edu]} & \hyperlink{mod-TRAIN}{\tg{TRAIN}} \hyperlink{mod-EVAL}{\tg{EVAL}}\\ \midrule
\hypertarget{asg-UVA-9}{}\textbf{Final project (55\%)} & Team project from a list of 20+ ideas (e.g., distillation, mini coding agent with sub-agents, MemGPT-style memory, debate agents, NL-to-PDDL planner, Lean/SAT assistants, misalignment replication); video, presentation, report & End-to-end agent design and evaluation & -- & \hyperlink{mod-ARCH}{\tg{ARCH}} \hyperlink{mod-EVAL}{\tg{EVAL}} \hyperlink{mod-APP}{\tg{APP}}\\ \midrule
\end{longtable}}
\chapter{UW-Madison: CS 839 AI Agents}\hypertarget{course-UW}{}
{\small\begin{tabularx}{\linewidth}{@{}>{\bfseries}l X@{}}\toprule
Institution & University of Wisconsin--Madison\\
Term & Spring 2026\\
Instructors & Aws Albarghouthi\\
Website & \href{https://github.com/barghouthi/ai-agents}{https://github.com/barghouthi/ai-agents}\\
Level & Graduate research seminar\\
Archetype & Research seminar (paper reading)\\
Public materials & Slides: none; video: none; code: Demo notebook; MIT license\\
Best for & A curated, critical reading list from ReAct to self-improving and evolutionary agents, plus a no-frameworks coding-agent assignment.\\
\bottomrule\end{tabularx}}

A paper seminar with two main papers and several supplementary papers per session.
It runs from CoT, ReAct, Toolformer and ToT to coding agents, deep research, agents for science, prompt-injection security, agent optimization (DSPy, TextGrad, ADAS, GEPA), self-improving and evolutionary agents, search, multi-agent systems and memory.
The buy-in assignment, a CLI coding agent built on raw API calls with no agent frameworks, is one of the best single exercises in the collection.
The repository is MIT-licensed.
\section{Lectures}
{\scriptsize\setlength{\tabcolsep}{2.5pt}\begin{longtable}{P{0.45cm} P{1.2cm} P{2.5cm} P{3.6cm} P{3.6cm} P{3.3cm} P{1.5cm}}
\toprule \# & Date & Title & Key topics & Readings & Materials & Tags\\ \midrule\endhead
\hypertarget{lec-UW-1}{}1 & Tue 20 Jan 2026 & \textbf{Class overview} & Course logistics, overview of the agents space & -- & -- & \tg{ORG}\\ \midrule
\hypertarget{lec-UW-2}{}2 & Thu 22 Jan 2026 & \textbf{Building agents} & Instructor intro to building LLM agents & Building Effective Agents (Anthropic) \href{https://www.anthropic.com/engineering/building-effective-agents}{[anthropic.com]} (listed as "see also" in week 1) & Demo notebook \href{https://github.com/barghouthi/ai-agents/blob/main/demos/agents.ipynb}{[GitHub]} (session association assumed) & \hyperlink{mod-ARCH}{\tg{ARCH}} \hyperlink{mod-TOOL}{\tg{TOOL}}\\ \midrule
\hypertarget{lec-UW-3}{}3 & Tue 27 Jan 2026 & \textbf{Reasoning and acting I} & Chain-of-thought, ReAct & CoT \href{https://arxiv.org/abs/2201.11903}{[arXiv:2201.11903]} ; ReAct \href{https://arxiv.org/abs/2210.03629}{[arXiv:2210.03629]} ; supp.: Self-Consistency \href{https://arxiv.org/abs/2203.11171}{[arXiv:2203.11171]} , Kimi K2 \href{https://arxiv.org/abs/2507.20534}{[arXiv:2507.20534]} , Llama 3 Herd \href{https://arxiv.org/abs/2407.21783}{[arXiv:2407.21783]} , Constitutional AI \href{https://arxiv.org/abs/2212.08073}{[arXiv:2212.08073]} & -- & \hyperlink{mod-RSN}{\tg{RSN}} \hyperlink{mod-FND}{\tg{FND}}\\ \midrule
\hypertarget{lec-UW-4}{}4 & Thu 29 Jan 2026 & \textbf{Reasoning and acting II} & Tool learning, Tree/Graph of Thoughts & Toolformer \href{https://arxiv.org/abs/2302.04761}{[arXiv:2302.04761]} ; Tree of Thoughts \href{https://arxiv.org/abs/2305.10601}{[arXiv:2305.10601]} ; supp.: LLMs as Tool Makers \href{https://arxiv.org/abs/2305.17126}{[arXiv:2305.17126]} , ART \href{https://arxiv.org/abs/2303.09014}{[arXiv:2303.09014]} , Graph of Thoughts \href{https://ojs.aaai.org/index.php/AAAI/article/view/29720}{[ojs.aaai.org]} & -- & \hyperlink{mod-TOOL}{\tg{TOOL}} \hyperlink{mod-RSN}{\tg{RSN}}\\ \midrule
\hypertarget{lec-UW-5}{}5 & Tue 03 Feb 2026 & \textbf{Coding agents I} & SWE benchmarks, SWE-agent ACI & SWE-bench \href{https://arxiv.org/abs/2310.06770}{[arXiv:2310.06770]} ; SWE-agent \href{https://arxiv.org/abs/2405.15793}{[arXiv:2405.15793]} ; supp.: Terminal-Bench \href{https://openreview.net/forum?id=a7Qa4CcHak}{[openreview.net]} , Agentless \href{https://arxiv.org/abs/2407.01489}{[arXiv:2407.01489]} & -- & \hyperlink{mod-CODE}{\tg{CODE}} \hyperlink{mod-EVAL}{\tg{EVAL}}\\ \midrule
\hypertarget{lec-UW-6}{}6 & Thu 05 Feb 2026 & \textbf{Coding agents II} & Open agent platforms, secure code generation benchmarks & OpenHands \href{https://arxiv.org/abs/2407.16741}{[arXiv:2407.16741]} ; BaxBench \href{https://arxiv.org/pdf/2502.11844}{[arXiv:2502.11844]} ; supp.: RepairAgent \href{https://arxiv.org/abs/2403.17134}{[arXiv:2403.17134]} , RedCode \href{https://arxiv.org/abs/2411.07781}{[arXiv:2411.07781]} , WebArena \href{https://arxiv.org/abs/2307.13854}{[arXiv:2307.13854]} & -- & \hyperlink{mod-CODE}{\tg{CODE}} \hyperlink{mod-SAFE}{\tg{SAFE}} \hyperlink{mod-EVAL}{\tg{EVAL}}\\ \midrule
\hypertarget{lec-UW-7}{}7 & Tue 10 Feb 2026 & \textbf{Deep research agents I} & Article writing, autonomous research labs & STORM \href{https://arxiv.org/abs/2402.14207}{[arXiv:2402.14207]} ; Agent Laboratory \href{https://arxiv.org/abs/2501.04227}{[arXiv:2501.04227]} ; supp.: ResearchAgent \href{https://arxiv.org/abs/2404.07738}{[arXiv:2404.07738]} , Deep research agents roadmap \href{https://arxiv.org/pdf/2506.18096}{[arXiv:2506.18096]} , Anthropic multi-agent research system \href{https://www.anthropic.com/engineering/multi-agent-research-system}{[anthropic.com]} & -- & \hyperlink{mod-APP}{\tg{APP}} \hyperlink{mod-MULTI}{\tg{MULTI}}\\ \midrule
\hypertarget{lec-UW-8}{}8 & Thu 12 Feb 2026 & \textbf{Deep research agents II} & Literature QA agents, research benchmarks & PaperQA \href{https://arxiv.org/abs/2312.07559}{[arXiv:2312.07559]} ; LiveResearchBench \href{https://arxiv.org/abs/2510.14240}{[arXiv:2510.14240]} ; supp.: Can LLMs Generate Novel Research Ideas? \href{https://arxiv.org/abs/2409.04109}{[arXiv:2409.04109]} , DR Tulu \href{https://arxiv.org/pdf/2511.19399}{[arXiv:2511.19399]} & -- & \hyperlink{mod-APP}{\tg{APP}} \hyperlink{mod-MEM}{\tg{MEM}} \hyperlink{mod-EVAL}{\tg{EVAL}}\\ \midrule
\hypertarget{lec-UW-9}{}9 & Tue 17 Feb 2026 & \textbf{Agents for science I} & AI Scientist, autonomous chemistry & The AI Scientist \href{https://arxiv.org/abs/2408.06292}{[arXiv:2408.06292]} ; Autonomous chemical research with LLMs \href{https://www.nature.com/articles/s41586-023-06792-0}{[nature.com]} ; supp.: ChemCrow \href{https://arxiv.org/abs/2304.05376}{[arXiv:2304.05376]} & -- & \hyperlink{mod-APP}{\tg{APP}}\\ \midrule
\hypertarget{lec-UW-10}{}10 & Thu 19 Feb 2026 & \textbf{Agents for science II} & AI co-scientist, hypothesis falsification & Towards an AI Co-Scientist \href{https://arxiv.org/abs/2502.18864}{[arXiv:2502.18864]} ; Agentic Sequential Falsifications \href{https://arxiv.org/abs/2502.09858}{[arXiv:2502.09858]} ; supp.: SciAgents \href{https://arxiv.org/abs/2409.05556}{[arXiv:2409.05556]} & -- & \hyperlink{mod-APP}{\tg{APP}} \hyperlink{mod-MULTI}{\tg{MULTI}}\\ \midrule
\hypertarget{lec-UW-11}{}11 & Tue 24 Feb 2026 & \textbf{Security \& prompt injection I} & Indirect prompt injection, AgentDojo benchmark & Indirect Prompt Injection \href{https://arxiv.org/abs/2302.12173}{[arXiv:2302.12173]} ; AgentDojo \href{https://arxiv.org/abs/2406.13352}{[arXiv:2406.13352]} ; supp.: Jailbroken \href{https://arxiv.org/abs/2307.02483}{[arXiv:2307.02483]} & -- & \hyperlink{mod-SAFE}{\tg{SAFE}} \hyperlink{mod-EVAL}{\tg{EVAL}}\\ \midrule
\hypertarget{lec-UW-12}{}12 & Thu 26 Feb 2026 & \textbf{Security \& prompt injection II} & Capability-based defenses, browser agent security & CaMeL \href{https://arxiv.org/abs/2503.18813}{[arXiv:2503.18813]} ; ceLLMate \href{https://arxiv.org/abs/2512.12594}{[arXiv:2512.12594]} ; supp.: Lethal trifecta \href{https://simonwillison.net/2025/Jun/16/the-lethal-trifecta/}{[simonwillison.net]} , Meta practical agent security \href{https://ai.meta.com/blog/practical-ai-agent-security/}{[ai.meta.com]} & -- & \hyperlink{mod-SAFE}{\tg{SAFE}}\\ \midrule
\hypertarget{lec-UW-13}{}13 & Tue 03 Mar 2026 & \textbf{Optimizing agents I} & Declarative LM programs, textual gradients & DSPy \href{https://arxiv.org/abs/2310.03714}{[arXiv:2310.03714]} ; TextGrad \href{https://arxiv.org/abs/2406.07496}{[arXiv:2406.07496]} ; supp.: LLMs as Optimizers \href{https://arxiv.org/abs/2309.03409}{[arXiv:2309.03409]} & -- & \hyperlink{mod-ARCH}{\tg{ARCH}} \hyperlink{mod-TRAIN}{\tg{TRAIN}}\\ \midrule
\hypertarget{lec-UW-14}{}14 & Thu 05 Mar 2026 & \textbf{Optimizing agents II} & Automated agent design, reflective prompt evolution & ADAS \href{https://arxiv.org/abs/2408.08435}{[arXiv:2408.08435]} ; GEPA \href{https://arxiv.org/abs/2507.19457}{[arXiv:2507.19457]} ; supp.: EvoPrompt \href{https://arxiv.org/abs/2309.08532}{[arXiv:2309.08532]} , Trace \href{https://arxiv.org/abs/2406.16218}{[arXiv:2406.16218]} & -- & \hyperlink{mod-ARCH}{\tg{ARCH}} \hyperlink{mod-TRAIN}{\tg{TRAIN}}\\ \midrule
\hypertarget{lec-UW-15}{}15 & Tue 10 Mar 2026 & \textbf{Self-improving agents I} & Lifelong skill libraries, self-referential agents & Voyager \href{https://arxiv.org/abs/2305.16291}{[arXiv:2305.16291]} ; Gödel Agent \href{https://arxiv.org/abs/2410.04444}{[arXiv:2410.04444]} & -- & \hyperlink{mod-ARCH}{\tg{ARCH}} \hyperlink{mod-WEB}{\tg{WEB}}\\ \midrule
\hypertarget{lec-UW-16}{}16 & Thu 12 Mar 2026 & \textbf{Self-improving agents II} & Evolutionary coding agents, live self-evolving SWE agent & AlphaEvolve \href{https://deepmind.google/discover/blog/alphaevolve-a-gemini-powered-coding-agent-for-designing-advanced-algorithms/}{[deepmind.google]} ; Live-SWE-agent \href{https://arxiv.org/abs/2511.13646}{[arXiv:2511.13646]} ; supp.: LATS \href{https://arxiv.org/abs/2310.04406}{[arXiv:2310.04406]} & -- & \hyperlink{mod-CODE}{\tg{CODE}} \hyperlink{mod-RSN}{\tg{RSN}}\\ \midrule
\hypertarget{lec-UW-17}{}17 & Tue 17 Mar 2026 & \textbf{Searching the solution space I} & Large-scale sampling and search for code and math & AlphaCode \href{https://arxiv.org/abs/2203.07814}{[arXiv:2203.07814]} ; AlphaGeometry \href{https://www.nature.com/articles/s41586-023-06747-5}{[nature.com]} ; supp.: AlphaProof \href{https://www.nature.com/articles/s41586-025-09833-y}{[nature.com]} , AlphaZero \href{https://arxiv.org/abs/1712.01815}{[arXiv:1712.01815]} & -- & \hyperlink{mod-RSN}{\tg{RSN}} \hyperlink{mod-CODE}{\tg{CODE}}\\ \midrule
\hypertarget{lec-UW-18}{}18 & Thu 19 Mar 2026 & \textbf{Searching the solution space II} & LLM-driven evolutionary search & ShinkaEvolve \href{https://arxiv.org/abs/2509.19349}{[arXiv:2509.19349]} ; Evolution through Large Models \href{https://arxiv.org/abs/2206.08896}{[arXiv:2206.08896]} ; supp.: Digital Red Queen \href{https://arxiv.org/abs/2601.03335}{[arXiv:2601.03335]} , Promptbreeder \href{https://arxiv.org/abs/2309.16797}{[arXiv:2309.16797]} & -- & \hyperlink{mod-RSN}{\tg{RSN}}\\ \midrule
\hypertarget{lec-UW-19}{}19 & Tue 24 Mar 2026 & \textbf{Multi-agent systems I} & Generative agents, role-playing agents & Generative Agents \href{https://arxiv.org/abs/2304.03442}{[arXiv:2304.03442]} ; CAMEL \href{https://arxiv.org/abs/2303.17760}{[arXiv:2303.17760]} ; supp.: Multiagent Debate \href{https://arxiv.org/abs/2305.14325}{[arXiv:2305.14325]} & -- & \hyperlink{mod-MULTI}{\tg{MULTI}}\\ \midrule
\hypertarget{lec-UW-20}{}20 & Thu 26 Mar 2026 & \textbf{Multi-agent systems II} & Multi-agent frameworks, scaling laws for agent systems, single vs multi-agent & AutoGen \href{https://arxiv.org/abs/2308.08155}{[arXiv:2308.08155]} ; Science of Scaling Agent Systems \href{https://arxiv.org/abs/2512.08296}{[arXiv:2512.08296]} ; supp.: Single agent systems \href{https://openhands.dev/blog/dont-sleep-on-single-agent-systems}{[openhands.dev]} , Prompt Infection \href{https://arxiv.org/abs/2410.07283}{[arXiv:2410.07283]} , MetaGPT \href{https://arxiv.org/abs/2308.00352}{[arXiv:2308.00352]} , AgentVerse \href{https://arxiv.org/abs/2308.10848}{[arXiv:2308.10848]} & -- & \hyperlink{mod-MULTI}{\tg{MULTI}} \hyperlink{mod-ARCH}{\tg{ARCH}}\\ \midrule
\hypertarget{lec-UW-21}{}21 & Tue 07 Apr 2026 & \textbf{Memory I} & OS-style memory hierarchy, cognitive architectures & MemGPT \href{https://arxiv.org/abs/2310.08560}{[arXiv:2310.08560]} ; CoALA \href{https://arxiv.org/abs/2309.02427}{[arXiv:2309.02427]} ; supp.: Think-in-Memory \href{https://arxiv.org/abs/2311.08719}{[arXiv:2311.08719]} , A-Mem \href{https://arxiv.org/pdf/2502.12110}{[arXiv:2502.12110]} , Recursive language models \href{https://arxiv.org/abs/2512.24601}{[arXiv:2512.24601]} & -- & \hyperlink{mod-MEM}{\tg{MEM}} \hyperlink{mod-ARCH}{\tg{ARCH}}\\ \midrule
\hypertarget{lec-UW-22}{}22 & Thu 09 Apr 2026 & \textbf{Memory \& other topics II} & Mixture-of-agents, context folding, long-horizon task measurement & Mixture-of-Agents \href{https://arxiv.org/abs/2406.04692}{[arXiv:2406.04692]} ; Context-Folding \href{https://arxiv.org/pdf/2510.11967}{[arXiv:2510.11967]} ; supp.: More Agents Is All You Need \href{https://arxiv.org/pdf/2402.05120}{[arXiv:2402.05120]} , Measuring AI Ability to Complete Long Tasks \href{https://arxiv.org/pdf/2503.14499}{[arXiv:2503.14499]} & -- & \hyperlink{mod-MEM}{\tg{MEM}} \hyperlink{mod-MULTI}{\tg{MULTI}} \hyperlink{mod-EVAL}{\tg{EVAL}}\\ \midrule
\hypertarget{lec-UW-23}{}23 & Week of 14 Apr 2026 & \textbf{Project work} & -- & -- & -- & \tg{ORG}\\ \midrule
\hypertarget{lec-UW-24}{}24 & Week of 21 Apr 2026 & \textbf{Project presentations} & -- & -- & -- & \tg{ORG}\\ \midrule
\hypertarget{lec-UW-25}{}25 & Week of 28 Apr 2026 & \textbf{Project presentations} & -- & -- & -- & \tg{ORG}\\ \midrule
\end{longtable}}
\section{Assignments and projects}
{\scriptsize\setlength{\tabcolsep}{2.5pt}\begin{longtable}{P{2.8cm} P{5.4cm} P{3.0cm} P{3.5cm} P{1.5cm}}
\toprule Name & What you build & Skills & Starter code / resources & Tags\\ \midrule\endhead
\hypertarget{asg-UW-1}{}\textbf{Build a Coding Agent (buy-in, individual, 15\%)} & A CLI coding agent like Claude Code / Cursor CLI / Gemini CLI, using raw LLM API calls with no agent frameworks (LangChain, Strands, CrewAI banned; thin wrappers like litellm allowed). It must parse and execute tool calls, have a polished UI, a README, and one non-trivial program the agent built & Agent loop, function calling, tool execution, CLI UX & None. Spec: \href{https://github.com/barghouthi/ai-agents/blob/main/assignment.md}{[GitHub]} & \hyperlink{mod-CODE}{\tg{CODE}} \hyperlink{mod-TOOL}{\tg{TOOL}} \hyperlink{mod-ARCH}{\tg{ARCH}}\\ \midrule
\hypertarget{asg-UW-2}{}\textbf{Paper presentation (10\%)} & 25-minute presentation of assigned paper(s) with literature context and critique & Critical reading & -- & \tg{ORG}\\ \midrule
\hypertarget{asg-UW-3}{}\textbf{Paper reviews (10\%)} & Written reviews submitted on Canvas & Critical reading & -- & \tg{ORG}\\ \midrule
\hypertarget{asg-UW-4}{}\textbf{Final research project (50\%)} & Research project on AI agents (pairs encouraged); one-page proposal due Feb 24; write-up and presentation & Research design, agent implementation, evaluation & -- & \hyperlink{mod-ARCH}{\tg{ARCH}} \hyperlink{mod-EVAL}{\tg{EVAL}}\\ \midrule
\end{longtable}}
\chapter{JHU: EN.601.727 Machine Programming}\hypertarget{course-JHU}{}
{\small\begin{tabularx}{\linewidth}{@{}>{\bfseries}l X@{}}\toprule
Institution & Johns Hopkins University\\
Term & Fall 2026 (in progress)\\
Instructors & Ziyang Li\\
Website & \href{https://machine-programming.github.io/}{https://machine-programming.github.io/}\\
Level & Graduate (PL + ML)\\
Archetype & Program synthesis to coding agents\\
Public materials & Slides: L1--12; video: none; code: A1--A3 public\\
Best for & Ties classic program synthesis to LLM coding agents, with two MCP lectures and a strong web-app agent assignment.\\
\bottomrule\end{tabularx}}

Treats code generation from both the symbolic side (enumerative search, types, specifications) and the LLM side (decoding, prompting, iterative refinement, agentic frameworks, MCP, pre- and post-training of code models).
The assignments are the main attraction: an enumerative synthesizer with an LLM step (A1), a rigorous pass@k evaluation of coding LLMs across three languages and four prompting strategies (A2), and a tool-calling web-app agent with secure file tools (A3).
The second half covers formal verification, security, HPC, neurosymbolic programming and student presentations.
\section{Lectures}
{\scriptsize\setlength{\tabcolsep}{2.5pt}\begin{longtable}{P{0.45cm} P{1.2cm} P{2.5cm} P{3.6cm} P{3.6cm} P{3.3cm} P{1.5cm}}
\toprule \# & Date & Title & Key topics & Readings & Materials & Tags\\ \midrule\endhead
\hypertarget{lec-JHU-1}{}1 & Tue Sep 1 & \textbf{Overview \& Introduction to Machine Programming} & Course intro; what machine programming is & -- & \href{https://machine-programming.github.io/lectures/lecture-1.pdf}{[PDF]} & \tg{ORG}\\ \midrule
\hypertarget{lec-JHU-2}{}2 & Thu Sep 3 & \textbf{Syntax, Semantics, and Bottom-up Inductive Synthesis} & DSL grammars/semantics; bottom-up enumeration from examples & -- & \href{https://machine-programming.github.io/lectures/lecture-2.pdf}{[PDF]} & \hyperlink{mod-RSN}{\tg{RSN}}\\ \midrule
\hypertarget{lec-JHU-3}{}3 & Tue Sep 8 & \textbf{Type Systems and Top-down Enumerative Synthesis} & Types for pruning; top-down search & -- & \href{https://machine-programming.github.io/lectures/lecture-3.pdf}{[PDF]} & \hyperlink{mod-RSN}{\tg{RSN}}\\ \midrule
\hypertarget{lec-JHU-4}{}4 & Thu Sep 10 & \textbf{Functional Specifications and Synthesis} & Synthesis from logical/functional specs & -- & \href{https://machine-programming.github.io/lectures/lecture-4.pdf}{[PDF]} & \hyperlink{mod-RSN}{\tg{RSN}}\\ \midrule
\hypertarget{lec-JHU-5}{}5 & Tue Sep 15 & \textbf{Language Modeling for Synthesis} & LMs as program generators & Code LLM readings (readings page) & \href{https://machine-programming.github.io/lectures/lecture-5.pdf}{[PDF]} & \hyperlink{mod-FND}{\tg{FND}} \hyperlink{mod-CODE}{\tg{CODE}}\\ \midrule
\hypertarget{lec-JHU-6}{}6 & Thu Sep 17 & \textbf{Decoding and Prompting for Synthesis} & Decoding strategies; prompting for code & -- & \href{https://machine-programming.github.io/lectures/lecture-6.pdf}{[PDF]} & \hyperlink{mod-FND}{\tg{FND}} \hyperlink{mod-CODE}{\tg{CODE}}\\ \midrule
\hypertarget{lec-JHU-7}{}7 & Tue Sep 22 & \textbf{Iterative Refinement and Agentic Framework} & Feedback/refinement loops; agent framework (A1 due) & -- & \href{https://machine-programming.github.io/lectures/lecture-7.pdf}{[PDF]} & \hyperlink{mod-ARCH}{\tg{ARCH}} \hyperlink{mod-CODE}{\tg{CODE}}\\ \midrule
\hypertarget{lec-JHU-8}{}8 & Thu Sep 24 & \textbf{Agentic Frameworks: Model Context Protocol} & MCP & -- & \href{https://machine-programming.github.io/lectures/lecture-8.pdf}{[PDF]} & \hyperlink{mod-TOOL}{\tg{TOOL}} \hyperlink{mod-ARCH}{\tg{ARCH}}\\ \midrule
\hypertarget{lec-JHU-9}{}9 & Tue Sep 29 & \textbf{Agentic Frameworks: Model Context Protocol (2)} & MCP, continued & -- & \href{https://machine-programming.github.io/lectures/lecture-9.pdf}{[PDF]} & \hyperlink{mod-TOOL}{\tg{TOOL}} \hyperlink{mod-ARCH}{\tg{ARCH}}\\ \midrule
\hypertarget{lec-JHU-10}{}10 & Thu Oct 1 & \textbf{Agentic Frameworks: Selected Topics} & Agent framework topics (unspecified) & Agents readings: MASAI, AutoCodeRover & \href{https://machine-programming.github.io/lectures/lecture-10.pdf}{[PDF]} & \hyperlink{mod-ARCH}{\tg{ARCH}} \hyperlink{mod-CODE}{\tg{CODE}}\\ \midrule
\hypertarget{lec-JHU-11}{}11 & Tue Oct 6 & \textbf{Pre-training of Coding Language Models (1)} & Code-LLM pre-training (data, tokenization, objectives) & e.g., StarCoder, Code Llama, DeepSeek Coder, The Stack, FIM (readings page) & \href{https://machine-programming.github.io/lectures/lecture-11.pdf}{[PDF]} & \hyperlink{mod-TRAIN}{\tg{TRAIN}} \hyperlink{mod-CODE}{\tg{CODE}}\\ \midrule
\hypertarget{lec-JHU-12}{}12 & Thu Oct 8 & \textbf{Pre-training of Coding Language Models (2)} & Pre-training, continued (A2 due Fri Oct 9) & as above & \href{https://machine-programming.github.io/lectures/lecture-12.pdf}{[PDF]} & \hyperlink{mod-TRAIN}{\tg{TRAIN}} \hyperlink{mod-CODE}{\tg{CODE}}\\ \midrule
\hypertarget{lec-JHU-13}{}13 & Tue Oct 13 & \textbf{Post-training of Coding Language Models} & Post-training of code LLMs & e.g., CodeDPO, LintSeq (readings page) & n/a (not yet posted) & \hyperlink{mod-TRAIN}{\tg{TRAIN}} \hyperlink{mod-CODE}{\tg{CODE}}\\ \midrule
\hypertarget{lec-JHU-14}{}14 & Thu Oct 15 & \textbf{Special Topics in Training Language Models} & -- & -- & -- & \hyperlink{mod-TRAIN}{\tg{TRAIN}}\\ \midrule
\hypertarget{lec-JHU-15}{}15 & Tue Oct 20 & \textbf{PL Kaleidoscope: Formal Verification and Theorem Proving} & Verification, theorem proving (e.g., Lean) & -- & -- & \hyperlink{mod-APP}{\tg{APP}} \hyperlink{mod-CODE}{\tg{CODE}}\\ \midrule
\hypertarget{lec-JHU-16}{}16 & Tue Oct 27 & \textbf{PL Kaleidoscope: Differentiable Programming and Machine Learning} & Differentiable programming & -- & -- & \hyperlink{mod-APP}{\tg{APP}}\\ \midrule
\hypertarget{lec-JHU-17}{}17 & Thu Oct 29 & \textbf{PL Kaleidoscope: Security and Privacy} & PL for security/privacy (e.g., CodeQL) & -- & -- & \hyperlink{mod-SAFE}{\tg{SAFE}} \hyperlink{mod-APP}{\tg{APP}}\\ \midrule
\hypertarget{lec-JHU-18}{}18 & Tue Nov 3 & \textbf{PL Kaleidoscope: Hardware and High-Performance Computing} & HW/HPC programming & -- & -- & \hyperlink{mod-APP}{\tg{APP}}\\ \midrule
\hypertarget{lec-JHU-19}{}19 & Thu Nov 5 & \textbf{PL Kaleidoscope: Design, Interactive, and Visual Programming} & Design, interactive and visual programming & -- & -- & \hyperlink{mod-APP}{\tg{APP}}\\ \midrule
\hypertarget{lec-JHU-20}{}20 & Tue Nov 10 & \textbf{Oral Presentations (1)} & Student paper presentations & -- & -- & \tg{ORG}\\ \midrule
\hypertarget{lec-JHU-21}{}21 & Thu Nov 12 & \textbf{Oral Presentations (2)} & Student paper presentations & -- & -- & \tg{ORG}\\ \midrule
\hypertarget{lec-JHU-22}{}22 & Tue Nov 17 & \textbf{Oral Presentations (3)} & Student paper presentations & -- & -- & \tg{ORG}\\ \midrule
\hypertarget{lec-JHU-23}{}23 & Thu Nov 19 & \textbf{Oral Presentations (4)} & Student paper presentations & -- & -- & \tg{ORG}\\ \midrule
\hypertarget{lec-JHU-24}{}24 & Tue Dec 1 & \textbf{Oral Presentations (5)} & Student paper presentations & -- & -- & \tg{ORG}\\ \midrule
\hypertarget{lec-JHU-25}{}25 & Thu Dec 3 & \textbf{Oral Presentations (6)} & Student paper presentations & -- & -- & \tg{ORG}\\ \midrule
\hypertarget{lec-JHU-26}{}26 & Tue Dec 8 & \textbf{Special Topics on Neurosymbolic Programming (1)} & Neurosymbolic programming (e.g., Scallop) & -- & -- & \hyperlink{mod-RSN}{\tg{RSN}} \hyperlink{mod-APP}{\tg{APP}}\\ \midrule
\hypertarget{lec-JHU-27}{}27 & Thu Dec 10 & \textbf{Special Topics on Neurosymbolic Programming (2)} & Neurosymbolic programming, continued & -- & -- & \hyperlink{mod-RSN}{\tg{RSN}} \hyperlink{mod-APP}{\tg{APP}}\\ \midrule
\hypertarget{lec-JHU-28}{}28 & TBD (Dec 15-23) & \textbf{Final Presentation Session} & Final project presentations & -- & -- & \tg{ORG}\\ \midrule
\end{longtable}}
\section{Assignments and projects}
{\scriptsize\setlength{\tabcolsep}{2.5pt}\begin{longtable}{P{2.8cm} P{5.4cm} P{3.0cm} P{3.5cm} P{1.5cm}}
\toprule Name & What you build & Skills & Starter code / resources & Tags\\ \midrule\endhead
\hypertarget{asg-JHU-1}{}\textbf{A1: Inductive Synthesis (due Sep 22)} & A bottom-up enumerative synthesizer for a geometric-shape DSL (grow / eliminate observational equivalents / search); your own string DSL and synthesizer; then prompts plus output parsing so Claude Sonnet 4.6 (via Bedrock) solves at least 80\% of string tasks & Enumerative search, pruning by observational equivalence, DSL design, prompt design, parsing LLM output & \href{https://github.com/machine-programming/assignment-1}{[GitHub]} & \hyperlink{mod-RSN}{\tg{RSN}} \hyperlink{mod-FND}{\tg{FND}} \hyperlink{mod-CODE}{\tg{CODE}}\\ \midrule
\hypertarget{asg-JHU-2}{}\textbf{A2: Evaluating Coding LLMs (due Oct 9)} & Execution harnesses for Python/Rust/OCaml; 4 prompting strategies (zero-shot, two-step CoT, iterative refinement with feedback up to 3 rounds, a custom method of about 5 calls or fewer); pass@1/2/3 pipeline over 24 experiments (Claude Haiku vs Sonnet x 3 languages x 4 strategies) on a 16-item xCodeEval/MBPP set; plots & Sandboxed execution with timeouts, error classification, pass@k, experiment scripting, visualization & \href{https://github.com/machine-programming/assignment-2}{[GitHub]} & \hyperlink{mod-EVAL}{\tg{EVAL}} \hyperlink{mod-CODE}{\tg{CODE}} \hyperlink{mod-RSN}{\tg{RSN}}\\ \midrule
\hypertarget{asg-JHU-3}{}\textbf{A3: Coding LLM Agents (due TBD; repo public)} & "WAA" (Web-App Agent): a Python tool-calling loop (tool registry, system prompt, LLM query, parse, execute, terminate or hit the turn limit), file-system tools with path validation and protected files, TODO tools; the agent then builds a personal site, an Express chat room and a creative app; extra credit adds a database plus DB tools & Agent loop design, tool and termination protocols, secure file tools, error recovery, mock-LLM testing, Jest/Supertest/Playwright & \href{https://github.com/machine-programming/assignment-3}{[GitHub]} & \hyperlink{mod-ARCH}{\tg{ARCH}} \hyperlink{mod-TOOL}{\tg{TOOL}} \hyperlink{mod-CODE}{\tg{CODE}}\\ \midrule
\hypertarget{asg-JHU-4}{}\textbf{Final project (35\%)} & Research project in program synthesis/analysis or related (teams of up to 3): 2-page proposal, implementation, report of at least 4 pages (NeurIPS template), 10-min talk & Research design, implementation, writing & n/a (\href{https://machine-programming.github.io/project.html}{[machine-programming.github.io]}) & \tg{ORG} \hyperlink{mod-CODE}{\tg{CODE}}\\ \midrule
\hypertarget{asg-JHU-5}{}\textbf{Oral presentation (10\%)} & Paper presentation in the Nov-Dec sessions & Paper reading/presenting & -- & \tg{ORG}\\ \midrule
\end{longtable}}
\chapter{CU Boulder: CSCI 7000-005 Building AI Agents}\hypertarget{course-CU}{}
{\small\begin{tabularx}{\linewidth}{@{}>{\bfseries}l X@{}}\toprule
Institution & University of Colorado Boulder\\
Term & Spring 2026\\
Instructors & Danny Dig\\
Website & \href{https://danny.cs.colorado.edu/courses/csci7000-005_Sp26/index.html}{https://danny.cs.colorado.edu/courses/csci7000-005\_Sp26/index.html}\\
Level & Graduate seminar; competent SE assumed\\
Archetype & Research seminar (agentic software engineering)\\
Public materials & Slides: 4 instructor decks; video: none; code: none\\
Best for & Research reading on agentic software engineering: refactoring agents, SWE-agent, skills, memory, MCP design choices.\\
\bottomrule\end{tabularx}}

A seminar with two papers per meeting, one-page critiques and a team research project with four milestones.
Its reading list leans toward agentic software engineering, from agent surveys and an agentic-SE roadmap through refactoring agents and the classic patterns (ReAct, Reflexion, Toolformer, Gorilla, PAL, AutoGen, MetaGPT, debate) to 2025--2026 work on skills, memory, safety, coding-agent configuration and MCP design.
The ``tool experience report'' (run a paper's tool on your own case and critique it) is a good habit to borrow.
\section{Lectures}
{\scriptsize\setlength{\tabcolsep}{2.5pt}\begin{longtable}{P{0.45cm} P{1.2cm} P{2.5cm} P{3.6cm} P{3.6cm} P{3.3cm} P{1.5cm}}
\toprule \# & Date & Title & Key topics & Readings & Materials & Tags\\ \midrule\endhead
\hypertarget{lec-CU-1}{}1 & Mon 01/12 & \textbf{Course introduction (Dig)} & Logistics; how to read research papers & Griswold, "How to Read an Engineering Research Paper" \href{http://cseweb.ucsd.edu/\%7Ewgg/CSE210/howtoread.html}{[cseweb.ucsd.edu]} & \href{https://danny.cs.colorado.edu/courses/csci7000-005_Sp26/pdfs/CSCI_7000_005_W1L1.pdf}{[PDF]} & \tg{ORG}\\ \midrule
\hypertarget{lec-CU-2}{}2 & Wed 01/14 & \textbf{Survey of agents (Dig)} & LLM agent surveys: general and SE & Wang et al., LLM-based autonomous agents survey \href{https://link.springer.com/article/10.1007/s11704-024-40231-1}{[link.springer.com]} ; Liu et al., LLM-Based Agents for SE survey \href{https://arxiv.org/pdf/2409.02977}{[arXiv:2409.02977]} & \href{https://danny.cs.colorado.edu/courses/csci7000-005_Sp26/pdfs/W1_L2_SurveyOfAgents.pdf}{[PDF]} & \hyperlink{mod-ARCH}{\tg{ARCH}} \hyperlink{mod-CODE}{\tg{CODE}}\\ \midrule
\hypertarget{lec-CU-3}{}3 & Wed 01/21 & \textbf{Agentic SE roadmap (Dig)} & Foundational pillars and research roadmap for agentic SE & Hassan et al., Agentic Software Engineering \href{https://arxiv.org/pdf/2509.06216}{[arXiv:2509.06216]} & \href{https://danny.cs.colorado.edu/courses/csci7000-005_Sp26/pdfs/W2_L2_Roadmap.pdf}{[PDF]} & \hyperlink{mod-CODE}{\tg{CODE}} \hyperlink{mod-ARCH}{\tg{ARCH}}\\ \midrule
\hypertarget{lec-CU-4}{}4 & Mon 01/26 & \textbf{Refactoring agents (students)} & Multi-agent rename refactoring; refactoring benchmark & Bellur et al., Multi-Agent Coordinated Rename Refactoring \href{https://arxiv.org/abs/2601.00482}{[arXiv:2601.00482]} ; Gautam et al., RefactorBench \href{https://openreview.net/pdf?id=NiNIthntx7}{[openreview.net]} & -- & \hyperlink{mod-CODE}{\tg{CODE}} \hyperlink{mod-MULTI}{\tg{MULTI}} \hyperlink{mod-EVAL}{\tg{EVAL}}\\ \midrule
\hypertarget{lec-CU-5}{}5 & Wed 01/28 & \textbf{MM-Assist \& MANTRA (Dig)} & LLM-assisted refactoring (move method); multi-agent refactoring & Bellur et al., Together We Are Better \href{https://danny.cs.colorado.edu/papers/ICSME25_MM-Assist.pdf}{[PDF]} ; Xu et al., MANTRA \href{https://arxiv.org/abs/2503.14340}{[arXiv:2503.14340]} & \href{https://danny.cs.colorado.edu/courses/csci7000-005_Sp26/pdfs/W3_L2_MM-Mantra.pdf}{[PDF]} & \hyperlink{mod-CODE}{\tg{CODE}} \hyperlink{mod-MULTI}{\tg{MULTI}}\\ \midrule
\hypertarget{lec-CU-6}{}6 & Mon 02/02 & \textbf{Project idea pitches} & Students pitch projects & -- & \href{https://danny.cs.colorado.edu/courses/csci7000-005_Sp26/projectProposal.html}{[danny.cs.colorado.edu]} & \tg{ORG}\\ \midrule
\hypertarget{lec-CU-7}{}7 & Wed 02/04 & \textbf{Project Bazaar (1), in person} & Team formation & -- & \href{https://danny.cs.colorado.edu/courses/csci7000-005_Sp26/projectTeams.html}{[danny.cs.colorado.edu]} & \tg{ORG}\\ \midrule
\hypertarget{lec-CU-8}{}8 & Mon 02/09 & \textbf{Project Bazaar (2), in person} & Team formation & -- & (same as above) & \tg{ORG}\\ \midrule
\hypertarget{lec-CU-9}{}9 & Wed 02/11 & \textbf{ReAct \& Web API code gen (students)} & Reasoning+acting loop; API integration code benchmark & Yao et al., ReAct \href{https://arxiv.org/abs/2210.03629}{[arXiv:2210.03629]} ; Maninger et al., Benchmarking Web API Integration Code Generation \href{https://arxiv.org/pdf/2509.20172v4}{[arXiv:2509.20172]} & -- & \hyperlink{mod-RSN}{\tg{RSN}} \hyperlink{mod-TOOL}{\tg{TOOL}} \hyperlink{mod-EVAL}{\tg{EVAL}}\\ \midrule
\hypertarget{lec-CU-10}{}10 & Mon 02/16 & \textbf{Milestone 1 presentations} & Research questions, challenges, approach, related work & -- & \href{https://danny.cs.colorado.edu/courses/csci7000-005_Sp26/milestone1.html}{[danny.cs.colorado.edu]} & \tg{ORG}\\ \midrule
\hypertarget{lec-CU-11}{}11 & Wed 02/18 & \textbf{Reflexion \& Toolformer (students)} & Verbal self-reflection; self-taught tool use & Shinn et al., Reflexion \href{https://arxiv.org/abs/2303.11366}{[arXiv:2303.11366]} ; Schick et al., Toolformer \href{https://arxiv.org/pdf/2302.04761}{[arXiv:2302.04761]} & -- & \hyperlink{mod-RSN}{\tg{RSN}} \hyperlink{mod-TOOL}{\tg{TOOL}} \hyperlink{mod-TRAIN}{\tg{TRAIN}}\\ \midrule
\hypertarget{lec-CU-12}{}12 & Mon 02/23 & \textbf{Gorilla \& PAL (students)} & API-call LLMs; program-aided reasoning & Patil et al., Gorilla \href{https://arxiv.org/pdf/2305.15334}{[arXiv:2305.15334]} ; Gao et al., PAL \href{https://arxiv.org/pdf/2211.10435}{[arXiv:2211.10435]} & -- & \hyperlink{mod-TOOL}{\tg{TOOL}} \hyperlink{mod-RSN}{\tg{RSN}}\\ \midrule
\hypertarget{lec-CU-13}{}13 & Wed 02/25 & \textbf{SWE-agent \& AutoCodeRover (students)} & Agent-computer interfaces; repo-level program repair & Yang et al., SWE-agent \href{https://arxiv.org/pdf/2405.15793}{[arXiv:2405.15793]} ; Zhang et al., AutoCodeRover \href{https://arxiv.org/pdf/2404.05427}{[arXiv:2404.05427]} & -- & \hyperlink{mod-CODE}{\tg{CODE}} \hyperlink{mod-ARCH}{\tg{ARCH}}\\ \midrule
\hypertarget{lec-CU-14}{}14 & Mon 03/02 & \textbf{AutoGen \& MetaGPT (students)} & Multi-agent conversation frameworks; SOP-driven role agents & Wu et al., AutoGen \href{https://arxiv.org/pdf/2308.08155}{[arXiv:2308.08155]} ; Hong et al., MetaGPT \href{https://arxiv.org/pdf/2308.00352}{[arXiv:2308.00352]} & -- & \hyperlink{mod-MULTI}{\tg{MULTI}} \hyperlink{mod-ARCH}{\tg{ARCH}}\\ \midrule
\hypertarget{lec-CU-15}{}15 & Wed 03/04 & \textbf{Multiagent debate \& Agentless (students)} & Debate for factuality/reasoning; demystifying SE agents & Du et al., Multiagent Debate \href{https://openreview.net/pdf?id=zj7YuTE4t8}{[openreview.net]} ; Xia et al., Demystifying LLM-Based SE Agents \href{https://dl.acm.org/doi/epdf/10.1145/3715754}{[dl.acm.org]} & -- & \hyperlink{mod-MULTI}{\tg{MULTI}} \hyperlink{mod-RSN}{\tg{RSN}} \hyperlink{mod-CODE}{\tg{CODE}}\\ \midrule
\hypertarget{lec-CU-16}{}16 & Mon 03/09 & \textbf{LangGraph tutorial \& OpenHands} & Hands-on LangGraph; open coding-agent platform; tool experience report & Wang et al., OpenHands \href{https://arxiv.org/pdf/2407.16741}{[arXiv:2407.16741]} & -- & \hyperlink{mod-ARCH}{\tg{ARCH}} \hyperlink{mod-CODE}{\tg{CODE}}\\ \midrule
\hypertarget{lec-CU-17}{}17 & Wed 03/11 & \textbf{Milestone 2 presentations} & Project progress & -- & \href{https://danny.cs.colorado.edu/courses/csci7000-005_Sp26/milestone2.html}{[danny.cs.colorado.edu]} & \tg{ORG}\\ \midrule
\hypertarget{lec-CU-18}{}18 & Mon 03/23 & \textbf{RefAgent \& Agentic Skills (students)} & Multi-agent refactoring; SoK on agent skills & Oueslati et al., RefAgent \href{https://arxiv.org/pdf/2511.03153}{[arXiv:2511.03153]} ; Jiang et al., SoK: Agentic Skills \href{https://arxiv.org/pdf/2602.20867}{[arXiv:2602.20867]} & -- & \hyperlink{mod-CODE}{\tg{CODE}} \hyperlink{mod-ARCH}{\tg{ARCH}} \hyperlink{mod-TOOL}{\tg{TOOL}}\\ \midrule
\hypertarget{lec-CU-19}{}19 & Wed 03/25 & \textbf{AppAgent \& coding-agent configuration (students)} & Smartphone GUI agent; config files of AI coding agents & Zhang et al., AppAgent \href{https://arxiv.org/pdf/2312.13771}{[arXiv:2312.13771]} ; Santos et al., Decoding the Configuration of AI Coding Agents \href{https://arxiv.org/pdf/2511.09268}{[arXiv:2511.09268]} & -- & \hyperlink{mod-WEB}{\tg{WEB}} \hyperlink{mod-CODE}{\tg{CODE}}\\ \midrule
\hypertarget{lec-CU-20}{}20 & Mon 03/30 & \textbf{Milestone 3 presentations} & Project progress & -- & \href{https://danny.cs.colorado.edu/courses/csci7000-005_Sp26/milestone3.html}{[danny.cs.colorado.edu]} & \tg{ORG}\\ \midrule
\hypertarget{lec-CU-21}{}21 & Wed 04/01 & \textbf{MemSkill \& More with Less (students)} & Memory skills for agents; efficiency & Zhang et al., MemSkill \href{https://arxiv.org/abs/2602.02474}{[arXiv:2602.02474]} ; Gao \& Peng, More with Less \href{https://arxiv.org/pdf/2510.16786}{[arXiv:2510.16786]} & -- & \hyperlink{mod-MEM}{\tg{MEM}} \hyperlink{mod-DEPLOY}{\tg{DEPLOY}}\\ \midrule
\hypertarget{lec-CU-22}{}22 & Mon 04/06 & \textbf{MetaClaw \& code smells (students)} & MetaClaw; causal analysis of smells in LLM code & Xia et al., MetaClaw \href{https://arxiv.org/pdf/2603.17187}{[arXiv:2603.17187]} ; Velasco et al., Causal Perspective on Smells in LLM-Generated Code \href{https://arxiv.org/abs/2511.15817}{[arXiv:2511.15817]} & -- & \hyperlink{mod-CODE}{\tg{CODE}} \hyperlink{mod-EVAL}{\tg{EVAL}}\\ \midrule
\hypertarget{lec-CU-23}{}23 & Wed 04/08 & \textbf{SafePro \& ToolDreamer (students)} & Agent safety; tool retrieval/selection & Zhou et al., SafePro \href{https://arxiv.org/pdf/2601.06663}{[arXiv:2601.06663]} ; Sengupta et al., ToolDreamer \href{https://arxiv.org/pdf/2510.19791v2}{[arXiv:2510.19791]} & -- & \hyperlink{mod-SAFE}{\tg{SAFE}} \hyperlink{mod-TOOL}{\tg{TOOL}}\\ \midrule
\hypertarget{lec-CU-24}{}24 & Mon 04/13 & \textbf{Swiss Cheese safety \& MCP design (students)} & Layered AI-safety model; tool orchestration vs code execution in MCP & Shamsujjoha et al., Swiss Cheese Model for AI Safety \href{https://arxiv.org/abs/2408.02205}{[arXiv:2408.02205]} ; Felendler et al., From Tool Orchestration to Code Execution: MCP Design Choices \href{https://arxiv.org/abs/2602.15945}{[arXiv:2602.15945]} & -- & \hyperlink{mod-SAFE}{\tg{SAFE}} \hyperlink{mod-TOOL}{\tg{TOOL}}\\ \midrule
\hypertarget{lec-CU-25}{}25 & Wed 04/15 & \textbf{Guest talk: Ameya Ketkar (Uber), building AI agents in industry} & Industry agents; multi-agent SRE incident analysis & Kataria, Intelligent Site Reliability Engineering \href{https://www.researchgate.net/publication/399119176_Intelligent_Site_Reliability_Engineering_A_Multi-agent_LLM_Framework_for_Automated_Incident_Analysis_and_Root_Cause_Determination}{[researchgate.net]} & -- & \hyperlink{mod-DEPLOY}{\tg{DEPLOY}} \hyperlink{mod-MULTI}{\tg{MULTI}} \hyperlink{mod-APP}{\tg{APP}}\\ \midrule
\hypertarget{lec-CU-26}{}26 & Mon 04/20 & \textbf{Class summary and insights (Dig), in person} & Course wrap-up & -- & -- & \tg{ORG}\\ \midrule
\hypertarget{lec-CU-27}{}27 & Wed 04/22 & \textbf{Milestone 4 presentations} & Empirical evaluation results & -- & \href{https://danny.cs.colorado.edu/courses/csci7000-005_Sp26/milestone4.html}{[danny.cs.colorado.edu]} & \tg{ORG}\\ \midrule
\hypertarget{lec-CU-28}{}28 & Fri 04/24 & \textbf{Project Clinic (Dig)} & Project help & -- & -- & \tg{ORG}\\ \midrule
\end{longtable}}
\section{Assignments and projects}
{\scriptsize\setlength{\tabcolsep}{2.5pt}\begin{longtable}{P{2.8cm} P{5.4cm} P{3.0cm} P{3.5cm} P{1.5cm}}
\toprule Name & What you build & Skills & Starter code / resources & Tags\\ \midrule\endhead
\hypertarget{asg-CU-1}{}\textbf{Paper critiques (20\%)} & One-page critique per paper (liked / disliked / future directions / open questions / discussion points / summary), due 12pm on the day of class & Critical reading of research & -- & \tg{ORG}\\ \midrule
\hypertarget{asg-CU-2}{}\textbf{Tool Experience Report (critique option \#2)} & Download the tool a paper introduces, run it on your own use case, and write a one-page reflection (experience, strengths, weaknesses, practical impact, recommendations) & Hands-on evaluation of agent tools & -- & \hyperlink{mod-EVAL}{\tg{EVAL}} \hyperlink{mod-CODE}{\tg{CODE}}\\ \midrule
\hypertarget{asg-CU-3}{}\textbf{Paper presentation (20\%)} & Lead the discussion of an assigned paper & Presenting research & -- & \tg{ORG}\\ \midrule
\hypertarget{asg-CU-4}{}\textbf{Research project (50\%, teams of 3-4)} & A research-oriented or novel-industrial project on GenAI/agents in SE: proposal (02/02), M1 2-page ACM paper plus 10-min talk (RQs, challenges, approach, 5 or more related works), M2, M3, M4 (empirical evaluation), final conference-style paper (abstract, intro, 10 or more refs, conclusions, lessons learned) & Research design, agent implementation, empirical evaluation, scientific writing & -- & \hyperlink{mod-CODE}{\tg{CODE}} \hyperlink{mod-EVAL}{\tg{EVAL}} \tg{ORG}\\ \midrule
\end{longtable}}
\chapter{Mannheim: IE 685/686 Large Language Models and Agents}\hypertarget{course-MANN}{}
{\small\begin{tabularx}{\linewidth}{@{}>{\bfseries}l X@{}}\toprule
Institution & Universit\"at Mannheim (DWS)\\
Term & Spring 2026\\
Instructors & Ralph Peeters, Christian Bizer, Aaron Steiner\\
Website & \href{https://www.uni-mannheim.de/dws/teaching/course-details/courses-for-master-candidates/ie-685-large-language-models-and-agents/}{https://www.uni-mannheim.de/dws/teaching/course-details/courses-for-master-candidates/ie-685-large-language-models-and-agents/}\\
Level & Master's; Python + basic ML\\
Archetype & Applied master's course + team project\\
Public materials & Slides: 8 PDF decks; video: none; code: 5 exercise notebooks (zip)\\
Best for & Compact slide decks from LLM basics to context engineering, plus nine well-specified team-project briefs.\\
\bottomrule\end{tabularx}}

A short lecture block (introduction, post-training, prompting, agents and tool use with MCP, RAG, context engineering with deep agents, safety), each with an exercise notebook, followed by a team project.
The nine project briefs (multi-agent gaming, QA agents, job hunting, a shopping assistant, a tutor, a personal assistant, news-reading personas, Wikipedia generation, skill mining) are ready-made portfolio ideas, and each comes with suggested benchmarks.
\section{Lectures}
{\scriptsize\setlength{\tabcolsep}{2.5pt}\begin{longtable}{P{0.45cm} P{1.2cm} P{2.5cm} P{3.6cm} P{3.6cm} P{3.3cm} P{1.5cm}}
\toprule \# & Date & Title & Key topics & Readings & Materials & Tags\\ \midrule\endhead
\hypertarget{lec-MANN-1}{}1 & Tue 10 Feb 2026 & \textbf{Introduction (to LLMs)} & Introduction to LLMs, course organization & Zhao et al., A Survey of LLMs \href{https://arxiv.org/abs/2303.18223}{[arXiv:2303.18223]} ; Zhou et al., Survey on Pretrained Foundation Models \href{https://link.springer.com/article/10.1007/s13042-024-02443-6}{[link.springer.com]} & Slides \href{https://www.uni-mannheim.de/media/Einrichtungen/dws/Files_Teaching/Large_Language_Models_and_Agents/FSS2026/IE685_LA_01_Introduction.pdf}{[PDF]} & \hyperlink{mod-FND}{\tg{FND}} \tg{ORG}\\ \midrule
\hypertarget{lec-MANN-2}{}2 & Tue 17 Feb 2026 & \textbf{Post-training and Efficient Adaptation} & Post-training (instruction tuning / alignment), parameter-efficient adaptation & -- & Slides \href{https://www.uni-mannheim.de/media/Einrichtungen/dws/Files_Teaching/Large_Language_Models_and_Agents/FSS2026/IE685_LA_02_PosttrainingAndEfficientAdaptation.pdf}{[PDF]} ; exercise \href{https://www.uni-mannheim.de/media/Einrichtungen/dws/Files_Teaching/Large_Language_Models_and_Agents/FSS2026/IE685_LA_EX_LLM_Introduction.ipynb.zip}{[notebook zip]} & \hyperlink{mod-TRAIN}{\tg{TRAIN}} \hyperlink{mod-FND}{\tg{FND}}\\ \midrule
\hypertarget{lec-MANN-3}{}3 & Tue 24 Feb 2026 & \textbf{Prompt Engineering} & Prompting techniques & -- & Slides \href{https://www.uni-mannheim.de/media/Einrichtungen/dws/Files_Teaching/Large_Language_Models_and_Agents/FSS2026/IE685_LA_03_PromptEngineering.pdf}{[PDF]} ; exercise \href{https://www.uni-mannheim.de/media/Einrichtungen/dws/Files_Teaching/Large_Language_Models_and_Agents/FSS2026/IE686_LA_EX_Langchain_Introduction.ipynb.zip}{[notebook zip]} & \hyperlink{mod-FND}{\tg{FND}} \hyperlink{mod-RSN}{\tg{RSN}}\\ \midrule
\hypertarget{lec-MANN-4}{}4 & Tue 03 Mar 2026 & \textbf{LLM Agents and Tool Use} & LLM-based agents, tool use and environment interaction, MCP & Wang et al., Survey on LLM-based Autonomous Agents \href{https://link.springer.com/article/10.1007/s11704-024-40231-1}{[link.springer.com]} & Slides (78 pp.) \href{https://www.uni-mannheim.de/media/Einrichtungen/dws/Files_Teaching/Large_Language_Models_and_Agents/FSS2026/IE685_LA_04_LLMAgentsAndToolUse.pdf}{[PDF]} ; exercise \href{https://www.uni-mannheim.de/media/Einrichtungen/dws/Files_Teaching/Large_Language_Models_and_Agents/FSS2026/IE685_LA_EX_Agent_Programming_LangChain_MCP.ipynb.zip}{[notebook zip]} & \hyperlink{mod-ARCH}{\tg{ARCH}} \hyperlink{mod-TOOL}{\tg{TOOL}}\\ \midrule
\hypertarget{lec-MANN-5}{}5 & Tue 10 Mar 2026 & \textbf{Retrieval Augmented Generation} & RAG & -- & Slides \href{https://www.uni-mannheim.de/media/Einrichtungen/dws/Files_Teaching/Large_Language_Models_and_Agents/FSS2026/IE685_LA_05_RetrievalAugmentedGeneration.pdf}{[PDF]} ; exercise \href{https://www.uni-mannheim.de/media/Einrichtungen/dws/Files_Teaching/Large_Language_Models_and_Agents/FSS2026/IE685_LA_EX_Retrieval_Augmented_Generation.zip}{[zip]} & \hyperlink{mod-MEM}{\tg{MEM}}\\ \midrule
\hypertarget{lec-MANN-6}{}6 & Tue 17 Mar 2026 & \textbf{Context Engineering} & Context engineering for LLM agents; deep agents & -- & Slides \href{https://www.uni-mannheim.de/media/Einrichtungen/dws/Files_Teaching/Large_Language_Models_and_Agents/FSS2026/IE685_LA_06_ContextEngineering.pdf}{[PDF]} ; exercise \href{https://www.uni-mannheim.de/media/Einrichtungen/dws/Files_Teaching/Large_Language_Models_and_Agents/FSS2026/deep_agent.zip}{[zip]} & \hyperlink{mod-MEM}{\tg{MEM}} \hyperlink{mod-ARCH}{\tg{ARCH}}\\ \midrule
\hypertarget{lec-MANN-7}{}7 & Tue 24 Mar 2026 & \textbf{Introduction to Student Projects (IE686)} & Project goals, procedure, 9 topic areas + own idea, schedule & -- & Slides \href{https://www.uni-mannheim.de/media/Einrichtungen/dws/Files_Teaching/Large_Language_Models_and_Agents/FSS2026/IE686_LA_07_IntroStudentProjects.pdf}{[PDF]} & \tg{ORG}\\ \midrule
\hypertarget{lec-MANN-8}{}8 & Tue 14 Apr 2026 & \textbf{Project feedback and coaching} & -- & -- & -- & \tg{ORG}\\ \midrule
\hypertarget{lec-MANN-9}{}9 & Tue 21 Apr 2026 & \textbf{LLM Security and Safety (+ project coaching)} & Safety and security of LLM agents & Wang et al., Full Stack Safety survey \href{https://arxiv.org/abs/2504.15585}{[arXiv:2504.15585]} & Slides \href{https://www.uni-mannheim.de/media/Einrichtungen/dws/Files_Teaching/Large_Language_Models_and_Agents/FSS2026/IE685_LA_08_SecurityAndSafety.pdf}{[PDF]} & \hyperlink{mod-SAFE}{\tg{SAFE}}\\ \midrule
\hypertarget{lec-MANN-10}{}10 & 28 Apr, 05 May, 12 May 2026 & \textbf{Project work and coaching} & n/a (page lists "28.05.2026", likely a typo for 28.04) & -- & -- & \tg{ORG}\\ \midrule
\hypertarget{lec-MANN-11}{}11 & 19-20 May 2026 & \textbf{Presentation of project results} & -- & -- & -- & \tg{ORG}\\ \midrule
\hypertarget{lec-MANN-x1}{}x1 & (topic only, no slide deck found) & \textbf{Evaluation of LLMs and LLM agents} & Listed as syllabus topic 8 & Mohammadi et al., Evaluation and benchmarking of LLM agents (KDD 2025) \href{https://dl.acm.org/doi/10.1145/3711896.3736570}{[dl.acm.org]} & -- & \hyperlink{mod-EVAL}{\tg{EVAL}}\\ \midrule
\end{longtable}}
\section{Assignments and projects}
{\scriptsize\setlength{\tabcolsep}{2.5pt}\begin{longtable}{P{2.8cm} P{5.4cm} P{3.0cm} P{3.5cm} P{1.5cm}}
\toprule Name & What you build & Skills & Starter code / resources & Tags\\ \midrule\endhead
\hypertarget{asg-MANN-1}{}\textbf{Exercise: LLM Introduction} & Intro notebook for working with LLMs & LLM APIs/basics & \href{https://www.uni-mannheim.de/media/Einrichtungen/dws/Files_Teaching/Large_Language_Models_and_Agents/FSS2026/IE685_LA_EX_LLM_Introduction.ipynb.zip}{[notebook zip]} & \hyperlink{mod-FND}{\tg{FND}}\\ \midrule
\hypertarget{asg-MANN-2}{}\textbf{Exercise: LangChain Introduction} & LangChain basics notebook & LangChain & \href{https://www.uni-mannheim.de/media/Einrichtungen/dws/Files_Teaching/Large_Language_Models_and_Agents/FSS2026/IE686_LA_EX_Langchain_Introduction.ipynb.zip}{[notebook zip]} & \hyperlink{mod-FND}{\tg{FND}} \hyperlink{mod-ARCH}{\tg{ARCH}}\\ \midrule
\hypertarget{asg-MANN-3}{}\textbf{Exercise: Agent Programming with LangChain + MCP} & Tool-using agents with LangChain and MCP & Agents, tool calling, MCP & \href{https://www.uni-mannheim.de/media/Einrichtungen/dws/Files_Teaching/Large_Language_Models_and_Agents/FSS2026/IE685_LA_EX_Agent_Programming_LangChain_MCP.ipynb.zip}{[notebook zip]} & \hyperlink{mod-TOOL}{\tg{TOOL}} \hyperlink{mod-ARCH}{\tg{ARCH}}\\ \midrule
\hypertarget{asg-MANN-4}{}\textbf{Exercise: RAG} & RAG pipeline notebook & Retrieval, RAG & \href{https://www.uni-mannheim.de/media/Einrichtungen/dws/Files_Teaching/Large_Language_Models_and_Agents/FSS2026/IE685_LA_EX_Retrieval_Augmented_Generation.zip}{[zip]} & \hyperlink{mod-MEM}{\tg{MEM}}\\ \midrule
\hypertarget{asg-MANN-5}{}\textbf{Exercise: Deep Agent} & Deep-agent / context-engineering exercise & Context engineering, planning agents & \href{https://www.uni-mannheim.de/media/Einrichtungen/dws/Files_Teaching/Large_Language_Models_and_Agents/FSS2026/deep_agent.zip}{[zip]} & \hyperlink{mod-MEM}{\tg{MEM}} \hyperlink{mod-ARCH}{\tg{ARCH}}\\ \midrule
\hypertarget{asg-MANN-6}{}\textbf{Project topic 1: Multi-Agent Gaming} & Agent teams that cooperate in, or compete at, games (Codenames worked well before) & Multi-agent coordination, eval by win rate/ELO & AgentBench \href{https://github.com/THUDM/AgentBench}{[GitHub]} ; SmartPlay \href{https://github.com/microsoft/SmartPlay}{[GitHub]} & \hyperlink{mod-MULTI}{\tg{MULTI}} \hyperlink{mod-EVAL}{\tg{EVAL}}\\ \midrule
\hypertarget{asg-MANN-7}{}\textbf{Project topic 2: QA Agents} & RAG (multi-)agent QA over web search that explains its evidence & RAG, web search, benchmark eval & BrowseComp \href{https://openai.com/index/browsecomp/}{[openai.com]} ; Meta CRAG KDD Cup 2024 \href{https://www.aicrowd.com/challenges/meta-comprehensive-rag-benchmark-kdd-cup-2024}{[aicrowd.com]} & \hyperlink{mod-MEM}{\tg{MEM}} \hyperlink{mod-WEB}{\tg{WEB}} \hyperlink{mod-EVAL}{\tg{EVAL}}\\ \midrule
\hypertarget{asg-MANN-8}{}\textbf{Project topic 3: Job Hunting Agents} & From a query plus CV, search job APIs, extract facts, present structured results & Tool/API use, extraction & \href{https://apislist.com/category/28/jobs}{[apislist.com]} & \hyperlink{mod-TOOL}{\tg{TOOL}} \hyperlink{mod-APP}{\tg{APP}}\\ \midrule
\hypertarget{asg-MANN-9}{}\textbf{Project topic 4: Online Shopping Assistant} & Product search/recommendation agent with structured feature comparison & Web agents, benchmarks & AgentBench ; WebArena \href{https://webarena.dev/}{[webarena.dev]} ; WebShop \href{https://webshop-pnlp.github.io/}{[webshop-pnlp.github.io]} & \hyperlink{mod-WEB}{\tg{WEB}} \hyperlink{mod-APP}{\tg{APP}} \hyperlink{mod-EVAL}{\tg{EVAL}}\\ \midrule
\hypertarget{asg-MANN-10}{}\textbf{Project topic 5: Personal Tutor} & Agent that explains concepts, finds knowledge gaps, tracks progress & Memory, personalization & -- & \hyperlink{mod-APP}{\tg{APP}} \hyperlink{mod-MEM}{\tg{MEM}}\\ \midrule
\hypertarget{asg-MANN-11}{}\textbf{Project topic 6: Personal Assistant} & Agent that manages calendar/mail/chat & Tool use, personal-data integration & -- & \hyperlink{mod-TOOL}{\tg{TOOL}} \hyperlink{mod-APP}{\tg{APP}}\\ \midrule
\hypertarget{asg-MANN-12}{}\textbf{Project topic 7: Personas Reading the News} & Persona agents browsing political news; study effects on memory/beliefs & Agent simulation, memory & Personas provided by instructors (not public) & \hyperlink{mod-MULTI}{\tg{MULTI}} \hyperlink{mod-MEM}{\tg{MEM}} \hyperlink{mod-APP}{\tg{APP}}\\ \midrule
\hypertarget{asg-MANN-13}{}\textbf{Project topic 8: Wikipedia Article Generation} & RAG/agent system to generate Wikipedia articles, evaluated against real articles with rubrics & RAG, LLM-as-judge & WikiGenBench, WikiAutoGen (named, no URLs) & \hyperlink{mod-MEM}{\tg{MEM}} \hyperlink{mod-EVAL}{\tg{EVAL}} \hyperlink{mod-APP}{\tg{APP}}\\ \midrule
\hypertarget{asg-MANN-14}{}\textbf{Project topic 9: Skill-Mining from Public Data} & RAG system that summarizes developer skills/roles from GitHub profiles & RAG, extraction & GitHub profiles & \hyperlink{mod-MEM}{\tg{MEM}} \hyperlink{mod-APP}{\tg{APP}}\\ \midrule
\hypertarget{asg-MANN-15}{}\textbf{Project topic X: Own idea} & Student-proposed topic, subject to approval & -- & -- & \hyperlink{mod-APP}{\tg{APP}}\\ \midrule
\end{longtable}}
\chapter{UChicago: AI Agents for Social Science and Society}\hypertarget{course-UCHI}{}
{\small\begin{tabularx}{\linewidth}{@{}>{\bfseries}l X@{}}\toprule
Institution & University of Chicago (KnowledgeLab)\\
Term & Winter 2026\\
Instructors & James A. Evans\\
Website & \href{https://github.com/KnowledgeLab/AI-Agents-for-Social-Science-and-Society-2026}{https://github.com/KnowledgeLab/AI-Agents-for-Social-Science-and-Society-2026}\\
Level & Grad / advanced undergrad (computational social science)\\
Archetype & Applied: agents as research instruments\\
Public materials & Slides: none; video: none; code: 9 weekly notebooks\\
Best for & Agents as simulated subjects and research assistants; notebooks on multi-agent simulation, LoRA, RL, steering and red-teaming.\\
\bottomrule\end{tabularx}}

Treats LLM agents as tools for social science: simulated subjects and ``digital doubles'', research assistants and automated scientists.
Agent-engineering content is concentrated in week~3 (prompted multi-agent simulation with AutoGen and Concordia, RAG) and week~5 (LoRA fine-tuning, benchmarking, tools and MCP). Other weeks cover deep-learning foundations, causal designs, interpretability and steering, RL, multimodal agents, and alignment.
Two notebooks (weeks 6 and 8) were too large to inspect.
\section{Lectures}
{\scriptsize\setlength{\tabcolsep}{2.5pt}\begin{longtable}{P{0.45cm} P{1.2cm} P{2.5cm} P{3.6cm} P{3.6cm} P{3.3cm} P{1.5cm}}
\toprule \# & Date & Title & Key topics & Readings & Materials & Tags\\ \midrule\endhead
\hypertarget{lec-UCHI-0}{}0 & Self-paced & \textbf{Introduction / Review} & Python, math review, Claude Code setup & Python tutorial \href{https://docs.python.org/3/tutorial/}{[docs.python.org]} ; Yang et al. "SWE-agent" (2024); *Deep Learning* ch. 2-3; Claude Code setup tutorial (Canvas) & -- & \tg{ORG} \hyperlink{mod-CODE}{\tg{CODE}}\\ \midrule
\hypertarget{lec-UCHI-1}{}1 & 2026-01-09 & \textbf{Deep Learning and Social Agents} & Deep vs. shallow nets; links between DL, statistical learning, and social-science agent models & Evans *Thinking with Deep Learning* (preface, ch. 1, 3); Cheng et al. "Exploring LLM-Based Intelligent Agents"; Bubeck et al. "Sparks of AGI"; Farrell, Gopnik, Shalizi, Evans "Large AI Models are Cultural and Social Technologies" (Science 2025); Park et al. "Generative Agents" (UIST 2023) & Notebook \href{https://github.com/KnowledgeLab/AI-Agents-for-Social-Science-and-Society-2026/blob/main/Tutorials-Homework_Notebooks/Week_1_Intro_NNs.ipynb}{[GitHub]} & \hyperlink{mod-FND}{\tg{FND}} \hyperlink{mod-APP}{\tg{APP}}\\ \midrule
\hypertarget{lec-UCHI-2}{}2 & 2026-01-16 & \textbf{Text Learning, Transformers, and Diffusion Models} & word2vec to transformers to diffusion; CoT, ToT, graph-of-thought & Vaswani et al. "Attention Is All You Need"; Jurafsky \& Martin ch. 10; Lu et al. "Cultural tendencies in generative AI"; Zhang \& Evans "LM Perplexity Predicts Scientific Surprise". Supplemental: Wei et al. CoT; Lightman et al. "Let's Verify Step by Step"; Yao et al. "Tree of Thoughts" & Notebook \href{https://github.com/KnowledgeLab/AI-Agents-for-Social-Science-and-Society-2026/blob/main/Tutorials-Homework_Notebooks/Week_2.ipynb}{[GitHub]} & \hyperlink{mod-FND}{\tg{FND}} \hyperlink{mod-RSN}{\tg{RSN}}\\ \midrule
\hypertarget{lec-UCHI-3}{}3 & 2026-01-23 & \textbf{LLMs Prompted for Multi-Agent Simulation} & LLM digital doubles; in-silico social simulation; RAG & Kozlowski \& Evans "Simulating Subjects"; Leibo et al. "From Predictive Pattern Completion to Social and Cultural Norms"; Cao, Pan, Evans; Xie et al. (PNAS 2025); Swanson et al. "The Virtual Lab of AI agents" (Nature 2025). Supplemental: RAG (Lewis et al.), Concordia, AutoGen, AFlow, LangGraph, Faiss, AI Scientist v1/v2, "Why do multi-agent LLM systems fail?", Argyle et al. "Out of One, Many" & Notebook \href{https://github.com/KnowledgeLab/AI-Agents-for-Social-Science-and-Society-2026/blob/main/Tutorials-Homework_Notebooks/Week_3.ipynb}{[GitHub]} & \hyperlink{mod-MULTI}{\tg{MULTI}} \hyperlink{mod-MEM}{\tg{MEM}} \hyperlink{mod-APP}{\tg{APP}}\\ \midrule
\hypertarget{lec-UCHI-4}{}4 & 2026-01-30 & \textbf{Experimental Designs with AI Agent and Human Interactions} & Causal identification; mixed human/AI subjects; persuasion experiments & Broska et al. "Mixed Subjects Design"; Costello et al. "Durably reducing conspiracy beliefs" (Science 2024); Potter et al. "Hidden Persuaders" (EMNLP 2024); Lai et al. "Biased AI improves human decision-making"; Horton "Homo Silicus"; Koch et al. "Deep learning for causal inference" & Notebook \href{https://github.com/KnowledgeLab/AI-Agents-for-Social-Science-and-Society-2026/blob/main/Tutorials-Homework_Notebooks/week_4_2026.ipynb}{[GitHub]} & \hyperlink{mod-APP}{\tg{APP}} \hyperlink{mod-EVAL}{\tg{EVAL}}\\ \midrule
\hypertarget{lec-UCHI-5}{}5 & 2026-02-06 & \textbf{Sampling, Fine Tuning, Benchmarking, and Tools} & Sampling bias; LoRA/QLoRA; agent benchmarks; tools and MCP & Hu et al. "LoRA"; Mohammadi et al. "Evaluation and benchmarking of LLM agents" (KDD 2025); Anthropic "What is MCP?"; Shumailov et al. "Model collapse" (Nature 2024); Chu et al. "media diets predict public opinion". Supplemental: Kapoor et al. "AI Agents That Matter"; QLoRA; DPO; AgentBench; DSPy; AIOS & Notebook \href{https://github.com/KnowledgeLab/AI-Agents-for-Social-Science-and-Society-2026/blob/main/Tutorials-Homework_Notebooks/Week_5_2026.ipynb}{[GitHub]} & \hyperlink{mod-TRAIN}{\tg{TRAIN}} \hyperlink{mod-EVAL}{\tg{EVAL}} \hyperlink{mod-TOOL}{\tg{TOOL}}\\ \midrule
\hypertarget{lec-UCHI-6}{}6 & 2026-02-13 & \textbf{Auto-Encoders, Interpretability, and Steering Agents} & Sparse autoencoders; personas; activation steering & Bricken et al. "Towards Monosemanticity"; Cunningham et al. "Sparse Autoencoders Find Highly Interpretable Features"; Kim, Evans, Schein "Linear representations of political perspective" (ICLR 2025); Chen et al. "Persona vectors". Supplemental: Scaling Monosemanticity, ActAdd, ITI, In-Context Vectors, Transcoders & Notebook \href{https://github.com/KnowledgeLab/AI-Agents-for-Social-Science-and-Society-2026/blob/main/Tutorials-Homework_Notebooks/week6_2026.ipynb}{[GitHub]} & \hyperlink{mod-SAFE}{\tg{SAFE}} \hyperlink{mod-FND}{\tg{FND}}\\ \midrule
\hypertarget{lec-UCHI-7}{}7 & 2026-02-20 & \textbf{Reinforcement Learning to Optimize AI Agents and Institutions} & RL basics; RLHF; RL for reasoning; multi-agent RL & Sutton-style RL ch. 1; Géron ch. 18; Mnih et al. DQN (Nature 2015); Bai et al. "Constitutional AI"; Ouyang et al. InstructGPT; Guo et al. DeepSeek-R1 (Nature 2025); Kim et al. "Reasoning Models Generate Societies of Thought"; Zhang, Neubig, Yue (pre/mid-training vs. RL); "Hypothetical Minds"; "RATE". Supplemental: Christiano et al.; Agent Lightning; OpenPipe ART (GRPO) & Notebook \href{https://github.com/KnowledgeLab/AI-Agents-for-Social-Science-and-Society-2026/blob/main/Tutorials-Homework_Notebooks/Week_7_RL.ipynb}{[GitHub]} & \hyperlink{mod-TRAIN}{\tg{TRAIN}} \hyperlink{mod-RSN}{\tg{RSN}} \hyperlink{mod-MULTI}{\tg{MULTI}}\\ \midrule
\hypertarget{lec-UCHI-8}{}8 & 2026-02-27 & \textbf{Multi-Modal and Embodied Agents} & Audio, video, and image agents; diffusion; robot embodiment & Bishop ch. 10, 20 (CNNs, diffusion); Yao et al. "ReAct"; OpenAI GPT-4V system card; Song et al. "Embodied Task Planning with LLMs"; Ludwig \& Mullainathan "ML as a Tool for Hypothesis Generation" (QJE 2024); Guilbeault et al. "Age and gender distortion" & Notebook \href{https://github.com/KnowledgeLab/AI-Agents-for-Social-Science-and-Society-2026/blob/main/Tutorials-Homework_Notebooks/Week_8_Multimodal.ipynb}{[GitHub]} & \hyperlink{mod-WEB}{\tg{WEB}} \hyperlink{mod-RSN}{\tg{RSN}} \hyperlink{mod-APP}{\tg{APP}}\\ \midrule
\hypertarget{lec-UCHI-9}{}9 & 2026-03-06 & \textbf{Alignment, Ethics, Safety, and Novelty} & Value alignment; deception; guardrails; effects on the space of ideas & Russell "Human-Compatible AI"; Hubinger et al. "Sleeper Agents"; Ngo et al. "Alignment Problem from a DL Perspective"; Tomasev et al. "Distributional AGI Safety"; Zou et al. "Representation Engineering"; Gabriel, Keeling, Manzini, Evans "We need a new ethics for a world of AI agents" (Nature 2025). Supplemental: AGrail; Si et al. "Can LLMs Generate Novel Research Ideas?"; Moral Machine & Notebook \href{https://github.com/KnowledgeLab/AI-Agents-for-Social-Science-and-Society-2026/blob/main/Tutorials-Homework_Notebooks/Week_9.ipynb}{[GitHub]} & \hyperlink{mod-SAFE}{\tg{SAFE}} \hyperlink{mod-APP}{\tg{APP}}\\ \midrule
\hypertarget{lec-UCHI-10}{}10 & 2026-03-13 & \textbf{Presentation of Projects (via Zoom)} & Final project presentations & -- & -- & \tg{ORG}\\ \midrule
\end{longtable}}
\section{Assignments and projects}
{\scriptsize\setlength{\tabcolsep}{2.5pt}\begin{longtable}{P{2.8cm} P{5.4cm} P{3.0cm} P{3.5cm} P{1.5cm}}
\toprule Name & What you build & Skills & Starter code / resources & Tags\\ \midrule\endhead
\hypertarget{asg-UCHI-1}{}\textbf{Week 1 notebook (Intro NNs)} & Feed-forward nets, wide \& deep and deep \& cross networks on your own dataset; compare activations, dropout, optimizers; error analysis & PyTorch, scikit-learn MLP & \href{https://github.com/KnowledgeLab/AI-Agents-for-Social-Science-and-Society-2026/blob/main/Tutorials-Homework_Notebooks/Week_1_Intro_NNs.ipynb}{[GitHub]} & \hyperlink{mod-FND}{\tg{FND}}\\ \midrule
\hypertarget{asg-UCHI-2}{}\textbf{Week 2 notebook} & 3 of 5 modules: word2vec on your own corpus, seq2seq/attention, fine-tuning a transformer with SHAP, BERT similarity and perplexity, diffusion hyperparameters; memo pilot & gensim, spaCy, HF transformers, shap, umap & \href{https://github.com/KnowledgeLab/AI-Agents-for-Social-Science-and-Society-2026/blob/main/Tutorials-Homework_Notebooks/Week_2.ipynb}{[GitHub]} & \hyperlink{mod-FND}{\tg{FND}}\\ \midrule
\hypertarget{asg-UCHI-3}{}\textbf{Week 3 notebook} & Prompting (0-shot, few-shot, chain-of-reasoning) via the OpenAI API; simulated survey and social-media respondents compared with human data; AutoGen multi-agent dialogues with reflection; Concordia situated simulation; RAG over 5-10 documents vs. a vanilla model & openai, LangChain, AutoGen, Concordia, FAISS/ChromaDB & \href{https://github.com/KnowledgeLab/AI-Agents-for-Social-Science-and-Society-2026/blob/main/Tutorials-Homework_Notebooks/Week_3.ipynb}{[GitHub]} & \hyperlink{mod-MULTI}{\tg{MULTI}} \hyperlink{mod-MEM}{\tg{MEM}} \hyperlink{mod-FND}{\tg{FND}}\\ \midrule
\hypertarget{asg-UCHI-4}{}\textbf{Week 4 notebook} & Deep causal inference (S/T-learner, TARNet on IHDP), causal mediation, double/debiased ML, prediction-powered inference on your own question & TensorFlow/Keras, causal ML & \href{https://github.com/KnowledgeLab/AI-Agents-for-Social-Science-and-Society-2026/blob/main/Tutorials-Homework_Notebooks/week_4_2026.ipynb}{[GitHub]} & \hyperlink{mod-APP}{\tg{APP}} \hyperlink{mod-EVAL}{\tg{EVAL}}\\ \midrule
\hypertarget{asg-UCHI-5}{}\textbf{Week 5 notebook} & 3 of 4 modules: sampling; LoRA/QLoRA fine-tuning with varied ranks plus CKA comparison; designing an agent benchmark; a tool-using agent and an MCP server & LoRA, benchmarking, MCP & \href{https://github.com/KnowledgeLab/AI-Agents-for-Social-Science-and-Society-2026/blob/main/Tutorials-Homework_Notebooks/Week_5_2026.ipynb}{[GitHub]} & \hyperlink{mod-TRAIN}{\tg{TRAIN}} \hyperlink{mod-EVAL}{\tg{EVAL}} \hyperlink{mod-TOOL}{\tg{TOOL}}\\ \midrule
\hypertarget{asg-UCHI-6}{}\textbf{Week 6 notebook} & Interpretability and steering (contents not read: file over 10 MB) & -- & \href{https://github.com/KnowledgeLab/AI-Agents-for-Social-Science-and-Society-2026/blob/main/Tutorials-Homework_Notebooks/week6_2026.ipynb}{[GitHub]} & \hyperlink{mod-SAFE}{\tg{SAFE}} \hyperlink{mod-FND}{\tg{FND}}\\ \midrule
\hypertarget{asg-UCHI-7}{}\textbf{Week 7 notebook (RL)} & Bandits (epsilon-greedy, annealing, UCB), RL for alignment, then RL for reasoning or multi-agent RL; memo pilot & gymnasium, stable-baselines3, PPO/DQN/GRPO & \href{https://github.com/KnowledgeLab/AI-Agents-for-Social-Science-and-Society-2026/blob/main/Tutorials-Homework_Notebooks/Week_7_RL.ipynb}{[GitHub]} & \hyperlink{mod-TRAIN}{\tg{TRAIN}} \hyperlink{mod-MULTI}{\tg{MULTI}}\\ \midrule
\hypertarget{asg-UCHI-8}{}\textbf{Week 8 notebook (Multimodal)} & Multimodal and embodied agents (contents not read: file over 10 MB) & -- & \href{https://github.com/KnowledgeLab/AI-Agents-for-Social-Science-and-Society-2026/blob/main/Tutorials-Homework_Notebooks/Week_8_Multimodal.ipynb}{[GitHub]} & \hyperlink{mod-WEB}{\tg{WEB}} \hyperlink{mod-APP}{\tg{APP}}\\ \midrule
\hypertarget{asg-UCHI-9}{}\textbf{Week 9 notebook} & Any 4 of 7 modules: safety judge, Constitutional AI, deception detection, moral geometry, automated red-teaming, AGrail guardrails; 150-300-word reflections (not AI-delegable) & Qwen2.5-7B-Instruct (4-bit), sentence-transformers, AGrail & \href{https://github.com/KnowledgeLab/AI-Agents-for-Social-Science-and-Society-2026/blob/main/Tutorials-Homework_Notebooks/Week_9.ipynb}{[GitHub]} & \hyperlink{mod-SAFE}{\tg{SAFE}} \hyperlink{mod-EVAL}{\tg{EVAL}}\\ \midrule
\hypertarget{asg-UCHI-10}{}\textbf{Weekly memos and Ignite presentations} & 300-500-word reading memos (from week 2) with peer upvoting; 5-minute code presentations & Critical reading, communication & -- & \tg{ORG}\\ \midrule
\hypertarget{asg-UCHI-11}{}\textbf{Final project} & Agent project for social insight, solo or teams of up to 4: at least 3 data types and 3 model types from 3 weeks, qualitatively validated; Substack/Medium post (3,600 words and 13 visuals solo; 5,000 and 17 for groups) with an open repo, plus an Ignite talk & End-to-end applied agent research & -- & \hyperlink{mod-APP}{\tg{APP}} \tg{ORG}\\ \midrule
\end{longtable}}
\chapter{Stanford 329A: CS329A Self-Improving AI Agents}\hypertarget{course-S329A}{}
{\small\begin{tabularx}{\linewidth}{@{}>{\bfseries}l X@{}}\toprule
Institution & Stanford University\\
Term & Autumn 2025\\
Instructors & Aakanksha Chowdhery, Azalia Mirhoseini\\
Website & \href{https://cs329a.stanford.edu/}{https://cs329a.stanford.edu/}\\
Level & Graduate; CS224N or CS229S\\
Archetype & Research course (self-improvement levers) + homeworks\\
Public materials & Slides: none linked; video: 9 recordings (Stanford Online); code: HW1--HW3 public repos (HW1 MIT, HW3 Apache-2.0)\\
Best for & The clearest path through test-time compute, verifiers, RL and open-ended self-improvement, with three public homeworks.\\
\bottomrule\end{tabularx}}

Asks how an agent improves through interaction with itself and its environment, and takes the levers in order: repeated sampling and test-time compute, verifiers and process reward models, learning from tool and execution feedback, planning and tree search, and train-time RL (STaR, GRPO, DAPO).
The second half covers open-ended self-improvement (ADAS, AI Scientist, AlphaEvolve), deep-research and software-engineering agents, memory, and long-horizon evaluation, with five guest talks.
The homeworks follow the same arc: test-time scaling on AIME25, code verification with LLM judges and synthetic unit tests, then a self-improving agentic workflow.
No slide decks are posted. The video links were matched to lectures through a third-party index and then checked on YouTube.
\section{Lectures}
{\scriptsize\setlength{\tabcolsep}{2.5pt}\begin{longtable}{P{0.45cm} P{1.2cm} P{2.5cm} P{3.6cm} P{3.6cm} P{3.3cm} P{1.5cm}}
\toprule \# & Date & Title & Key topics & Readings & Materials & Tags\\ \midrule\endhead
\hypertarget{lec-S329A-1}{}1 & Mon Sep 22 & \textbf{Course Overview} & Course logistics; framing of self-improving agents & -- & \href{https://www.youtube.com/watch?v=6YnLB0XbTnI}{[video]} & \tg{ORG}\\ \midrule
\hypertarget{lec-S329A-2}{}2 & Fri Sep 26 & \textbf{Test-time Compute Scaling} & Repeated sampling, inference scaling laws, architecture search over inference techniques & Large Language Monkeys; Archon; Scaling LLM Test-Time Compute Optimally (Snell et al.); How Do Large Language Monkeys Get Their Power (Laws)? & \href{https://www.youtube.com/watch?v=-Ggc37xLj_Y}{[video]} & \hyperlink{mod-RSN}{\tg{RSN}} \hyperlink{mod-EVAL}{\tg{EVAL}}\\ \midrule
\hypertarget{lec-S329A-3}{}3 & Mon Sep 29 & \textbf{Robust Verification} & Outcome/process verifiers, weak verifiers, generation-verification gap & Shrinking the Generation-Verification Gap with Weak Verifiers; Training Verifiers to Solve Math Word Problems; Let's Verify Step by Step; Math-Shepherd & \href{https://www.youtube.com/watch?v=p7TdPUcPoik}{[video]} & \hyperlink{mod-RSN}{\tg{RSN}} \hyperlink{mod-EVAL}{\tg{EVAL}}\\ \midrule
\hypertarget{lec-S329A-4}{}4 & Fri Oct 3 & \textbf{Learning from Feedback with Tools/Code} & Reason+act loops, RL from execution feedback, AI feedback (HW1 out) & ReAct; RLEF; Constitutional AI & \href{https://www.youtube.com/watch?v=Lxh9RF5S-K0}{[video]} & \hyperlink{mod-TOOL}{\tg{TOOL}} \hyperlink{mod-CODE}{\tg{CODE}} \hyperlink{mod-TRAIN}{\tg{TRAIN}}\\ \midrule
\hypertarget{lec-S329A-5}{}5 & Mon Oct 6 & \textbf{Multi-step Reasoning/Planning} & Multi-step RL with synthetic data, tree search, decomposition, parallel planning & SWiRL; Language Agent Tree Search; SPRINT; ADaPT; Wider or Deeper? (adaptive branching tree search) & \href{https://www.youtube.com/watch?v=Ml_fp9XkB8Y}{[video]} & \hyperlink{mod-RSN}{\tg{RSN}} \hyperlink{mod-TRAIN}{\tg{TRAIN}}\\ \midrule
\hypertarget{lec-S329A-6}{}6 & Fri Oct 10 & \textbf{Train Time Scaling/Scaling RL} & Self-taught reasoning, GRPO, large-scale RL recipes (project proposal due) & STaR; DeepSeekMath; DAPO & \href{https://www.youtube.com/watch?v=yVnmHSAy3ck}{[video]} & \hyperlink{mod-TRAIN}{\tg{TRAIN}} \hyperlink{mod-RSN}{\tg{RSN}}\\ \midrule
\hypertarget{lec-S329A-7}{}7 & Mon Oct 13 & \textbf{Open-Ended Evolution of Self-Improving Agents} & Automated agent design, automated science, evolutionary coding agents (HW2 out Oct 14) & Automated Design of Agentic Systems; The AI Scientist; AlphaEvolve & -- & \hyperlink{mod-ARCH}{\tg{ARCH}} \hyperlink{mod-CODE}{\tg{CODE}} \hyperlink{mod-APP}{\tg{APP}}\\ \midrule
\hypertarget{lec-S329A-8}{}8 & Fri Oct 17 & \textbf{Self-Improvement with Search \& Deep Research Agents} & Search-based code generation, agentic search for reasoning models, deep research & AlphaCode; AlphaCode 2 Tech Report; Search-o1 & \href{https://www.youtube.com/watch?v=Uni9dqyuuDM}{[video]} & \hyperlink{mod-WEB}{\tg{WEB}} \hyperlink{mod-RSN}{\tg{RSN}} \hyperlink{mod-CODE}{\tg{CODE}}\\ \midrule
\hypertarget{lec-S329A-9}{}9 & Mon Oct 20 & \textbf{Guest: Melvin Johnson (Google DeepMind), Evolution of Post-training from Chatbots to Agents} & Post-training for agents (from title only) & -- & -- & \tg{ORG} \hyperlink{mod-TRAIN}{\tg{TRAIN}}\\ \midrule
\hypertarget{lec-S329A-10}{}10 & Fri Oct 24 & \textbf{Midterm presentations (1)} & 5-10 min project presentations (HW3 out) & -- & -- & \tg{ORG}\\ \midrule
\hypertarget{lec-S329A-11}{}11 & Mon Oct 27 & \textbf{Midterm presentations (2)} & Project presentations & -- & -- & \tg{ORG}\\ \midrule
\hypertarget{lec-S329A-12}{}12 & Fri Oct 31 & \textbf{Midterm presentations (3)} & Project presentations & -- & -- & \tg{ORG}\\ \midrule
\hypertarget{lec-S329A-13}{}13 & Mon Nov 3 & \textbf{Agentic Frameworks for Software Engineering} & Test-time scaling for SWE, GPU-kernel generation, agent-system interfaces & CodeMonkeys; KernelBench; Improving Parallel Program Performance with LLM Optimizers via Agent-System Interfaces & -- & \hyperlink{mod-CODE}{\tg{CODE}} \hyperlink{mod-ARCH}{\tg{ARCH}}\\ \midrule
\hypertarget{lec-S329A-14}{}14 & Fri Nov 7 & \textbf{Guest: Augmenting Agents with Memory (Junchen Jiang, LMCache/UChicago)} & Long-context representations, OS-style memory, KV-cache reuse (HW3 due) & Cartridges; MemGPT; CacheBlend & -- & \hyperlink{mod-MEM}{\tg{MEM}} \hyperlink{mod-DEPLOY}{\tg{DEPLOY}}\\ \midrule
\hypertarget{lec-S329A-15}{}15 & Mon Nov 10 & \textbf{Guest: Denny Zhou (Google DeepMind), LLM Reasoning} & LLM reasoning (from title only) & -- & -- & \tg{ORG} \hyperlink{mod-RSN}{\tg{RSN}}\\ \midrule
\hypertarget{lec-S329A-16}{}16 & Fri Nov 14 & \textbf{Guest: Thang Luong (Google DeepMind), Towards AI Superhuman Reasoning} & Superhuman reasoning (from title only) & -- & -- & \tg{ORG} \hyperlink{mod-RSN}{\tg{RSN}}\\ \midrule
\hypertarget{lec-S329A-17}{}17 & Mon Nov 17 & \textbf{Agentic Evaluations \& Long-Horizon Tasks} & Task-length metrics, economically valuable tasks, research-synthesis benchmarks & Measuring AI Ability to Complete Long Tasks; GDPVal; DeepScholar-Bench & \href{https://www.youtube.com/watch?v=8JAqLnTaZu4}{[video]} & \hyperlink{mod-EVAL}{\tg{EVAL}}\\ \midrule
\hypertarget{lec-S329A-18}{}18 & Fri Nov 21 & \textbf{Guest: Misha Laskin (Reflection AI), Building Agentic Systems for Autonomy} & Autonomous agentic systems (from title only) & -- & -- & \tg{ORG} \hyperlink{mod-ARCH}{\tg{ARCH}}\\ \midrule
\hypertarget{lec-S329A-19}{}19 & Mon Dec 1 & \textbf{Guest: Danny Driess (Physical Intelligence), Multimodal AI Agents in Robotics} & Multimodal robot agents (from title only) & -- & -- & \tg{ORG} \hyperlink{mod-APP}{\tg{APP}}\\ \midrule
\hypertarget{lec-S329A-20}{}20 & Fri Dec 5 & \textbf{Future Research Areas} & Open problems in self-improving agents (from title only) & -- & \href{https://www.youtube.com/watch?v=AyO6wyu4DEg}{[video]} & \hyperlink{mod-FND}{\tg{FND}}\\ \midrule
\end{longtable}}
\section{Assignments and projects}
{\scriptsize\setlength{\tabcolsep}{2.5pt}\begin{longtable}{P{2.8cm} P{5.4cm} P{3.0cm} P{3.5cm} P{1.5cm}}
\toprule Name & What you build & Skills & Starter code / resources & Tags\\ \midrule\endhead
\hypertarget{asg-S329A-1}{}\textbf{HW1: Test-time compute scaling (out Oct 3, due Oct 13; 130 pts + 2 bonus)} & On AIME25 (30 problems, \texttt{math-ai/\allowbreak{}aime25}) with Qwen3-Next-80B-A3B-Instruct via Together AI: zero-shot baseline; \texttt{MajorityVoting} over budgets 1-16; best-of-N with an LLM aggregator (\texttt{LLMVoting}, generative reward model); an RLEF-style \texttt{SelfImprovementSystem} (generate, critique, regenerate); plots plus a cost-per-API-call analysis & Repeated sampling, majority voting, LLM-as-judge, critique-and-refine loops, accuracy-vs-compute analysis & \href{https://github.com/stanford-cs329a/cs329a-homework1-fall2025-public}{[GitHub]} & \hyperlink{mod-RSN}{\tg{RSN}} \hyperlink{mod-EVAL}{\tg{EVAL}}\\ \midrule
\hypertarget{asg-S329A-2}{}\textbf{HW2: Code generation and verification (out Oct 14, due Oct 22 per schedule / Oct 23 per grading; 110 pts)} & On HumanEval with Llama-3.2-3B-Instruct-Turbo, code run in a Docker sandbox: \texttt{calculate\_\allowbreak{}accuracy} and failure analysis; pass@k for k = 3, 6, 9; an \texttt{LLMJudge} that picks among candidates via JSON output; an \texttt{LLMUnitTestGenerator} (Qwen2.5-72B) that writes 5 tests from the docstring; checks how often ground truth passes the synthetic tests; pass@3 with synthetic vs ground-truth tests & Sandboxed execution, pass@k, LLM-as-judge, synthetic unit tests as weak verifiers, error analysis & \href{https://github.com/stanford-cs329a/cs329a-homework2-fall2025-public}{[GitHub]} & \hyperlink{mod-CODE}{\tg{CODE}} \hyperlink{mod-EVAL}{\tg{EVAL}} \hyperlink{mod-RSN}{\tg{RSN}}\\ \midrule
\hypertarget{asg-S329A-3}{}\textbf{HW3: Self-improving agentic workflow (out Oct 24, due Nov 7; 120 pts)} & An \texttt{APIManager} with Google Custom Search, Polygon stock data, Wolfram Alpha and weather (Nominatim + Open-Meteo) tools; LLM-based \texttt{route\_\allowbreak{}query} with Pydantic/JSON mode; a \texttt{MultiLMAgent} with query decomposition, multi-model fusion, iterative self-refinement and an adaptive \texttt{run\_\allowbreak{}pipeline} (full credit at 70\% or more QA accuracy); a \texttt{DeepResearchAgent} that writes a cited 4-5 paragraph report & Tool routing, structured outputs, decomposition, ensembling/fusion, self-refinement loops, deep-research agents, LLM-judged QA eval & \href{https://github.com/stanford-cs329a/cs329a-homework3-fall2025-public}{[GitHub]} & \hyperlink{mod-TOOL}{\tg{TOOL}} \hyperlink{mod-WEB}{\tg{WEB}} \hyperlink{mod-ARCH}{\tg{ARCH}} \hyperlink{mod-EVAL}{\tg{EVAL}}\\ \midrule
\hypertarget{asg-S329A-4}{}\textbf{Project proposal (2.5\%, due Oct 10)} & Proposal for the team research project & Research design & -- & \tg{ORG}\\ \midrule
\hypertarget{asg-S329A-5}{}\textbf{Midterm presentation + report (10\%)} & 5-10 min presentation (Oct 24, 27, 31) plus report & Presenting, writing & -- & \tg{ORG}\\ \midrule
\hypertarget{asg-S329A-6}{}\textbf{Final project (35\%, due Dec 10) + poster (2.5\%, Dec 12)} & Original research in teams of 2-4 on course topics or benchmarks for agentic workflows (API credits provided) & Research, implementation, writing & n/a (past examples: \href{https://cs329a.stanford.edu/pastprojects.html}{[cs329a.stanford.edu]}) & \tg{ORG} \hyperlink{mod-TRAIN}{\tg{TRAIN}} \hyperlink{mod-EVAL}{\tg{EVAL}}\\ \midrule
\end{longtable}}
\chapter{Berkeley '25: CS194/294-196 Agentic AI}\hypertarget{course-BRKAG}{}
{\small\begin{tabularx}{\linewidth}{@{}>{\bfseries}l X@{}}\toprule
Institution & UC Berkeley (RDI)\\
Term & Fall 2025\\
Instructors & Dawn Song (with guest speakers)\\
Website & \href{https://rdi.berkeley.edu/agentic-ai/f25}{https://rdi.berkeley.edu/agentic-ai/f25}\\
Level & Undergraduate and graduate; ML/DL background\\
Archetype & Guest-lecture series + AgentBeats project (also a MOOC)\\
Public materials & Slides: 9 of 13 sessions; video: 12 of 13 sessions; code: AgentBeats tutorial repo (no license stated)\\
Best for & Industry and lab speakers on post-training, agent evaluation and deployment, and an evaluator-agent project (AgentBeats).\\
\bottomrule\end{tabularx}}

A guest-speaker series (OpenAI, NVIDIA, Microsoft, Meta, Sierra, Google DeepMind and others) bracketed by Dawn Song's introduction and closing lecture on agent safety and security.
Recurring themes are post-training for verifiable agentic tasks, agent evaluation and its statistics, multi-agent systems, and lessons from production deployment.
The project is the distinctive part. In phase~1 you build a \emph{green} agent, an evaluator that turns a benchmark into tasks, an environment and scoring. In phase~2 you build the agent under test, called \emph{white} in the course and \emph{purple} in the current AgentBeats competition.
The same lectures ran as a MOOC at agenticai-learning.org/f25.
\section{Lectures}
{\scriptsize\setlength{\tabcolsep}{2.5pt}\begin{longtable}{P{0.45cm} P{1.2cm} P{2.5cm} P{3.6cm} P{3.6cm} P{3.3cm} P{1.5cm}}
\toprule \# & Date & Title & Key topics & Readings & Materials & Tags\\ \midrule\endhead
\hypertarget{lec-BRKAG-1}{}1 & Mon Sep 8 & \textbf{Introduction} & Speaker: Dawn Song (UC Berkeley). Course intro and logistics & -- & \href{https://rdi.berkeley.edu/agentic-ai/slides/introduction_25.pdf}{slides} & \tg{ORG} \hyperlink{mod-FND}{\tg{FND}}\\ \midrule
\hypertarget{lec-BRKAG-2}{}2 & Mon Sep 15 & \textbf{LLM Agents Overview} & Speaker: Yann Dubois (OpenAI). LLM foundations and agent overview & \href{https://rdi.berkeley.edu/agentic-ai/slides/d11.pdf}{Kimi K2: Open Agentic Intelligence}; \href{https://rdi.berkeley.edu/agentic-ai/slides/d22.pdf}{DeepSeek-V3 Technical Report} & \href{https://rdi.berkeley.edu/agentic-ai/slides/lecture1.pdf}{slides} \href{https://www.youtube.com/watch?v=r1qZpYAmqmg}{video} & \hyperlink{mod-FND}{\tg{FND}} \hyperlink{mod-TRAIN}{\tg{TRAIN}}\\ \midrule
\hypertarget{lec-BRKAG-3}{}3 & Mon Sep 22 & \textbf{Evolution of system designs from an AI engineer perspective} & Speaker: Yangqing Jia (NVIDIA). AI system and infrastructure design & -- & \href{https://rdi.berkeley.edu/agentic-ai/slides/Yangqing.pdf}{slides} \href{https://www.youtube.com/watch?v=xqRAS6rAouo\&list=PLS01nW3RtgoqGkm4UeqNeZLccW-OGc1fJ\&index=2}{video} & \hyperlink{mod-ARCH}{\tg{ARCH}} \hyperlink{mod-DEPLOY}{\tg{DEPLOY}}\\ \midrule
\hypertarget{lec-BRKAG-4}{}4 & Mon Sep 29 & \textbf{Post-Training Verifiable Agents} & Speaker: Jiantao Jiao (NVIDIA). Post-training agents on verifiable tasks & \href{https://openai.com/index/introducing-swe-bench-verified/}{Introducing SWE-bench Verified}; \href{https://openai.com/index/browsecomp/}{BrowseComp} & \href{https://rdi.berkeley.edu/agentic-ai/slides/AgenticAI_Jiantao\%20Jiao.pdf}{slides} \href{https://www.youtube.com/watch?v=3l0Zxus34es\&list=PLS01nW3RtgoqGkm4UeqNeZLccW-OGc1fJ\&index=3}{video} & \hyperlink{mod-TRAIN}{\tg{TRAIN}} \hyperlink{mod-EVAL}{\tg{EVAL}} \hyperlink{mod-CODE}{\tg{CODE}} \hyperlink{mod-WEB}{\tg{WEB}}\\ \midrule
\hypertarget{lec-BRKAG-5}{}5 & Mon Oct 6 & \textbf{Agent Evaluation \& Project Overview} & Speaker: n/a (not listed). Agent evaluation; AgentBeats project overview & \href{https://arxiv.org/pdf/2503.16416}{Survey on Evaluation of LLM-based Agents} & \href{https://rdi.berkeley.edu/agentic-ai/slides/LLM\%20Agent\%20Evaluations\%20\%26\%20Project\%20Overview.pdf}{slides} \href{https://www.youtube.com/watch?v=VfOA2a0dj4w}{video} & \hyperlink{mod-EVAL}{\tg{EVAL}} \tg{ORG}\\ \midrule
\hypertarget{lec-BRKAG-6}{}6 & Mon Oct 13 & \textbf{Some Challenges and Lessons from Training Agentic Models} & Speaker: Weizhu Chen (Microsoft). Training agentic models & -- & \href{https://rdi.berkeley.edu/agentic-ai/slides/weizhu.pdf}{slides} \href{https://www.youtube.com/watch?v=xNxrBHZPDvM}{video} & \hyperlink{mod-TRAIN}{\tg{TRAIN}}\\ \midrule
\hypertarget{lec-BRKAG-7}{}7 & Mon Oct 20 & \textbf{Multi-Agent AI} & Speaker: Noam Brown (OpenAI). Multi-agent AI & -- & \href{https://rdi.berkeley.edu/agentic-ai/slides/multiagent_berkeley_2025.pptx.pdf}{slides} \href{https://www.youtube.com/watch?v=SrLcGdVOb9w}{video} & \hyperlink{mod-MULTI}{\tg{MULTI}} \hyperlink{mod-RSN}{\tg{RSN}}\\ \midrule
\hypertarget{lec-BRKAG-8}{}8 & Mon Oct 27 & \textbf{Predictable Noise in LLM} & Speaker: Sida Wang (Meta). Noise and error bars in LLM evals & \href{https://arxiv.org/pdf/2411.00640}{Adding Error Bars to Evals} & \href{https://rdi.berkeley.edu/agentic-ai/slides/PredEval.pdf}{slides} \href{https://www.youtube.com/watch?v=HV8pugcFVO0}{video} & \hyperlink{mod-EVAL}{\tg{EVAL}}\\ \midrule
\hypertarget{lec-BRKAG-9}{}9 & Mon Nov 3 & \textbf{AI Agents to Automate Scientific Discoveries} & Speaker: James Zou (Stanford). Agents for science (Virtual Lab, Paper2Agent) & \href{https://www.nature.com/articles/s41586-025-09442-9}{The Virtual Lab of AI agents designs new SARS-CoV-2 nanobodies}; \href{https://arxiv.org/abs/2509.06917}{Paper2Agent} & \href{https://www.youtube.com/watch?v=yqPIsTTdUkc}{video} & \hyperlink{mod-APP}{\tg{APP}} \hyperlink{mod-MULTI}{\tg{MULTI}}\\ \midrule
\hypertarget{lec-BRKAG-10}{}10 & Mon Nov 10 & \textbf{Practical Lessons from Deploying Real-World AI Agents} & Speaker: Clay Bavor (Sierra). Production customer-facing agents; simulation testing & \href{https://arxiv.org/pdf/2506.07982}{$\tau$2-Bench}; \href{https://sierra.ai/blog/voice-sims-test-agents-in-real-world-conditions-before-they-talk-to-your-customers}{Sierra Voice Sims blog} & \href{https://rdi.berkeley.edu/agentic-ai/slides/Practical_Lessons_AI_Agents.pdf}{slides} \href{https://www.youtube.com/watch?v=sfJM4LaiYsM}{video} & \hyperlink{mod-DEPLOY}{\tg{DEPLOY}} \hyperlink{mod-EVAL}{\tg{EVAL}} \hyperlink{mod-APP}{\tg{APP}}\\ \midrule
\hypertarget{lec-BRKAG-11}{}11 & Mon Nov 17 & \textbf{Multi-Agent Systems in the Era of LLMs} & Speaker: Oriol Vinyals (Google DeepMind). Multi-agent systems with LLMs & -- & \href{https://www.youtube.com/watch?v=ntjOxjZMaac\&t=2s}{video} & \hyperlink{mod-MULTI}{\tg{MULTI}}\\ \midrule
\hypertarget{lec-BRKAG-12}{}12 & Mon Dec 1 (2:00-3:30pm PT) & \textbf{Autonomous Agents: Embodiment, Interaction, and Learning} & Speaker: Peter Stone (UT Austin / Sony AI). Embodied agents and RL (Gran Turismo, SLAC) & \href{https://www.cs.utexas.edu/\%7Epstone/Papers/bib2html/b2hd-nature22.html}{Outracing Champion Gran Turismo Drivers with Deep RL}; \href{https://www.cs.utexas.edu/\%7Epstone/Papers/bib2html/b2hd-jiaheng_hu_2025.html}{SLAC} & \href{https://www.youtube.com/watch?v=iDhzzugMOLA}{video} & \hyperlink{mod-APP}{\tg{APP}} \hyperlink{mod-TRAIN}{\tg{TRAIN}}\\ \midrule
\hypertarget{lec-BRKAG-13}{}13 & Mon Dec 8 (11:30am-1:30pm PT) & \textbf{Agentic AI Safety \& Security} & Speaker: Dawn Song (UC Berkeley). Safety and security of agentic AI & -- & \href{https://youtube.com/live/CvZDJxd4LKM?feature=share}{video} (livestream) & \hyperlink{mod-SAFE}{\tg{SAFE}}\\ \midrule
\end{longtable}}
\section{Assignments and projects}
{\scriptsize\setlength{\tabcolsep}{2.5pt}\begin{longtable}{P{2.8cm} P{5.4cm} P{3.0cm} P{3.5cm} P{1.5cm}}
\toprule Name & What you build & Skills & Starter code / resources & Tags\\ \midrule\endhead
\hypertarget{asg-BRKAG-1}{}\textbf{Project Phase 1: Green agent (2 units 45\%, 3/4 units 50\%)} & An evaluator ("green") agent on AgentBeats that wraps a benchmark or game as an agent. It defines tasks, environment and scoring, receives participant endpoints plus a config over A2A, runs the assessment and returns JSON results as artifacts. Deliverables: group formation (due 9/22), proposal (5\%, due 10/8), early demo of a 2-minute video plus milestone report (10\%, due 10/20), final green agent (30\% or 35\%, released 11/7, due 12/17) & Benchmark design, environment and scoring design, A2A agent servers, Docker packaging, reproducible evaluation & \href{https://github.com/RDI-Foundation/agentbeats-tutorial}{[GitHub]} & \hyperlink{mod-EVAL}{\tg{EVAL}} \hyperlink{mod-MULTI}{\tg{MULTI}} \hyperlink{mod-ARCH}{\tg{ARCH}}\\ \midrule
\hypertarget{asg-BRKAG-2}{}\textbf{Project Phase 2: White agent (2 units 15\%, 3/4 units 20\%)} & An agent under test ("white"; "purple" in current AgentBeats docs) that competes on green-agent benchmarks. Implementation (10\% or 15\%) plus a 1-2 page report (5\%); released 11/24, due 12/17 & Agent design, tool use, A2A interoperability, performance on evaluator agents & \href{https://github.com/RDI-Foundation/agentbeats-tutorial}{[GitHub]} & \hyperlink{mod-ARCH}{\tg{ARCH}} \hyperlink{mod-TOOL}{\tg{TOOL}} \hyperlink{mod-EVAL}{\tg{EVAL}}\\ \midrule
\hypertarget{asg-BRKAG-3}{}\textbf{Article (1 unit, 40\%)} & A written article, due 12/7. 1-unit students may instead do Phase 1 of the "Agent Track" (minimal coding, e.g. a simple prompt-engineered agent), also 40\% & Technical writing & -- & \tg{ORG}\\ \midrule
\hypertarget{asg-BRKAG-4}{}\textbf{Quizzes (10-20\%)} & Lecture quizzes (MOOC: completion-graded, usually 5 multiple-choice questions) & Lecture comprehension & -- & \tg{ORG}\\ \midrule
\end{longtable}}
\chapter{Berkeley Adv.: CS294/194-280 Advanced Large Language Model Agents}\hypertarget{course-BRKADV}{}
{\small\begin{tabularx}{\linewidth}{@{}>{\bfseries}l X@{}}\toprule
Institution & UC Berkeley (RDI)\\
Term & Spring 2025\\
Instructors & Dawn Song, Xinyun Chen, Kaiyu Yang\\
Website & \href{https://rdi.berkeley.edu/adv-llm-agents/sp25}{https://rdi.berkeley.edu/adv-llm-agents/sp25}\\
Level & Upper-division undergraduate and graduate\\
Archetype & Guest-lecture series + Lean labs (also a MOOC)\\
Public materials & Slides: All 12 lectures; video: All 12 lectures; code: 2 Lean labs in a public Drive folder (no license)\\
Best for & Research talks on reasoning, post-training, coding and web agents and formal math, plus a compiler-verified agent lab.\\
\bottomrule\end{tabularx}}

Picks up where Berkeley's Fall 2024 LLM Agents course left off: inference-time reasoning, iterative preference optimization and open post-training recipes (Tulu~3), language-agent memory and planning, coding agents and vulnerability detection, and multimodal web and computer-use agents.
A four-lecture block on formal mathematics (AlphaProof, autoformalization, theorem proving, abstraction discovery) is followed by a closing lecture on agent safety and security.
Lab~2 is the main draw for agent builders: a verify-and-retry agent loop in which the Lean compiler is the verifier, and grading rejects 	exttt{sorry}. Lab~1, writing two Lean~4 problems with formal specifications, is dataset work rather than an agent task.
\section{Lectures}
{\scriptsize\setlength{\tabcolsep}{2.5pt}\begin{longtable}{P{0.45cm} P{1.2cm} P{2.5cm} P{3.6cm} P{3.6cm} P{3.3cm} P{1.5cm}}
\toprule \# & Date & Title & Key topics & Readings & Materials & Tags\\ \midrule\endhead
\hypertarget{lec-BRKADV-1}{}1 & Mon Jan 27, 2025 & \textbf{Inference-Time Techniques for LLM Reasoning} & Speaker: Xinyun Chen (Google DeepMind). Inference-time reasoning, self-correction, self-debugging; course intro deck & LLMs as Optimizers \href{https://arxiv.org/abs/2309.03409}{[arXiv:2309.03409]} ; LLMs Cannot Self-Correct Reasoning Yet \href{https://arxiv.org/abs/2310.01798}{[arXiv:2310.01798]} ; Teaching LLMs to Self-Debug \href{https://arxiv.org/abs/2304.05128}{[arXiv:2304.05128]} & \href{https://rdi.berkeley.edu/adv-llm-agents/slides/inference_time_techniques_lecture_sp25.pdf}{slides} \href{https://rdi.berkeley.edu/adv-llm-agents/slides/llm-agents-berkeley-intro-sp25.pdf}{intro} \href{https://www.youtube.com/live/g0Dwtf3BH-0}{video} & \hyperlink{mod-RSN}{\tg{RSN}} \tg{ORG}\\ \midrule
\hypertarget{lec-BRKADV-2}{}2 & Mon Feb 3, 2025 & \textbf{Learning to reason with LLMs} & Speaker: Jason Weston (Meta). Reasoning via preference optimization and self-verification & DPO \href{https://arxiv.org/abs/2305.18290}{[arXiv:2305.18290]} ; Iterative Reasoning Preference Optimization \href{https://arxiv.org/abs/2404.19733}{[arXiv:2404.19733]} ; Chain-of-Verification \href{https://arxiv.org/abs/2309.11495}{[arXiv:2309.11495]} & \href{https://rdi.berkeley.edu/adv-llm-agents/slides/Jason-Weston-Reasoning-Alignment-Berkeley-Talk.pdf}{slides} \href{https://www.youtube.com/live/_MNlLhU33H0}{video} & \hyperlink{mod-TRAIN}{\tg{TRAIN}} \hyperlink{mod-RSN}{\tg{RSN}}\\ \midrule
\hypertarget{lec-BRKADV-3}{}3 & Mon Feb 10, 2025 & \textbf{On Reasoning, Memory, and Planning of Language Agents} & Speaker: Yu Su (Ohio State University). Implicit reasoning, long-term memory, planning with world models & Grokked Transformers \href{https://arxiv.org/abs/2405.15071}{[arXiv:2405.15071]} ; HippoRAG \href{https://arxiv.org/abs/2405.14831}{[arXiv:2405.14831]} ; Is Your LLM Secretly a World Model of the Internet? \href{https://arxiv.org/abs/2411.06559}{[arXiv:2411.06559]} & \href{https://rdi.berkeley.edu/adv-llm-agents/slides/language_agents_YuSu_Berkeley.pdf}{slides} \href{https://www.youtube.com/live/zvI4UN2_i-w}{video} & \hyperlink{mod-RSN}{\tg{RSN}} \hyperlink{mod-MEM}{\tg{MEM}} \hyperlink{mod-ARCH}{\tg{ARCH}}\\ \midrule
\hypertarget{lec-BRKADV-4}{}4 & Mon Feb 24, 2025 & \textbf{Open Training Recipes for Reasoning in Language Models} & Speaker: Hanna Hajishirzi (University of Washington). Open post-training (OLMo/Tulu), DPO vs PPO & Tulu 3 \href{https://arxiv.org/abs/2411.15124}{[arXiv:2411.15124]} ; Unpacking DPO and PPO \href{https://arxiv.org/abs/2406.09279}{[arXiv:2406.09279]} ; OpenScholar \href{https://arxiv.org/abs/2411.14199}{[arXiv:2411.14199]} & \href{https://rdi.berkeley.edu/adv-llm-agents/slides/OLMo-Tulu-Reasoning-Hanna.pdf}{slides} \href{https://www.youtube.com/live/cMiu3A7YBks}{video} & \hyperlink{mod-TRAIN}{\tg{TRAIN}} \hyperlink{mod-RSN}{\tg{RSN}}\\ \midrule
\hypertarget{lec-BRKADV-5}{}5 & Mon Mar 3, 2025 & \textbf{Coding Agents and AI for Vulnerability Detection} & Speaker: Charles Sutton (Google DeepMind). Coding agents, interactive tools, vulnerability finding (Big Sleep) & Interactive Tools for LM Agents \href{https://arxiv.org/abs/2409.16165}{[arXiv:2409.16165]} ; From Naptime to Big Sleep \href{https://googleprojectzero.blogspot.com/2024/10/from-naptime-to-big-sleep.html}{[googleprojectzero.blogspot.com]} & \href{https://rdi.berkeley.edu/adv-llm-agents/slides/Code\%20Agents\%20and\%20AI\%20for\%20Vulnerability\%20Detection.pdf}{slides} \href{https://www.youtube.com/live/JCk6qJtaCSU}{video} & \hyperlink{mod-CODE}{\tg{CODE}} \hyperlink{mod-TOOL}{\tg{TOOL}} \hyperlink{mod-SAFE}{\tg{SAFE}}\\ \midrule
\hypertarget{lec-BRKADV-6}{}6 & Mon Mar 10, 2025 & \textbf{Multimodal Autonomous AI Agents} & Speaker: Ruslan Salakhutdinov (CMU/Meta). Web agents, visual web benchmarks, tree search & Mind2Web \href{https://arxiv.org/abs/2306.06070}{[arXiv:2306.06070]} ; WebArena \href{https://arxiv.org/abs/2307.13854}{[arXiv:2307.13854]} ; VisualWebArena \href{https://jykoh.com/vwa}{[jykoh.com]} ; Tree Search for LM Agents \href{https://jykoh.com/search-agents}{[jykoh.com]} & \href{https://rdi.berkeley.edu/adv-llm-agents/slides/ruslan-multimodal.pdf}{slides} \href{https://www.youtube.com/live/RPINOYM12RU}{video} & \hyperlink{mod-WEB}{\tg{WEB}} \hyperlink{mod-EVAL}{\tg{EVAL}} \hyperlink{mod-RSN}{\tg{RSN}}\\ \midrule
\hypertarget{lec-BRKADV-7}{}7 & Mon Mar 17, 2025 & \textbf{Multimodal Agents -- From Perception to Action} & Speaker: Caiming Xiong (Salesforce AI Research). Computer-use/GUI agents, OS benchmarks & OSWorld \href{https://arxiv.org/pdf/2404.07972}{[arXiv:2404.07972]} ; AGUVIS \href{https://arxiv.org/pdf/2412.04454}{[arXiv:2412.04454]} & \href{https://rdi.berkeley.edu/adv-llm-agents/slides/Multimodal_Agent_caiming.pdf}{slides} \href{https://www.youtube.com/live/n__Tim8K2IY}{video} & \hyperlink{mod-WEB}{\tg{WEB}} \hyperlink{mod-EVAL}{\tg{EVAL}}\\ \midrule
\hypertarget{lec-BRKADV-8}{}8 & Mon Mar 31, 2025 (10am-noon PT) & \textbf{AlphaProof: when reinforcement learning meets formal mathematics} & Speaker: Thomas Hubert (Google DeepMind). RL plus Lean for IMO-level proofs & IMO silver-medal blog \href{https://deepmind.google/discover/blog/ai-solves-imo-problems-at-silver-medal-level/}{[deepmind.google]} ; Mastering Chess and Shogi \href{https://arxiv.org/pdf/1712.01815}{[arXiv:1712.01815]} ; The Future of Mathematics? \href{https://www.youtube.com/watch?v=Dp-mQ3HxgDE}{[video]} ; Building the Mathematical Library of the Future \href{https://www.quantamagazine.org/building-the-mathematical-library-of-the-future-20201001/}{[quantamagazine.org]} & \href{https://rdi.berkeley.edu/adv-llm-agents/slides/alphaproof.pdf}{slides} \href{https://www.youtube.com/live/3gaEMscOMAU}{video} & \hyperlink{mod-RSN}{\tg{RSN}} \hyperlink{mod-TRAIN}{\tg{TRAIN}} \hyperlink{mod-APP}{\tg{APP}}\\ \midrule
\hypertarget{lec-BRKADV-9}{}9 & Mon Apr 7, 2025 & \textbf{Language models for autoformalization and theorem proving} & Speaker: Kaiyu Yang (Meta FAIR). Autoformalization, retrieval-augmented Lean proving & LeanDojo \href{https://arxiv.org/abs/2306.15626}{[arXiv:2306.15626]} ; Autoformalization with LLMs \href{https://arxiv.org/abs/2205.12615}{[arXiv:2205.12615]} ; Autoformalizing Euclidean Geometry \href{https://arxiv.org/abs/2405.17216}{[arXiv:2405.17216]} & \href{https://rdi.berkeley.edu/adv-llm-agents/slides/mathverification.pdf}{slides} \href{https://www.youtube.com/live/cLhWEyMQ4mQ}{video} & \hyperlink{mod-RSN}{\tg{RSN}} \hyperlink{mod-APP}{\tg{APP}} \hyperlink{mod-TOOL}{\tg{TOOL}}\\ \midrule
\hypertarget{lec-BRKADV-10}{}10 & Mon Apr 14, 2025 & \textbf{Bridging Informal and Formal Mathematical Reasoning (Berkeley page: "Advanced topics in theorem proving")} & Speaker: Sean Welleck (CMU). Informal-to-formal proving, context, self-taught Lean reasoning & Draft, Sketch, and Prove \href{https://arxiv.org/abs/2210.12283}{[arXiv:2210.12283]} ; miniCTX \href{https://www.arxiv.org/pdf/2408.03350}{[arXiv:2408.03350]} ; Lean-STaR \href{https://arxiv.org/abs/2407.10040}{[arXiv:2407.10040]} ; ImProver \href{https://arxiv.org/abs/2410.04753}{[arXiv:2410.04753]} & \href{https://rdi.berkeley.edu/adv-llm-agents/slides/welleck2025_berkeley_bridging.pdf}{slides} \href{https://www.youtube.com/live/Gy5Nm17l9oo}{video} & \hyperlink{mod-RSN}{\tg{RSN}} \hyperlink{mod-APP}{\tg{APP}}\\ \midrule
\hypertarget{lec-BRKADV-11}{}11 & Mon Apr 21, 2025 (10am-noon PT) & \textbf{Abstraction and Discovery with LLM Agents} & Speaker: Swarat Chaudhuri (UT Austin). In-context theorem-proving agents, learned concept libraries & An In-Context Learning Agent for Formal Theorem-Proving \href{https://arxiv.org/abs/2310.04353}{[arXiv:2310.04353]} ; Symbolic Regression with a Learned Concept Library \href{https://arxiv.org/abs/2409.09359}{[arXiv:2409.09359]} & \href{https://rdi.berkeley.edu/adv-llm-agents/slides/swarat.pdf}{slides} \href{https://www.youtube.com/live/IHc0TEMrEdY}{video} & \hyperlink{mod-RSN}{\tg{RSN}} \hyperlink{mod-ARCH}{\tg{ARCH}} \hyperlink{mod-APP}{\tg{APP}}\\ \midrule
\hypertarget{lec-BRKADV-12}{}12 & Mon Apr 28, 2025 & \textbf{Towards building safe and secure agentic AI} & Speaker: Dawn Song (UC Berkeley). Agent security: prompt-injection detection, memory/RAG poisoning, privilege control & Privtrans \href{https://dawnsong.io/papers/privtrans.pdf}{[PDF]} ; DataSentinel \href{https://arxiv.org/abs/2504.11358}{[arXiv:2504.11358]} ; AgentPoison \href{https://arxiv.org/abs/2407.12784}{[arXiv:2407.12784]} ; Progent \href{https://arxiv.org/html/2504.11703v1}{[arXiv:2504.11703]} & \href{https://rdi.berkeley.edu/adv-llm-agents/slides/dawn-agentic-ai.pdf}{slides} \href{https://www.youtube.com/live/ti6yPE2VPZc}{video} & \hyperlink{mod-SAFE}{\tg{SAFE}} \hyperlink{mod-DEPLOY}{\tg{DEPLOY}}\\ \midrule
\end{longtable}}
\section{Assignments and projects}
{\scriptsize\setlength{\tabcolsep}{2.5pt}\begin{longtable}{P{2.8cm} P{5.4cm} P{3.0cm} P{3.5cm} P{1.5cm}}
\toprule Name & What you build & Skills & Starter code / resources & Tags\\ \midrule\endhead
\hypertarget{asg-BRKADV-1}{}\textbf{Lab 1: Dataset contribution (Lean)} & 2 verifiable coding problems, each with \texttt{description.\allowbreak{}txt}, \texttt{signature.\allowbreak{}json}, \texttt{test.\allowbreak{}json} (at least 5 tests with expected and "unexpected" values) and a \texttt{task.\allowbreak{}lean} holding a \texttt{\_\allowbreak{}code} implementation and \texttt{\_\allowbreak{}spec} theorems between fixed CODE/SPEC markers; implementations should be non-trivial (about 30+ lines) and the specs should fully constrain behavior & Lean 4 functional programming and tactic proofs (induction, simp, rw, cases), writing formal specs, test design & \href{https://drive.google.com/file/d/1e5Q47rFHzrKS2jHK5WUCfUTSboa-jHtU/view}{[drive.google.com]} (instructions); \href{https://drive.google.com/file/d/1_GlMZgKbDq56Tsw6kCe-vLU6s2e2IF0q/view}{[drive.google.com]} (lab1-datapoint.zip, not opened) & \hyperlink{mod-CODE}{\tg{CODE}} \hyperlink{mod-RSN}{\tg{RSN}} \hyperlink{mod-EVAL}{\tg{EVAL}}\\ \midrule
\hypertarget{asg-BRKADV-2}{}\textbf{Lab 2: Verified Lean coding agent} & \texttt{main\_\allowbreak{}workflow} in \texttt{src/\allowbreak{}main.\allowbreak{}py}: from a task description and Lean template, return \texttt{"code"} and \texttt{"proof"}; suggested planning / generation / verification agents with Lean-error feedback loops and optional RAG over \texttt{documents/\allowbreak{}} (\texttt{src/\allowbreak{}embedding\_\allowbreak{}db.\allowbreak{}py}); graded by compiling the code with the proof as \texttt{sorry}, then with the proof; \texttt{sorry}, trivial theorems and hard-coding get no credit; submit with \texttt{make zip}, plus an optional 2-page design PDF & Agent loop design, compiler-in-the-loop verification, retry/feedback, RAG, multi-agent decomposition, Lean proofs & \href{https://drive.google.com/file/d/1_8ASqnVR99BTiKd-I7g8Ov7-UiXsJeHW/view}{[drive.google.com]} (instructions); \href{https://drive.google.com/file/d/1Yr99_WrQoqoO9DzwwPMmbyfwQaRPihD8/view}{[drive.google.com]} (lab2-starter-code.zip, not opened) & \hyperlink{mod-ARCH}{\tg{ARCH}} \hyperlink{mod-CODE}{\tg{CODE}} \hyperlink{mod-RSN}{\tg{RSN}} \hyperlink{mod-MULTI}{\tg{MULTI}} \hyperlink{mod-MEM}{\tg{MEM}}\\ \midrule
\hypertarget{asg-BRKADV-3}{}\textbf{Course project (Berkeley, 2+ units)} & Team project, applications track (3-4) or research track (2-3, by application): proposal (due 2/24), milestone (due 3/31), poster (5/5), recording and report (5/16); 3/4 units also graded on implementation & Research/product design, implementation, writing, presenting & -- & \tg{ORG} \hyperlink{mod-APP}{\tg{APP}}\\ \midrule
\hypertarget{asg-BRKADV-4}{}\textbf{AgentX competition (MOOC Ninja/Legendary tiers)} & Entrepreneurship track (pitch deck, demo video, live product) or Research track (7-8 page paper, video, public GitHub repo); build period Mar 31 - May 31, 2025; Demo Day Aug 2, 2025 at the Agentic AI Summit & Building and pitching LLM agent products or research & n/a (\href{https://rdi.berkeley.edu/agentx/}{[rdi.berkeley.edu]}) & \hyperlink{mod-APP}{\tg{APP}} \tg{ORG}\\ \midrule
\hypertarget{asg-BRKADV-5}{}\textbf{Written article} & About 500-word public post (lecture summary or AgentX write-up), pass/no pass & Technical writing & -- & \tg{ORG}\\ \midrule
\hypertarget{asg-BRKADV-6}{}\textbf{Quizzes (12)} & One completion-graded quiz per lecture, 5 multiple-choice questions each (Google Forms) & Lecture review & -- & \tg{ORG}\\ \midrule
\end{longtable}}
\chapter{Stanford 329T: CS329T Trustworthy ML: Building and Evaluating Agentic Systems}\hypertarget{course-S329T}{}
{\small\begin{tabularx}{\linewidth}{@{}>{\bfseries}l X@{}}\toprule
Institution & Stanford University\\
Term & Fall 2025\\
Instructors & Anupam Datta, John Mitchell\\
Website & \href{https://web.stanford.edu/class/cs329t/syllabus.html}{https://web.stanford.edu/class/cs329t/syllabus.html}\\
Level & Graduate; CS229 + deep learning\\
Archetype & Project course on agent evaluation\\
Public materials & Slides: 11 decks (10 sessions); video: none (not recorded); code: HW1 handout PDF; starter branch 404\\
Best for & Evaluation-first agent building: Goal-Plan-Action evaluation, LLM judges, MCP tool descriptions, GEPA and red-teaming.\\
\bottomrule\end{tabularx}}

A project course on building agentic applications and then evaluating them so they can be improved step by step.
After LLM background, prompting and RAG, it introduces its main evaluation idea, checking whether an agent's goals, plans and actions line up (\emph{Agent GPA}).
Guest talks cover systematic agent evaluation and refining LLM judges, MCP and writing tool descriptions, data curation and RL for agents, GEPA, and agent security.
The one homework traces a LangGraph multi-agent research system in TruLens, has you label its failure modes by hand, and compares your labels with an LLM judge's.
The starter branch linked from the handout no longer exists; an updated version of the notebook is on the TruLens main branch.
\section{Lectures}
{\scriptsize\setlength{\tabcolsep}{2.5pt}\begin{longtable}{P{0.45cm} P{1.2cm} P{2.5cm} P{3.6cm} P{3.6cm} P{3.3cm} P{1.5cm}}
\toprule \# & Date & Title & Key topics & Readings & Materials & Tags\\ \midrule\endhead
\hypertarget{lec-S329T-1}{}1 & Tue Sep 23 & \textbf{Course introduction and overview} & Logistics; 2023/2024 project examples; LLM background; prompting, fine-tuning, agentic AI & LangGraph agentic concepts \href{https://langchain-ai.github.io/langgraph/concepts/agentic_concepts/}{[langchain-ai.github.io]} ; OpenAI "A practical guide to building agents" \href{https://cdn.openai.com/business-guides-and-resources/a-practical-guide-to-building-agents.pdf}{[PDF]} & \href{https://web.stanford.edu/class/cs329t/slides/Lecture\%201\%20-\%20CS\%20329T\%20Fall\%202025.pdf}{slides} & \tg{ORG} \hyperlink{mod-FND}{\tg{FND}} \hyperlink{mod-ARCH}{\tg{ARCH}}\\ \midrule
\hypertarget{lec-S329T-2}{}2 & Thu Sep 25 & \textbf{Models, Prompting and RAG} & LLM strengths and limits; prompting techniques; retrieval-augmented generation (HW1 released) & \href{https://arxiv.org/abs/2407.12858}{[arXiv:2407.12858]} ; \href{https://arxiv.org/abs/2406.06608}{[arXiv:2406.06608]} ; \href{https://arxiv.org/pdf/2202.03629}{[arXiv:2202.03629]} ; supplement: DeepLearning.AI "Building and Evaluating Data Agents" \href{https://www.deeplearning.ai/short-courses/building-and-evaluating-data-agents/}{[deeplearning.ai]} & \href{https://web.stanford.edu/class/cs329t/slides/Lecture\%202\%20-\%20CS\%20329T\%20Fall\%202025.pdf}{slides}, \href{https://web.stanford.edu/class/cs329t/homeworks/CS329T_HW1_Fall2025.pdf}{HW1} & \hyperlink{mod-FND}{\tg{FND}} \hyperlink{mod-MEM}{\tg{MEM}}\\ \midrule
\hypertarget{lec-S329T-3}{}3 & Tue Sep 30 & \textbf{Agentic AI} & Agent foundations; sample agent systems; evaluating goals, plans and actions & Anthropic "Building effective agents" \href{https://www.anthropic.com/engineering/building-effective-agents}{[anthropic.com]} ; "What is your Agent's GPA? Evaluating and improving agent trustworthiness" (blog; no URL given) & \href{https://web.stanford.edu/class/cs329t/slides/Lecture\%203\%20--\%20CS\%20329T\%20Fall\%202025.pdf}{slides} & \hyperlink{mod-ARCH}{\tg{ARCH}} \hyperlink{mod-EVAL}{\tg{EVAL}}\\ \midrule
\hypertarget{lec-S329T-4}{}4 & Thu Oct 2 & \textbf{Recap, further discussion, and planning for projects} & Review of lectures 2-3; project directions & Together "Open Deep Research" cookbook \href{https://github.com/togethercomputer/together-cookbook/blob/main/Agents/Together_Open_Deep_Research_CookBook.ipynb}{[GitHub]} & \href{https://web.stanford.edu/class/cs329t/slides/Lecture\%204\%20--\%20CS\%20329T\%20Fall\%202025.pdf}{slides} & \tg{ORG} \hyperlink{mod-WEB}{\tg{WEB}}\\ \midrule
\hypertarget{lec-S329T-5}{}5 & Tue Oct 7 & \textbf{Project proposals and feedback - Group 1} & Student proposals & -- & -- & \tg{ORG}\\ \midrule
\hypertarget{lec-S329T-6}{}6 & Thu Oct 9 & \textbf{Project proposals and feedback - Group 2} & Student proposals & -- & -- & \tg{ORG}\\ \midrule
\hypertarget{lec-S329T-7}{}7 & Tue Oct 14 & \textbf{Systematic evaluation of agentic systems (guest: Allison Jia)} & Goal/plan/action evaluation; refining the LLM judge; evaluation with and without ground truth (HW1 due 4:30pm) & \href{https://arxiv.org/abs/2507.21504}{[arXiv:2507.21504]} ; \href{https://arxiv.org/abs/2410.10934}{[arXiv:2410.10934]} ; \href{https://arxiv.org/abs/2510.08847v1}{[arXiv:2510.08847]} & \href{https://web.stanford.edu/class/cs329t/slides/Allison\%20Jia\%20-\%20slides.pdf}{slides} & \hyperlink{mod-EVAL}{\tg{EVAL}}\\ \midrule
\hypertarget{lec-S329T-8}{}8 & Thu Oct 16 & \textbf{MCP, effective tool descriptions, evals (guests: Yusuf Ozuysal, Ishita Gupta)} & Model Context Protocol; writing tools for agents & Anthropic MCP announcement \href{https://www.anthropic.com/news/model-context-protocol}{[anthropic.com]} ; MCP intro \href{http://modelcontextprotocol.io/docs/getting-started/intro}{[modelcontextprotocol.io]} ; Anthropic "Writing tools for agents" \href{https://www.anthropic.com/engineering/writing-tools-for-agents}{[anthropic.com]} & \href{https://web.stanford.edu/class/cs329t/slides/Model\%20Context\%20Protocol\%20\%28\%20Yusuf\%20Ozuysal\%29.pdf}{slides (Ozuysal)}, \href{https://web.stanford.edu/class/cs329t/slides/ishita-mcp.pdf}{slides (Gupta)} & \hyperlink{mod-TOOL}{\tg{TOOL}} \hyperlink{mod-EVAL}{\tg{EVAL}}\\ \midrule
\hypertarget{lec-S329T-9}{}9 & Tue Oct 21 & \textbf{Deep dive into LLMs} & LLM internals and training overview & Andrej Karpathy video \href{https://www.youtube.com/watch?v=7xTGNNLPyMI}{[video]} & \href{https://web.stanford.edu/class/cs329t/slides/Lecture\%207\%20--\%20CS\%20329T\%20Fall\%202025.pdf}{slides} & \hyperlink{mod-FND}{\tg{FND}} \hyperlink{mod-TRAIN}{\tg{TRAIN}}\\ \midrule
\hypertarget{lec-S329T-10}{}10 & Thu Oct 23 & \textbf{Data curation and RL to optimize agents (guest: Trung Vu, Bespoke Labs)} & Synthetic data; multi-turn tool-use RL; agent environments and benchmarks & Bespoke Labs blog \href{https://www.bespokelabs.ai/blog/improving-multi-turn-tool-use-with-reinforcement-learning}{[bespokelabs.ai]} ; \href{https://arxiv.org/abs/2504.04736}{[arXiv:2504.04736]} ; environments: SWE-bench \href{https://arxiv.org/abs/2310.06770}{[arXiv:2310.06770]} , SWE-smith \href{https://arxiv.org/abs/2504.21798}{[arXiv:2504.21798]} , ARE \href{https://arxiv.org/html/2509.17158v1}{[arXiv:2509.17158]} (GAIA2 \href{https://huggingface.co/blog/gaia2}{[huggingface.co]}), OSWorld \href{https://arxiv.org/abs/2404.07972}{[arXiv:2404.07972]} , Terminal-Bench \href{https://www.tbench.ai/}{[tbench.ai]} ; RL: DAPO \href{https://arxiv.org/abs/2503.14476}{[arXiv:2503.14476]} , ScaleRL \href{https://arxiv.org/pdf/2510.13786}{[arXiv:2510.13786]} ; optional: \href{https://arxiv.org/abs/2506.04178}{[arXiv:2506.04178]} , \href{https://arxiv.org/abs/2507.01921}{[arXiv:2507.01921]} & \href{https://web.stanford.edu/class/cs329t/slides/Lecture\%208\%20-\%20\%5B10_23_2025\%5D\%20Reinforcement\%20Learning\%20\%20for\%20Agents.pdf}{slides} & \hyperlink{mod-TRAIN}{\tg{TRAIN}} \hyperlink{mod-TOOL}{\tg{TOOL}} \hyperlink{mod-EVAL}{\tg{EVAL}}\\ \midrule
\hypertarget{lec-S329T-11}{}11 & Tue Oct 28 & \textbf{GEPA (guest: Lakshya Agrawal)} & Reflective prompt evolution vs. RL & GEPA \href{https://arxiv.org/abs/2507.19457}{[arXiv:2507.19457]} ; DSPy \href{https://dspy.ai/}{[dspy.ai]} ; Colab \href{https://colab.research.google.com/drive/1-Ed72XYr_XXJx0127fTxneBY4CUO0Bmv}{[colab.research.google.com]} ; ARC-AGI notebook \href{https://github.com/gepa-ai/gepa/blob/arc_agi/src/gepa/examples/dspy_full_program_evolution/arc_agi.ipynb}{[GitHub]} & \href{https://web.stanford.edu/class/cs329t/slides/Lecture\%209\%20-\%2020251028\%20GEPA_Talk.pdf}{slides} & \hyperlink{mod-TRAIN}{\tg{TRAIN}} \hyperlink{mod-RSN}{\tg{RSN}}\\ \midrule
\hypertarget{lec-S329T-12}{}12 & Thu Oct 30 & \textbf{Security for agents: adversarial attacks and guardrails (guests: Matt Fredrikson, Andy Zou)} & Attacks, defenses and benchmarks for agents & Offense: \href{https://arxiv.org/abs/2307.15043}{[arXiv:2307.15043]} , \href{https://arxiv.org/abs/2507.20526}{[arXiv:2507.20526]} ; defense: \href{https://arxiv.org/abs/2310.01405}{[arXiv:2310.01405]} , \href{https://arxiv.org/abs/2406.04313}{[arXiv:2406.04313]} ; benchmarks: \href{https://arxiv.org/abs/2402.04249}{[arXiv:2402.04249]} , \href{https://arxiv.org/abs/2410.09024}{[arXiv:2410.09024]} , Gray Swan Arena \href{https://app.grayswan.ai/arena}{[app.grayswan.ai]} & \href{https://web.stanford.edu/class/cs329t/slides/lecture\%2010\%20-\%2020251030\%20-\%20security\%20for\%20agents\%20-\%20andy\%20zou.pdf}{slides} & \hyperlink{mod-SAFE}{\tg{SAFE}} \hyperlink{mod-EVAL}{\tg{EVAL}}\\ \midrule
\hypertarget{lec-S329T-13}{}13 & Thu Nov 6 & \textbf{Mid-term project presentations and feedback - Group 1} & Student presentations & -- & -- & \tg{ORG}\\ \midrule
\hypertarget{lec-S329T-14}{}14 & Tue Nov 11 & \textbf{Mid-term project presentations and feedback - Group 2} & Student presentations & -- & -- & \tg{ORG}\\ \midrule
\hypertarget{lec-S329T-15}{}15 & Thu Nov 13 & \textbf{Multi-modal frontier models with agentic capabilities (guest: Deepak Ramachandran, Google DeepMind)} & Multimodal agents & Reading: TBD & n/a (no slides posted) & \hyperlink{mod-FND}{\tg{FND}} \hyperlink{mod-ARCH}{\tg{ARCH}}\\ \midrule
\hypertarget{lec-S329T-16}{}16 & Tue Nov 18 & \textbf{Project semi-final presentation, with feedback - Group 1} & Student presentations & -- & -- & \tg{ORG}\\ \midrule
\hypertarget{lec-S329T-17}{}17 & Thu Nov 20 & \textbf{Project semi-final presentation, with feedback - Group 2} & Student presentations & -- & -- & \tg{ORG}\\ \midrule
\hypertarget{lec-S329T-18}{}18 & Dec 2 or 4 & \textbf{Final Project Fair} & 2.5-hour session, 3:30-6:00pm: 2-3 minute team openers, then poster and demo session & -- & -- & \tg{ORG} \hyperlink{mod-APP}{\tg{APP}}\\ \midrule
\end{longtable}}
\section{Assignments and projects}
{\scriptsize\setlength{\tabcolsep}{2.5pt}\begin{longtable}{P{2.8cm} P{5.4cm} P{3.0cm} P{3.5cm} P{1.5cm}}
\toprule Name & What you build & Skills & Starter code / resources & Tags\\ \midrule\endhead
\hypertarget{asg-S329T-1}{}\textbf{HW1: Measuring and Improving Agent GPA (released Sep 25, due Oct 14 4:30pm; 15\% of grade)} & Work through a provided notebook with a LangGraph multi-agent web-research system (planner and executor nodes plus other agents; Tavily web search, OpenAI models), traced in TruLens. Annotate traces by hand with Goal-Plan-Action failure modes, rerun with LLM-judge metrics (\texttt{session.\allowbreak{}get\_\allowbreak{}leaderboard()}), compare judge scores with your annotations, then change prompts or architecture and rerun. Prerequisite: DeepLearning.AI "Building and Evaluating Data Agents" (about 2 hours). Deliverable: one-page summary (annotation rationale, metrics analysis, justified improvements) plus a notebook printout, as one PDF on Gradescope. Rubric: completion 40\%, annotation quality 20\%, metrics analysis 30\%, improvements 10\% & LangGraph agent orchestration, web-search tools, tracing, human annotation of traces, LLM-as-judge validation against human labels, failure-mode analysis, prompt/architecture iteration & Linked in the PDF (404 as of 2026-10-09): \href{https://github.com/truera/trulens/blob/josh/stanford-course/examples/experimental/web-search-agent-human-annotation-flow.ipynb}{[GitHub]} . A notebook with the same filename exists on main: \href{https://github.com/truera/trulens/blob/main/examples/experimental/web-search-agent-human-annotation-flow.ipynb}{[GitHub]} (added to main 2025-10-30 in PR \#2237, after HW1 was due; since changed for LangChain 1.0 (\#2279) and the Feedback-to-Metric API migration (\#2513), so it is not the version students used. Branch josh/stanford-course no longer exists; checked via GitHub API 2026-10-09) & \hyperlink{mod-EVAL}{\tg{EVAL}} \hyperlink{mod-MULTI}{\tg{MULTI}} \hyperlink{mod-WEB}{\tg{WEB}} \hyperlink{mod-TOOL}{\tg{TOOL}} \hyperlink{mod-ARCH}{\tg{ARCH}}\\ \midrule
\hypertarget{asg-S329T-2}{}\textbf{Course project (70\%, teams of 2-3)} & Team-built agentic AI application with evaluation: proposal presentation (week 3), mid-term presentations (weeks 7-8), semi-final dry-runs (week 9), final presentation and poster/demo fair (week 10), written report due at final exam time & Project scoping, agent building, evaluation design, presenting & -- & \hyperlink{mod-APP}{\tg{APP}} \hyperlink{mod-EVAL}{\tg{EVAL}} \tg{ORG}\\ \midrule
\end{longtable}}
\chapter{Supplementary Resources}\hypertarget{supplements}{}
These non-course resources (from the bookmarks folder, plus the MIT agent-security lecture added later) are referenced throughout the modules.
\section{Model Context Protocol docs (version 2026-07-28)}\label{sup-mcp}
{\small \href{https://modelcontextprotocol.io/docs/2026-07-28/getting-started/intro}{https://modelcontextprotocol.io/docs/2026-07-28/getting-started/intro} \quad \hyperlink{mod-TOOL}{\tg{TOOL}} \hyperlink{mod-ARCH}{\tg{ARCH}}}\par
MCP is the open standard for connecting an AI application (the \emph{host}) to external context and tools. The host creates one \emph{client} per \emph{server}. Servers expose \emph{tools} (callable functions with JSON-Schema inputs), \emph{resources} (contextual data) and \emph{prompts} (templates), over \emph{stdio} for local servers or \emph{Streamable HTTP} for remote ones, with OAuth recommended.
\textbf{Version caveat:} the 2026-07-28 revision makes the protocol stateless. Every request carries its version and capabilities in \texttt{\_meta}, servers must support a new \texttt{server/discover} method that replaces the old initialize handshake (clients may call it, but can also send requests without it), change notifications move to an opt-in \texttt{subscriptions/listen} stream, and \emph{sampling} and \emph{logging} are deprecated. Course material written before mid-2026 (including several lectures in this guide) teaches the older stateful handshake. The concepts still carry over, but check the current docs before you write code.
\section{Anthropic Engineering: Code execution with MCP}\label{sup-codeexec}
{\small \href{https://www.anthropic.com/engineering/code-execution-with-mcp}{https://www.anthropic.com/engineering/code-execution-with-mcp} \quad \hyperlink{mod-TOOL}{\tg{TOOL}} \hyperlink{mod-MEM}{\tg{MEM}} \hyperlink{mod-DEPLOY}{\tg{DEPLOY}}}\par
Argues that agents connected to many MCP servers should call tools \emph{through code} instead of through direct tool calls. Each server is exposed as a file tree of small code APIs. The agent reads only the definitions it needs (progressive disclosure) and writes a script that runs in a sandbox, so intermediate data never passes through the context window. In the worked example, context drops from about 150{,}000 to 2{,}000 tokens. The cost is running a secure sandbox. This ties tool use to context engineering and to Skills, and CU Boulder assigns a paper that studies the same design choice (\hyperlink{lec-CU-24}{CU Boulder~24}).
\section{OpenClaw documentation}\label{sup-openclaw}
{\small \href{https://docs.openclaw.ai/}{https://docs.openclaw.ai/} \quad \hyperlink{mod-ARCH}{\tg{ARCH}} \hyperlink{mod-DEPLOY}{\tg{DEPLOY}} \hyperlink{mod-SAFE}{\tg{SAFE}}}\par
An MIT-licensed, self-hosted, always-on personal assistant. A single \emph{gateway} process owns sessions, routing and \emph{channels} (Discord, Signal, Slack, WhatsApp, \ldots). Agents get tools, cron jobs, hooks, webhooks and on-demand skills from the ClawHub registry. It is a useful reference architecture for long-running, multi-channel agents and a cautionary case study: CMU's safety lecture cites an OpenClaw incident report (\hyperlink{lec-CMU-13}{CMU~13}), and UVA discusses the OpenClaw family in its personal-agents lecture (\hyperlink{lec-UVA-18}{UVA~18}). \emph{Catalogued only.} Under your install policy, vet it before installing anything.
\section{Berkeley BAIR: TinyAgent, Function Calling at the Edge}\label{sup-tinyagent}
{\small \href{https://bair.berkeley.edu/blog/2024/05/29/tiny-agent/}{https://bair.berkeley.edu/blog/2024/05/29/tiny-agent/} \quad \hyperlink{mod-TOOL}{\tg{TOOL}} \hyperlink{mod-TRAIN}{\tg{TRAIN}} \hyperlink{mod-DEPLOY}{\tg{DEPLOY}}}\par
A case study in specializing a \emph{small} model for agentic tool use. GPT-4-Turbo generates synthetic function-calling plans, which are validated programmatically. A 1.1B or 7B model is LoRA-fine-tuned to emit LLMCompiler-style DAG plans over 16 macOS tools. A small DeBERTa classifier replaces embedding retrieval for tool selection, roughly halving the prompt. The model is then quantized to 4-bit for llama.cpp. The 1.1B model (80.4\%) beats GPT-4-Turbo (79.1\%) on this fixed toolset while running locally.
Code: \href{https://github.com/SqueezeAILab/TinyAgent}{SqueezeAILab/TinyAgent} (MIT). Read it after \hyperlink{mod-TOOL}{Module~3 (Tool Use and Function Calling)} and the SFT lecture in \hyperlink{mod-TRAIN}{Module~11 (Training and Optimizing Agents)}.
\section{MIT 6.566 guest lecture: AI Agent Security (Anish Athalye, April 2026)}\label{sup-agentsec}
{\small \href{https://github.com/anishathalye/ai-agent-security-lecture}{https://github.com/anishathalye/ai-agent-security-lecture} \quad \hyperlink{mod-SAFE}{\tg{SAFE}} \hyperlink{mod-TOOL}{\tg{TOOL}} \hyperlink{mod-ARCH}{\tg{ARCH}}}\par
A single lecture, with a \href{https://www.youtube.com/watch?v=w0oGeKxD5Fc}{recording} and a repository of thirteen small Python scripts, that builds an agent from scratch and then attacks it. The scripts go from raw completion and chat templates, through tools, native tool calling, an agent loop and a CodeAct-style code agent, to a web-summarization prompt-injection demo. They end with two defenses: the dual-LLM pattern and CaMeL, which runs benign and adversarial scenarios with and without CaMeL. The assigned reading is the CaMeL paper, \emph{Defeating Prompt Injections by Design} (Debenedetti et al., 2025). Some scripts need Ollama with a local model and others need an API key. No license is stated. Run it after \hyperlink{mod-TOOL}{Module~3 (Tool Use and Function Calling)} and alongside \hyperlink{mod-SAFE}{Module~12 (Safety, Security and Oversight)}. Under your install policy, vet \texttt{uv} and the dependencies first.
\section{Anthropic: Building effective agent automations (Martin and Avilla, October 2026)}\label{sup-anthauto}
{\small \href{https://claude.dev/blog/building-effective-agent-automations/}{https://claude.dev/blog/building-effective-agent-automations/} \quad \hyperlink{mod-DEPLOY}{\tg{DEPLOY}} \hyperlink{mod-SAFE}{\tg{SAFE}} \hyperlink{mod-TOOL}{\tg{TOOL}}}\par
A worked example of an unattended, scheduled agent (``daily-brief'', on Claude Managed Agents) that reads sources, tracks what changed and posts a brief to Slack. It is organised as six components: sources, destination, agent, schedule, memory and guardrails. Its rules are the concrete form of the durable-state and idempotency advice in \hyperlink{mod-DEPLOY}{Module~13 (Deployment, Efficiency and Operations)}. Read each source from a stored bookmark, not a fixed time window. Report a failed read as unreadable, never as a quiet day. Count a post as sent only when Slack confirms it, and only then update the bookmarks and the ledger. Re-read preferences every run from a store the agent cannot edit. Use read-only credentials wherever the agent only reads, and cap each run's spending at three to five times a normal run at first. Posting requires no human approval. Read-only permissions, a host allowlist and spending limits constrain the automation, but malicious source content can still alter the brief. The post gives design reasoning and failure examples but no measurements. Our suggested companion reading: the \hyperref[sup-openclaw]{OpenClaw} documentation, for comparison with an always-on architecture.
\section{Notion: How we built a personal data scientist for every Notion employee (David Thomas, October 2026)}\label{sup-notion}
{\small \href{https://www.notion.com/blog/how-we-built-a-personal-data-scientist-for-every-notion-employee}{https://www.notion.com/blog/how-we-built-a-personal-data-scientist-for-every-notion-employee} \quad \hyperlink{mod-MEM}{\tg{MEM}} \hyperlink{mod-DEPLOY}{\tg{DEPLOY}} \hyperlink{mod-EVAL}{\tg{EVAL}} \hyperlink{mod-APP}{\tg{APP}}}\par
A production case study of Data Scout, an internal agent that answers plain-language questions from governed Snowflake data through MCP. Each user signs in with their own credentials, so existing access controls apply. A curated, owned library separates data context (metrics, sources, filters), business context and routing guidance. Queries carry structured metadata so traces show which definitions were used. A nightly review agent records missing, stale or conflicting context; a separate triage agent deduplicates findings, applies authorised fixes, and sends ambiguous or higher-risk changes to a person. The post's main lesson is that a query can run cleanly and still answer the wrong question, through wrong meaning, routing, measurement or interpretation. Notion reports that weekly requests to its data team fell by about 80\% within two months and that about half of its employees use the agent. These figures are self-reported, and the post says its evaluation is still maturing: usage and query success do not prove accuracy. Read it with \hyperlink{mod-MEM}{Module~5 (Memory, Retrieval and Context Engineering)} for context engineering and \hyperlink{mod-DEPLOY}{Module~13 (Deployment, Efficiency and Operations)} for operations.
\section{Google Colab}\label{sup-colab}
{\small \href{https://colab.research.google.com/}{https://colab.research.google.com/} \quad \hyperlink{mod-DEPLOY}{\tg{DEPLOY}}}\par
A hosted Jupyter runtime. Use it to run course notebooks (UChicago, Northeastern, Mannheim, UVA) without a local setup, and to keep experiments with downloaded code off your own machine.
\part{Appendices}
\chapter{Codebase and Tool Index}
Repositories, frameworks, benchmarks and tools referenced by each course (as recorded during research).
\section{CMU}
{\small\begin{itemize}
\item Course GitHub org: \href{https://github.com/cmu-agents}{[GitHub]} (assignment-1, assignment-2, assignment-3, chess-app)
\item Modal (sandboxes and compute): \href{https://modal.com}{[modal.com]}
\item uv: \href{https://docs.astral.sh/uv}{[docs.astral.sh]}
\item OpenHands (lecture 2 readings, Frameworks 1): \href{https://github.com/All-Hands-AI/OpenHands}{[GitHub]} (repo URL from general knowledge, not seen on the course site)
\item LangGraph (Frameworks 2): \href{https://github.com/langchain-ai/langgraph}{[GitHub]} (from general knowledge)
\item Model Context Protocol spec (lecture 2): \href{https://modelcontextprotocol.io}{[modelcontextprotocol.io]} (from general knowledge)
\item Agent Skills protocol (A1 reference): URL n/a (named in the A1 handout)
\item Harbor (A2 task packaging): URL n/a
\item Tinker 0.24.0 and SkyRL Tinker server (A3): URL n/a. SkyRL and Miles are cited as RL-systems examples.
\item vLLM / SGLang (served policies, A1 and A3)
\item Other harness code cited in lecture 3 readings: Codex, OpenCode, Pi, Hermes Agent (URLs n/a)
\end{itemize}}
\section{Stanford '25}
{\small\begin{itemize}
\item Assignments repo: \href{https://github.com/mihail911/modern-software-dev-assignments}{[GitHub]} (Python 3.12, Poetry, Conda)
\item Ollama (local models in Fall 2025 week 1): \href{https://ollama.com}{[ollama.com]}
\item FastMCP: \href{https://github.com/jlowin/fastmcp}{[GitHub]}
\item MCP spec: \href{https://modelcontextprotocol.io/}{[modelcontextprotocol.io]} ; reference servers: \href{https://github.com/modelcontextprotocol/servers}{[GitHub]} ; TypeScript SDK: \href{https://github.com/modelcontextprotocol/typescript-sdk}{[GitHub]}
\item Claude Code docs: \href{https://code.claude.com/docs/en/setup}{[code.claude.com]} , skills: \href{https://code.claude.com/docs/en/skills}{[code.claude.com]} , settings: \href{https://docs.claude.com/en/docs/claude-code/settings}{[docs.claude.com]} ; Messages API: \href{https://docs.claude.com/en/api/messages}{[docs.claude.com]}
\item Agent Skills standard: \href{https://agentskills.io}{[agentskills.io]} ; sample skills: \href{https://github.com/mattpocock/skills}{[GitHub]}
\item mitmproxy (week 1, Fall 2026), plus the tutorial "Intercepting Claude Code Requests": \href{https://www.ai.moda/en/blog/tutorial-intercepting-claude-code-requests}{[ai.moda]}
\item HumanLayer advanced context engineering: \href{https://github.com/humanlayer/advanced-context-engineering-for-coding-agents}{[GitHub]}
\item awesome-claude-agents: \href{https://github.com/vijaythecoder/awesome-claude-agents}{[GitHub]} ; SuperClaude: \href{https://github.com/SuperClaude-Org/SuperClaude_Framework}{[GitHub]}
\item Warp: \href{https://www.warp.dev/}{[warp.dev]} ; Semgrep: \href{https://semgrep.dev/}{[semgrep.dev]} ; Graphite: \href{https://graphite.dev/}{[graphite.dev]} ; Vercel: \href{https://vercel.com/}{[vercel.com]} ; bolt.new ; Cursor ; Resolve.ai: \href{https://resolve.ai/}{[resolve.ai]} ; Cognition/Devin: \href{https://cognition.ai/}{[cognition.ai]}
\item Fall 2026 partners: Browserbase \href{https://www.browserbase.com}{[browserbase.com]} , OpenHands \href{https://www.openhands.dev}{[openhands.dev]} , CrewAI \href{https://www.crewai.com}{[crewai.com]} , CopilotKit \href{https://www.copilotkit.ai}{[copilotkit.ai]} , Milvus \href{https://milvus.io}{[milvus.io]} , Marimo \href{https://marimo.io}{[marimo.io]} , Phoenix \href{https://phoenix.arize.com}{[phoenix.arize.com]} , Unsloth \href{https://unsloth.ai}{[unsloth.ai]} , Anyscale \href{https://www.anyscale.com}{[anyscale.com]} , cmux \href{https://cmux.dev}{[cmux.dev]} , Pi \href{https://pi.dev}{[pi.dev]} , HeyGen \href{https://www.heygen.com}{[heygen.com]}
\item Paid 4-week Maven spin-off course (from a third-party source): \href{https://maven.com/the-modern-software-developer/ai-course}{[maven.com]}
\end{itemize}}
\section{Northeastern}
{\small\begin{itemize}
\item Course repo: \href{https://github.com/agenticllms/lecture-files}{[GitHub]} (folders DADA1, W02b, W03a, W03b, W04a, W04b, W05a, W05b; requirements.txt; example.env)
\item requirements.txt: langchain, langchain-openai, langchain-anthropic, langchain-google-genai, langchain-tavily, langchain-ollama, langgraph, python-dotenv, ipykernel, requests, gradio, langchain-openrouter (W05a adds langchain-pinecone, numpy, optional langchain-nvidia-ai-endpoints)
\item LangChain / LangGraph docs: \href{https://docs.langchain.com/oss/python/langgraph/workflows-agents}{[docs.langchain.com]} ; LangSmith eval/cost docs: \href{https://docs.langchain.com/langsmith/evaluation-concepts}{[docs.langchain.com]}
\item Ollama: \href{https://ollama.com/download}{[ollama.com]} ; Qwen 3.5 models: \href{https://ollama.com/library/qwen3.5}{[ollama.com]}
\item uv: \href{https://astral.sh/uv/install.sh}{[astral.sh]}
\item Pinecone: \href{https://www.pinecone.io}{[pinecone.io]}
\item Hugging Face Spaces: \href{https://huggingface.co/new-space}{[huggingface.co]}
\item Render: \href{https://render.com}{[render.com]} ; GitHub CLI: \href{https://cli.github.com}{[GitHub]}
\item Model Context Protocol: \href{https://modelcontextprotocol.io/docs/getting-started/intro}{[modelcontextprotocol.io]}
\item APIs used: Gemini (default model, GOOGLE\_\allowbreak{}API\_\allowbreak{}KEY), Tavily, Ticketmaster Discovery, Open-Meteo, OpenRouter, NVIDIA embeddings
\end{itemize}}
\section{Buffalo}
{\small\begin{itemize}
\item Chip Huyen, "Agents": \href{https://huyenchip.com/2025/01/07/agents.html}{[huyenchip.com]}
\item Anthropic, "Building Effective Agents": \href{https://www.anthropic.com/engineering/building-effective-agents}{[anthropic.com]}
\item OpenAI Agents SDK: \href{https://openai.github.io/openai-agents-python/}{[openai.github.io]}
\item LangGraph: \href{https://docs.langchain.com/oss/python/langgraph/overview}{[docs.langchain.com]}
\item LlamaIndex: \href{https://docs.llamaindex.ai/}{[docs.llamaindex.ai]}
\item Also mentioned without URLs: OpenAI function/tool calling docs, Ollama, OpenAI-compatible APIs, vector databases/local vector stores, Jupyter, Git/GitHub
\item Books (no URLs): Russell \& Norvig, *AI: A Modern Approach*; Caldwell, *The Agentic AI Bible*; Huyen, *AI Engineering* (O'Reilly 2025); Sutton \& Barto, *Reinforcement Learning* 2nd ed.
\item Research references (no URLs given): ReAct, Toolformer, RAG (Lewis et al.), Generative Agents, SWE-bench, WebArena
\end{itemize}}
\section{UVA}
{\small\begin{itemize}
\item Hugging Face smolagents: \href{https://huggingface.co/docs/smolagents/index}{[huggingface.co]}
\item LangChain / LangGraph: \href{https://www.langchain.com}{[langchain.com]} ; LangGraph Graph API: \href{https://docs.langchain.com/oss/python/langgraph/graph-api}{[docs.langchain.com]}
\item Model Context Protocol: \href{https://modelcontextprotocol.io/}{[modelcontextprotocol.io]} ; HF MCP course: \href{https://huggingface.co/learn/mcp-course/en/unit0/introduction}{[huggingface.co]}
\item Unsloth: \href{https://unsloth.ai/}{[unsloth.ai]}
\item Tinker (Thinking Machines): \href{https://thinkingmachines.ai/tinker/}{[thinkingmachines.ai]} ; install docs \href{https://tinker-docs.thinkingmachines.ai/install}{[tinker-docs.thinkingmachines.ai]}
\item Ollama (llava model), Google Colab, OpenAI Platform: \href{https://platform.openai.com/}{[platform.openai.com]}
\item Henry Kautz repos: gpt\_\allowbreak{}with\_\allowbreak{}tools \href{https://gitlab.com/HenryKautz/gpt_with_tools}{[gitlab.com]} ; MoltReader \href{https://github.com/HenryKautz/MoltReader}{[GitHub]} ; pdf\_\allowbreak{}ocr \href{https://github.com/HenryKautz/pdf_ocr}{[GitHub]} ; clip\_\allowbreak{}test \href{https://github.com/HenryKautz/clip_test}{[GitHub]}
\item HEPTAPOD: \href{https://github.com/tonymenzo/heptapod}{[GitHub]} ; Orchestral AI: \href{https://orchestral-ai.com}{[orchestral-ai.com]}
\item OpenClaw \href{https://openclaw.ai/}{[openclaw.ai]} ; NullClaw \href{https://nullclaw.org/}{[nullclaw.org]} ; NVIDIA NemoClaw \href{https://www.nvidia.com/en-us/ai/nemoclaw/}{[nvidia.com]} ; Claude Cowork \href{https://www.anthropic.com/product/claude-cowork}{[anthropic.com]}
\item Brave Search API \href{https://brave.com/search/api/}{[brave.com]} ; requests ; beautifulsoup4 ; OpenCV ; gradio ; edge-tts
\item Background reading on the home page: Anthropic "Effective context engineering for AI agents" \href{https://www.anthropic.com/engineering/effective-context-engineering-for-ai-agents}{[anthropic.com]} ; Chip Huyen "Building LLM applications for production" \href{https://huyenchip.com/2023/04/11/llm-engineering.html}{[huyenchip.com]} ; Antonio Gulli "Agent Design Patterns" preview \href{https://docs.google.com/document/d/1rsaK53T3Lg5KoGwvf8ukOUvbELRtH-V0LnOIFDxBryE/preview}{[docs.google.com]}
\item Additional core papers on the home page: ReAct \href{https://arxiv.org/abs/2210.03629}{[arXiv:2210.03629]} ; Self-Refine \href{https://arxiv.org/abs/2303.17651}{[arXiv:2303.17651]} ; Tree of Thoughts \href{https://arxiv.org/abs/2305.10601}{[arXiv:2305.10601]} ; Generative Agents \href{https://arxiv.org/abs/2304.03442}{[arXiv:2304.03442]} ; Magentic-One \href{https://arxiv.org/html/2411.04468v1}{[arXiv:2411.04468]} ; ToolOrchestra \href{https://arxiv.org/abs/2511.21689}{[arXiv:2511.21689]} ; narrow fine-tuning misalignment (Nature 2026) \href{https://doi.org/10.1038/s41586-025-09937-5}{[doi.org]}
\end{itemize}}
\section{UW-Madison}
{\small\begin{itemize}
\item Course demo notebook: \href{https://github.com/barghouthi/ai-agents/blob/main/demos/agents.ipynb}{[GitHub]}
\item DSPy (codebase mentioned alongside the DSPy paper): \href{https://arxiv.org/abs/2310.03714}{[arXiv:2310.03714]}
\item litellm (allowed thin wrapper in the assignment); LangChain, Strands, CrewAI (named as disallowed)
\item Example coding-agent CLIs: Claude Code, Cursor CLI, Gemini CLI
\item Benchmarks/environments in the readings: SWE-bench, Terminal-Bench, AgentDojo, WebArena, BaxBench, RedCode, LiveResearchBench
\item Presentation advice: \href{https://www.youtube.com/watch?v=aFT79TmffPk}{[video]} ; \href{https://pages.graphics.cs.wisc.edu/777-S11/2011/04/08/advice-on-presentations/}{[pages.graphics.cs.wisc.edu]}
\end{itemize}}
\section{JHU}
{\small\begin{itemize}
\item Course site source: \href{https://github.com/machine-programming/machine-programming.github.io}{[GitHub]}
\item Assignment repos: \href{https://github.com/machine-programming/assignment-1}{[GitHub]} , \href{https://github.com/machine-programming/assignment-2}{[GitHub]} , \href{https://github.com/machine-programming/assignment-3}{[GitHub]} (all Python, no license listed)
\item Amazon Bedrock (Converse API) with Claude Haiku / Claude Sonnet 4.6 (A1, A2); Google Gemini API (\texttt{gemini-2.\allowbreak{}5-pro} in A3's example config)
\item A3 stack: Node.js/npm, Express, WebSockets (optional), Jest, Supertest, Playwright
\item A2 toolchains: Python \texttt{subprocess}/\texttt{py\_\allowbreak{}compile}, Rust \texttt{cargo}, OCaml \texttt{ocaml}/\texttt{ocamlc} (dune optional)
\item Benchmarks/harnesses on the readings page: LiveCodeBench \href{https://livecodebench.github.io/}{[livecodebench.github.io]} , SWE-bench \href{https://arxiv.org/abs/2310.06770}{[arXiv:2310.06770]} , BigCode Eval Harness \href{https://github.com/bigcode-project/bigcode-evaluation-harness}{[GitHub]} , BigCodeBench \href{https://huggingface.co/blog/leaderboard-bigcodebench}{[huggingface.co]} , LMArena \href{https://lmarena.ai/?leaderboard}{[lmarena.ai]}
\item Agent papers on the readings page: MASAI \href{https://arxiv.org/pdf/2406.11638}{[arXiv:2406.11638]} , AutoCodeRover \href{https://arxiv.org/abs/2404.05427}{[arXiv:2404.05427]}
\item Project-idea tools: Scallop \href{https://scallop.github.io/}{[scallop.github.io]} , plus Lean, Coq, Z3, CodeQL, Datalog, Prolog, PDDL (Wikipedia links on the project page)
\item Report templates: \href{https://www.overleaf.com/latex/templates/neurips-2024/tpsbbrdqcmsh}{[overleaf.com]} , \href{https://typst.app/universe/package/bloated-neurips/0.8.0/}{[typst.app]}
\end{itemize}}
\section{CU Boulder}
{\small\begin{itemize}
\item LangGraph (in-class tutorial on 03/09). The site gives no URL.
\item OpenHands (paper): \href{https://arxiv.org/pdf/2407.16741}{[arXiv:2407.16741]}
\item SWE-agent (paper): \href{https://arxiv.org/pdf/2405.15793}{[arXiv:2405.15793]} ; AutoCodeRover (paper): \href{https://arxiv.org/pdf/2404.05427}{[arXiv:2404.05427]}
\item AutoGen (paper): \href{https://arxiv.org/pdf/2308.08155}{[arXiv:2308.08155]} ; MetaGPT (paper): \href{https://arxiv.org/pdf/2308.00352}{[arXiv:2308.00352]}
\item ACM proceedings template (milestones/final): \href{https://www.acm.org/publications/proceedings-template}{[acm.org]}
\item Piazza: \href{https://piazza.com/colorado/spring2026/csci7000005}{[piazza.com]} ; Canvas (submissions, login required)
\end{itemize}}
\section{Mannheim}
{\small\begin{itemize}
\item LangChain: \href{https://www.langchain.com/langchain}{[langchain.com]}
\item Tavily (search API): \href{https://www.tavily.com/}{[tavily.com]}
\item MCP (named in the exercise filename)
\item AgentBench: \href{https://github.com/THUDM/AgentBench}{[GitHub]}
\item SmartPlay: \href{https://github.com/microsoft/SmartPlay}{[GitHub]}
\item WebArena: \href{https://webarena.dev/}{[webarena.dev]} ; WebShop: \href{https://webshop-pnlp.github.io/}{[webshop-pnlp.github.io]}
\item BrowseComp: \href{https://openai.com/index/browsecomp/}{[openai.com]} ; Meta CRAG (KDD Cup 2024): \href{https://www.aicrowd.com/challenges/meta-comprehensive-rag-benchmark-kdd-cup-2024}{[aicrowd.com]}
\item Course evaluation PDFs: \href{https://www.uni-mannheim.de/media/Einrichtungen/dws/Files_Teaching/Large_Language_Models_and_Agents/FSS2026/IE_685_Large_Language_Models_and_Agents_Lecture_Evaluation_FSS2026.pdf}{[PDF]} ; \href{https://www.uni-mannheim.de/media/Einrichtungen/dws/Files_Teaching/Large_Language_Models_and_Agents/FSS2026/IE_686_Large_Language_Models_and_Agents_Project_Evaluation_FSS2026.pdf}{[PDF]}
\end{itemize}}
\section{UChicago}
{\small\begin{itemize}
\item Course repo: \href{https://github.com/KnowledgeLab/AI-Agents-for-Social-Science-and-Society-2026}{[GitHub]} (notebooks plus \texttt{Tutorials-Homework\_\allowbreak{}Notebooks/\allowbreak{}data/\allowbreak{}})
\item Python tutorial: \href{https://docs.python.org/3/tutorial/}{[docs.python.org]}
\item Named in the syllabus or notebooks without URLs: Claude Code, LangGraph, LangChain, AutoGen, Concordia (DeepMind), Faiss, ChromaDB, DSPy, AIOS, AFlow, Kubeflow, MLflow, OpenPipe ART, Agent Lightning, Neuronpedia, AGrail, PyTorch, Keras/TensorFlow, scikit-learn, gensim, spaCy, Hugging Face transformers, shap, gymnasium, stable-baselines3, Petals (Llama-2), OpenAI API, SWE-agent
\item Google Colab (Pro recommended); UChicago RCC; Canvas
\end{itemize}}
\section{Stanford 329A}
{\small\begin{itemize}
\item Homework repos: \href{https://github.com/stanford-cs329a/cs329a-homework1-fall2025-public}{[GitHub]} (MIT), \href{https://github.com/stanford-cs329a/cs329a-homework2-fall2025-public}{[GitHub]} (no license shown), \href{https://github.com/stanford-cs329a/cs329a-homework3-fall2025-public}{[GitHub]} (Apache-2.0). All Python notebooks plus a course package (\texttt{cs329\_\allowbreak{}hw1}, \texttt{cs329\_\allowbreak{}hw}, \texttt{cs329a\_\allowbreak{}hw3})
\item Together AI API (all homeworks): Qwen3-Next-80B-A3B-Instruct, Llama-3.2-3B-Instruct-Turbo, Qwen2.5-72B-Instruct-Turbo, Meta-Llama-3.1-8B-Instruct-Turbo, Llama-3.3-70B-Instruct-Turbo-Free, gemma-3n-E4B-it, gpt-oss-20B
\item Benchmarks/datasets: AIME25 (Hugging Face \texttt{math-ai/\allowbreak{}aime25}), HumanEval (\href{https://arxiv.org/abs/2107.03374}{[arXiv:2107.03374]}); a QA dataset bundled with HW3 (not described)
\item HW3 external APIs: Google Custom Search, Polygon, Wolfram Alpha, Geopy Nominatim, Open-Meteo; Pydantic for structured outputs
\item Docker (HW2 \texttt{humaneval-sandbox} image), conda, python-dotenv, tenacity, ThreadPoolExecutor
\item Key papers with code/systems from readings: Archon \href{https://www.arxiv.org/abs/2409.15254}{[arXiv:2409.15254]}, MemGPT \href{https://arxiv.org/abs/2310.08560}{[arXiv:2310.08560]}, CacheBlend \href{http://arxiv.org/abs/2405.16444}{[arXiv:2405.16444]}, KernelBench \href{https://arxiv.org/pdf/2502.10517}{[arXiv:2502.10517]}, CodeMonkeys \href{https://arxiv.org/abs/2501.14723}{[arXiv:2501.14723]}, DAPO \href{https://arxiv.org/abs/2503.14476}{[arXiv:2503.14476]}, AlphaEvolve \href{https://storage.googleapis.com/deepmind-media/DeepMind.com/Blog/alphaevolve-a-gemini-powered-coding-agent-for-designing-advanced-algorithms/AlphaEvolve.pdf}{[PDF]}
\end{itemize}}
\section{Berkeley '25}
{\small\begin{itemize}
\item AgentBeats tutorial: \href{https://github.com/RDI-Foundation/agentbeats-tutorial}{[GitHub]} (examples plus a local assessment runner, \texttt{uv run agentbeats-run scenarios/\allowbreak{}debate/\allowbreak{}scenario.\allowbreak{}toml}; scenarios: debate with an LLM-as-judge green agent, and tau2; Python/uv, Docker images for linux/amd64, published to GitHub Container Registry; no license listed). The README points to separate agent-template and green-agent-template repos.
\item The current competition page links starter code at \href{https://github.com/agentbeats/tutorial}{[GitHub]} and docs at \href{https://docs.agentbeats.dev/tutorial/}{[docs.agentbeats.dev]} (not opened)
\item AgentBeats platform: \href{https://agentbeats.dev}{[agentbeats.dev]} (not opened)
\item A2A (Agent-to-Agent) protocol: \href{https://a2a-protocol.org}{[a2a-protocol.org]} ; A2A SDK or Google ADK for building agents
\item LLM providers named in the tutorial: Google Gemini, OpenAI, Anthropic, OpenRouter, Ollama
\item Benchmarks in the readings: SWE-bench Verified, BrowseComp, $\tau$2-Bench
\item AgentX-AgentBeats competition: \href{https://rdi.berkeley.edu/agentx-agentbeats}{[rdi.berkeley.edu]}
\end{itemize}}
\section{Berkeley Adv.}
{\small\begin{itemize}
\item Lab materials folder (4 files, modified 24 Jun 2025): \href{https://drive.google.com/drive/folders/1DsXvOqfvFs4yBc2ZrLuOCXAFXBb5q1kg}{[drive.google.com]} --- Lab1\_\allowbreak{}Instructions.md (12 KB), lab1-datapoint.zip (2 KB), Lab2\_\allowbreak{}Instructions.md (10 KB), lab2-starter-code.zip (66 KB)
\item Lab FAQ (Google Doc, did not render): \href{https://docs.google.com/document/d/1AmVd4tU8wvYlINDmUuEuw6TIvxl95yjgui9JZAZ8P80/edit}{[docs.google.com]}
\item Lab 2 stack, per the instructions: Lean 4 with Mathlib and Aesop, built with Lake (\texttt{lake lean Lean4CodeGenerator.\allowbreak{}lean}); Python harness with \texttt{src/\allowbreak{}agents.\allowbreak{}py} (GPT-4o for generation, o3-mini for light reasoning), \texttt{src/\allowbreak{}embedding\_\allowbreak{}models.\allowbreak{}py} (OpenAI and lightweight embeddings), \texttt{src/\allowbreak{}embedding\_\allowbreak{}db.\allowbreak{}py}; \texttt{make test} / \texttt{make zip}; approved models only (o1-pro named as not approved); up to 60 s of extended thinking per question
\item Lean learning resources named in Lab 1: Functional Programming in Lean, Theorem Proving in Lean 4, the mathlib4 "All Tactics in Lean 4" list, Natural Number Game (NNG4)
\item Problem sources suggested in Lab 1: LiveCodeBench, APPS, MBPP, LeetCode, HackerRank, Codeforces
\item Systems and benchmarks in the readings: LeanDojo, WebArena, VisualWebArena, Mind2Web, OSWorld, AGUVIS, Tulu 3/OLMo, HippoRAG, Big Sleep (Project Zero)
\item Community: Edstem \href{https://edstem.org/us/join/QMmJkA}{[edstem.org]} (Berkeley), Discord \href{https://discord.gg/NWVpQ9rBvd}{[discord.gg]} (MOOC); Agentic AI Summit \href{https://rdi.berkeley.edu/events/agentic-ai-summit}{[rdi.berkeley.edu]}
\end{itemize}}
\section{Stanford 329T}
{\small\begin{itemize}
\item TruLens (open-source LLM app tracing and evaluation): \href{https://github.com/truera/trulens}{[GitHub]} . HW1 uses \texttt{run\_\allowbreak{}dashboard()}, human feedback on traces, LLM-judge metrics and \texttt{session.\allowbreak{}get\_\allowbreak{}leaderboard()}
\item LangGraph: agentic concepts page \href{https://langchain-ai.github.io/langgraph/concepts/agentic_concepts/}{[langchain-ai.github.io]} ; the main-branch notebook pins \texttt{langgraph==0.\allowbreak{}5.\allowbreak{}4}
\item Tavily search API and the OpenAI API (HW1 needs both keys; the handout estimates under \$2.00 in OpenAI cost)
\item Substitute HW1 notebook (main branch): \href{https://github.com/truera/trulens/blob/main/examples/experimental/web-search-agent-human-annotation-flow.ipynb}{[GitHub]} . It uses \texttt{trulens}, \texttt{trulens-apps-langgraph}, \texttt{trulens-providers-openai}, \texttt{langchain\_\allowbreak{}tavily}, \texttt{langchain\_\allowbreak{}openai} and \texttt{langchain\_\allowbreak{}experimental}, with nodes planner, executor, web\_\allowbreak{}researcher, chart\_\allowbreak{}generator, chart\_\allowbreak{}summarizer and synthesizer. Judge metrics: Plan Quality, Plan Adherence, Execution Efficiency, Logical Consistency, plus the RAG triad (Groundedness, Answer Relevance, Context Relevance)
\item DeepLearning.AI short course "Building and Evaluating Data Agents": \href{https://www.deeplearning.ai/short-courses/building-and-evaluating-data-agents/}{[deeplearning.ai]}
\item Model Context Protocol: \href{https://modelcontextprotocol.io/docs/getting-started/intro}{[modelcontextprotocol.io]}
\item DSPy \href{https://dspy.ai/}{[dspy.ai]} and GEPA \href{https://github.com/gepa-ai/gepa}{[GitHub]}
\item Together cookbook, Open Deep Research: \href{https://github.com/togethercomputer/together-cookbook/blob/main/Agents/Together_Open_Deep_Research_CookBook.ipynb}{[GitHub]}
\item Agent benchmarks/environments: SWE-bench, SWE-smith, ARE/GAIA2, OSWorld, Terminal-Bench \href{https://www.tbench.ai/}{[tbench.ai]} , Gray Swan Arena \href{https://app.grayswan.ai/arena}{[app.grayswan.ai]}
\item AI coding assistants (Cursor, Claude Code) are explicitly allowed and encouraged for HW1
\end{itemize}}
\section{Supplements}
{\small\begin{itemize}
\item \hyperref[sup-mcp]{Model Context Protocol docs (version 2026-07-28)}: \href{https://modelcontextprotocol.io/docs/2026-07-28/getting-started/intro}{https://modelcontextprotocol.io/docs/2026-07-28/getting-started/intro}
\item \hyperref[sup-codeexec]{Anthropic Engineering: Code execution with MCP}: \href{https://www.anthropic.com/engineering/code-execution-with-mcp}{https://www.anthropic.com/engineering/code-execution-with-mcp}
\item \hyperref[sup-openclaw]{OpenClaw documentation}: \href{https://docs.openclaw.ai/}{https://docs.openclaw.ai/}
\item \hyperref[sup-tinyagent]{Berkeley BAIR: TinyAgent, Function Calling at the Edge}: \href{https://bair.berkeley.edu/blog/2024/05/29/tiny-agent/}{https://bair.berkeley.edu/blog/2024/05/29/tiny-agent/}
\item \hyperref[sup-agentsec]{MIT 6.566 guest lecture: AI Agent Security (Anish Athalye, April 2026)}: \href{https://github.com/anishathalye/ai-agent-security-lecture}{https://github.com/anishathalye/ai-agent-security-lecture}
\item \hyperref[sup-anthauto]{Anthropic: Building effective agent automations (Martin and Avilla, October 2026)}: \href{https://claude.dev/blog/building-effective-agent-automations/}{https://claude.dev/blog/building-effective-agent-automations/}
\item \hyperref[sup-notion]{Notion: How we built a personal data scientist for every Notion employee (David Thomas, October 2026)}: \href{https://www.notion.com/blog/how-we-built-a-personal-data-scientist-for-every-notion-employee}{https://www.notion.com/blog/how-we-built-a-personal-data-scientist-for-every-notion-employee}
\item \hyperref[sup-colab]{Google Colab}: \href{https://colab.research.google.com/}{https://colab.research.google.com/}
\item TinyAgent code: \href{https://github.com/SqueezeAILab/TinyAgent}{github.com/SqueezeAILab/TinyAgent}
\end{itemize}}
\chapter{Canonical Reading Checklist}
The key readings from every module in one checklist. Lecture-level readings are in the catalogue tables.
\begin{multicols}{2}\small
\par\medskip\noindent\textbf{\hyperlink{mod-FND}{M1 LLM Foundations and Prompting}}\par
\begin{itemize}[label=$\square$]
\item Vaswani et al., \href{https://arxiv.org/abs/1706.03762}{\emph{Attention Is All You Need}} (2017)
\item Ouyang et al., \href{https://arxiv.org/abs/2203.02155}{\emph{InstructGPT}} (2022)
\item Wei et al., \href{https://arxiv.org/abs/2201.11903}{\emph{Chain-of-Thought Prompting}} (2022)
\item Belcak et al., \href{https://arxiv.org/abs/2506.02153}{\emph{Small Language Models are the Future of Agentic AI}} (2025)
\item Karpathy, \href{https://www.youtube.com/watch?v=7xTGNNLPyMI}{\emph{Deep Dive into LLMs}} (video)
\end{itemize}
\par\medskip\noindent\textbf{\hyperlink{mod-ARCH}{M2 The Agent Loop and Agent Architectures}}\par
\begin{itemize}[label=$\square$]
\item Anthropic, \href{https://www.anthropic.com/engineering/building-effective-agents}{\emph{Building Effective Agents}}
\item Yao et al., \href{https://arxiv.org/abs/2210.03629}{\emph{ReAct}} (2022)
\item Sumers et al., \href{https://arxiv.org/abs/2309.02427}{\emph{CoALA: Cognitive Architectures for Language Agents}} (2023)
\item Chip Huyen, \href{https://huyenchip.com/2025/01/07/agents.html}{\emph{Agents}} (2025)
\item Wang et al., \href{https://arxiv.org/abs/2308.11432}{\emph{A Survey on LLM-based Autonomous Agents}} (2024)
\end{itemize}
\par\medskip\noindent\textbf{\hyperlink{mod-TOOL}{M3 Tool Use and Function Calling}}\par
\begin{itemize}[label=$\square$]
\item Schick et al., \href{https://arxiv.org/abs/2302.04761}{\emph{Toolformer}} (2023)
\item Patil et al., \href{https://arxiv.org/abs/2305.15334}{\emph{Gorilla}} (2023)
\item Wang et al., \href{https://arxiv.org/abs/2402.01030}{\emph{CodeAct: Executable Code Actions}} (2024)
\item Gao et al., \href{https://arxiv.org/abs/2211.10435}{\emph{PAL}} (2022)
\item \href{https://arxiv.org/abs/2403.15452}{\emph{What Are Tools Anyway?}} (2024)
\item Anthropic, \href{https://www.anthropic.com/engineering/writing-tools-for-agents}{\emph{Writing effective tools for agents}}
\end{itemize}
\par\medskip\noindent\textbf{\hyperlink{mod-MCP}{M4 MCP and Agent Interoperability}}\par
\begin{itemize}[label=$\square$]
\item MCP specification and \href{https://modelcontextprotocol.io/docs/2026-07-28/tutorials/security/security_best_practices}{security best practices}
\item Anthropic, \href{https://www.anthropic.com/engineering/code-execution-with-mcp}{\emph{Code execution with MCP}} (2025)
\item Felendler et al., \href{https://arxiv.org/abs/2602.15945}{\emph{From Tool Orchestration to Code Execution: MCP Design Choices}} (2026)
\item \href{https://arxiv.org/abs/2503.23278}{\emph{MCP Landscape \& Security}} (arXiv 2503.23278)
\end{itemize}
\par\medskip\noindent\textbf{\hyperlink{mod-MEM}{M5 Memory, Retrieval and Context Engineering}}\par
\begin{itemize}[label=$\square$]
\item Lewis et al., \href{https://arxiv.org/abs/2005.11401}{\emph{Retrieval-Augmented Generation}} (2020)
\item Liu et al., \href{https://arxiv.org/abs/2307.03172}{\emph{Lost in the Middle}} (2023)
\item Packer et al., \href{https://arxiv.org/abs/2310.08560}{\emph{MemGPT}} (2023)
\item Park et al., \href{https://arxiv.org/abs/2304.03442}{\emph{Generative Agents}} (2023)
\item Anthropic, \href{https://www.anthropic.com/engineering/effective-context-engineering-for-ai-agents}{\emph{Effective context engineering for AI agents}}
\item Asai et al., \href{https://arxiv.org/abs/2310.11511}{\emph{Self-RAG}} (2023); Yan et al., \href{https://arxiv.org/abs/2401.15884}{\emph{Corrective RAG}} (2024)
\end{itemize}
\par\medskip\noindent\textbf{\hyperlink{mod-RSN}{M6 Reasoning, Planning, Reflection and Search}}\par
\begin{itemize}[label=$\square$]
\item Wang et al., \href{https://arxiv.org/abs/2203.11171}{\emph{Self-Consistency}} (2022)
\item Yao et al., \href{https://arxiv.org/abs/2305.10601}{\emph{Tree of Thoughts}} (2023)
\item Shinn et al., \href{https://arxiv.org/abs/2303.11366}{\emph{Reflexion}} (2023); Madaan et al., \href{https://arxiv.org/abs/2303.17651}{\emph{Self-Refine}} (2023)
\item Zhou et al., \href{https://arxiv.org/abs/2205.10625}{\emph{Least-to-Most Prompting}} (2022)
\item Zhou et al., \href{https://arxiv.org/abs/2310.04406}{\emph{LATS}} (2023)
\item Kim et al., \href{https://arxiv.org/abs/2312.04511}{\emph{LLMCompiler}} (2023)
\end{itemize}
\par\medskip\noindent\textbf{\hyperlink{mod-MULTI}{M7 Multi-Agent Systems}}\par
\begin{itemize}[label=$\square$]
\item Wu et al., \href{https://arxiv.org/abs/2308.08155}{\emph{AutoGen}} (2023)
\item Hong et al., \href{https://arxiv.org/abs/2308.00352}{\emph{MetaGPT}} (2023)
\item Li et al., \href{https://arxiv.org/abs/2303.17760}{\emph{CAMEL}} (2023)
\item Du et al., \href{https://arxiv.org/abs/2305.14325}{\emph{Multiagent Debate}} (2023)
\item Cemri et al., \href{https://arxiv.org/abs/2503.13657}{\emph{Why Do Multi-Agent LLM Systems Fail?}} (2025)
\item Anthropic, \href{https://www.anthropic.com/engineering/multi-agent-research-system}{\emph{How we built our multi-agent research system}}
\end{itemize}
\par\medskip\noindent\textbf{\hyperlink{mod-CODE}{M8 Coding Agents}}\par
\begin{itemize}[label=$\square$]
\item Jimenez et al., \href{https://arxiv.org/abs/2310.06770}{\emph{SWE-bench}} (2023)
\item Yang et al., \href{https://arxiv.org/abs/2405.15793}{\emph{SWE-agent}} (2024)
\item Wang et al., \href{https://arxiv.org/abs/2407.16741}{\emph{OpenHands}} (2024)
\item Xia et al., \href{https://arxiv.org/abs/2407.01489}{\emph{Agentless}} (2024)
\item Zhang et al., \href{https://arxiv.org/abs/2404.05427}{\emph{AutoCodeRover}} (2024)
\item Hassan et al., \href{https://arxiv.org/abs/2509.06216}{\emph{Agentic Software Engineering}} roadmap (2025)
\item Anthropic, \href{https://www.anthropic.com/engineering/claude-code-best-practices}{\emph{Claude Code best practices}}
\end{itemize}
\par\medskip\noindent\textbf{\hyperlink{mod-WEB}{M9 Web, GUI, Computer-Use and Multimodal Agents}}\par
\begin{itemize}[label=$\square$]
\item Zhou et al., \href{https://arxiv.org/abs/2307.13854}{\emph{WebArena}} (2023); Koh et al., \href{https://arxiv.org/abs/2401.13649}{\emph{VisualWebArena}} (2024)
\item Xie et al., \href{https://arxiv.org/abs/2404.07972}{\emph{OSWorld}} (2024)
\item Deng et al., \href{https://arxiv.org/abs/2306.06070}{\emph{Mind2Web}} (2023)
\item Yao et al., \href{https://arxiv.org/abs/2207.01206}{\emph{WebShop}} (2022)
\item Zhang et al., \href{https://arxiv.org/abs/2312.13771}{\emph{AppAgent}} (2023)
\item Wang et al., \href{https://arxiv.org/abs/2305.16291}{\emph{Voyager}} (2023)
\end{itemize}
\par\medskip\noindent\textbf{\hyperlink{mod-EVAL}{M10 Evaluating Agents}}\par
\begin{itemize}[label=$\square$]
\item Zheng et al., \href{https://arxiv.org/abs/2306.05685}{\emph{Judging LLM-as-a-Judge}} (2023)
\item Yao et al., \href{https://arxiv.org/abs/2406.12045}{\emph{$\tau$-bench}} (2024)
\item Liu et al., \href{https://arxiv.org/abs/2308.03688}{\emph{AgentBench}} (2023)
\item Kapoor et al., \href{https://arxiv.org/abs/2407.01502}{\emph{AI Agents That Matter}} (2024)
\item Mohammadi et al., \href{https://arxiv.org/abs/2507.21504}{\emph{Evaluation and Benchmarking of LLM Agents}} (KDD 2025)
\item METR, \href{https://arxiv.org/abs/2503.14499}{\emph{Measuring AI Ability to Complete Long Tasks}} (2025)
\end{itemize}
\par\medskip\noindent\textbf{\hyperlink{mod-TRAIN}{M11 Training and Optimizing Agents}}\par
\begin{itemize}[label=$\square$]
\item Ouyang et al., \href{https://arxiv.org/abs/2203.02155}{\emph{InstructGPT}} (2022)
\item Shao et al., \href{https://arxiv.org/abs/2402.03300}{\emph{DeepSeekMath (GRPO)}} (2024)
\item DeepSeek-AI, \href{https://arxiv.org/abs/2501.12948}{\emph{DeepSeek-R1}} (2025)
\item Singh et al., \href{https://arxiv.org/abs/2312.06585}{\emph{ReST-EM}} (2023)
\item Hu et al., \href{https://arxiv.org/abs/2106.09685}{\emph{LoRA}} (2021); Dettmers et al., \href{https://arxiv.org/abs/2305.14314}{\emph{QLoRA}} (2023)
\item Khattab et al., \href{https://arxiv.org/abs/2310.03714}{\emph{DSPy}} (2023); Agrawal et al., \href{https://arxiv.org/abs/2507.19457}{\emph{GEPA}} (2025)
\end{itemize}
\par\medskip\noindent\textbf{\hyperlink{mod-SAFE}{M12 Safety, Security and Oversight}}\par
\begin{itemize}[label=$\square$]
\item Greshake et al., \href{https://arxiv.org/abs/2302.12173}{\emph{Indirect Prompt Injection}} (2023)
\item Debenedetti et al., \href{https://arxiv.org/abs/2406.13352}{\emph{AgentDojo}} (2024)
\item Debenedetti et al., \href{https://arxiv.org/abs/2503.18813}{\emph{CaMeL: Defeating Prompt Injections by Design}} (2025)
\item Willison, \href{https://simonwillison.net/2025/Jun/16/the-lethal-trifecta/}{\emph{The lethal trifecta}}
\item Wang et al., \href{https://arxiv.org/abs/2504.15585}{\emph{Full Stack Safety survey}} (2025)
\item Hubinger et al., \href{https://arxiv.org/abs/2401.05566}{\emph{Sleeper Agents}} (2024)
\end{itemize}
\par\medskip\noindent\textbf{\hyperlink{mod-DEPLOY}{M13 Deployment, Efficiency and Operations}}\par
\begin{itemize}[label=$\square$]
\item Google, \href{https://sre.google/sre-book/introduction/}{\emph{SRE book, introduction}}
\item Anthropic, \href{https://www.anthropic.com/engineering/code-execution-with-mcp}{\emph{Code execution with MCP}} (2025)
\item Erdogan et al., \href{https://arxiv.org/abs/2409.00608}{\emph{TinyAgent}} (2024)
\item Gao \& Peng, \href{https://arxiv.org/abs/2510.16786}{\emph{More with Less}} (2025)
\item Anthropic, \href{https://claude.dev/blog/building-effective-agent-automations/}{\emph{Building effective agent automations}} (2026)
\item Notion, \href{https://www.notion.com/blog/how-we-built-a-personal-data-scientist-for-every-notion-employee}{\emph{How we built a personal data scientist for every Notion employee}} (2026)
\end{itemize}
\par\medskip\noindent\textbf{\hyperlink{mod-APP}{M14 Applications: Deep Research, Science and Society}}\par
\begin{itemize}[label=$\square$]
\item Shao et al., \href{https://arxiv.org/abs/2402.14207}{\emph{STORM}} (2024)
\item L\'ala et al., \href{https://arxiv.org/abs/2312.07559}{\emph{PaperQA}} (2023)
\item Asai et al., \href{https://arxiv.org/abs/2411.14199}{\emph{OpenScholar}} (2024)
\item Lu et al., \href{https://arxiv.org/abs/2408.06292}{\emph{The AI Scientist}} (2024)
\item Gottweis et al., \href{https://arxiv.org/abs/2502.18864}{\emph{Towards an AI Co-Scientist}} (2025)
\item Park et al., \href{https://arxiv.org/abs/2304.03442}{\emph{Generative Agents}} (2023)
\end{itemize}
\end{multicols}
\chapter{Tag Matrix}
Raw tag counts per course (a lecture can carry several tags), including logistics (\tg{ORG}).
{\small\setlength{\tabcolsep}{3.5pt}\begin{longtable}{lcccccccccccccc}
\toprule Course & \hyperlink{mod-FND}{\tg{FND}} & \hyperlink{mod-ARCH}{\tg{ARCH}} & \hyperlink{mod-TOOL}{\tg{TOOL}} & \hyperlink{mod-MEM}{\tg{MEM}} & \hyperlink{mod-RSN}{\tg{RSN}} & \hyperlink{mod-MULTI}{\tg{MULTI}} & \hyperlink{mod-CODE}{\tg{CODE}} & \hyperlink{mod-WEB}{\tg{WEB}} & \hyperlink{mod-EVAL}{\tg{EVAL}} & \hyperlink{mod-TRAIN}{\tg{TRAIN}} & \hyperlink{mod-SAFE}{\tg{SAFE}} & \hyperlink{mod-DEPLOY}{\tg{DEPLOY}} & \hyperlink{mod-APP}{\tg{APP}} & \tg{ORG}\\ \midrule
\hyperlink{course-CMU}{CMU} & \textcolor{black!25}{0} & 3 & 1 & 3 & 3 & 2 & 2 & 1 & 2 & 4 & 2 & 2 & 3 & 1\\
\hyperlink{course-STAN26}{Stanford '26} & 2 & 3 & 2 & 2 & \textcolor{black!25}{0} & 1 & 8 & \textcolor{black!25}{0} & 2 & \textcolor{black!25}{0} & 2 & 4 & 3 & \textcolor{black!25}{0}\\
\hyperlink{course-STAN25}{Stanford '25} & 2 & 1 & 2 & 1 & 1 & 2 & 13 & \textcolor{black!25}{0} & 2 & \textcolor{black!25}{0} & 2 & 3 & 3 & 3\\
\hyperlink{course-NEU}{Northeastern} & 3 & 6 & 4 & 6 & 4 & 3 & \textcolor{black!25}{0} & \textcolor{black!25}{0} & 1 & 1 & 2 & 5 & \textcolor{black!25}{0} & 2\\
\hyperlink{course-UB}{Buffalo} & 1 & 2 & 1 & 3 & 2 & 1 & 1 & \textcolor{black!25}{0} & 1 & \textcolor{black!25}{0} & 2 & 2 & 1 & 3\\
\hyperlink{course-UVA}{UVA} & 2 & 4 & 4 & 2 & 4 & 4 & 1 & 2 & 3 & 2 & 1 & 1 & 3 & 3\\
\hyperlink{course-UW}{UW-Madison} & 1 & 6 & 2 & 3 & 5 & 5 & 4 & 1 & 5 & 2 & 3 & \textcolor{black!25}{0} & 4 & 4\\
\hyperlink{course-JHU}{JHU} & 2 & 4 & 2 & \textcolor{black!25}{0} & 5 & \textcolor{black!25}{0} & 8 & \textcolor{black!25}{0} & \textcolor{black!25}{0} & 4 & 1 & \textcolor{black!25}{0} & 7 & 8\\
\hyperlink{course-CU}{CU Boulder} & \textcolor{black!25}{0} & 6 & 6 & 1 & 4 & 5 & 10 & 1 & 3 & 1 & 2 & 2 & 1 & 10\\
\hyperlink{course-MANN}{Mannheim} & 3 & 2 & 1 & 2 & 1 & \textcolor{black!25}{0} & \textcolor{black!25}{0} & \textcolor{black!25}{0} & 1 & 1 & 1 & \textcolor{black!25}{0} & \textcolor{black!25}{0} & 5\\
\hyperlink{course-UCHI}{UChicago} & 3 & \textcolor{black!25}{0} & 1 & 1 & 3 & 2 & 1 & 1 & 2 & 2 & 2 & \textcolor{black!25}{0} & 5 & 2\\
\hyperlink{course-S329A}{Stanford 329A} & 1 & 3 & 1 & 1 & 7 & \textcolor{black!25}{0} & 4 & 1 & 3 & 4 & \textcolor{black!25}{0} & 1 & 2 & 9\\
\hyperlink{course-BRKAG}{Berkeley '25} & 2 & 1 & \textcolor{black!25}{0} & \textcolor{black!25}{0} & 1 & 3 & 1 & 1 & 4 & 4 & 1 & 2 & 3 & 2\\
\hyperlink{course-BRKADV}{Berkeley Adv.} & \textcolor{black!25}{0} & 2 & 2 & 1 & 9 & \textcolor{black!25}{0} & 1 & 2 & 2 & 3 & 2 & 1 & 4 & 1\\
\hyperlink{course-S329T}{Stanford 329T} & 4 & 3 & 2 & 1 & 1 & \textcolor{black!25}{0} & \textcolor{black!25}{0} & 1 & 5 & 3 & 1 & \textcolor{black!25}{0} & 1 & 9\\
\bottomrule\end{longtable}}
\chapter{Provenance and Caveats}
\begin{itemize}
\item \textbf{Method.} The 20 links in the Chrome bookmarks folder \emph{CMU-Agentic} were each visited. Static pages were fetched directly. The JavaScript-rendered sites (CMU, Stanford) were read in a live browser. Course repositories were read through GitHub. Four more courses (Stanford CS329A, Berkeley Agentic AI Fall 2025, Berkeley Advanced LLM Agents Spring 2025, Stanford CS329T) and the MIT agent-security lecture were added on the same date from a targeted search for open 2025--2026 courses, with candidates found by a Codex consultant and each one then read directly. Each course was summarised into a Markdown table in \texttt{research/}, and this PDF is generated from those tables by \texttt{guide/build\_guide.py}.
\item \textbf{Classification.} Tags were assigned from session titles, listed topics and readings. They are a navigation aid, not ground truth, and some sessions could reasonably be tagged differently.
\item \textbf{CMU.} The lecture table comes from the rendered schedule (October 9, 2026). The milestone dates on that page (A1 Sep~14, A2 Oct~1, A3 Oct~29) take precedence over the due dates in the assignment table, which come from the handouts. Recordings are linked from the schedule page because individual video URLs could not be extracted, and the YouTube playlist ID was not verified.
\item \textbf{Stanford.} Fall 2026 is current: only weeks 1--3 have slides and assignments so far, and most Thursday sessions from week 3 on are guest talks. Fall 2025 dates come from a translated copy of the syllabus and may be off by a day.
\item \textbf{In-progress courses} (CMU, Stanford '26, Northeastern, JHU, Buffalo) will publish more material after this snapshot. Re-run the research step to refresh.
\item \textbf{Unread items.} Some large slide decks could not be text-extracted (several Mannheim PDFs, a few UVA PPTX files), so their topics are inferred from titles. The UChicago week 6 and 8 notebooks were too large to inspect.
\item \textbf{Licences.} Only some repositories state a licence (UW CS 839: MIT; TinyAgent: MIT; OpenClaw: MIT). Before redistributing course material in your own repository, check each licence. Linking is always safe.
\end{itemize}
\section{Per-course research notes}
These are the notes recorded while each course was being researched.
\subsection{CMU}
{\small\begin{itemize}
\item \textbf{Lecture schedule verified.} The lecture rows in the guide come from cmu-11768-verified.md, read from the official schedule (\href{https://www.cmu-agents.com/\#/schedule}{[cmu-agents.com]}) rendered in Chrome on 2026-10-09. It lists 23 lectures with dates, modules, lecturers where named, and readings. The lecture table in this file was reconstructed earlier from secondary sources (a blog post, beri.net and YouTube titles) and is superseded wherever the two disagree.
\item \textbf{Materials.} Slide PDFs are posted for lectures 1-13 (\href{https://www.cmu-agents.com/slides/}{[cmu-agents.com]}). Lectures 14-23 had no slides yet on 2026-10-09. The schedule marks recordings for lectures 1-9; the guide links them through the schedule page because the individual YouTube links were not extracted.
\item \textbf{Guest lectures.} Lectures 22 and 23 are guest lectures (Karthik Narasimhan, Sasha Rush) with no announced topic.
\item \textbf{A2/A3 visibility.} A2 and A3 repos are now public, last updated Sep 20 and Oct 7, 2026 respectively. Earlier summaries called them enrollment-only.
\item Neubig's X announcement (\href{https://x.com/gneubig/status/2072730570304430183}{[x.com]}) could not be fetched (HTTP 402). Its text is known only from search snippets.
\end{itemize}}
\subsection{Stanford '25}
{\small\begin{itemize}
\item \textbf{The official syllabus did not load.} The "Syllabus" tab on themodernsoftware.dev and /fall2025 is rendered client-side, so WebFetch returned only the overview, team, and logistics. /syllabus returns 404. The Fall 2025 lecture table comes from the authorized Korean translation (discuss.pytorch.kr). The slide URLs are copied from that page; I did not open them individually.
\item \textbf{Day-of-week conflict:} the Korean source labels the sessions Mon/Fri (e.g. 9/22 Mon, 9/26 Fri), while the official site lists Tue/Thu 5:30-6:20 PM. The dates may belong to a different meeting pattern or be translation artifacts, so treat them as approximate.
\item \textbf{Fall 2026 schedule.} The Fall 2026 session list was later read from the rendered Syllabus tab in Chrome (2026-10-09); see stanford-cs146s-fall2026-verified.md. Slides are posted for weeks 1-3, and assignments for weeks 1-3 are on \texttt{master}.
\item \textbf{Assignment URL drift:} the Fall 2025 syllabus links to \texttt{master/\allowbreak{}week4.\allowbreak{}.\allowbreak{}8}, but \texttt{master} now holds only the Fall 2026 weeks 1-3. Use the \texttt{fall2025} branch URLs given above, which I confirmed via the GitHub tree API. Names in parentheses are the syllabus names, which differ from the titles in the assignment.md files.
\item Fall 2025 week 1 also included a \texttt{week1/\allowbreak{}README.\allowbreak{}md}, and week 5 and week 7 had README.md files plus \texttt{docs/\allowbreak{}TASKS.\allowbreak{}md}. I did not read those READMEs in full.
\item The Unit42 reading link on the source page is wrapped in a Google redirect; I extracted the underlying URL.
\item The Fall 2026 office-hours listing includes Vijay Daita, but the team grid lists one TA as "TBD".
\end{itemize}}
\subsection{Northeastern}
{\small\begin{itemize}
\item Session-page-to-lecture mapping is my inference. The Week 3a page ("Memory \& Structured Output") sits on T 9/22, whose lecture title is "Simple Agent Architecture". Its notebooks follow the memory/structured-output docs listed under F 9/25. Tuesday notebooks may run one session ahead of or behind the lecture topic.
\item Only week3b-slides.html and week4a-slides.html were found. week5a-slides.html returned 404. No slides are linked for the other sessions.
\item "DADA" and "ICA" are not expanded anywhere I saw. Week 1 says "no ICA this week", so ICA probably means in-class activity.
\item The DADA1 due date and the course grading breakdown were not found on any page I saw.
\item All pages loaded successfully.
\end{itemize}}
\subsection{Buffalo}
{\small\begin{itemize}
\item The page loaded fully. It is the only page: there are no separate schedule, syllabus, lab or slide subpages.
\item Readings are not mapped to weeks. Reading Notes 1-4 have no specified readings.
\item No public slides, videos or starter code. The tags come from the week topics only.
\item The page shows no 410 vs 510 differences in workload.
\end{itemize}}
\subsection{UVA}
{\small\begin{itemize}
\item The home page and the calendar loaded fully, as did topic pages 1-15. Slide (.pptx/.pdf), lesson-plan and code files are relative links; I did not fetch them one by one, so the filenames are as shown on the pages.
\item The URL path is \texttt{\textasciitilde{}rmw7my}, but the page names Henry Kautz as instructor. The home page title is "CS 6501: Workshop on Building AI Agents".
\item Thursday sessions with no calendar text (15 Jan, 22 Jan, 5 Feb, 19 Mar) and Tue 17 Feb are probably lab continuations. They are listed as placeholders, not as separate lectures. Counted substantive topic sessions: 15 topics across 14 dated lecture slots, plus 3 presentation days.
\item The home page topic list mentions items with no dedicated calendar session (smolagents/HF pipelines, Self-Refine, search-augmented generation, multi-agent systems, AgentFold). These were presumably covered inside the topics above.
\item The fetch tool summarized pages, so exercise details were paraphrased, not copied verbatim.
\end{itemize}}
\subsection{UW-Madison}
{\small\begin{itemize}
\item README.md, assignment.md, course-info.md, LICENSE, and the repo contents (via the GitHub API) all loaded. The repo holds only those files plus \texttt{demos/\allowbreak{}agents.\allowbreak{}ipynb}.
\item The README gives the schedule by week with Tue/Thu splits. Session dates were derived from the week start dates (Jan 20 = Tue). Week 3 is listed as "Feb 10" and Week 9 as "Mar 24".
\item The instructor's name is not on any course page. It was inferred from the GitHub owner and LICENSE (Aws Albarghouthi). A web search found no separate UW course page.
\item The demo notebook's contents were not inspected, and its link to session 2 is an assumption.
\end{itemize}}
\subsection{JHU}
{\small\begin{itemize}
\item The course is in progress (Fall 2026). As of 2026-10-09 the schedule links slides only for lectures 1-12. Lectures 13 onward have titles only.
\item The assignments page says A3's link "will be posted when released", but \href{https://github.com/machine-programming/assignment-3}{[GitHub]} is already public with a full README.
\item Lecture PDFs 7, 8 and 10 are over 10 MB and could not be read through WebFetch, so their topics are inferred from the titles. Slide URLs for the other lectures come from the relative links and were not fetched one by one.
\item The readings page (\href{https://machine-programming.github.io/readings.html}{[machine-programming.github.io]}) is organized by topic (Code LLMs, Evaluation, Agents, Improving Code Generation, Interpretability), not by lecture. The readings listed in the table are topical matches, not official per-lecture assignments. The "Agents" section is thin (MASAI, AutoCodeRover, one TBD).
\item The readings page credits Columbia COMS 6998, MIT 6.S981, UCSD CSE 291 and UW CS294 as inspiration.
\item The presentation page (presentation.html) was not fetched.
\end{itemize}}
\subsection{CU Boulder}
{\small\begin{itemize}
\item All of the index page loaded. Subpages fetched: syllabus, researchProject, paperCritiques, milestone1, finalProject. Not fetched: paperPresentation.html, projectProposal.html, projectTeams.html, milestone2-4.html.
\item The course is finished (Spring 2026). Only 4 instructor slide decks are public (lectures 1, 2, 3, 5). Slide URLs were resolved from relative links and not fetched one by one.
\item Date inconsistency: the schedule says "Sun 04/26 ... submit final project", while the research-project and final-project pages give 04/27 12pm.
\item The location is listed as ECES 114 on the index page but DLC 1B20 for the in-person sessions.
\item Lecture titles for student-led sessions are my labels built from the paper names. The site lists only presenter names and papers.
\item The syllabus has no AI-use policy and no formal learning objectives.
\end{itemize}}
\subsection{Mannheim}
{\small\begin{itemize}
\item The course page loaded fully. Slide and exercise links on the page are relative and resolve under \texttt{https:/\allowbreak{}/\allowbreak{}www.\allowbreak{}uni-mannheim.\allowbreak{}de/\allowbreak{}media/\allowbreak{}.\allowbreak{}.\allowbreak{}.\allowbreak{}}; three PDFs (lectures 04, 06, 07) were fetched and downloaded successfully. Only the project-intro deck (07) was actually read. Key topics for lectures 2-6 and 8 are inferred from their titles and the syllabus list, not from the slide contents.
\item The WebFetch tool could not extract text from the slide PDFs (binary). Lecture 4 is a 78-page deck tagged German-language (\texttt{/\allowbreak{}Lang(de)}), but the visible slides of deck 07 are in English.
\item Syllabus topic 8, "Evaluation of LLMs and LLM agents", has no dedicated slide deck or dated session on the page. It may be folded into other lectures or the projects.
\item Exercise zips were not downloaded or opened (per instructions; no code was run).
\item The page lists one session date as "28.05.2026", which is out of order and probably means 28.04.2026. The IE686 project deck's own schedule uses Wednesday coaching dates (15.04-13.05) and presentations on 19.05 and 20.05.
\item No separate public GitHub repo or videos were found for this course.
\end{itemize}}
\subsection{UChicago}
{\small\begin{itemize}
\item The README loaded fully. It gives citations only; apart from the Python tutorial, no reading URLs are provided.
\item Notebooks for weeks 6 and 8 exceed WebFetch's 10 MB limit and were not read. Their rows rely on the filename and the syllabus topic. The other notebooks were skimmed from their first \textasciitilde{}100K characters only, so the descriptions may miss later modules.
\item The source gives conflicting final-project dates: Mar 14, 5pm for the blog post and Mar 11 for presentations (Section III), but Week 10 says presentations on Mar 13 via Zoom with the project due Mar 13, 5pm.
\item The README prints "REACT ... (Yao et al., 2025)"; the ReAct paper itself is from 2022/ICLR 2023. Two week-7 readings ("Hypothetical Minds", "RATE") are listed without authors.
\item TA name spelling varies in the README ("Jesse Zhou" vs. "Jesse Zhu").
\item This is mostly a computational-social-science course. Agent-engineering content (tool use, MCP, multi-agent frameworks, RAG) is concentrated in weeks 3 and 5. Weeks 1-2 and 4 are DL foundations and causal methods.
\end{itemize}}
\subsection{Stanford 329A}
{\small\begin{itemize}
\item I confirmed the consultant's claim of 9 recordings and no slides. The course site links no videos, slides, Canvas or Panopto; its only PDFs are reading papers and past-project files. The recordings are on Stanford Online's YouTube channel (playlist \href{https://www.youtube.com/playlist?list=PLangBM27OtEA}{[video]}, titled "CS329A Self-Improving AI Agents"). The playlist page did not render a video list through WebFetch. The nine video IDs came from a third-party index (\href{https://github.com/flowioo/stanford-cs329a-self-improving-ai-agents}{[GitHub]}), and each watch page was then opened to confirm its official "Part N" title. Third-party sources date the uploads to Aug 3, 2026. I did not check the YouTube license.
\item Recording-to-lecture mapping is by title: Parts 1-6 are lectures 1-6, Part 7 ("Self-Improvement and Deep Research Agents") is lecture 8, Part 8 is lecture 17, and Part 9 is lecture 20. Lecture 7 (Open-Ended Evolution), lecture 13 (SWE agents), lecture 14 (memory) and all guest talks have no public recording.
\item The schedule is static content on the home page, which WebFetch could read. WebFetch returns converted text, so I could not confirm the HTML markup; the page shows columns \#, Date, Description, Paper Readings*, Deadlines. The footnote says readings "may be updated closer to the class date." The site numbers the midterm sessions as one row ("10-12"), split here into lectures 10, 11 and 12.
\item The HW2 due date conflicts: the schedule says Oct 22, the grading section says Oct 23. The notebooks give no due dates; the dates in the table come from the course site. HW2 and HW3 point totals are summed from the notebooks' part values.
\item I read the homework details from the READMEs and the raw notebooks (\texttt{student\_\allowbreak{}homework1.\allowbreak{}ipynb}, \texttt{homework.\allowbreak{}ipynb}, \texttt{homework3.\allowbreak{}ipynb}). I did not open the package source folders. The HW3 notebook text names \texttt{deep\_\allowbreak{}research\_\allowbreak{}agent.\allowbreak{}py}, but its code imports \texttt{cs329a\_\allowbreak{}hw3.\allowbreak{}DeepResearchAgent}.
\item "Key topics" for guest talks and lecture 20 come from titles only. For other lectures they come from titles and listed readings, not from watching the videos.
\item Prerequisites come from the Stanford bulletin, not the course site. The Stanford Online course page (\href{https://online.stanford.edu/courses/cs329a-self-improving-ai-agents}{[online.stanford.edu]}) returned no content.
\item Past-project papers and posters are relative links under \texttt{winter2025projects/\allowbreak{}} on the course site. I did not open them.
\end{itemize}}
\subsection{Berkeley '25}
{\small\begin{itemize}
\item Opened: the course page \href{https://rdi.berkeley.edu/agentic-ai/f25}{[rdi.berkeley.edu]} (twice, the second time for exact link targets), the MOOC page \href{https://agenticai-learning.org/f25}{[agenticai-learning.org]}, the competition page \href{https://rdi.berkeley.edu/agentx-agentbeats}{[rdi.berkeley.edu]}, and the tutorial repo README. I tried the Oct 6 slide PDF, but its text was not machine-readable (Google Slides export). No other slide PDFs and no videos were opened. WebFetch kept a cached copy of that PDF in the Claude Code tool-results folder automatically; I did not save it on purpose.
\item The schedule was read through WebFetch's text conversion, not raw HTML, so I could not confirm the underlying markup. Its columns are Date, Lecture, and Supplemental Readings. Slide links are relative; the absolute URLs here are resolved by me, with spaces and \texttt{\&} percent-encoded.
\item The consultant's counts check out: 9 sessions have slides (Sep 8, 15, 22, 29; Oct 6, 13, 20, 27; Nov 10) and 12 have video (11 recordings plus the Dec 8 livestream). There is one project with two phases (green, then white) on AgentBeats. Both the tutorial repo and the competition page exist.
\item No speaker is listed for Oct 6 (Agent Evaluation \& Project Overview) on either the course page or the MOOC page. The slide metadata gives only the title and date (10.6.2025). The speaker is n/a.
\item The competition page has changed since Fall 2025. It now says "purple" agents, with Phase 1 (green) running Oct 16, 2025 to Jan 31, 2026, Phase 2 (purple) running Mar 2 to Jun 2, 2026 in four sprints, and different starter-code links. None of these current deadlines are used here. The project framing and dates in the table come from the course page's grading and timeline (in-course final deadline 12/17/2025). The MOOC's deadline was Jan 31, 2026.
\item The MOOC page says lectures started Sep 15. The Sep 8 Introduction appears only on the in-person course page.
\item Readings are the course page's "Supplemental Readings" for each session. Kimi K2 and DeepSeek-V3 are hosted on the course site as d11.pdf and d22.pdf.
\item The project phase descriptions (A2A messages, JSON artifacts, Docker) come from the tutorial README and the current competition page. The course page gives only weights and dates for the phases, so how much of that detail applied in Fall 2025 is inferred.
\item Teaching staff listings differ: a third-party search snippet of the intro slides mentions a reader, John Wang, who is not on the course page. I did not verify this.
\end{itemize}}
\subsection{Berkeley Adv.}
{\small\begin{itemize}
\item The consultant's report mostly holds up. There are 12 lectures, each with slides and a YouTube video, and the topics match. There are two Lean lab instruction files plus two ZIPs in the stated Drive folder. Neither site links an official GitHub repo, and a web search found none. Corrections: the ZIPs are a Lab 1 datapoint example (2 KB) and the Lab 2 starter code (66 KB), not a single starter ZIP; Lab 1 is a dataset-contribution task (write problems and specs), not an agent; the Berkeley grading table lists a single "Lab" component.
\item Course number CS294/194-280 is confirmed on the Berkeley page. The MOOC page gives no course number, but the Lab FAQ title reads "CS 194/294 Advanced LLM Agents".
\item Lecture 10's title differs between pages: "Bridging Informal and Formal Mathematical Reasoning" (MOOC page, which matches the slide filename) and "Advanced topics in theorem proving" (Berkeley page).
\item Slide URLs come from the Berkeley page (\texttt{/\allowbreak{}adv-llm-agents/\allowbreak{}slides/\allowbreak{}}); the MOOC page links the same filenames under \texttt{/\allowbreak{}llm-agents-mooc/\allowbreak{}slides/\allowbreak{}}. Neither set was fetched one by one, so the links are as published and were not checked for liveness. The Lecture 5 filename has literal spaces, encoded here as \%20.
\item Readings: the Berkeley page gives titles, and the MOOC page gives URLs in the same order. They were paired by position, and the arXiv IDs match the titles as far as I know, but the abstracts were not fetched.
\item Lab specs were read through WebFetch, which displays text. The ZIPs were not opened. Neither instruction file gives a Lean version, a due date or a numeric pass threshold. The Berkeley page gives the lab window as released 3/31, due 4/28; the MOOC deadline was May 31, 2025. The Lab FAQ Google Doc did not render without sign-in.
\item Quiz links (12 Google Forms), the quiz archive doc and certificate forms are on the MOOC page and are not listed here.
\item No license is stated anywhere. Both sites are GitHub Pages builds, but no source repo is linked.
\end{itemize}}
\subsection{Stanford 329T}
{\small\begin{itemize}
\item Verification (2026-10-09): the base page \href{https://web.stanford.edu/class/cs329t/}{[web.stanford.edu]} shows no sessions and links only to the Autumn 2024 archive (and lists earlier offerings), which is why it can look like there is no Fall 2025 offering. The Fall 2025 schedule is on a separate page, \href{https://web.stanford.edu/class/cs329t/syllabus.html}{[web.stanford.edu]} , which the home page refers to but whose URL the fetched text did not show. The same syllabus is in the archive at cs329t.1262/syllabus.html. \href{https://web.stanford.edu/class/cs329t/schedule.html}{[web.stanford.edu]} returns 404. Fall 2025 is confirmed by the slide filenames ("CS 329T Fall 2025"), 2025-dated guest decks, "Dec 2 or 4, 2025" on the syllabus, and the HW1 header "CS329T: Trustworthy Machine Learning (Fall 2025)".
\item The full title "Trustworthy Machine Learning: Building and Evaluating Agentic Systems" comes from the HW1 handout. The website header says only "CS 329T". The home page mentions guest tutorials on "fairness, ML explainability, and privacy", which looks carried over from older offerings; the Fall 2025 schedule itself is agent-focused.
\item Counts: 18 class meetings in the table (10 lecture/guest sessions with slides, 1 guest session without slides, 7 student-presentation or project-fair sessions); 11 slide decks (session 8 has two); 1 homework; no course videos.
\item The HW1 notebook link was taken from the PDF's embedded hyperlink, because the PDF text says only "Here is the notebook". It points to the \texttt{josh/\allowbreak{}stanford-course} branch of truera/trulens, which returns 404. The main-branch file with the same name exists, but it is not confirmed to be the assigned version: it has no numbered "step 9" human-annotation step (its step 9 is TruLens tracing setup, and annotations appear unnumbered in the Example 1-4 sections as an "Annotated Logical Consistency" feedback), and it does not mention Stanford or CS329T.
\item Snowflake is not named in the course text; it appears only in the contact emails for Datta and Reini. TruLens is not mentioned on the home page; it appears in HW1.
\item The second Sep 30 reading ("What is your Agent's GPA?") has no URL on the syllabus and is not linked here.
\item Slide contents were not read (PDFs were checked only for loading), so per-session topics are paraphrased from the syllabus.
\end{itemize}}
\end{document}
